\documentclass[11pt]{book}

%---------------------------------------
% Packages
%---------------------------------------
% === Document Setup and Language ===
\usepackage[english]{babel}        % English language
\usepackage[utf8]{inputenc}        % Unicode support
\usepackage[T1]{fontenc}           % Font encoding
\usepackage{subfiles}              % Allow for subfiles which can be compiled standalone
\setcounter{secnumdepth}{2}

% === Math Packages ===
\usepackage{amsfonts, amsmath, amssymb, amsthm}  % AMS math packages
\usepackage{dsfont}                % For \mathds{1}
\usepackage{mathtools}             % Enhanced math tools
\usepackage{accents}               % Mathematical accents

% === Typography and Fonts ===
%\usepackage{mathptmx}              % Times New Roman font
\usepackage{newtxtext,newtxmath}   % Enhanced Times New Roman with math support
\usepackage{csquotes}              % Smart quotes and quotation management
\linespread{1.05}                  % Line spacing

% === Graphics and Diagrams ===
\usepackage{tikz-cd, tikz}         % Diagrams and drawings
\usetikzlibrary{
	positioning, 	% For relative positioning (below=of, right=of)
	calc,        	% For coordinate calculations (finding midpoints)
	fit,			% Needed for the bounding box (fit)
	patterns,
	decorations.markings,
	decorations.pathmorphing,	% Needed for squiggly arrows
	backgrounds
}
\usepackage{graphicx}              % Image handling
\usepackage{booktabs}              % \toprule, \midrule, \bottomrule for tables

% === Page Layout ===
\usepackage{geometry}
\geometry{left=1.3in,right=1.2in}
\setlength{\parskip}{\smallskipamount}
\setlength{\parindent}{0pt}

% === Lists and Enumeration ===
\usepackage{enumerate}             % Enhanced enumeration
\renewcommand{\labelenumi}{\normalfont{(\arabic{enumi})}} % Default to [(1)]
\usepackage{etoolbox}              % Programming tools for LaTeX

% === Hyperlinks and References ===
\usepackage[hyphens]{url}
\urlstyle{same}
\usepackage{bookmark}
\usepackage{hyperref}              % Must be loaded after most other packages
\hypersetup{
	colorlinks=true,
	linkcolor=black,
	filecolor=black,
	urlcolor=black,
	citecolor=black,
	bookmarksdepth=subsubsection
}

% === Bibliography ===
\usepackage[
backend=biber,
natbib,
style=alphabetic,
maxcitenames=4,
maxalphanames=4,
maxbibnames=99
]{biblatex}
\renewbibmacro{in:}{}              % Suppresses "In: " before journal names
\addbibresource{C:/Users/LocalAdmin/OneDrive/Documenten OneDrive/A - Regensburg/Research and notes/Main_Bibliography.bib}
%\addbibresource{Bibliography.bib}

% === Optional Settings (uncomment if needed) ===
% \setcounter{secnumdepth}{2}
% \setcounter{tocdepth}{2}

%---------------------------------------
% Gestalten
%---------------------------------------
\DeclareMathOperator{\fpqc}{fpqc}

%---------------------------------------
% Algebraic geometry / motivic cohomology
%---------------------------------------
\DeclareMathOperator{\KGL}{KGL}
\DeclareMathOperator{\Nis}{Nis}
\DeclareMathOperator{\Et}{Et}
\DeclareMathOperator{\cdh}{cdh}
\DeclareMathOperator{\vdim}{vdim}
\DeclareMathOperator{\Krdim}{Kr.dim}
\DeclareMathOperator{\qcqs}{qcqs}
\DeclareMathOperator{\Sm}{Sm}
\DeclareMathOperator{\SSeq}{SSeq}
\DeclareMathOperator{\Sym}{Sym}
\DeclareMathOperator{\KH}{KH}

%---------------------------------------
% Topos theory
%---------------------------------------
\newcommand{\ssquare}{\,\square\,}
\DeclareMathOperator{\Aff}{Aff}
\DeclareMathOperator{\Mono}{Mono}
\DeclareMathOperator{\Gerb}{Gerb}
\DeclareMathOperator{\hyp}{hyp}
\DeclareMathOperator{\mono}{mono}
\DeclareMathOperator{\cons}{cons}
\DeclareMathOperator{\wcons}{wcons}
\DeclareMathOperator{\quottext}{quot}
\DeclareMathOperator{\epi}{epi}
\DeclareMathOperator{\MonoCong}{MonoCong}
\DeclareMathOperator{\EpiCong}{EpiCong}
\DeclareMathOperator{\Cong}{Cong}
\DeclareMathOperator{\HypCong}{HypCong}
\DeclareMathOperator{\Mdl}{Mdl}
\DeclareMathOperator{\Acyc}{Acyc}
\DeclareMathOperator{\Sat}{Sat}
\DeclareMathOperator{\SSat}{SSat}
\DeclareMathOperator{\Conn}{Conn}
\DeclareMathOperator{\EffEpi}{EffEpi}
\DeclareMathOperator{\Epi}{Epi}
\DeclareMathOperator{\Cotr}{Cotr}
\DeclareMathOperator{\All}{All}
\DeclareMathOperator{\Pos}{Pos}
\DeclareMathOperator{\post}{post}
\DeclareMathOperator{\El}{El}
\DeclareMathOperator{\cocompl}{cocompl}
\newcommand{\Topos}{\catname{Topos}}
\newcommand{\Logos}{\catname{Logos}}
\newcommand{\bbCat}{\bbC \mathrm{at}}
\newcommand{\bbLog}{\bbL \mathrm{ogos}}
\newcommand{\bbTop}{\bbT \mathrm{opos}}
%\renewcommand{\top}{\mathrm{top}}
\DeclareMathOperator{\Geom}{Geom}
\DeclareMathOperator{\Pt}{Pt}
\DeclareMathOperator{\im}{im}				% Image
\DeclareMathOperator{\coim}{coim}			% Coimage
\DeclareMathOperator{\Act}{Act}
\DeclareMathOperator{\Bun}{Bun}
\DeclareMathOperator{\Site}{Site}
\DeclareMathOperator{\cont}{cont}
\DeclareMathOperator{\cocont}{cocont}
\DeclareMathOperator{\wccolim}{w.c.colim}

%---------------------------------------
% Miscellaneous
%---------------------------------------
\DeclareMathOperator{\Atl}{Atl}
\DeclareMathOperator{\DiffStk}{DiffStk}
\DeclareMathOperator{\SepStk}{SepStk}
\renewcommand{\H}{\mathrm{H}}
\newcommand{\SH}{\mathrm{SH}}
\DeclareMathOperator{\htp}{htp}
\DeclareMathOperator{\rep}{rep}
\DeclareMathOperator{\Loc}{Loc}
\DeclareMathOperator{\modtext}{mod}
\newcommand{\absmod}[1]{\abs{#1}_{\modtext}}
\DeclareMathOperator{\sep}{sep}
\DeclareMathOperator{\PB}{PB}
\DeclareMathOperator{\PF}{PF}
\DeclareMathOperator{\NerveDiff}{N^{\mathrm{Diff}}}
\newcommand{\Deltaalg}[1]{\Delta^{#1}_{\mathrm{alg}}}
%\newcommand{\Act}[1]{#1 \textup{Act}}
\DeclareMathOperator{\shear}{shear}
\DeclareMathOperator{\Sph}{Sph}
\DeclareMathOperator{\Sub}{Sub}
\DeclareMathOperator{\Open}{Open}
\DeclareMathOperator{\open}{open}
\DeclareMathOperator{\Corr}{Corr}
\DeclareMathOperator{\bbMap}{\mathbb{M} ap}
\DeclareMathOperator{\bbHom}{\mathbb{H} om}
\DeclareMathOperator{\gen}{gen}
\DeclareMathOperator{\Bgl}{B_{gl}}
\newcommand{\PrnBdl}[2]{\textup{PrnBdl}_{#1}(#2)}
\newcommand{\catnameb}[1]{{\normalfont\textbf{#1}}}
\DeclareMathOperator{\bLieGrpd}{\catnameb{LieGrpd}}
\newcommand{\cFun}{\Ff \textup{un}}
\DeclareMathOperator{\Tr}{Tr}
\DeclareMathOperator{\tr}{tr}
\DeclareMathOperator{\Shp}{Shp}
\DeclareMathOperator{\THH}{THH}

%---------------------------------------
% Higher Algebra
%---------------------------------------

\newcommand{\Assoc}{\mathcal{A}\mathit{ssoc}}
\newcommand{\Comm}{\mathcal{C}\!\,\mathit{omm}}
\newcommand{\Lie}{\mathcal{L}\mathit{ie}}
\newcommand{\Triv}{\mathcal{T}\!\!\mathit{riv}}
\newcommand{\oTens}{\mathcal{T}\!\!\mathit{ens}}
\newcommand{\oLMod}{\mathcal{L}\!\mathcal{M}\mathit{od}}
\newcommand{\oRMod}{\mathcal{R}\!\mathcal{M}\mathit{od}}
\newcommand{\oBMod}{\mathcal{B}\!\mathcal{M}\mathit{od}}
\newcommand{\oMod}{\mathcal{M}\mathit{od}}
\newcommand{\oAlg}{\mathcal{A}\mathit{lg}}
\newcommand{\oEnd}{\mathcal{E}\mathit{nd}}
\newcommand{\oCMon}{\mathcal{C}\!\mathcal{M}\mathit{on}}
\newcommand{\oCGrp}{\mathcal{C}\mathcal{G}\mathit{rp}}
\newcommand{\oSp}{\mathcal{S}\mathit{p}}
\newcommand{\Mult}{\mathcal{M}}
\DeclareMathOperator{\oCat}{\mathcal{C}\!\mathit{at}}
\DeclareMathOperator{\oFun}{\mathcal{F}\!\mathit{un}}
%\newcommand{\Mult}{\mathcal{M}\mathit{ult}}
\DeclareMathOperator{\Env}{Env}
\DeclareMathOperator{\Op}{Op}
\DeclareMathOperator{\AdTrip}{AdTrip}
\DeclareMathOperator{\rev}{rev}
\DeclareMathOperator{\tiso}{iso}
\DeclareMathOperator{\oTw}{\overline{\Tw}}
\DeclareMathOperator{\lax}{lax}
\DeclareMathOperator{\NTw}{NTw}
\newcommand{\OpCocart}{\Oo^{\amalg}}
\newcommand{\OpCart}{\Oo^{\times}}
\DeclareMathOperator{\Day}{Day}
\newcommand{\oDay}{\mathcal{D}\mathit{ay}}
\newcommand{\dlax}{\textup{-lax}}
\newcommand{\textred}{\textup{red}}
\newcommand{\exc}{\textup{exc}}
\DeclareMathOperator{\BiFun}{BiFun}
\DeclareMathOperator{\BV}{BV}
\renewcommand{\Bar}{\mathrm{Bar}}
\newcommand{\dnil}{\textup{-nil}}
\DeclareMathOperator{\Ext}{Ext}
\DeclareMathOperator{\Ex}{Ex}
\DeclareMathOperator{\Tor}{Tor}
\newcommand{\enr}{\textup{enr}}
\newcommand{\heart}{\heartsuit}
\newcommand{\dsm}{\textup{-sm}}
\DeclareMathOperator{\Free}{Free}

%---------------------------------------
% Stable Homotopy Theory
%---------------------------------------
\DeclareMathOperator{\ku}{ku}
\DeclareMathOperator{\KU}{KU}
\DeclareMathOperator{\MU}{MU}
\DeclareMathOperator{\MO}{MO}
\DeclareMathOperator{\MSO}{MSO}
\DeclareMathOperator{\BO}{BO}
\DeclareMathOperator{\bGr}{\mathbf{Gr}}
\DeclareMathOperator{\BUP}{BUP}
\DeclareMathOperator{\bBUP}{\mathbf{BUP}}
\DeclareMathOperator{\Th}{Th}
\DeclareMathOperator{\Pic}{Pic}
\DeclareMathOperator{\pic}{pic}
\DeclareMathOperator{\pre}{pre}
\DeclareMathOperator{\Line}{Line}
%\DeclareMathOperator{\Triv}{Triv}
\DeclareMathOperator{\PT}{PT}
\newcommand{\interiorsymb}[1]{ {\kern0pt#1}^{\mathrm{o}}}
\newcommand{\PSp}{\catname{PSp}}
\DeclareMathOperator{\Emb}{Emb}
\newcommand{\Ntop}{N^{\textup{top}}}
\newcommand{\Homtop}{\Hom^{\textup{top}}}
\DeclareMathOperator{\fw}{fw}
\DeclareMathOperator{\Shift}{Shift}
\DeclareMathOperator{\Der}{Der}

%---------------------------------------
% Category Theory Operators
%---------------------------------------

% Maps and functors
\DeclareMathOperator{\Hom}{Hom}
\DeclareMathOperator{\Aut}{Aut}
\DeclareMathOperator{\End}{End}
\DeclareMathOperator{\iHom}{\underline{\operatorname{Hom}}}			% Internal Hom
\DeclareMathOperator{\ihom}{\underline{\operatorname{hom}}}			% Internal Hom
\DeclareMathOperator{\uHom}{\underline{\operatorname{Hom}}}			% Internal Hom
\DeclareMathOperator{\Map}{Map}
\DeclareMathOperator{\map}{map}
\DeclareMathOperator{\Path}{Path}
\DeclareMathOperator{\Fun}{Fun}
\newcommand{\FunL}{\Fun^{\textup{L}}}
\newcommand{\FunR}{\Fun^{\textup R}}
\newcommand{\ulFun}{\ul{\Fun}}
\DeclareMathOperator{\Nat}{Nat}

% Fibrations and Related Concepts
\DeclareMathOperator{\Lift}{Lift}
\DeclareMathOperator{\lift}{lift}
\DeclareMathOperator{\cocart}{cocart}
\DeclareMathOperator{\CoCart}{CoCart}
\DeclareMathOperator{\Cocart}{Cocart}
\DeclareMathOperator{\isCocart}{\texttt{isCocart}}
\DeclareMathOperator{\Cart}{Cart}
\DeclareMathOperator{\RFib}{RFib}
\DeclareMathOperator{\LFib}{LFib}
\DeclareMathOperator{\St}{St}
\DeclareMathOperator{\Un}{Un}
\DeclareMathOperator{\cart}{cart}
\newcommand{\dcart}{\textup{-cart}}
\DeclareMathOperator{\cc}{cc}
\DeclareMathOperator{\ct}{ct}
\DeclareMathOperator{\Str}{Str}
\DeclareMathOperator{\el}{el}

%---------------------------------------
% Algebraic Geometry
%---------------------------------------
\DeclareMathOperator{\Spec}{Spec}
\DeclareMathOperator{\QCoh}{QCoh}
\DeclareMathOperator{\IndCoh}{IndCoh}
\DeclareMathOperator{\Perf}{Perf}
\DeclareMathOperator{\Proj}{Proj}
\DeclareMathOperator{\Zar}{Zar}
\DeclareMathOperator{\Sch}{Sch}
\newcommand{\ad}{\mathrm{ad}}
\newcommand{\corp}{\mathrm{corp}}
\newcommand{\loc}{\mathrm{loc}}
\newcommand{\mot}{\mathrm{mot}}
\newcommand{\ultra}{\mathrm{ultra}}

%---------------------------------------
% Other math-operators
%---------------------------------------
\DeclareMathOperator{\Ar}{Ar}
\DeclareMathOperator{\BC}{BC}
\DeclareMathOperator{\GL}{GL}
\DeclareMathOperator{\Gr}{Gr}
\DeclareMathOperator{\Ho}{Ho}
\let\Im\relax
\DeclareMathOperator{\Im}{Im}
\DeclareMathOperator{\Ind}{Ind}
\DeclareMathOperator{\Iso}{Iso}
\DeclareMathOperator{\LConst}{LConst}
\DeclareMathOperator{\Nm}{Nm}
\DeclareMathOperator{\Pro}{Pro}
\DeclareMathOperator{\Lan}{Lan}
\DeclareMathOperator{\Ran}{Ran}
\DeclareMathOperator{\Rex}{Rex}
\DeclareMathOperator{\PSh}{PSh}
\DeclareMathOperator{\Shape}{Shape}
\DeclareMathOperator{\Shv}{Shv}
\usepackage{relsize}
\newcommand{\simp}{\mathlarger{\boldsymbol{\Delta}}}
\newcommand{\bOmega}{{\boldsymbol{\Omega}}}
\DeclareMathOperator{\Span}{Span}
\newcommand{\SpanTwo}{{\textup{\textsc{Span}}}_2}
\newcommand{\SpanTwoP}[4]{\SpanTwo(#1,#2)_{#4,#3}}
\DeclareMathOperator{\NSpan}{NSpan}
\DeclareMathOperator{\Tw}{Tw}
\DeclareMathOperator{\LM}{LM}
\let\Re\relax
\DeclareMathOperator{\Re}{Re}
\DeclareMathOperator{\Fact}{Fact}

\DeclareMathOperator{\SD}{SD}				% Suave dual
\DeclareMathOperator{\PD}{PD}				% Prim dual
\DeclareMathOperator{\Suave}{Suave}			% Suave objects
\DeclareMathOperator{\Prim}{Prim}			% Prim objects

\DeclareMathOperator{\shift}{shift}
\DeclareMathOperator{\act}{act}
\DeclareMathOperator{\supp}{supp}
\DeclareMathOperator{\ex}{ex}
\DeclareMathOperator{\lex}{lex}
\DeclareMathOperator{\rex}{rex}
\DeclareMathOperator{\chain}{chain}
\DeclareMathOperator{\sd}{sd}
\DeclareMathOperator{\qc}{qc}
\DeclareMathOperator{\Poly}{Poly}
\DeclareMathOperator{\Mot}{Mot}

%---------------------------------------
% Abbreviations
%---------------------------------------

\newcommand{\can}{\mathrm{can}}             % Canonical
\newcommand{\ho}{\mathrm{ho}}               % Homotopy category
\newcommand{\nondeg}{\mathrm{nondeg}}       % Non-degenerate
\newcommand{\univ}{\mathrm{univ}}           % Universal
\newcommand{\prop}{\textup{pr}}
\newcommand{\inv}{\textup{inv}}
\DeclareMathOperator{\ac}{ac}				% Associated category
\DeclareMathOperator{\pb}{pb}				% Pullback

\DeclareMathOperator{\fgt}{fgt}				% Forgetful functor
\DeclareMathOperator{\id}{id}				% Identity
\DeclareMathOperator{\pr}{pr}				% Projection
\DeclareMathOperator{\incl}{incl}				% Inclusion
\DeclareMathOperator{\const}{const}			% Forgetful
\DeclareMathOperator{\sk}{sk}				% Skeleton
\DeclareMathOperator{\cosk}{cosk}			% Coskeleton
\DeclareMathOperator{\gl}{gl}               % Global
\DeclareMathOperator{\grp}{grp}             % Group completion
\DeclareMathOperator{\qis}{qis}             % Quasi-isomorphism
\DeclareMathOperator{\cyl}{cyl}             % Cylinder
\DeclareMathOperator{\surj}{surj}             % Surjective
\newcommand{\et}{\textup{ét}}				% Étale
\newcommand{\coh}{\mathrm{coh}}				% Coherent
\newcommand{\acc}{\mathrm{acc}}				% Accessible
\renewcommand{\small}{\mathrm{small}}

\DeclareMathOperator{\free}{free}
\DeclareMathOperator{\fin}{fin}
\DeclareMathOperator{\pt}{pt}
\DeclareMathOperator{\aug}{{aug}}
\DeclareMathOperator{\st}{st}

\DeclareMathOperator{\cofib}{cofib}
\DeclareMathOperator{\fib}{fib}
\DeclareMathOperator{\colim}{colim}
\let\lim\relax
\DeclareMathOperator{\lim}{lim}
\DeclareMathOperator{\coeq}{coeq}
\DeclareMathOperator{\coker}{coker}
\newcommand{\ev}{\textup{ev}}
\newcommand{\coev}{\textup{coev}}
\DeclareMathOperator{\res}{res}
\DeclareMathOperator{\ind}{ind}
\DeclareMathOperator{\coind}{coind}

%---------------------------------------
% Categories
%---------------------------------------
\newcommand{\catname}[1]{{\textup{#1}}}
\newcommand{\Set}{\catname{Set}}
\newcommand{\An}{\catname{An}}
\newcommand{\CondAn}{\catname{CondAn}}
\newcommand{\Ab}{\catname{Ab}}
\newcommand{\Alg}{\catname{Alg}}
\newcommand{\CAlg}{\catname{CAlg}}
\newcommand{\Cat}{\catname{Cat}}
\newcommand{\CAT}{\catname{CAT}}
\newcommand{\CRing}{\catname{CRing}}
\newcommand{\qCat}{\catname{qCat}}
\newcommand{\CGrp}{\catname{CGrp}}
\newcommand{\CMon}{\catname{CMon}}
\newcommand{\Fin}{\catname{Fin}}
\newcommand{\FinGrpd}{\catname{FinGrpd}}
\newcommand{\Grp}{\catname{Grp}}
\newcommand{\Grpd}{\catname{Grpd}}
\newcommand{\Mod}{\catname{Mod}}
\newcommand{\coMod}{\catname{coMod}}
\newcommand{\Mon}{\catname{Mon}}
\newcommand{\LMod}{\catname{LMod}}
\newcommand{\RMod}{\catname{RMod}}
\newcommand{\BMod}{\catname{BMod}}
\newcommand{\Rep}{\catname{Rep}}
\renewcommand{\Top}{\catname{Top}}
\newcommand{\TopGrpd}{\catname{TopGrpd}}
\newcommand{\bTopGrpd}{\mathbf{TopGrpd}}
\newcommand{\Spc}{\catname{Spc}}
\newcommand{\Sp}{\catname{Sp}}
\newcommand{\Vect}{\catname{Vect}}
\newcommand{\Orb}{\textup{Orb}}
\newcommand{\Orbit}{\textup{Orbit}}
\newcommand{\Glo}{\textup{Glo}}

\newcommand{\sAn}{\s\An}
\newcommand{\dAn}{\mathrm{d}\An}
\newcommand{\bVect}{\mathbf{V}\catname{ect}}
\newcommand{\LCHaus}{\catname{LCHaus}}
\newcommand{\CHaus}{\catname{CHaus}}
\newcommand{\ProFin}{\catname{ProFin}}
\newcommand{\AnRing}{\catname{AnRing}}
\newcommand{\AnStck}{\catname{AnStck}}
\newcommand{\Ring}{\catname{Ring}}
\newcommand{\twoCAlg}{\catname{2CAlg}}
\newcommand{\Idem}{\catname{Idem}}
\newcommand{\coIdem}{\catname{coIdem}}
\newcommand{\StRing}{\catname{StRing}}
\newcommand{\PreStRing}{\catname{PreStRing}}
\newcommand{\Gest}{\catname{Gest}}
\newcommand{\RingGest}{\catname{RingGest}}
\newcommand{\LieGrpd}{\catname{LieGrpd}}
\newcommand{\sTop}{\catname{sTop}}
\newcommand{\Diff}{\catname{Diff}}
\newcommand{\sSet}{\catname{sSet}}
\newcommand{\Poset}{\catname{Poset}}

\renewcommand{\Pr}{\textup{Pr}}
\newcommand{\PrL}{\textup{Pr}^{\textup{L}}}
\newcommand{\PrR}{\textup{Pr}^{\textup{R}}}

\def\ulSp{\smash{\ul{{\setbox0=\hbox{\Sp}\dp0=1pt \ht0=0pt \box0\relax}}}}
\def\ulSpc{\smash{\ul{{\setbox0=\hbox{\Spc}\dp0=1pt \ht0=0pt \box0\relax}}}}
\def\ulOrbSp{\smash{\ul{{\setbox0=\hbox{\Orb\Sp}\dp0=1pt \ht0=0pt \box0\relax}}}}
\def\ulOrbSpc{\smash{\ul{{\setbox0=\hbox{\Orb\Spc}\dp0=1pt \ht0=0pt \box0\relax}}}}

%---------------------------------------
% Miscellaneous
%---------------------------------------

\newcommand{\h}{\textup{h}}
\newcommand{\CW}{\mathrm{CW}}
\DeclareMathOperator{\hocolim}{hocolim}
\DeclareMathOperator{\holim}{holim}
\DeclareMathOperator{\Exc}{Exc}
\newcommand{\s}{\textup{s}}
\DeclareMathOperator{\Sing}{Sing}
\DeclareMathOperator{\Kan}{Kan}
\newcommand{\bKan}{\mathbf{Kan}}
\newcommand{\bqCat}{\mathbf{qCat}}
\DeclareMathOperator{\Ch}{\catname{Ch}}
\DeclareMathOperator{\inj}{inj}
\DeclareMathOperator{\all}{all}
\DeclareMathOperator{\isotext}{iso}
\DeclareMathOperator{\Cut}{Cut}
\DeclareMathOperator{\LCut}{LCut}
\DeclareMathOperator{\BCut}{BCut}
\DeclareMathOperator{\Red}{Red}
\DeclareMathOperator{\Segal}{Segal}
\DeclareMathOperator{\Seg}{Seg}
\DeclareMathOperator{\CSeg}{CSeg}
\DeclareMathOperator{\add}{add}
\DeclareMathOperator{\sadd}{sadd}
\newcommand{\Piinfty}[1]{\Pi_{\infty}(#1)}
\DeclareMathOperator{\oplax}{oplax}

\newcommand{\qednow}{\pushQED{\qed}\qedhere\popQED} % Use this 

\newcommand{\quadtext}[1]{\quad\textrm{#1}\quad}
\newcommand{\qquadtext}[1]{\qquad\textrm{#1}\qquad}
\newcommand{\qin}{\quad\in\quad}
\newcommand{\ul}[1]{\underline{#1}}
\newcommand{\abs}[1]{\lvert #1 \rvert}
\newcommand{\geom}[1]{\| #1 \|}
\newcommand{\floor}[1]{\lfloor #1 \rfloor}

\newcommand{\catop}{^{\textup{op}}}
\let\op\relax
\DeclareMathOperator{\op}{op}

\DeclareMathOperator{\unit}{\mathds{1}}    % unit for \otimes
\newcommand{\blank}{{\textup{ --}}}
\newcommand{\ceil}[1]{\lceil #1 \rceil}
\newcommand{\lra}[1]{\langle #1 \rangle}

\newcommand{\iso}{\xrightarrow{\;\smash{\raisebox{-0.5ex}{\ensuremath{\scriptstyle\sim}}}\;}}
% \newcommand{\mapsfrom}{\mathrel{\reflectbox{$\mapsto$}}}
\newcommand{\xrightrightarrows}[2]{\begin{tikzcd}[ampersand replacement=\&, cramped] \vphantom{x}\rar[shift left=0.45ex, "#1"] \rar[shift right=0.45ex, swap, "#2"]\& \vphantom{x} \end{tikzcd}}
\newcommand{\xadjto}[2]{\begin{tikzcd}[ampersand replacement=\&, cramped] \vphantom{x}\rar[shift left=0.45ex, "#1"]\& \lar[shift left=0.45ex, "#2"] \vphantom{x} \end{tikzcd}}

%---------------------------------------
% Alphabets
%---------------------------------------
\let\S\relax
\let\P\relax
\DeclareMathOperator{\C}{\mathbb C}
\DeclareMathOperator{\F}{\mathbb F}
\DeclareMathOperator{\N}{\mathbb N}
\DeclareMathOperator{\P}{\mathbb P}
\DeclareMathOperator{\Q}{\mathbb Q}
\DeclareMathOperator{\R}{\mathbb R}
\DeclareMathOperator{\S}{\mathbb S}
\DeclareMathOperator{\Z}{\mathbb Z}
\DeclareMathOperator{\D}{\mathrm{D}}
\renewcommand{\phi}{\varphi}
\renewcommand{\epsilon}{\varepsilon}
\renewcommand{\aa}{\mathfrak{a}}
\newcommand{\pp}{\mathfrak{p}}
\newcommand{\fa}{\mathfrak{a}}
\newcommand{\fb}{\mathfrak{b}}
\newcommand{\fc}{\mathfrak{c}}
\newcommand{\fm}{\mathfrak{m}}
\newcommand{\mm}{\mathfrak{m}}
\newcommand{\fn}{\mathfrak{n}}
\newcommand{\fp}{\mathfrak{p}}

% Caligraphic font
\newcommand{\Aa}{\mathcal{A}}
\newcommand{\Bb}{\mathcal{B}}
\newcommand{\Cc}{\mathcal{C}}
\newcommand{\Dd}{\mathcal{D}}
\newcommand{\Ee}{\mathcal{E}}
\newcommand{\Ff}{\mathcal{F}}
\newcommand{\Gg}{\mathcal{G}}
\newcommand{\Hh}{\mathcal{H}}
\newcommand{\Ii}{\mathcal{I}}
\newcommand{\Jj}{\mathcal{J}}
\newcommand{\Kk}{\mathcal{K}}
\newcommand{\Ll}{\mathcal{L}}
\newcommand{\Mm}{\mathcal{M}}
\newcommand{\Nn}{\mathcal{N}}
\newcommand{\Oo}{\mathcal{O}}
\newcommand{\Pp}{\mathcal{P}}
\newcommand{\Qq}{\mathcal{Q}}
\newcommand{\Rr}{\mathcal{R}}
\newcommand{\Ss}{\mathcal{S}}
\newcommand{\Tt}{\mathcal{T}}
\newcommand{\Uu}{\mathcal{U}}
\newcommand{\Vv}{\mathcal{V}}
\newcommand{\Ww}{\mathcal{W}}
\newcommand{\Xx}{\mathcal{X}}
\newcommand{\Yy}{\mathcal{Y}}
\newcommand{\Zz}{\mathcal{Z}}

% Bold font
\newcommand{\bA}{\mathbf{A}}
\newcommand{\bB}{\mathbf{B}}
\newcommand{\bC}{\mathbf{C}}
\newcommand{\bD}{\mathbf{D}}
\newcommand{\bE}{\mathbf{E}}
\newcommand{\bF}{\mathbf{F}}
\newcommand{\bG}{\mathbf{G}}
\newcommand{\bH}{\mathbf{H}}
\newcommand{\bI}{\mathbf{I}}
\newcommand{\bJ}{\mathbf{J}}
\newcommand{\bK}{\mathbf{K}}
\newcommand{\bL}{\mathbf{L}}
\newcommand{\bM}{\mathbf{M}}
\newcommand{\bN}{\mathbf{N}}
\newcommand{\bO}{\mathbf{O}}
\newcommand{\bP}{\mathbf{P}}
\newcommand{\bQ}{\mathbf{Q}}
\newcommand{\bR}{\mathbf{R}}
\newcommand{\bS}{\mathbf{S}}
\newcommand{\bT}{\mathbf{T}}
\newcommand{\bU}{\mathbf{U}}
\newcommand{\bV}{\mathbf{V}}
\newcommand{\bW}{\mathbf{W}}
\newcommand{\bX}{\mathbf{X}}
\newcommand{\bY}{\mathbf{Y}}
\newcommand{\bZ}{\mathbf{Z}}

% Blackboard bold
\newcommand{\bbA}{\mathbb{A}}
\newcommand{\bbB}{\mathbb{B}}
\newcommand{\bbC}{\mathbb{C}}
\newcommand{\bbD}{\mathbb{D}}
\newcommand{\bbE}{\mathbb{E}}
\newcommand{\bbF}{\mathbb{F}}
\newcommand{\bbG}{\mathbb{G}}
\newcommand{\bbH}{\mathbb{H}}
\newcommand{\bbI}{\mathbb{I}}
\newcommand{\bbJ}{\mathbb{J}}
\newcommand{\bbK}{\mathbb{K}}
\newcommand{\bbL}{\mathbb{L}}
\newcommand{\bbM}{\mathbb{M}}
\newcommand{\bbN}{\mathbb{N}}
\newcommand{\bbO}{\mathbb{O}}
\newcommand{\bbP}{\mathbb{P}}
\newcommand{\bbQ}{\mathbb{Q}}
\newcommand{\bbR}{\mathbb{R}}
\newcommand{\bbS}{\mathbb{S}}
\newcommand{\bbT}{\mathbb{T}}
\newcommand{\bbU}{\mathbb{U}}
\newcommand{\bbV}{\mathbb{V}}
\newcommand{\bbW}{\mathbb{W}}
\newcommand{\bbX}{\mathbb{X}}
\newcommand{\bbY}{\mathbb{Y}}
\newcommand{\bbZ}{\mathbb{Z}}

%---------------------------------------
% Environments
%---------------------------------------

% Theorem-style environments
\theoremstyle{plain}
\newtheorem{theorem}{Theorem}[chapter]

\newtheorem{conjecture}[theorem]{Conjecture}
\newtheorem{corollary}[theorem]{Corollary}
\newtheorem{lemma}[theorem]{Lemma}
\newtheorem{proposition}[theorem]{Proposition}
\newtheorem{thmx}{Theorem}
\renewcommand{\thethmx}{\Alph{thmx}}

% Definition-style environments
\theoremstyle{definition}
\newtheorem{axiom}{Axiom}
\newtheorem{claim}[theorem]{Claim}
\newtheorem{construction}[theorem]{Construction}
\newtheorem{convention}[theorem]{Convention}
\newtheorem{definition}[theorem]{Definition}
\newtheorem{discussion}[theorem]{Discussion}
\newtheorem{example}[theorem]{Example}
\newtheorem{exercise}[theorem]{Exercise}
\newtheorem{fact}[theorem]{Fact}
\newtheorem{notation}[theorem]{Notation}
\newtheorem{observation}[theorem]{Observation}
\newtheorem{question}[theorem]{Question}
\newtheorem{remark}[theorem]{Remark}
\newtheorem{terminology}[theorem]{Terminology}
\newtheorem{warning}[theorem]{Warning}

% Unnumbered environments
\newtheorem*{remark*}{Remark}
\newtheorem*{example*}{Example}
\newtheorem*{theorem*}{Theorem}
\newtheorem*{recognition*}{Recognition Principle for Infinite Loop Spaces}
\newtheorem*{claim*}{Claim}
\newtheorem*{convention*}{Convention}
\newtheorem*{definition*}{Definition}
\newtheorem*{exercise*}{Exercise}
\newtheorem*{fact*}{Fact}
\newtheorem*{question*}{Question}
\newtheorem*{warning*}{Warning}
\renewcommand{\theaxiom}{\Alph{axiom}}
\newtheorem{subaxiom}{Axiom}
\renewcommand{\thesubaxiom}{\Alph{axiom}.\arabic{subaxiom}}
\renewcommand{\theHsubaxiom}{\theaxiom.\thesubaxiom}  % Makes sure the hyperlinks remember in which axiom the subaxiom appears
\newtheorem{subaxiomprime}{Axiom}
\renewcommand{\thesubaxiomprime}{\Alph{axiom}.\arabic{subaxiomprime}'}
\renewcommand{\theHsubaxiomprime}{\theaxiom.\thesubaxiomprime}  % Makes sure the hyperlinks remember in which axiom the subaxiom appears

%---------------------------------------
% Margin notes
%---------------------------------------
\newcommand{\red}[1]{{\color{red}{\textbf #1}}}
\setlength{\parindent}{0cm}
\setlength{\parskip}{0.8ex}
\setlength\marginparwidth{57pt}
\newcommand{\notehelper}[3]{\textcolor{#3}{$\blacksquare$}\marginpar{\ifodd\thepage\raggedright\else\raggedleft\fi\color{#3}\tiny \textbf{#2:} #1}}
\newcommand\bastiaan[1]{\notehelper{#1}{Bastiaan}{blue}}
\newcommand\bastiaanlong[1]{{\itshape \color{blue} #1}}

%---------------------------------------
% Tikzcd
%---------------------------------------

\tikzcdset{pullback/.style = {phantom, "\lrcorner"{very near start}}}
\tikzcdset{pullbackdl/.style = {phantom, "{\reflectbox{$\lrcorner$}}"{very near start}}}
\tikzcdset{pushout/.style = {phantom, "{\rotatebox{180}{$\lrcorner$}}"{very near end}}}

\tikzset{curve/.style={settings={#1},to path={(\tikztostart)
			.. controls ($(\tikztostart)!\pv{pos}!(\tikztotarget)!\pv{height}!270:(\tikztotarget)$)
			and ($(\tikztostart)!1-\pv{pos}!(\tikztotarget)!\pv{height}!270:(\tikztotarget)$)
			.. (\tikztotarget)\tikztonodes}}}
\tikzcdset{scale cd/.style={every label/.append style={scale=#1},cells={nodes={scale=#1}}}}
\newcommand{\pullback}{\mathrel{\text{\tikz{
				\draw[line width=0.6pt] (0,0) -- (0.5em,0);
				\draw[line width=0.6pt] (0.5em,0) -- (0.5em,0.6em);
		}}\,}}
\newcommand{\pullbackto}{\mathrel{\text{\tikz{
				\draw[line width=0.6pt] (0,0) -- (0.5em,0);
				\draw[line width=0.6pt] (0.5em,0) -- (0.5em,0.6em);
		}}\kern0.5em\to}}
\newcommand{\pushout}{\mathrel{\text{\tikz{
				\draw[line width=0.4pt] (0,0) -- (0,0.5em);
				\draw[line width=0.4pt] (0,0.5em) -- (0.5em,0.5em);
		}}\,}}
\newcommand{\pushoutto}{\mathrel{\text{\tikz{
				\draw[line width=0.4pt] (0,0) -- (0,0.5em);
				\draw[line width=0.4pt] (0,0.5em) -- (0.5em,0.5em);
		}}\kern0.5em\to}}

%---------------------------------------
% Other
%---------------------------------------

\newcommand{\quot}{/\!\!/}

% Use this when you want to write a \colon but the arrow goes from right to left, as in $Y \leftarrow X\noloc f$.
\newcommand\noloc{%
	\nobreak
	\mspace{6mu plus 1mu}
	{:}
	\nonscript\mkern-\thinmuskip
	\mathpunct{}
	\mspace{2mu}
}

\newcommand{\pvto}[2]{{%
		\kern1pt
		\big(
		\begin{matrix}
			\text{\tiny $#1$} \\[-7pt]
			\scriptstyle\downarrow \\[-7pt]
			\text{\tiny $#2$}
		\end{matrix}
		\big)
}}
\newcommand{\puto}[3]{{%
		\kern1pt
		\big(
		\begin{matrix}
			\text{\tiny $#1$ \hspace{-6pt}\phantom{$#3$}} \\[-9pt]
			\scriptstyle\downarrow \raisebox{2pt}{\text{\tiny $#3$}} \\[-9pt]
			\text{\tiny $#2$ \hspace{-6pt}\phantom{$#3$}}
		\end{matrix}
		\big)
}}
\newcommand{\apvto}[2]{\raisebox{0pt}[1.3\ht\strutbox][1.3\dp\strutbox]{$\smash{\pvto{#1}{#2}}$}}
\newcommand{\pwto}[2]{{%
		\kern1pt
		\big(
		\begin{matrix}
			\text{\tiny $#1$} \\[-5pt]
			\text{\tiny $#2$}
		\end{matrix}
		\big)
}}
\newcommand{\rbox}[3]{\raisebox{0pt}[#1\ht\strutbox][#2\dp\strutbox]{$#3$}}

%---------------------------------------
% Cleverref
%---------------------------------------
\usepackage[capitalise]{cleveref}  % Smart cross-referencing

\crefname{axiom}{axiom}{axioms}
\crefname{claim}{claim}{claims}
\crefname{conjecture}{conjecture}{conjectures}
\crefname{construction}{construction}{constructions}
\crefname{convention}{convention}{conventions}
\crefname{corollary}{corollary}{corollaries}
\crefname{definition}{definition}{definitions}
\crefname{discussion}{discussion}{discussions}
\crefname{example}{example}{examples}
\crefname{exercise}{exercise}{exercises}
\crefname{fact}{fact}{facts}
\crefname{lemma}{lemma}{lemmas}
\crefname{notation}{notation}{notations}
\crefname{observation}{observation}{observations}
\crefname{proposition}{proposition}{propositions}
\crefname{question}{question}{questions}
\crefname{remark}{remark}{remarks}
\crefname{subaxiom}{axiom}{axioms}
\crefname{terminology}{terminology}{terminologies}
\crefname{theorem}{theorem}{theorems}
\crefname{warning}{warning}{warnings}

\Crefname{axiom}{Axiom}{Axioms}
\Crefname{claim}{Claim}{Claims}
\Crefname{conjecture}{Conjecture}{Conjectures}
\Crefname{construction}{Construction}{Constructions}
\Crefname{convention}{Convention}{Conventions}
\Crefname{corollary}{Corollary}{Corollaries}
\Crefname{definition}{Definition}{Definitions}
\Crefname{discussion}{Discussion}{Discussions}
\Crefname{example}{Example}{Examples}
\Crefname{exercise}{Exercise}{Exercises}
\Crefname{fact}{Fact}{Facts}
\Crefname{lemma}{Lemma}{Lemmas}
\Crefname{notation}{Notation}{Notations}
\Crefname{observation}{Observation}{Observations}
\Crefname{proposition}{Proposition}{Propositions}
\Crefname{question}{Question}{Questions}
\Crefname{remark}{Remark}{Remarks}
\Crefname{subaxiom}{Axiom}{Axioms}
\Crefname{terminology}{Terminology}{Terminologies}
\Crefname{theorem}{Theorem}{Theorems}
\Crefname{warning}{Warning}{Warnings}

% Appendix chapters are labelled \label[appendix]{...} so that \Cref calls them
% ``Appendix'' rather than ``Chapter'' (main-matter chapters keep ``Chapter'').
\crefname{appendix}{appendix}{appendices}
\Crefname{appendix}{Appendix}{Appendices}


\usepackage{float}
\usepackage{etoolbox}
\usepackage[T1]{fontenc}
\usepackage[utf8]{inputenc}

\begin{document}
	\pagenumbering{roman}

	\title{Gestalten\\ \vspace{10pt}
	\Large University of Regensburg, Summer term 2026}
	\author{\large Organizers: Bastiaan Cnossen and Sil Linskens \\ Note-taker: Bastiaan Cnossen}

	\maketitle

	%	\newgeometry{top=1in,bottom=1in,left=1.5in,right=1.5in,headsep=1.5cm,marginparwidth=1.25in}
	\newgeometry{left=1.3in,right=1.2in, top=3.8cm, headsep=1.3cm}


	\addto\captionsenglish{\renewcommand{\chaptername}{Talk}}

	\begin{KeepFromToc}
		\tableofcontents
	\end{KeepFromToc}

	\titlespacing*{\chapter}{0pt}{20pt}{15pt}
	
	\chapter*{Acknowledgments}
	% \addcontentsline{toc}{chapter}{Acknowledgments}
	% \addtocontents{toc}{\vskip 15pt}

	All mathematical material in these notes is taken from \citet{Scholze_Gestalten}, which is joint work with Stefanich. These notes are based on notes taken live during the talks, and we thank all participants of the seminar for their contributions. We also thank Peter Scholze for explaining the arguments in \Cref{chap:Compact_Objects} to us. 
	
	\begingroup
	\let\clearpage\relax
	\let\cleardoublepage\relax
	\chapter*{Use of generative AI}
	This document has been written with significant involvement of generative AI:
	\begin{itemize}
		\item The notes taken live during the seminar were turned into proper TeX by Claude Opus. This process also corrected various typos / unfinished sentences but otherwise did not greatly affect the mathematical contents.
		\item Certain material from \cite{Scholze_Gestalten} not covered during the seminar has been added to the notes by Claude Opus or Codex. This happened for the following chapters:
		\begin{itemize}
			\item The end of \Cref{chap:Properties_Of_Maps};
			\item The material at the end of \Cref{chap:Injections_Surjections_Gestalten};
			\item The last three sections of \Cref{chap:Transmutations};
			\item The material in \Cref{chap:Cartier_Duality}.
		\end{itemize}
		\item The background material on six-functor formalisms in \Cref{chap:Six_Functor_Formalisms} was originally added by Claude Opus, taken primarily from \cite{Heyer_Mann,CLL_Span2,CLL_Span}. It has subsequently been edited.
		\item Some mathematical content has been generated by Codex and Claude Fable:
		\begin{itemize}
			\item The content of \Cref{chap:Compact_Objects} was mostly written by Codex, though with significant prompting and manual editing;
			\item The content of the exploratory \Cref{chap:Zariski_Gestalten} was almost entirely written autonomously by Claude Fable (originally to test its capabilities, then added to the notes as it seems genuinely interesting).
		\end{itemize}
	\end{itemize}
	\endgroup

	\titlespacing*{\chapter}{0pt}{50pt}{40pt}

	\restoregeometry
	\openany

	\mainmatter
	\pagestyle{plain}
	\setcounter{secnumdepth}{2}

	\chapter{Presentable categories}
	\label[chapter]{chap:Presentable_Categories}

    \emph{Talk given by Sil.}

	\section{Preliminaries}

	Concepts we should know:
	\begin{enumerate}
		\item[(-2)] \emph{Regular cardinal $\kappa$}: any union of sets of cardinality $< \kappa$ indexed by a set of cardinality $< \kappa$ has cardinality $< \kappa$.
		\item[(-1)] \emph{Essentially $\kappa$-small ($\infty$-)category $C$}: (If $\kappa$ uncountable) The set of isomorphism classes of $C$ has cardinality $< \kappa$, and for all $x,y \in C$ the homotopy groups $\pi_i\hom_C(x,y)$ have cardinality $<\kappa$ (For $\kappa = \aleph_0$ one would take the smallest subcategory of $\Cat_{\infty}$ generated under finite colimits by $[0]$ and $[1]$);
		\item[(0)] \emph{$\kappa$-filtered category $I$}: Every $\kappa$-small diagram in $I$ admits a cocone;
		\item[(1)] \emph{$\kappa$-compact object $X \in C$}: $\Hom_C(X,-)\colon C \to \An$ preserves $\kappa$-filtered colimits;
		\item[(2)] \emph{$\kappa$-presentable categories}: A category $C$ which has small colimits and such that every object is a $\kappa$-filtered colimit of $\kappa$-compact objects;
		\item[(3)] \emph{Presentable categories}: A category $C$ which is $\kappa$-presentable for some $\kappa$.
	\end{enumerate}

	\begin{remark}
		$C$ is presentable iff it ``admits a presentation:'' There is a small category $S$ such that $C \simeq \lra{S}[R^{-1}]$, where $\lra{S} = \PSh(S) = \Fun(S\catop,\An)$ is the free cocompletion of $S$, where $R$ is a (small) set of morphisms in $\lra{S}$, and where $\lra{S}[R^{-1}] \subseteq \lra{S}$ is the full subcategory spanned by those $X \in \lra{S}$ such that $\Hom_{\lra{S}}(f,X)\colon \Hom_{\lra{S}}(B,X) \to \Hom_{\lra{S}}(A,X)$ is an isomorphism for every morphism $f\colon A \to B$ in $R$.

		Similarly, $C$ is $\kappa$-presentable if we can choose $S$ and $R$ above such that $R$ consists of maps between $\kappa$-compact objects.
	\end{remark}

	\begin{remark}
		The description of $\lra{S}[R^{-1}]$ as a full subcategory of $\lra{S}$ is an instance of the principle that ``free constructions in $\PrL$ correspond to cofree constructions in $\Cat$, via passing to right adjoints.'' Indeed, the localization functor $\lra{S} \to \lra{S}[R^{-1}]$ is a morphism in $\PrL$, hence preserves colimits and admits a right adjoint; this right adjoint is automatically fully faithful and identifies $\lra{S}[R^{-1}]$ with the subcategory described above. See \cite[Proposition 5.5.4.20]{lurie2009HTT}.
	\end{remark}

	\begin{definition}
		We denote by $\PrL$ the category of presentable categories and cocontinuous functors. We let $\PrL_{\kappa} \subseteq \PrL$ denote the (non-full) subcategory of $\kappa$-presentable categories and those cocontinuous functors which preserve $\kappa$-compact objects.
	\end{definition}

	\begin{example}
		The categories $\An$, $\Sp$, $D(R)$ are all in $\PrL_{\aleph_0}$. For $X$ a locally compact Hausdorff space, $\Shv(X) \in \PrL_{\aleph_1}$. Also the category of almost modules is $\aleph_1$-presentable.
	\end{example}

	\begin{lemma}
		\label[lemma]{lem:Functor_Category_Presentable}
		Let $\kappa$ be an uncountable regular cardinal. If $D$ is
		$\kappa$-presentable and $I$ is a $\kappa$-small category, then also
		$\Fun(I,D)$ is $\kappa$-presentable. Its $\kappa$-compact generators
		are given by $\Fun(I,D^{\kappa})$.
	\end{lemma}
	\begin{proof}
		This follows from the argument given in \cite[Proposition 5.4.4.3]{lurie2009HTT}; note that no
		enlargement of the regular cardinal is needed: since $I$ is
		$\kappa$-small, the argument only uses that $\kappa$-filtered colimits
		of anima commute with $\kappa$-small limits. This commutation holds for
		every regular cardinal; we assume $\kappa$ uncountable only to stay
		within our standing conventions, under which the $\kappa$-compact objects
		of $\Fun(I,D)$ are genuine $\kappa$-small colimits of objects of
		$\Fun(I,D^{\kappa})$ rather than mere retracts thereof, so that no
		separate idempotent-completion step (as would be required for
		$\kappa = \aleph_0$) intervenes.
	\end{proof}

	\begin{definition}
		Let $\Rex_{\kappa}$ be the category of categories with $\kappa$-small colimits. (In the case $\kappa = \aleph_0$ we assume the categories are additionally idempotent complete).
	\end{definition}

	\begin{proposition}
		Let $C$ be a $\kappa$-presentable category. Then the subcategory $C^{\kappa} \subseteq C$ is essentially small, and the inclusion $C^{\kappa} \hookrightarrow C$ uniquely extends to an equivalence $\Ind_{\kappa}(C^{\kappa}) \iso C$. Moreover, the assignments $C \mapsto C^{\kappa}$ and $D \mapsto \Ind_{\kappa}(D)$ define an equivalence
		\[
			\Rex_{\kappa} \rightleftarrows \PrL_{\kappa}.
		\]
	\end{proposition}

	\begin{theorem}
		\label[theorem]{thm:PrL_Kappa_Presenable}
		The category $\PrL_{\kappa}$ is $\kappa$-presentable. (In other words: $\PrL_{\kappa} \in \PrL_{\kappa}$.)
	\end{theorem}

	\begin{proof}
		By the previous proposition, $\PrL_{\kappa}$ is equivalent to $\Rex_{\kappa}$, so we need to show $\Rex_{\kappa}$ is $\kappa$-presentable. To this end, we consider the free-forgetful adjunction
		\[
			\lra{-}^{\kappa}\colon \Cat \rightleftarrows \Rex_{\kappa}\colon \fgt.
		\]
		We will first show that this adjunction is monadic. By the Barr--Beck theorem, we need to show:
		\begin{itemize}
			\item The functor $\fgt$ is conservative. This is clear.
			\item The category $\Rex_{\kappa}$ admits $\fgt$-split simplicial colimits, and $\fgt\colon \Rex_{\kappa} \to \Cat$ preserves these. Here we say that a functor $C_{\bullet}\colon \simp\catop \to \Rex_{\kappa}$ is \emph{$\fgt$-split} if its image $\fgt(C_{\bullet})\colon \simp\catop \to \Cat$ admits a splitting.
		\end{itemize}
		For the second claim, consider such a diagram $C_{\bullet}$, and let $C_{-1} := \colim_{[n] \in \simp\catop} C_n$ be its colimit in $\Cat$. Since split simplicial colimits are absolute, the colimit functors for $\kappa$-small diagrams in the $C_n$ descend to colimit functors in $C_{-1}$. Thus $C_{-1}$ admits $\kappa$-small colimits, and each structure map $\iota_n \colon C_n \to C_{-1}$ preserves them. We claim that the augmented diagram is a colimit in $\Rex_{\kappa}$; that is, for every $D \in \Rex_{\kappa}$ the canonical map
		\[
			\Hom_{\Rex_{\kappa}}(C_{-1}, D) \to \lim_{[n] \in \simp} \Hom_{\Rex_{\kappa}}(C_n,D)
		\]
		is an equivalence. Since $C_{-1}$ is the colimit in $\Cat$, we have
		\[
			\Fun(C_{-1},D) \iso \lim_{[n] \in \simp} \Fun(C_n,D).
		\]
		It remains only to identify the colimit-preserving functors on both sides. This property is detected after precomposition with $\iota_0$: one direction uses that $\iota_0$ preserves $\kappa$-small colimits, while the converse uses a splitting $s \colon C_{-1} \to C_0$ with $\iota_0s \simeq \id_{C_{-1}}$ to write every $\kappa$-small diagram in $C_{-1}$ as the image under $\iota_0$ of a diagram in $C_0$. Hence the displayed equivalence restricts to the desired equivalence on the full subcategories of $\kappa$-small-colimit-preserving functors.

		We now know that $\Rex_{\kappa}$ is monadic over $\Cat$. Moreover, $\Cat$ is $\aleph_0$-presentable, as it is a Bousfield localization of $\sAn = \Fun(\simp\catop,\An)$ at a small set of morphisms between compact objects. The free $\kappa$-cocompletion monad is $\kappa$-accessible, so $\Rex_{\kappa}$ admits all colimits and is generated by the objects $\lra{[n]}^{\kappa}$. These generators are $\kappa$-compact, because $[n]$ is $\kappa$-compact in $\Cat$ and a $\kappa$-small colimit in a $\kappa$-filtered colimit of $\kappa$-cocomplete categories is computed after some stage; equivalently, $\fgt\colon \Rex_{\kappa} \to \Cat$ preserves $\kappa$-filtered colimits.
	\end{proof}

	\begin{proposition}
		The object $\PrL_{\kappa}$ of $\PrL_{\kappa}$ is $\kappa$-compact if and only if $\kappa$ is not a limit cardinal. (In particular, this applies to $\kappa = \aleph_1$.)
	\end{proposition}

	\begin{lemma}\label[lemma]{lem:Fully_Faithful_On_Aleph1_Compacts}
		Let $F\colon C \to D$ be a morphism in $\PrL_{\kappa}$, and let $G$ denote its right adjoint. Then the following are equivalent:
		\begin{enumerate}[(1)]
			\item The functor $F$ is fully faithful;
			\item The induced functor
			\[
				F^{\kappa}\colon C^{\kappa} \to D^{\kappa}
			\]
			is fully faithful;
			\item The unit $\eta_X\colon X \to GFX$ is an isomorphism for every $X \in C^{\kappa}$.
		\end{enumerate}
	\end{lemma}
	\begin{proof}
		The implication $(1) \Rightarrow (3)$ is the usual adjunction criterion for a functor to be fully faithful. The implication $(3) \Rightarrow (2)$ follows from the adjunction: for $X,Y \in C^{\kappa}$, we have
		\[
			\Hom_D(FX,FY) \simeq \Hom_C(X,GFY) \simeq \Hom_C(X,Y).
		\]
		Conversely, assume that $F^{\kappa}$ is fully faithful. For $Y \in C^{\kappa}$, the map $\eta_Y\colon Y \to GFY$ induces, for every $X \in C^{\kappa}$, the map
		\[
			\Hom_C(X,Y) \to \Hom_C(X,GFY) \simeq \Hom_D(FX,FY),
		\]
		which is an isomorphism by assumption. Since the $\kappa$-compact objects generate $C$ under colimits, this implies that $\eta_Y$ is an isomorphism. Thus $(2) \Rightarrow (3)$.

		It remains to see that $(2)$ implies $(1)$. Consider the subcategory $C' \subseteq C$ consisting of those objects $X$ such that the map $\Hom_C(X,Y) \to \Hom_D(FX,FY)$ is an isomorphism for each $Y \in C$. Note that we have $C^{\kappa} \subseteq C'$: if $X$ is $\kappa$-compact, then both $\Hom_C(X,-)$ and $\Hom_D(FX,F(-))$ preserve $\kappa$-filtered colimits, and every $Y \in C$ may be written as a $\kappa$-filtered colimit of $\kappa$-compact objects. Also note that $C'$ is closed under arbitrary colimits, since both $\Hom_C(-,Y)$ and $\Hom_D(F(-),FY)$ send colimits in $X$ to limits. It follows that $C' = C$, showing that $F$ is fully faithful.
	\end{proof}

	\section{Basic properties of \texorpdfstring{$\PrL$}{Pr\textasciicircum L}}

	\begin{theorem}
		Let $F\colon C \to D$ be a functor between $C,D \in \PrL$.
		\begin{enumerate}
			\item Then $F$ admits a right adjoint if and only if it preserves colimits.
			\item Moreover, $F$ admits a left adjoint if and only if it preserves limits and it is $\kappa$-accessible for some regular cardinal $\kappa$ (i.e.\ it preserves $\kappa$-filtered colimits).
		\end{enumerate}
	\end{theorem}

	\begin{corollary}
		There is an equivalence $(\PrL)\catop \simeq \PrR$.
	\end{corollary}

	\begin{proposition}
		\begin{enumerate}
			\item The category $\PrL$ admits small limits, which are preserved by $\fgt\colon \PrL \to \Cat$.
			\item The category $\PrL_{\kappa}$ admits limits, preserved by the composite
			\[
				\PrL_{\kappa} \simeq \Rex_{\kappa} \to \Cat.
			\]
			\item If $\kappa$ is uncountable, then the inclusion $\PrL_{\kappa} \hookrightarrow \PrL$ preserves $\kappa$-small limits.
		\end{enumerate}
	\end{proposition}

	\begin{proof}
		For (1) and (3) we refer to Lurie's HTT. For (2), one needs to use that $\Rex_{\kappa} \to \Cat$ is monadic.
	\end{proof}

	\begin{proposition}
		The categories $\PrR$ and $\PrR_{\kappa}$ admit small limits, and the functors $\PrR_{\kappa} \to \PrR \to \Cat$ preserve small limits.
	\end{proposition}

	\begin{corollary}
		The categories $\PrL$ and $\PrL_{\kappa}$ admit colimits (which are computed by passing to adjoints and computing the limit in $\PrR$ resp.\ $\PrR_{\kappa}$.)
	\end{corollary}

	\begin{example}
		Let $I$ be a set.
		\begin{enumerate}
			\item The product $\prod_{i \in I} C_i$ in $\PrL$ is just the product of underlying categories. But the product in $\PrL_{\kappa}$ would be $\Ind_{\kappa}(\prod_{i \in I} C_i^{\kappa})$.
			\item The coproduct $\coprod_{i \in I} C_i$ in $\PrL$ is also the product $\prod_{i \in I} C_i$ in $\Cat$. So is the coproduct in $\PrL_{\kappa}$.
			\item The pushout of $*/H$ and $*/G$ over $*$ is $*/(G \star H)$, where $G \star H$ is the free product of $G$ and $H$, and where $*/G$ denotes the classifying anima of $G$. But the pushout $\lra{*/H} \sqcup_{\lra{*}} \lra{*/G}$ is given by the pullback $\lra{*/H} \times_{\An} \lra{*/G}$ in $\Cat$.
		\end{enumerate}
	\end{example}

	\begin{remark}
		Item (2) of the previous example is striking: the natural map
		\[
			\coprod_{i \in I} C_i \to \prod_{i \in I} C_i
		\]
		in $\PrL$ is an isomorphism for \emph{arbitrary} indexing sets $I$. Adapting the notion of \emph{semiadditivity} from \cite[Definition 6.1.6.13]{lurie2016HA} to arbitrary cardinals, this says that $\PrL$ is $\kappa$-semiadditive for every cardinal $\kappa$. This is an extremely unusual property of $\PrL$.

		The analogous statement in $\PrL_{\kappa}$ requires $|I| < \kappa$: the coproduct $\coprod_{i \in I}^{\PrL_{\kappa}} C_i$ is always the naive product $\prod_i C_i$, whereas the product $\prod_{i \in I}^{\PrL_{\kappa}} C_i = \Ind_{\kappa}(\prod_i C_i^{\kappa})$ generally differs from it; the two agree precisely when $|I| < \kappa$.
	\end{remark}

	\section{Tensor products in \texorpdfstring{$\PrL$}{Pr\textasciicircum L}}

	\begin{proposition}
		Let $C,D \in \PrL$. There exists a universal presentable category $C \otimes D$ with a functor $C \times D \to C \otimes D$ which preserves colimits in both variables. If $C,D \in \PrL_{\kappa}$, then $C \otimes D \in \PrL_{\kappa}$, and the functor $C \times D \to C \otimes D$ preserves $\kappa$-compact objects.
	\end{proposition}

	\begin{proof}
		We may set $\lra{S} \otimes \lra{T} = \lra{S \times T}$. Moreover, we may set $\lra{S}[R^{-1}] \otimes \lra{T}[Q^{-1}] = \lra{S \times T}[W^{-1}]$, where $W$ consists of the morphisms $f \otimes \id_t$ for $f \in R$ and $t \in T$, together with $\id_s \otimes g$ for $g \in Q$ and $s \in S$ (here $s$ and $t$ range over the representables).
	\end{proof}

	\begin{corollary}
		The categories $\PrL$ and $\PrL_{\kappa}$ admit symmetric monoidal structures $\Pr^{L,\otimes}$, $\Pr^{L,\otimes}_{\kappa}$.
	\end{corollary}

	\begin{proof}
		We define these as suboperads of $\CAT^{\times}$.
	\end{proof}

	\begin{example}
		\begin{enumerate}
			\item For $C = \lra{S}$, we get $C \otimes D \simeq \Fun(S\catop,D)$.
			\item For $C \in \PrL_{\kappa}$ and $D \in \PrL$, we have $C \otimes D \simeq \Fun^{\kappa\text{-}\lex}((C^{\kappa})\catop,D)$.
			\item Let $A$ be a (static) ring. Then $D_{\geq 0}(A)$ has compact projective generators $A^n$ for $n \geq 0$, so $D_{\geq 0}(A) \simeq \Fun^{\times}(\{A^n\}\catop,\An)$. It follows that
			\[
				D_{\geq 0}(A) \otimes C \simeq \Fun^{\times}(\{A^n\}\catop,C).
			\]
			(Objects of the RHS are also known as \emph{$A$-module objects in $C$}.)
		\end{enumerate}
	\end{example}

	\begin{proposition}
		The symmetric monoidal category $\Pr^{L,\otimes}$ is closed monoidal, with internal hom given by
		\[
			[C,D] = \FunL(C,D),
		\]
		the full subcategory of $\Fun(C,D)$ spanned by the colimit-preserving functors.

		Moreover, also $\Pr^{L,\otimes}_{\kappa}$ is closed monoidal, with internal hom
		\[
			[C,D] = \Ind_{\kappa}(\Fun^{\kappa\text{-}\rex}(C^{\kappa},D^{\kappa})).
		\]
	\end{proposition}

	\begin{proof}
		A colimit-preserving functor $E \to \FunL(C,D)$ is the same as a functor $E \times C \to D$ which preserves colimits in both variables, which by universality is the same as a colimit-preserving functor $E \otimes C \to D$. It remains to show that $\FunL(C,D)$ is presentable. Since $\kappa$-small limits in $\PrL$ and $\PrL_{\kappa}$ agree, we may reduce to the case $C = \lra{[n]}^{\kappa}$. Unwinding the statement in this case, we must prove that for every $D \in \PrL_{\kappa}$, the functor
		\[
			\Ind_{\kappa}(\Fun([n],D^{\kappa})) \to \Fun([n], \Ind_{\kappa}(D^{\kappa}))
		\]
		is an equivalence. This is HTT.5.3.5.15.
	\end{proof}

	\begin{proposition}
		We have $(\PrL_{\kappa})^{\otimes} \in \CAlg(\PrL_{\kappa})$.
	\end{proposition}

	\begin{proof}
		Since $(\PrL_{\kappa})^{\otimes} \simeq \Rex_{\kappa}^{\otimes}$, it suffices to prove the claim for the latter. The only thing that remains is to show that the functor
		\[
			- \otimes -\colon \Rex_{\kappa} \times \Rex_{\kappa} \to \Rex_{\kappa}
		\]
		preserves $\kappa$-compact objects. This in turn reduces to the claim that the tensor product of $\lra{[n]}^{\kappa}$ and $\lra{[m]}^{\kappa}$ is $\kappa$-compact. But this tensor product is $\lra{[n] \times [m]}^{\kappa}$, which is indeed $\kappa$-compact.
	\end{proof}

	\section{\texorpdfstring{$C$}{C}-linear presentable categories}

	Consider $C \in \CAlg(\PrL_{\kappa})$. For $A \in \CAlg(C)$, we may consider the module category $\Mod_A(C) \in \CAlg(\PrL_{\kappa})$. Moreover, the forgetful functor $\Mod_A(C) \to C$ admits a left adjoint
	\[
		\free\colon C \to \Mod_A(C), \qquad X \mapsto A \otimes X,
	\]
	which defines a functor in $\CAlg(\PrL_{\kappa})$.

	Given two $A$-modules $M,M'$, we obtain a \emph{relative tensor product over $A$}
	\[
		M \otimes_A M' = \colim_{\simp\catop}\left(\begin{tikzcd}[ampersand replacement=\&, cramped, arrow style=tikz]
			\dots \ar[r, shift left=3] \ar[r, shift left=1] \ar[r, shift left=-1] \ar[r, shift left=-3] \& M \otimes A \otimes A \otimes M' \ar[r, shift left=2] \ar[r] \ar[r, shift left=-2] \& M \otimes A \otimes M' \ar[r, shift left] \ar[r, shift right] \& M \otimes M'
		\end{tikzcd}\right).
	\]

	\begin{definition}
		Let $C \in \CAlg(\PrL_{\kappa})$. Then we may form the module category $\Mod_C(\PrL_{\kappa})$ (where we are using that $\PrL_{\kappa} \in \CAlg(\PrL_{\kappa})$). We refer to the objects of $\Mod_C(\PrL_{\kappa})$ as \emph{$C$-linear $\kappa$-presentable categories} and to its morphisms as \emph{$C$-linear functors}.
	\end{definition}

	Given $D \in \CAlg(\PrL_{\kappa})$ and $B \in \CAlg(D)$, there is always an equivalence
	\[
		\CAlg(D)_{B/} \simeq \CAlg(\Mod_B(D)).
	\]

	\begin{theorem}
		\label[theorem]{thm:Embeddings_CAlg(C)_Into_CAlg(Mod_C)}
		Given $A \in \CAlg(C)$ as before, the relative tensor product defines a symmetric monoidal structure on $\Mod_A(C)$, turning the module category into an object
		\[
			\Mod_A(C) \in \CAlg(\Mod_C(\PrL_{\kappa})) \simeq \CAlg(\PrL_{\kappa})_{C/}.
		\]
		Moreover, the assignment $A \mapsto \Mod_A(C)$ defines a fully faithful functor
		\[
			\Mod_{(-)}(C)\colon \CAlg(C) \hookrightarrow \CAlg(\Mod_C(\PrL_{\kappa}))
		\]
		which is moreover a morphism in $\CAlg(\PrL_{\kappa})$. Its right adjoint
		\[
			\End_{(-)}(\unit) \colon \CAlg(\Mod_C(\PrL_{\kappa})) \to \CAlg(C)
		\]
		sends an object $D \in \CAlg(\PrL_{\kappa})_{C/}$ to the endomorphism object of its monoidal unit. This functor preserves $\kappa$-filtered colimits.
	\end{theorem}

	\begin{notation}\label[notation]{nota:Module_Endomorphism_Adjunction}
		For $C \in \CAlg(\PrL_{\kappa})$, we refer to the adjunction of \Cref{thm:Embeddings_CAlg(C)_Into_CAlg(Mod_C)} as the \emph{module-endomorphism adjunction}
		\[
			\begin{tikzcd}
				\CAlg(C) \rar[hook, shift left]{\Mod_{(-)}(C)} & \CAlg(\Mod_C(\PrL_{\kappa})), \lar[shift left]{\End_{(-)}(\unit)}
			\end{tikzcd}
		\]
		Thus $R \in \CAlg(C)$ is sent to $\Mod_R(C)$, and the right adjoint sends $D \in \CAlg(\Mod_C(\PrL_{\kappa}))$ to $\End_D(\unit_D) \in \CAlg(C)$.
	\end{notation}

	\begin{convention}
		From now on, we fix our cutoff cardinal to be $\kappa = \aleph_1$. This choice is somewhat arbitrary: we could have worked with any uncountable regular cardinal. The reason why we need uncountability is to make sure that various countable colimits of compact objects remain compact. For example, given two modules $M$ and $M'$ over an algebra $A$, the relative tensor product $M \otimes_A M'$ defined above is computed as the geometric realization of the bar complex, hence as a \emph{countable} colimit. Moreover, while $\kappa$-compact objects in a presheaf category are generally \emph{retracts} of $\kappa$-small colimits of representables, for $\kappa$ uncountable we may absorb the retract into the colimit. 
		
		For these reasons, working with an uncountable cutoff cardinal $\kappa$ gives better behavior of the $\kappa$-compact objects. We will record some basic facts about $\kappa$-compact objects in \Cref{chap:Compact_Objects}.
	\end{convention}

	\section{Stefanich rings}
	\label[section]{sec:Stefanich_Rings}

	We are now ready to define the main object of study in these notes.

	\begin{definition}
		Set
		\[
			0\Pr := \An, \qquad 1\Pr := \PrL_{\aleph_1}.
		\]
		For $n \geq 1$, define inductively
		\[
			(n+1)\Pr := \Mod_{n\Pr}(1\Pr) \in \CAlg(1\Pr).
		\]
	\end{definition}

	\begin{notation}\label[notation]{nota:Absolute_Transition_Adjunctions}
		For every $n \geq 1$, the module-endomorphism adjunction gives an adjunction
		\[
			\Mod^{\bbP}_n\colon \CAlg((n-1)\Pr) \rightleftarrows \CAlg(n\Pr) \noloc \End^{\bbP}_n,
		\]
		where
		\[
			\Mod^{\bbP}_n(R) := \Mod_R((n-1)\Pr), \qquad \End^{\bbP}_n(D) := \End_D(\unit_D).
		\]
	\end{notation}

	\begin{definition}
		We define the category of \emph{Stefanich rings} as the colimit
		\[
			\StRing \; := \; \colim_{1\Pr}\bigl(\CAlg(\An) \xrightarrow{\Mod^{\bbP}_1} \CAlg(1\Pr) \xrightarrow{\Mod^{\bbP}_2} \CAlg(2\Pr) \xrightarrow{\Mod^{\bbP}_3} \dots\bigr).
		\]
		Equivalently, passing to the right adjoints gives
		\[
			\StRing \simeq \lim_{\CAT}\bigl(\CAlg(\An) \xleftarrow{\End^{\bbP}_1} \CAlg(1\Pr) \xleftarrow{\End^{\bbP}_2} \CAlg(2\Pr) \xleftarrow{\End^{\bbP}_3} \dots\bigr).
		\]
	\end{definition}

	\begin{notation}
		The initial Stefanich ring is denoted by
		\[
			\bbP := (\ast,\An,1\Pr,2\Pr,\dots).
		\]
		The notations $\Mod^{\bbP}_n$ and $\End^{\bbP}_n$ are used in anticipation of later generalizations for arbitrary Stefanich rings in place of $\bbP$.
	\end{notation}

	\chapter{Quasicoherent sheaves of higher categories I}
	\label[chapter]{chap:QCoh_I}

	\emph{Talk given by Christoph.}

	The goal of the next three talks is to construct higher analogues of quasicoherent sheaves on fpqc-stacks. In this talk we treat the affine case, meaning stacks of the form $\Spec(A)$ for a commutative ring spectrum $A \in \CAlg(\Sp)$. We write
	\[
		\Aff_{\S} := \CAlg(\Sp)\catop
	\]
	for the category of \emph{affine spectral schemes}, and denote by $\Spec(A)$ the affine spectral scheme corresponding to $A$. A map $\Spec(B) \to \Spec(A)$ in $\Aff_{\S}$ is the same thing as a map $A \to B$ in $\CAlg(\Sp)$.

	For each such affine $\Spec(A)$ we will define presentable categories
	\[
		0\Pr_A,\; 1\Pr_A,\; 2\Pr_A,\ldots
	\]
	which should be thought of as higher versions of $\QCoh(\Spec(A))$. These categories assemble into a Stefanich ring $\Oo(\Spec(A))$, resulting in a functor
	\[
		\Oo\colon \Aff_{\S}\catop \hookrightarrow \StRing.
	\]
	The main goals of the talk are then the following:
	\begin{itemize}
		\item We show that the resulting assignments $\Spec(A) \mapsto n\Pr_A$ can be upgraded to a six-functor formalism
		\[
			\Span(\Aff_{\S}, \mathrm{all}, \textup{countably presented maps}) \to (n+1)\Pr, \qquad \Spec(A) \mapsto n\Pr_A.
		\]
		\item We show that these assignments satisfy \emph{ambidexterity} in sufficiently high degrees.
	\end{itemize}

	\section{Affine Stefanich rings}

	We start by constructing the categories $n\Pr_A$.
	
	\begin{definition}
		Consider a commutative ring spectrum $A \in \CAlg(\Sp)$. We denote its \emph{derived category} of modules by $\D(A)$, or by
		\[
			0\Pr_A := \D(A) := \Mod_A(\Sp) \in \CAlg(1\Pr)
		\]
		in the indexing introduced below.
		For $n \geq 1$, we inductively define the presentable category $n\Pr_A$ by
		\[
			n\Pr_A := \Mod_{(n-1)\Pr_A}(1\Pr) \qin \CAlg(1\Pr).
		\]
		Observe that $n\Pr_A$ comes equipped with a canonical functor $n\Pr \to n\Pr_A$. This is defined inductively: for $n = 0$, this is the unique symmetric monoidal left adjoint $\An \to \Mod_A(\Sp) = 0\Pr_A$, while for $n \geq 1$ it is defined as the base change functor
		\[
			n\Pr = \Mod_{(n-1)\Pr}(1\Pr) \xrightarrow{(n-1)\Pr_A \otimes_{(n-1)\Pr} -} \Mod_{(n-1)\Pr_A}(1\Pr) = n\Pr_A.
		\]
		We may thus regard $n\Pr_A$ as an object
		\[
			n\Pr_A \qin \CAlg(1\Pr)_{n\Pr/} \; \simeq \; \CAlg(\Mod_{n\Pr}(1\Pr)) \; \simeq \; \CAlg((n+1)\Pr).
		\]
	\end{definition}

	\begin{observation}
		Recall from \Cref{nota:Absolute_Transition_Adjunctions} the fully faithful functor
		\[
			\Mod^{\bbP}_n\colon \CAlg((n-1)\Pr) \hookrightarrow \CAlg(n\Pr), \qquad R \mapsto \Mod_R((n-1)\Pr),
		\]
		with right adjoint $\End^{\bbP}_n(D) := \End_D(\unit_D)$. Note that for $A \in \CAlg(\Sp)$, there is a canonical equivalence
		\[
			n\Pr_A = \Mod_{(n-1)\Pr_A}(1\Pr) \simeq \Mod_{(n-1)\Pr_A}(\Mod_{(n-1)\Pr}(1\Pr)) = \Mod_{(n-1)\Pr_A}(n\Pr) = \Mod^{\bbP}_{n+1}((n-1)\Pr_A).
		\]
		In particular, since $\Mod^{\bbP}_{n+1}$ is fully faithful, there is a canonical equivalence
		\begin{equation}
		\label{eq:Structure_Maps_St(A)}
			\End^{\bbP}_{n+1}(n\Pr_A) \simeq (n-1)\Pr_A.
		\end{equation}
	\end{observation}

	\begin{definition}
		Given a commutative ring spectrum $A$, we define the associated Stefanich ring
		\[
			\Oo(\Spec(A)) \qin \StRing \; = \; \lim_{\CAT}\bigl(\CAlg(\An) \xleftarrow{\End^{\bbP}_1} \CAlg(1\Pr) \xleftarrow{\End^{\bbP}_2} \CAlg(2\Pr) \xleftarrow{\End^{\bbP}_3} \dots\bigr)
		\]
		to be the sequence
		\[
			\Oo(\Spec(A)) := (\Omega^{\infty}(A), \D(A), 1\Pr_A, 2\Pr_A, \ldots),
		\]
		equipped with the structure maps from \eqref{eq:Structure_Maps_St(A)}.
	\end{definition}

	\begin{lemma}\label[lemma]{lem:Colimit_Of_Fully_Faithful_Left_Adjoints}
		Let $(C_n)_{n \geq n_0}$ be a sequence of objects of $1\Pr$, together with fully faithful left adjoints $F_n\colon C_n \to C_{n+1}$ with right adjoints $G_n\colon C_{n+1} \to C_n$. Let
		\[
			C := \colim_{1\Pr}(C_{n_0} \xhookrightarrow{F_{n_0}} C_{n_0+1} \xhookrightarrow{F_{n_0+1}} \dots).
		\]
		Then for every $n \geq n_0$ the canonical functor $\iota_n\colon C_n \to C$ is fully faithful.
	\end{lemma}
	\begin{proof}
		Since the $F_n$ are left adjoints, the colimit defining $C$ may equivalently be computed as the limit along the right adjoints,
		\[
			C \simeq \lim_{\CAT}( \dots \xrightarrow{G_{n+1}} C_{n+1} \xrightarrow{G_n} C_n),
		\]
		and under this identification $\iota_n$ is left adjoint to the projection $q_n\colon C \to C_n$. We must show that the unit map $R \to q_n\iota_n(R)$ is an isomorphism for each $R \in C_n$. Since the functors $F_m$ are fully faithful, the sequence of objects
		\[
			X_j = F_{j-1}F_{j-2}\cdots F_n(R) \qin C_j
		\]
		for $j \geq n$ defines an object $X := (X_j)_{j \geq n}$ in this limit. We claim that the identification $R = X_n = q_n(X)$ exhibits $X$ as a left-adjoint object to $R$ under $q_n$ (that is, it induces an isomorphism $X \cong \iota_n(R)$). Indeed, for any object $B = (B_j)_{j \geq n}$ of the limit, we have
		\[
			\begin{aligned}
				\Hom_{C}(X, B)
				&\simeq \lim_{j \geq n} \Hom_{C_j}(X_j,B_j) \\
				&\simeq \lim_{j \geq n} \Hom_{C_n}(R,G_n \cdots G_{j-1}(B_j)) \\
				&\simeq \Hom_{C_n}(R,B_n).
			\end{aligned}
		\]
		Here the second equivalence uses the adjunctions $F_m \dashv G_m$, and the last equivalence uses the compatibility equivalences $G_i(B_{i+1}) \simeq B_i$ defining the object $B$ of the limit.
		This shows that the unit of the adjunction $\iota_n \dashv q_n$ at $R \in C_n$ is the identity map $R \to R$, hence is an isomorphism, showing that $\iota_n$ is fully faithful.
	\end{proof}

	\begin{lemma}
		\label[lemma]{lem:Affine_Stefanich_Rings}
		The assignment $\Spec(A) \mapsto \Oo(\Spec(A))$ defines a fully faithful functor
		\[
			\Oo\colon \Aff_{\S}\catop \hookrightarrow \StRing.
		\]
	\end{lemma}
	\begin{proof}
		Since we have $n\Pr_A = \Mod^{\bbP}_{n+1}((n-1)\Pr_A)$ for all $n$, this assignment is equivalent to the following composite:
		\[
			\CAlg(\Sp) \xrightarrow{\Mod_{(-)}(\Sp)} \CAlg(1\Pr) \hookrightarrow \StRing,
		\]
		where the second functor is fully faithful by \Cref{lem:Colimit_Of_Fully_Faithful_Left_Adjoints}, being a colimit over the fully faithful left adjoints
		\[
			\Mod^{\bbP}_n\colon \CAlg((n-1)\Pr) \hookrightarrow \CAlg(n\Pr).
		\]
		But also the first functor is fully faithful, as it factors as
		\[
			\CAlg(\Sp) \xhookrightarrow{\Mod_{(-)}(\Sp)} \CAlg(\Mod_{\Sp}(1\Pr)) \hookrightarrow \CAlg(1\Pr),
		\]
		where we use that $\Sp \in \CAlg(1\Pr)$ is an idempotent algebra.
	\end{proof}
	
	\begin{notation}
		Consider a map $f\colon A \to B$ in $\CAlg(\Sp)$. By the functoriality of the associated Stefanich ring, this results in a map of Stefanich rings
		\[
			\Oo(\Spec(f))\colon \Oo(\Spec(A)) \to \Oo(\Spec(B)).
		\]
		The $n$-th component is a symmetric monoidal functor of the form
		\[
			f_n^*\colon n\Pr_A \to n\Pr_B,
		\]
		which we refer to as the \emph{$n$-th pullback functor}. Concretely, $f_0^*\colon \Mod_A(\Sp) \to \Mod_B(\Sp)$ is extension of scalars along $f$, and for $n > 0$ the functor $f_n^*\colon \Mod_{(n-1)\Pr_A}(1\Pr) \to \Mod_{(n-1)\Pr_B}(1\Pr)$ is extension of scalars along
		\[
			f_{n-1}^*\colon (n-1)\Pr_A \to (n-1)\Pr_B.
		\]
	\end{notation}

	\section{Extension and restriction of scalars}

	In preparation for the construction of the six-functor formalism $\Spec(A) \mapsto n\Pr_A$, we recall some basic properties of extension and restriction of scalars in presentably symmetric monoidal categories.

	Consider some $\Vv \in \CAlg(\PrL)$. Given a map $g\colon R \to S$ in $\CAlg(\Vv)$, we obtain an adjunction
	\[
		g^*\colon \Mod_R(\Vv) \rightleftarrows \Mod_S(\Vv) \colon g_*
	\]
	given by extension/restriction of scalars. The functor $g^*$ is given on underlying objects in $\Vv$ by $M \mapsto M \otimes_R S$. It is symmetric monoidal.

	\begin{remark}
		\label[remark]{rmk:Basic_Properties_Extension_Scalars}
		We recall some basic properties of this adjunction:
		\begin{enumerate}[(a)]
			\item The functor $g_*\colon \Mod_S(\Vv) \to \Mod_R(\Vv)$ commutes with the forgetful functors to $\Vv$. It is functorial in $g$: we have $(g'g)_* \simeq g_* g'_*$. Consequently, $g^*$ is functorial in $g$.
			\item The functor $g^*$ preserves free modules. Recall that an $R$-module is called \emph{free} if it lies in the image of extension of scalars $\Vv \to \Mod_R(\Vv), \; M \mapsto R \otimes M$ along the unit map $\unit \to R$, and similarly for $S$-modules. Since $g^*(R \otimes M) \simeq S \otimes M$ and the composite $\unit \to R \xrightarrow{g} S$ is the unit map of $S$, the functor $g^*$ sends the free $R$-module $R \otimes M$ to the free $S$-module $S \otimes M$.
			\item \label{it:Restriction_Scalars_Preserves_Colimits} The functor $g_*\colon \Mod_S(\Vv) \to \Mod_R(\Vv)$ preserves colimits. This holds because the map $\Mod_R(\Vv) \to \Vv$ creates colimits.
			\item \label{it:Restriction_Scalars_Projection_Formula} The functor $g_*$ satisfies the \textbf{projection formula}: for $N \in \Mod_S(\Vv)$ and $M \in \Mod_R(\Vv)$, the map
			\[
				M \otimes_R g_*N \to g_*(g^*M \otimes_S N)
			\]
			adjoint to $g^*(M \otimes_R g_*N) \simeq g^*M \otimes_S g^*g_*N \to g^*M \otimes_S N$ is an isomorphism. Since both sides preserve colimits in $M$, we may reduce the claim to the case of a free module $M = f^*M' = R \otimes M'$, in which case we see that both sides reduce to $M' \otimes g_*N$ by an explicit computation of the formula for the relative tensor product.
			\item \label{it:Restriction_Scalars_Base_Change} The functors $g_*$ satisfy \textbf{base change}: Consider a pushout square in $\CAlg(\Vv)$ of the form
			\[
				\begin{tikzcd}
					R \rar \dar & S \dar \\
					T \rar & U,
				\end{tikzcd}
			\]
			meaning that $U \simeq T \otimes_R S$ as commutative algebras. By passing to extension of scalars functors, we obtain a commutative diagram in $\CAlg(1\Pr)$ of the form
			\[
				\begin{tikzcd}
					\Mod_R(\Vv) \rar \dar & \Mod_S(\Vv) \dar \\
					\Mod_T(\Vv) \rar & \Mod_U(\Vv).
				\end{tikzcd}
			\]
			The horizontal functors have right adjoints, and we need to show that both composites $\Mod_S(\Vv) \to \Mod_T(\Vv)$ agree, along the canonical map. Computing both sides, we get $T \otimes_R M$ versus $(T \otimes_R S) \otimes_S M$, and these agree. (Details omitted.)
			\item \label{it:Restriction_Scalars_Internal_Right_Adjoint} Finally, assume that $\Vv \in \CAlg(1\Pr)$, and let $g\colon R \to S$ be a countably presented map, meaning that $S$ is $\aleph_1$-compact when regarded as an object of $\CAlg_R(\Vv)$. Then the functor
			\[
				g_*\colon \Mod_S(\Vv) \to \Mod_R(\Vv)
			\]
			preserves $\aleph_1$-compact objects by \Cref{cor:Restriction_Kappa_Compact_Algebra_Preserves_Compacts}, applied with $\kappa = \aleph_1$ and $C = \Vv$. In combination with \eqref{it:Restriction_Scalars_Preserves_Colimits}, we conclude that $g_*$ is a right adjoint to $g^*$ internal to the 2-category $1\Pr$.
		\end{enumerate}
	\end{remark}

	\begin{lemma}
		\label[lemma]{lem:Mod_Preserves_Countably_Presented}
		Let $\Vv \in \CAlg(1\Pr)$, and let $g\colon R \to S$ be a countably presented map in $\CAlg(\Vv)$. Then also the induced functor
		\[
			g^*\colon \Mod_R(\Vv) \to \Mod_S(\Vv)
		\]
		is countably presented in $\CAlg(\Mod_{\Vv}(1\Pr)) \simeq \CAlg_{\Vv}(1\Pr)$.
	\end{lemma}
	\begin{proof}
		Recall from \Cref{thm:Embeddings_CAlg(C)_Into_CAlg(Mod_C)} that the functor
		\[
			\CAlg(\Mod_{\Vv}(1\Pr)) \to \CAlg(\Vv), \qquad D \mapsto \End_{D}(\unit_D)
		\]
		preserves $\aleph_1$-filtered colimits. It follows that its left adjoint
		\[
			\CAlg(\Vv) \hookrightarrow \CAlg(\Mod_{\Vv}(1\Pr)), \qquad R \mapsto \Mod_R(\Vv)
		\]
		preserves $\aleph_1$-compact objects. Replacing $\Vv$ by $\Mod_R(\Vv)$ then gives the claim.
	\end{proof}

	\begin{corollary}
		\label[corollary]{cor:nPr_Preserves_Countably_Presented}
		Let $g\colon A \to B$ be a countably presented morphism of commutative ring spectra. Then for each $n \geq 0$, the functor
		\[
			g_n^*\colon n\Pr_A \to n\Pr_B
		\]
		is countably presented as a map in $\CAlg((n+1)\Pr)$.
	\end{corollary}
	\begin{proof}
		By induction this is immediate from \Cref{lem:Mod_Preserves_Countably_Presented}.
	\end{proof}

	\begin{corollary}
		\label[corollary]{cor:nPr_B_Aleph1_Compact}
		Let $f\colon A \to B$ be a countably presented map in $\CAlg(\Sp)$. Then, for every $n \geq 0$, the object $n\Pr_B \in (n+1)\Pr_A \simeq \Mod_{n\Pr_A}(1\Pr)$ is $\aleph_1$-compact.
	\end{corollary}
	\begin{proof}
		By the previous corollary, $n\Pr_B$ is $\aleph_1$-compact as an object of $\CAlg_{n\Pr_A}(1\Pr)$. By \Cref{lem:Kappa_Compact_Algebra_Iff_Module} this is equivalent to it being $\aleph_1$-compact as an object of $\Mod_{n\Pr_A}(1\Pr)$.
	\end{proof}

	\section{Construction of the six-functor formalism}

	We will now construct the desired six-functor formalism on $\Aff_{\S}$ we are after.

	\begin{definition}
		A map $\Spec(B) \to \Spec(A)$ in $\Aff_{\S}$ is called \emph{countably presented} if $B$ is countably presented as an $A$-algebra.
	\end{definition}

	This class of maps satisfies the usual closure properties.

	\begin{lemma}\label[lemma]{lem:Affine_Countably_Presented_Calculus}
		Let $A \in \CAlg(\Sp)$.
		\begin{enumerate}[(i)]
			\item Countably presented $A$-algebras are closed under countable colimits, retracts, and finite colimits in $\CAlg_A(\Sp)$.
			\item Countably presented maps of affine spectral schemes are stable under base change and composition.
			\item If $A \to B$ and $A \to C$ are countably presented and $B \to C$ is a map over $A$, then $B \to C$ is countably presented.
			\item If $\Spec(B) \to \Spec(A)$ is countably presented and admits a section $\Spec(A) \to \Spec(B)$, then the section is countably presented.
		\end{enumerate}
	\end{lemma}
	\begin{proof}
		Part (i) is \Cref{lem:Kappa_Compact_Algebra_Closure_Properties}(1), applied with $\kappa = \aleph_1$ and $C = \Sp$.

		For base change in (ii), suppose $A \to B$ is countably presented and $A \to A'$ is any map. The base change is $\Spec(A' \otimes_A B) \to \Spec(A')$, and $A' \to A' \otimes_A B$ is countably presented by \Cref{lem:Kappa_Compact_Algebra_Closure_Properties}(2). For composition, if $A \to B$ and $B \to C$ are countably presented, then $A \to C$ is countably presented by \Cref{lem:Kappa_Compact_Algebra_Closure_Properties}(3).

		Part (iii) is \Cref{lem:Kappa_Compact_Algebra_Closure_Properties}(4). Finally, a section $\Spec(A) \to \Spec(B)$ corresponds to an $A$-algebra retraction $B \to A$, so (iv) is exactly \Cref{lem:Kappa_Compact_Algebra_Section}.
	\end{proof}

	We now come to the main result of this talk:

	\begin{theorem}
		\label[theorem]{thm:Six_Functor_Formalism_Pr_On_Affines}
		For every $n \geq 0$, there exists a six-functor formalism
		\begin{align*}
			\Span(\Aff_{\S}, \mathrm{all}, \textup{countably presented maps}) \quad &\to \qquad (n+1)\Pr \\
			\Spec(A) \qquad \qquad\qquad \quad & \mapsto \qquad n\Pr_A
			\\
			\Spec(A) \xleftarrow{f} \Spec(B) \xrightarrow{g} \Spec(C)\qquad  &\mapsto \qquad n\Pr_A \xrightarrow{f_n^*} n\Pr_B \xrightarrow{g_{n,*}} n\Pr_C,
		\end{align*}
		where $g_{n,*}$ is right adjoint to $g_n^*$.
	\end{theorem}
	\begin{proof}
		We start with the case $n = 0$. The functoriality of module categories provides a functor
		\[
			\Aff_{\S}\catop \to 1\Pr, \qquad \Spec(A) \mapsto \Mod_A(\Sp), \qquad f \mapsto f^*.
		\]
		To extend this to the span category via the right adjoints, it suffices to check the following properties:
		\begin{itemize}
			\item For every countably presented map $g\colon \Spec(A) \to \Spec(B)$, the functor $g^*\colon \Mod_B(\Sp) \to \Mod_A(\Sp)$ admits a right adjoint $g_*$ internal to the 2-category $1\Pr$. This was verified in \Cref{rmk:Basic_Properties_Extension_Scalars}\eqref{it:Restriction_Scalars_Internal_Right_Adjoint}.
			\item These right adjoints $g_*$ satisfy base change. This was verified in \Cref{rmk:Basic_Properties_Extension_Scalars}\eqref{it:Restriction_Scalars_Base_Change}.
			\item For every such $f$, the right adjoint $f_*$ satisfies the projection formula. This was verified in \Cref{rmk:Basic_Properties_Extension_Scalars}\eqref{it:Restriction_Scalars_Projection_Formula}.
		\end{itemize}

		The proof for $n > 0$ is essentially identical. Functoriality of the assignment $A \mapsto n\Pr_A$ provides a functor
		\[
			\Aff_{\S}\catop \to (n+1)\Pr, \qquad \Spec(A) \mapsto n\Pr_A, \qquad f \mapsto f^*_n.
		\]
		If $g\colon \Spec(B) \to \Spec(A)$ is countably presented, then by \Cref{cor:nPr_Preserves_Countably_Presented} so is the functor
		\[
			g^*_{n-1}\colon (n-1)\Pr_A \to (n-1)\Pr_B
		\]
		as a morphism in $\CAlg(n\Pr)$. It then follows from \Cref{rmk:Basic_Properties_Extension_Scalars}\eqref{it:Restriction_Scalars_Internal_Right_Adjoint} that the associated restriction of scalars functor
		\[
			g_{n,*} = (g^*_{n-1})_* \colon n\Pr_B \to n\Pr_A
		\]
		is an internal right adjoint to $g^*_n$ in $(n+1)\Pr$. Moreover, the other parts of \Cref{rmk:Basic_Properties_Extension_Scalars} show that these right adjoints satisfy the required base change and projection formula conditions.
	\end{proof}

	\begin{remark}
		If we drop the countable presentation hypotheses, we still obtain a presentable six-functor formalism:
		\[
			\Span(\Aff_{\S}) \to \PrL, \qquad \Spec(A) \mapsto n\Pr_A
		\]
	\end{remark}

	\section{Ambidexterity}
	\label[section]{sec:Ambidexterity}

	It turns out that for $n \geq 1$ the symmetric monoidal pullback functor $f_n^*$ admits not only a well-behaved right adjoint $f_{n,*}$, but also a left adjoint $f_{n,\sharp}$, and the two coincide. This phenomenon, of a functor having identical left and right adjoints, is known as \emph{ambidexterity}, a term introduced by Hopkins--Lurie \cite{hopkinsLurie2013ambidexterity}. Before applying this to our functors $f_n^*$, we set up the general framework.

	\begin{definition}\label[definition]{def:Adjoint_In_2Cat}
		Let $\Cc$ be a 2-category, and let $F\colon X \to Y$ be a $1$-morphism in $\Cc$.
		\begin{itemize}
			\item We say $F$ is a \emph{left adjoint} if there exist a $1$-morphism $F^R\colon Y \to X$ together with $2$-morphisms
			\[
				\eta\colon \id_X \to F^R F \quad \text{in $\End_{\Cc}(X)$,} \qquad \epsilon\colon F F^R \to \id_Y \quad \text{in $\End_{\Cc}(Y)$,}
			\]
			such that the composites
			\[
				F \xrightarrow{F\eta} F F^R F \xrightarrow{\epsilon F} F \qquad \text{and} \qquad F^R \xrightarrow{\eta F^R} F^R F F^R \xrightarrow{F^R \epsilon} F^R
			\]
			are equivalent to the identity.
			\item Dually, $F$ is a \emph{right adjoint} if there exist a $1$-morphism $F^L\colon Y \to X$ and $2$-morphisms $\eta\colon \id_Y \to F F^L$, $\epsilon\colon F^L F \to \id_X$ satisfying the analogous triangle identities.
		\end{itemize}
	\end{definition}

	For the present discussion the reader may take $\Cc$ to be an object of $2\Pr$, though the definition makes sense for any reasonable notion of $(\infty,2)$-category. If $F$ is a left adjoint, the right adjoint $F^R$ is determined up to canonical isomorphism, and the left adjoint of $F^R$ is again $F$ (with the same $\eta$ and $\epsilon$).

	\begin{remark}
		In our formalism of presentable $(\infty,n)$-categories, there are forgetful functors $n\Pr \to k\Pr$ for any $k \leq n$. Indeed, the symmetric monoidal functor $\An \to 1\Pr$ (since $\An$ is the unit object) yields, by passing to modules, symmetric monoidal functors $n\Pr \to (n+1)\Pr$ inductively, whose right adjoints are the desired forgetful functors. In particular, \Cref{def:Adjoint_In_2Cat} applies to $1$-morphisms in $n\Pr$ for any $n \geq 2$.
	\end{remark}

	We can now state the general ambidexterity result.

	\begin{proposition}\label[proposition]{prop:General_Ambidexterity}
		Let $\Cc$ be an $(\infty,3)$-category, and let $F\colon X \to Y$ be a $1$-morphism in $\Cc$. The following are equivalent.
		\begin{enumerate}[(i)]
			\item The morphism $F$ admits a right adjoint $F^R\colon Y \to X$, and both
			\[
				\eta\colon \id_X \to F^R F \quad \text{in $\End_{\Cc}(X)$} \qquad \text{and} \qquad \epsilon\colon F F^R \to \id_Y \quad \text{in $\End_{\Cc}(Y)$}
			\]
			admit \emph{right} adjoints (in the respective $(\infty,2)$-categories).
			\item The morphism $F$ admits a left adjoint $F^L\colon Y \to X$, and both
			\[
				\eta\colon \id_Y \to F F^L \quad \text{in $\End_{\Cc}(Y)$} \qquad \text{and} \qquad \epsilon\colon F^L F \to \id_X \quad \text{in $\End_{\Cc}(X)$}
			\]
			admit \emph{left} adjoints.
		\end{enumerate}
		Moreover, when these conditions hold, there is a canonical identification $F^L \simeq F^R$.
	\end{proposition}

	\begin{remark}
		There is a dual version of \Cref{prop:General_Ambidexterity}, in which the words ``right'' and ``left'' are swapped in (i) and (ii) respectively; it again yields a canonical identification $F^L \simeq F^R$. However, in any given application, $\eta$ and $\epsilon$ must be assumed to admit the \emph{same} type of adjoint for the result to apply.
	\end{remark}

	\begin{proof}
		We show that (i) implies (ii) with $F^L = F^R$; the dual implication is similar. Given $(F^R, \eta, \epsilon)$ as in (i), define
		\[
			\eta'\colon \id_Y \to F F^R \qquad \text{and} \qquad \epsilon'\colon F^R F \to \id_X
		\]
		by $\eta' := \epsilon^R$ and $\epsilon' := \eta^R$. The triangle identities for $(F, F^R, \eta', \epsilon')$ follow from those of $(F, F^R, \eta, \epsilon)$ together with the fact that forming right adjoints is compatible with composition. By construction, $\eta'$ and $\epsilon'$ admit left adjoints (namely $\epsilon$ and $\eta$ themselves), giving (ii).
	\end{proof}

	We now establish an ambidexterity property for the functors $f_{n,*}$.

	\begin{proposition}
		\label[proposition]{prop:Affine_Ambidexterity}
		Let $n \geq 1$, and let $f\colon A \to B$ be a countably presented map in $\CAlg(\Sp)$.
		\begin{itemize}
			\item The pullback functor $f_n^*\colon n\Pr_A \to n\Pr_B$ admits a left adjoint $f_{n,\sharp}\colon n\Pr_B \to n\Pr_A$, which the canonical comparison identifies with the right adjoint, $f_{n,\sharp} \simeq f_{n,*}$. As a left adjoint, $f_{n,\sharp}$ satisfies base change and the projection formula.
			\item The category $n\Pr_B$ is self-dual as an object of $(n+1)\Pr_A \simeq \Mod_{n\Pr_A}(1\Pr)$.
		\end{itemize}
	\end{proposition}

	\begin{proof}
		Since we have the adjunction $f_n^* \dashv f_{n,*}$, we obtain a unit and counit
		\[
			\eta\colon \id \to f_{n,*} f_n^* \in \End_{(n+1)\Pr_A}(n\Pr_A) \qquad \qquad \epsilon \colon f_n^*f_{n,*} \to \id \in \End_{(n+1)\Pr_A}(n\Pr_B).
		\]
		We will deduce the proposition from \Cref{prop:General_Ambidexterity}, applied in the $(\infty,3)$-category $\Cc = (n+1)\Pr_A$ (which lies in $(n+2)\Pr$ and hence in $3\Pr$, since $n \geq 1$); concretely, we need to verify that $\eta$ and $\epsilon$ admit right adjoints. For $\eta$, note that evaluation at the unit induces an equivalence 
		\[
			\End_{(n+1)\Pr_A}(n\Pr_A) \iso n\Pr_A.
		\]
		Under this equivalence, the unit $\eta$ corresponds to the map
		\[
			f_{n-1}^*\colon (n-1)\Pr_A \to (n-1)\Pr_B,
		\]
		which indeed has a right adjoint $f_{n-1,*}$ in $n\Pr_A$.
		
		For $\epsilon$, note that there is an equivalence
		\begin{align*}
			\End_{(n+1)\Pr_A}(n\Pr_B) &= \End_{\Mod_{n\Pr_A}}(\Mod_{(n-1)\Pr_B}(n\Pr_A), \Mod_{(n-1)\Pr_B}(n\Pr_A)) \\
			&\simeq \Mod_{(n-1)\Pr_B \otimes_{(n-1)\Pr_A} (n-1)\Pr_B}(n\Pr_A).
		\end{align*}
		Under this equivalence, the counit $\epsilon$ corresponds to the `multiplication map'
		\[
			(n-1)\Pr_B \otimes_{(n-1)\Pr_A} (n-1)\Pr_B \to (n-1)\Pr_B.
		\]
		The left-hand side is equivalent to $(n-1)\Pr_{B \otimes_A B}$, and under this equivalence the map agrees with $\nabla^*_{n-1}$ induced by the codiagonal map $\nabla\colon B \otimes_A B \to B$. For this we also know it admits a $(n-1)\Pr_{B \otimes_A B}$-linear right adjoint.

		For the claim about base change, we apply the same result inside the 3-category $\Ar((n+1)\Pr_A)$, the arrow category of $(n+1)\Pr_A$. Since adjoints in this arrow category are precisely the adjointable squares, this gives the claim.

		Finally, for the duality data on $n\Pr_B$, we are using the following two maps:
		\[
			n\Pr_A \xrightarrow{f_n^*} n\Pr_B \xrightarrow{\nabla_{n,*}} n\Pr_B \otimes_{n \Pr_A} n\Pr_B
		\]
		\[
			n\Pr_B \otimes_{n\Pr_A} n\Pr_B \xrightarrow{\nabla_n^*} n\Pr_B \xrightarrow{f_{n,*}} n\Pr_A.
		\]
		The triangle identities are an immediate consequence of base change.
	\end{proof}

	\chapter{Descent for quasicoherent sheaves}
	\label[chapter]{chap:Descent_QCoh}

	\emph{Talk given by Melvin.}

	We use the following output of the previous talk:
	\begin{itemize}
		\item For $A \in \CAlg(\Sp)$, we defined categories $n\Pr_A$ by
		\[
			0\Pr_A := \D(A) := \Mod_A(\Sp), \qquad (n+1)\Pr_A := \Mod_{n\Pr_A}((n+1)\Pr) \simeq \Mod_{n\Pr_A}(1\Pr).
		\]
		\item The assignment $\Spec(A) \mapsto n\Pr_A$ defines a six-functor formalism on $\Aff_{\S}$ valued in $1\Pr$.
		\item For $n \geq 1$, we have ambidexterity: if $f\colon \Spec(A) \to \Spec(B)$ is countably presented, then the right adjoint $f_{n,*}\colon n\Pr_A \to n\Pr_B$ is also a left adjoint to $f_n^*$ in a canonical way.
	\end{itemize}

	\section{Hyperdescent}

	The goal of this chapter is the following hyperdescent statement:

	\begin{theorem}[{cf.\ \cite[Corollary 3.33, Proposition 3.45]{Mathew_Galois}}]\label[theorem]{thm:fpqc_Hyperdescent_nPr}
		The functor $n\Pr_{(-)}^!\colon \Aff_{\S}\catop \to 1\Pr$ (defined on maps by the right adjoints $f^!$) is a hypersheaf with respect to the fpqc topology for each $n \geq 0$.
	\end{theorem}

	\begin{remark}
		For $n \geq 1$, the maps $f^!$ agree with $f^*$ by ambidexterity.
	\end{remark}

	To make sense of this statement, we recall the following two definitions:

	\begin{definition}
		Let $C$ be a site with finite limits. A \emph{hypercover} of $A \in C$ is a simplicial object $A_{\bullet} \in \Fun(\simp\catop,C)$ together with a map $A_0 \to A$ which is a cover in $C$, such that if we right Kan extend $A_{\bullet}$ to $A(-) \in \Fun(\PSh_{\fin}(\simp)\catop,C)$, the map $A(\Delta^n) \to A(\partial \Delta^n)$ is a cover in $C$.
	\end{definition}

	\begin{definition}
		Recall the following basic definitions:
		\begin{itemize}
			\item For a static ring $B_0 \in \CAlg(\Ab)$ and a static $B_0$-module $M_0$, we say that $M_0$ is \emph{flat} if the functor $M_0 \otimes_{B_0} -\colon \Mod_{B_0}(\Ab) \to \Ab$ is exact. It is called \emph{faithfully flat} if, in addition, the functor $M_0 \otimes_{B_0} -$ is conservative, i.e.\ detects isomorphisms. (Equivalently it detects exact sequences.)
			\item A morphism $B_0 \to A_0$ in $\CAlg(\Ab)$ is called \emph{(faithfully) flat} if $A_0$ is (faithfully) flat as a $B_0$-module.
			\item A morphism $B \to A$ in $\CAlg(\Sp)$ is \emph{(faithfully) flat} if the map $\pi_0 B \to \pi_0 A$ in $\CAlg(\Ab)$ is (faithfully) flat and for each $i \in \Z$ the multiplication map $\pi_i B \otimes_{\pi_0 B} \pi_0 A \to \pi_i A$ is an isomorphism.
			\item A morphism $\Spec(A) \to \Spec(B)$ in $\Aff_{\S}$ is called (faithfully) flat if its associated map $B \to A$ in $\CAlg(\Sp)$ is.
		\end{itemize}
	\end{definition}

	\begin{definition}
		The \emph{fpqc topology} on $\Aff_{\S}$ is the Grothendieck topology in which a sieve over $\Spec(B)$ is declared to be a covering sieve if it contains a finite collection of morphisms $\Spec(A_i) \to \Spec(B)$ such that the induced map $\Spec(\prod_{i=1}^n A_i) \to \Spec(B)$ is faithfully flat.
	\end{definition}

	Observe that a functor $F\colon \Aff_{\S}\catop \to C$ is a hypersheaf with respect to the fpqc topology if and only if the following two properties are satisfied:
	\begin{itemize}
		\item The functor $F$ preserves finite products.
		\item For an fpqc-hypercover $\Spec(A_{\bullet}) \to \Spec(A)$, the induced map
		\[
			F(\Spec(A)) \to \lim_{[n] \in \simp} F(\Spec(A_n))
		\]
		is an isomorphism.
	\end{itemize}

	The proof has three parts; the first two treat countably presented fpqc-hypercovers:
	\begin{itemize}
		\item We first prove the theorem for $n = 0$.
		\item We then prove the general case by induction on $n$.
		\item Finally, we reduce arbitrary fpqc-hypercovers to the countably presented case.
	\end{itemize}

	\section{The base case \texorpdfstring{$n = 0$}{n = 0}}

	We will start by showing that the assignment $\Spec(A) \mapsto \D(A)$ satisfies fpqc-hyperdescent. Since we have $\D(\prod_i A_i) \iso \prod_i \D(A_i)$, it remains to show that for an fpqc-hypercover $\Spec(A_{\bullet}) \to \Spec(A)$ in $\Aff_{\S}$ the functor
	\[
		\D(A) \iso \lim_{[n] \in \simp}^! \D(A_{\bullet})
	\]
	is an equivalence. Here the limit on the right is taken along the functors $f^!$.

	\subsection{Step 1: The functor is fully faithful}

	We start by showing that this comparison functor is fully faithful.

	Note that the functor $\D(A) \iso \lim_{[n] \in \simp}^! \D(A_{\bullet})$ admits a left adjoint, given by sending an object $(M_n)_{[n] \in \simp}$ to the colimit $\colim_{[n] \in \simp\catop} (f_n)_* M_n$, where we denote by $f_n\colon \Spec(A_n) \to \Spec(A)$ the maps constituting the hypercover. So we must show that for any object $M \in \D(A)$, the counit map
	\[
		\colim_{[n] \in \simp\catop} (f_{n})_* f_n^! M \to M
	\]
	is an isomorphism of $A$-modules. Note that we may rewrite the left-hand side as
	\[
		\colim_{[n] \in \simp\catop} (f_{n})_* f_n^! M = \colim_{[n] \in \simp\catop} (f_{n})_* \iHom_{A_n}(A_n, f_n^! M) \simeq \colim_{[n] \in \simp\catop} (f_{n})_* \iHom_{A}(A_n, M).
	\]
	We may now filter $\simp$ by the subcategories $\simp_{\leq n}$ to rewrite this colimit in turn as
	\[
		\colim_{m \geq 0} \colim_{[n] \in (\simp_{\leq m})\catop} \iHom_A(A_n, M) \simeq \colim_{m \geq 0} \iHom_A(\lim_{[n] \in \simp_{\leq m}} A_\bullet, M),
	\]
	where we use that a colimit over $\simp_{\leq m}$ is a finite colimit and is thus preserved by the exact functor $\iHom_A(-, M)\colon \D(A) \to \D(A)$. If we similarly write $M \simeq \iHom_A(A, M) \simeq \colim_{m \geq 0} \iHom_A(A, M)$, then the comparison map is obtained by applying $\iHom_A(-, M)$ to the canonical map of pro-systems $A \to \{\lim_{\simp_{\leq m}} A_{\bullet}\}_{m \geq 0}$ in $\D(A)$. So we need to show that this map is a pro-isomorphism in $\D(A)$.

	We first justify the reduction to the connective case.

	\begin{lemma}\label[lemma]{lem:Connective_Truncation_fpqc_Hypercovers}
		Let $A \to A_{\bullet}$ be an fpqc-hypercover of commutative ring spectra. Put $A^+ := \tau_{\geq 0} A$ and $A_n^+ := \tau_{\geq 0} A_n$. Then $A^+ \to A_{\bullet}^+$ is an fpqc-hypercover. Moreover, for every $n \geq 0$ we have a canonical isomorphism
		\[
			A_n \simeq A \otimes_{A^+} A_n^+,
		\]
		and for every $m \geq 0$ the canonical map
		\[
			A \otimes_{A^+} \lim_{[n] \in \simp_{\leq m}} A_n^+ \longrightarrow \lim_{[n] \in \simp_{\leq m}} A_n
		\]
		is an isomorphism in $\D(A)$. Consequently, to show that
		\[
			A \to \{\lim_{[n] \in \simp_{\leq m}} A_n\}_{m \geq 0}
		\]
		is a pro-isomorphism in $\D(A)$, it suffices to prove the analogous statement for $A^+ \to A_{\bullet}^+$.
	\end{lemma}
	\begin{proof}
		We first record the basic flatness input. If $R \to S$ is flat, then $R^+ := \tau_{\geq 0}R \to S^+ := \tau_{\geq 0}S$ is flat and the canonical map
		\[
			R \otimes_{R^+} S^+ \to S
		\]
		is an isomorphism. More generally, for every $R$-algebra $T$ there is a canonical isomorphism
		\[
			T^+ \otimes_{R^+} S^+ \simeq \tau_{\geq 0}(T \otimes_R S).
		\]
		Indeed, both statements follow by comparing homotopy groups, using the defining condition for flatness:
		\[
			\pi_i S \simeq \pi_i R \otimes_{\pi_0 R} \pi_0 S.
		\]

		Let $K$ be a finite simplicial set, and write $A_K$ for the commutative ring spectrum corresponding to the value of the right Kan extension of $\Spec(A_{\bullet})$ on $K$; similarly, write $A^+_K$ for the corresponding object associated with $\Spec(A^+_{\bullet})$. We claim, by induction over the cells of $K$, that
		\[
			A^+_K \simeq \tau_{\geq 0} A_K,
		\]
		and that $A_K$ is flat over $A$.

		The claim is clear for $K = \emptyset$, where $A_K = A$. Suppose that $L$ is obtained from $K$ by attaching an $r$-simplex,
		\[
			L = K \cup_{\partial \Delta^r} \Delta^r.
		\]
		Then
		\[
			A_L \simeq A_K \otimes_{A_{\partial \Delta^r}} A_r.
		\]
		The matching map $A_{\partial \Delta^r} \to A_r$ is faithfully flat by the hypercover condition. Hence $A_K \to A_L$ is flat, and so $A_L$ is flat over $A$. Moreover, by the flatness input above and the induction hypothesis, we get
		\[
			\tau_{\geq 0}A_L
			\simeq
			\tau_{\geq 0}A_K \otimes_{\tau_{\geq 0}A_{\partial \Delta^r}} \tau_{\geq 0}A_r
			\simeq
			A^+_K \otimes_{A^+_{\partial \Delta^r}} A^+_r
			\simeq
			A^+_L.
		\]
		This proves the induction.

		Applying this to $K = \partial \Delta^r$, we see that the $r$-th matching map for $A^+ \to A^+_{\bullet}$ is the connective cover of the faithfully flat map
		\[
			A_{\partial \Delta^r} \to A_r.
		\]
		It is therefore faithfully flat, so $A^+ \to A^+_{\bullet}$ is again an fpqc-hypercover.

		Since $A_n$ is flat over $A$, the flatness input also gives $A_n \simeq A \otimes_{A^+} A_n^+$ for every $n$. Finally, the last displayed isomorphism follows because the functor
		\[
			A \otimes_{A^+} -\colon \D(A^+) \to \D(A)
		\]
		is exact and $\simp_{\leq m}$ is finite:
		\[
			A \otimes_{A^+} \lim_{[n] \in \simp_{\leq m}} A_n^+
			\simeq
			\lim_{[n] \in \simp_{\leq m}} (A \otimes_{A^+} A_n^+)
			\simeq
			\lim_{[n] \in \simp_{\leq m}} A_n.
		\]
		The final assertion follows by applying this exact base-change functor to the pro-isomorphism for the connective hypercover.
	\end{proof}

	By \Cref{lem:Connective_Truncation_fpqc_Hypercovers}, we may assume from now on that both $A$ and all of the $A_i$ are connective.

	Now, to show that the map $A \to \{\lim_{\simp_{\leq m}} A_{\bullet}\}_{m \geq 0}$ is an isomorphism in $\Pro(\D(A))$, we may equivalently show its cofiber is zero. The cofiber is given by the pro-system $\{C_m\}_{m \geq 0}$, where we set
	\[
		C_m := \cofib(A \to \lim_{\simp_{\leq m}} A_\bullet).
	\]
	It will thus suffice to show that the map $C_{m+2} \to C_m$ is the zero map in $\D(A)$ for each $m \geq 0$. For this, recall the following definition:

	\begin{definition}
		A morphism $X \to Y$ in $\D(A)$ is called a \emph{phantom map} if for every perfect $A$-module $Z$ and every map $Z \to X$, the composite $Z \to X \to Y$ is nullhomotopic. In particular, taking $Z = A[i]$ shows that a phantom map induces the zero map $\pi_i(X) \to \pi_i(Y)$ for each $i \in \Z$.
	\end{definition}

	\begin{observation}
		Every $A$-module $X$ may be written as a filtered colimit of perfect $A$-modules: simply consider the collection of maps $\{X' \to X\}$ with $X'$ perfect. If we write $X = \colim_{i \in I} X_i$, then we see that a map $X \to Y$ is phantom if and only if it lies in the kernel of the map
		\[
			\pi_0(\lim_i \hom_A(X_i, Y)) \cong \pi_0\hom_A(X, Y) \to \lim_i \pi_0 \hom_A(X_i, Y).
		\]
		More generally, the homotopy limit spectral sequence
		\[
			R^s\lim_i\pi_t\hom_A(X_i,Y) \Rightarrow \pi_{t-s}\hom_A(X,Y)
		\]
		identifies this kernel with the first positive piece of the induced filtration on $\pi_0\hom_A(X,Y)$; its associated graded pieces are subquotients of $R^s\lim_i\pi_s\hom_A(X_i,Y)$ for $s > 0$. In the sequential case, this reduces to the Milnor exact sequence for the tower of spectra $\hom_A(X_i,Y)$:
		\[
			0 \to \lim\nolimits^1_i \pi_{1}\hom_A(X_i,Y) \to \pi_0\hom_A(X,Y) \to \lim_i\pi_0\hom_A(X_i,Y) \to 0,
		\]
		so in that case the kernel is $\lim^1_i \pi_{1}\hom_A(X_i, Y)$.
	\end{observation}

	\begin{lemma}[{\cite[Theorem~4.2.5]{HPS1997Axiomatic}}]
		The composite of two phantom maps is nullhomotopic.
	\end{lemma}
	\begin{proof}[Proof sketch]
		In the phantom-tower proof of \citet{HPS1997Axiomatic}, one constructs a sequential tower by iteratively killing maps out of perfect objects. Phantom maps have filtration at least $1$ in this tower, detected by the $\lim^1$-term in the Milnor exact sequence. Composition is compatible with this filtration and adds filtration degrees. Thus a composite of two phantom maps lands in filtration at least $2$, i.e.\ in an $R^2\lim$-term. For a sequential tower, $R^s\lim$ vanishes for $s \geq 2$.
	\end{proof}

	In light of the lemma, it remains to show that each map $C_{m+1} \to C_m$ is a phantom map. We now make the following claim:

	\begin{claim}\label[claim]{cl:CountablyGeneratedFlat}
		For each $m \geq 0$, the $A$-module $C_m := \cofib(A \to \lim_{\simp_{\leq m}} A_\bullet)$ is a countably presented flat $A$-module in degree $m-1$, i.e.\ $C_m[-(m-1)]$ is a flat $A$-module.
	\end{claim}

	We will prove this claim below, but let us first explain why it implies that $C_{m+1} \to C_m$ is phantom. For this, recall the following spectral analogue of Lazard's theorem:

	\begin{theorem}[Lazard, HA 7.2.2.15]
		Let $A$ be a connective ring spectrum. Then every connective countably presented flat $A$-module $M$ is a sequential colimit of finite rank free $A$-modules:
		\[
			M \simeq \colim_{i \geq 0} A^{\oplus n_i}.
		\]
	\end{theorem}

	Combining this with \Cref{cl:CountablyGeneratedFlat}, we may thus write $C_{m+1} \simeq \colim_{i \geq 0} A^{\oplus n_i}[m]$. Now, if $Z$ is any perfect $A$-module and $Z \to C_{m+1}$ is any map, then it must factor through some finite stage $Z \to A^{\oplus n_i}[m]$. So it suffices to show that any map $A[m] \to C_m$ is nullhomotopic in $\D(A)$. But homotopy classes of such maps correspond to elements of $\pi_{m}(C_m)$, which is $0$ since by \Cref{cl:CountablyGeneratedFlat} $C_m$ is concentrated in degree $m-1$.

	We now move towards a proof of \Cref{cl:CountablyGeneratedFlat}. For this, recall that an $A$-module $M$ is flat if and only if $M \otimes_A \pi_0(A)$ is a flat static $\pi_0 A$-module. So we may check the claim after base changing to $\pi_0 A$. Observe that base change commutes with the cofiber and with the finite limit $\lim_{\leq m} A_{\bullet}$. So we may assume that $A$ and $A_{\bullet}$ are static.

	The augmented cosimplicial object $\{A_{\bullet}\}$ may be turned into a cochain complex
	\[
		A \to A_0 \to A_1 \to \dots
	\]
	by taking the alternating sums of the face maps (the unreduced version of Dold--Kan). We claim that this is an exact sequence. We may test this after base change along faithfully flat maps $A \to B$. Since our cosimplicial object is an fpqc-hypercover, we may then base change to objects after which the sequence becomes split, hence exact. Details are omitted here.

	It then follows (why?) that
	\[
		C_m = \cofib(A \to \lim_{\simp_{\leq m}} A_{\bullet}) \simeq \coker(A_{m-1} \to A_m)[m-1].
	\]
	Thus, \Cref{cl:CountablyGeneratedFlat} reduces to the claim that $\coker(A_{m-1} \to A_m)$ is a countably presented static flat $A$-module.

	Scholze says the following about this: ``Moreover, the same holds true after any base change in $A$, in particular along a quotient map $A \to A/I$; this shows that this cokernel is flat over $A$.''

	All in all, we have now shown that the functor $\D(A) \to \lim_{[n] \in \simp}^! \D(A_n)$ is fully faithful.

	\subsection{Step 2: Essential surjectivity}

	We now show the functor $\D(A) \to \lim_{[n] \in \simp}^! \D(A_n)$ is an equivalence.

	\textbf{Step 2a:} We first reduce to the case where $A_{\bullet}$ is a \v{C}ech cover. To this end, consider the following commutative square:
	\[
		\begin{tikzcd}
			A \rar \dar & A_i \dar \\
			A_0 \rar & A_0 \otimes_A A_i.
		\end{tikzcd}
	\]
	Since the cosimplicial object at the bottom is split, we have
	\[
		\D(A_0) \iso \lim^!_{[i] \in \simp} \D(A_0 \otimes_A A_i),
	\]
	and, more generally, $\D(A_0^{\otimes_A k}) \iso \lim^!_{[i] \in \simp} \D(A_0^{\otimes_A k} \otimes_A A_i)$ for all $k \geq 1$. So we are done if we can show that
	\[
		\D(A) \iso \lim^!_{[n] \in \simp} \D(A_0^{\otimes_A (n+1)}) \qquad \text{ and } \qquad \D(A_i) \iso \lim^!_{[n] \in \simp} \D(A_0^{\otimes_A (n+1)} \otimes_A A_i),
	\]
	which is the special case of \v{C}ech descent.

	\textbf{Step 2b:} Assume now that $A_n \cong A_0^{\otimes_A (n+1)}$ is the \v{C}ech cover of $A \to A_0$. To show that the map $\D(A) \to \lim_{[n] \in \simp}^! \D(A_n)$ is an isomorphism in $\PrR$, we may pass to left adjoints and show that
	\[
		\colim_{[n] \in \simp\catop} \D(A_n) \to \D(A)
	\]
	is an isomorphism in $\PrL$. Here the colimit is taken along the functors $(f_n)_*\colon \D(A_n) \to \D(A)$. Recall that each of these functors lives in $\Pr_{\D(A)} = \Mod_{\D(A)}(\Pr)$, and hence the colimit may be computed in $\Mod_{\D(A)}(\Pr)$. Also note that if we apply the functor $\D(A_0) \otimes_{\D(A)} -\colon \Mod_{\D(A)}(\Pr) \to \Mod_{\D(A_0)}(\Pr)$, the augmented simplicial diagram admits a splitting, since $\D(A_0) \otimes_{\D(A)} \D(A_n) \cong \D(A_{n+1})$. It follows that the map $\colim_{[n] \in \simp\catop} \D(A_n) \to \D(A)$ becomes an isomorphism after applying $\D(A_i) \otimes_{\D(A)} -$ for each $i$, and thus by passing to colimits also after applying $(\colim_i \D(A_i)) \otimes_{\D(A)} -$.

	Now, by Step 1 we know that the functor $\D(A) \to \lim^!_{\simp} \D(A_{\bullet}) \simeq \colim_* \D(A_{\bullet})$ is fully faithful. In particular, if we denote the image of the monoidal unit $\unit \in \D(A)$ by $M$, then we see that $M$ is sent to $\unit$ under the $\D(A)$-linear left adjoint $\colim_* \D(A_\bullet)$. If we consider the $\D(A)$-linear functor $\D(A) \to \colim_* \D(A_\bullet)$ determined by $M$, it thus follows that the composite
	\[
		\D(A) \to \colim_* \D(A_{\bullet}) \to \D(A)
	\]
	is the identity in $\Pr_{\D(A)}$, showing that $\D(A)$ is a retract of this colimit. It thus follows that the map $\colim_{[n] \in \simp\catop} \D(A_n) \to \D(A)$ becomes an isomorphism after applying $\D(A) \otimes_{\D(A)} -$. Since this is the identity functor, we are done.

	This finishes the proof in the countably presented case, for $n = 0$.

	\section{The induction step for \texorpdfstring{$n > 0$}{n > 0}}

	We now assume that $n > 0$, and that we have already proved the claim for $n - 1$. Observe that we again have an adjunction
	\[
		F_{n,\sharp}\colon \lim n\Pr_{A_{\bullet}} \rightleftarrows n\Pr_A \colon F_n^*.
	\]
	Also, in this case each of the individual adjunctions $n\Pr_A \rightleftarrows n\Pr_{A_{\bullet}}$ is an ambidextrous adjunction.

	We first show that $F_n^*$ is fully faithful, or equivalently that the counit map
	\[
		F_{n,\sharp} F_n^* \to \id
	\]
	is an isomorphism. Note that this adjunction is $n\Pr_A$-linear, so it suffices to show this claim on the monoidal unit: we need to show that the map $F_{n,\sharp} F_n^* (n-1)\Pr_A \to (n-1)\Pr_A$ is an isomorphism.

	For $n = 1$, this map is $F_{n,\sharp} F_n^* \D(A) \iso \D(A)$, which is the base case we just showed.

	For $n \geq 2$, we need to show that $F_{n,\sharp} F_n^* (n-1)\Pr_A \iso (n-1)\Pr_A$. One can show that the LHS is dualizable, so we may equivalently show that the dual map $(n-1)\Pr_A \to F_{n,\sharp} F_n^* (n-1)\Pr_A$ is an isomorphism. But this is a linear map, so it agrees with the map one level lower. This concludes the induction.

	Essential surjectivity is proved similarly to the previous situation, by reducing to the \v{C}ech-cover. (It's even a bit easier, as we have ambidexterity.)

	This finishes the proof.

	\section{Reduction to the countably presented case}

	Finally, we show that the claim for $n = 0$ can be reduced to the countably presented case. We use the following lemma:

	\begin{lemma}\label[lemma]{lem:Faithfully_Flat_Approximation}
		Every faithfully flat map $A \to B$ of commutative ring spectra is an $\aleph_1$-filtered colimit of countably presented such maps.
	\end{lemma}

	\begin{proof}
		First, if $A$ is static, then so is $B$. This case is due to Waterhouse;\footnote{W.\ C.\ Waterhouse, \emph{Basically bounded functors and flat sheaves}, Pacific J.\ Math.\ \textbf{57} (1975), no.~2, 597--610.} let us quickly recall the argument. By Lazard's theorem, the flat $A$-modules are the Ind-category of the finite projective $A$-modules, so we can write the underlying module of $B$ as an $\aleph_1$-filtered colimit of countably presented flat modules $M_i$. Since each $M_i$ is countably presented, it is $\aleph_1$-compact as an $A$-module, so $\Hom_A(M_i, -)$ commutes with $\aleph_1$-filtered colimits. We will use two consequences of this: any map from a countably presented $A$-module into $B = \colim_i M_i$ factors through some $M_i$, and any two such maps that agree after composing to $B$ already agree at some stage $M_i$.

		The tensor product $M_i \otimes_A M_i$ is again countably presented, so the composite $M_i \otimes_A M_i \to B \otimes_A B \to B$ of the canonical map with the multiplication of $B$ factors through some $M_j$, $j = j(i)$, compatibly with the maps to $B$; similarly the unit $A \to B$ factors through some $M_{i_0}$. Choosing inductively a sequence $i_0, i_1, \ldots$ with $i_{n+1} \geq j(i_n)$ yields compatible maps $M_{i_n} \otimes_A M_{i_n} \to M_{i_{n+1}}$, and the colimit $M := \colim(M_{i_0} \to M_{i_1} \to \ldots)$ acquires a multiplication
		\[
			M \otimes_A M = \colim_n ( M_{i_n} \otimes_A M_{i_n} ) \to \colim_n M_{i_{n+1}} = M
		\]
		and a unit $A \to M_{i_0} \to M$, both compatible with $M \to B$.

		The commutativity, associativity, and unitality of this multiplication each amount to an equality of two maps out of a countably presented module ($M_{i_n}^{\otimes 2}$, resp.\ $M_{i_n}^{\otimes 3}$, resp.\ $M_{i_n}$) that holds after composing into the commutative $A$-algebra $B$. By the second consequence above, each such identity holds after passing to a later stage; as there are only countably many of them, they can all be arranged to hold on a single cofinal subsystem. On such a subsystem $M$ is a countably presented flat $A$-subalgebra of $B$, and these subalgebras form an $\aleph_1$-filtered system with colimit $B$. This writes $B$ in the desired form.

		In general, write the map $A \to B$ as the $\aleph_1$-filtered colimit of all countably presented $A$-algebras $B_i$ mapping to $B$. We want to see that the set of $i$ for which $A \to B_i$ is faithfully flat is cofinal. We get a similar presentation of $\pi_0 A$ as the colimit of the $\pi_0 B_i$, and by the static case, for any $i$ there is some $j$ so that $\pi_0 B_i \to \pi_0 B_j$ factors over a faithfully flat $\pi_0 A$-algebra. Taking a sequence $i_0, i_1, \ldots$ so that each map $\pi_0 B_{i_0} \to \pi_0 B_{i_1} \to \ldots$ factors over a faithfully flat $\pi_0 A$-algebra, the colimit $B_j = \colim(B_{i_0} \to B_{i_1} \to \ldots)$ is still countably presented, and has the property that $\pi_0 A \to \pi_0 B_j$ is faithfully flat. Using a similar argument, we can arrange that $\pi_0 B_i \otimes_{\pi_0 A} \pi_n A \to \pi_n B_i$ is an isomorphism for any given $i$, and hence for all $i$ by a further countable colimit.
	\end{proof}

	As a consequence, every fpqc-hypercover is an $\aleph_1$-filtered colimit of countably presented fpqc-hypercovers.

	Now, given any hypercover $\Spec(A_{\bullet}) \to \Spec(A)$, we have
	\[
		\lim_{\simp}^! n\Pr_{A_\bullet} \simeq \lim^!_{\simp} \colim_r n\Pr_{A_{r,\bullet}} \simeq \colim_r \lim_{\simp}^! n\Pr_{A_{r,\bullet}} \simeq n\Pr_A,
	\]
	where we use that $1\Pr$ is $\aleph_1$-compactly generated so that $\aleph_1$-filtered colimits commute with countable limits. This finishes the proof of the general form of the theorem.

	\chapter{Quasicoherent sheaves of higher categories II}
	\label[chapter]{chap:QCoh_II}

	\emph{Talk given by Bastiaan.}

	We now pass from affine spectral schemes ($\Aff_{\S} := \CAlg(\Sp)\catop$) to fpqc-stacks.

	On affines, Talk 2 introduced the functors
	\[
		\Aff_{\S}\catop \to \CAlg((n+1)\Pr) = \CAlg(1\Pr)_{n\Pr/}, \qquad \Spec(A) \mapsto n\Pr_A
	\]
	for all $n \geq 0$, and showed they canonically extend to six-functor formalisms in which the countably presented maps are exactly the $!$-able maps, with $f_! = f_*$. Denote by
	\[
		\Shv_{\fpqc}(\Aff_{\S}) \; \subseteq \; \PSh(\Aff_{\S})
	\]	
	the subcategory of \emph{small} fpqc-sheaves. Here \emph{small} means that the sheaf is a small colimit of representable objects. (In fact, we will mostly restrict to sheaves which are countable colimits of affines.) By the main result of Talk 3, we obtain a unique limit-preserving extension
	\[
		\Shv_{\fpqc}(\Aff_{\S})\catop \to \CAlg(1\Pr)_{n\Pr/}, \qquad X \mapsto n\Pr_X.
	\]
	Since the endomorphism functor $\End\colon (n+1)\Pr \to n\Pr$ commutes with limits, this means that every fpqc-stack $X$ defines a Stefanich ring
	\[
		\Oo(X) := (\Omega^{\infty}\Gamma(X), \D(X), 1\Pr_X, 2\Pr_X, \ldots),
	\]
	with $n$-th component $\Oo(X)_n := (n-1)\Pr_X$ (using the low-degree convention $0\Pr_X := \D(X)$). Here $\Gamma(X) := \End_{\D(X)}(\unit_{\D(X)}) \in \CAlg(\Sp)$ denotes the spectrum of global functions on $X$; more generally, for a map $f\colon Y \to X$ we write
	\[
		\Oo(Y/X)_0 := f_*\unit_{\D(Y)} \in \D(X)
	\]
	for the level-zero relative structure object.\footnote{This notation is chosen to be compatible with the later notation $(\Oo(Y)/\Oo(X))_m$ for the relative objects of a map of Gestalten.} This extends the affine assignment $\Spec(A) \mapsto \Oo(\Spec(A))$ to a functor
	\[
		\Oo\colon \Shv_{\fpqc}(\Aff_{\S})\catop \to \StRing.
	\]

	For $n = 0$, the value $0\Pr_X = \D(X)$ is the usual derived $\infty$-category of quasicoherent sheaves on $X$, at least when $X$ is a countable colimit of affines.

	\begin{remark}
		In general, there is a difference between taking the limit defining $n\Pr_X$ inside $1\Pr = \PrL_{\aleph_1}$ or inside the larger category $\PrL$ of all presentable categories. Usually, $\QCoh(X)$ is defined using limits in $\PrL$. For countable colimits of affines the two limits agree.
	\end{remark}

	The rest of the chapter has four goals:
	\begin{itemize}
		\item For $n \geq 1$, we extend the assignment $X \mapsto n\Pr_X$ to a six-functor formalism, with $!$-able maps given by so-called `bi-countably presented maps';
		\item We show that ``every bi-countably presented map is étale'', and for $n \geq 2$ that in fact ``every bi-countably presented map is proper'';
		\item We deduce a Künneth formula for sheaves of $(n-1)$-presentable categories on a fiber product;
		\item We conclude that the Stefanich ring functor preserves countable colimits and finite limits.
	\end{itemize}

	\section{Bi-countably presented maps}

	We will now extend the assignment $X \mapsto n\Pr_X$ to a six-functor formalism. We may think of this as a `higher' version of the six-functor formalism on $X \mapsto \D(X)$. A first question is: what are the $!$-able maps? It turns out that at the higher categorical level, the class of $!$-able maps contains essentially \emph{all} maps. The relevant class of maps is the following:

	\begin{definition}\label[definition]{def:Bi_Countably_Presented}
		Let $A \in \CAlg(\Sp)$. A map $Z \to \Spec(A)$ of fpqc-sheaves is \emph{bi-countably presented over $\Spec(A)$} if $Z$ is a countable colimit, in $\Shv_{\fpqc}(\Aff_{\S})_{/\Spec(A)}$, of affines $\Spec(B_i)$ with $A \to B_i$ countably presented.

		A map $f\colon Y \to X$ of fpqc-stacks is \emph{bi-countably presented} if for every affine $\Spec(A) \to X$, the fiber product $Y \times_X \Spec(A)$ is bi-countably presented over $\Spec(A)$. Let $E$ denote the class of bi-countably presented maps.
	\end{definition}

	\begin{lemma}\label[lemma]{lem:Section_Of_Bicountably_Presented_Affine_Base}
		Let $S = \Spec(A)$, and let $p\colon Z \to S$ be bi-countably presented. Then every section $s\colon S \to Z$ of $p$ is bi-countably presented as a map to $Z$.
	\end{lemma}
	\begin{proof}
		We first record two elementary fpqc-local reductions. Let
		$S'=\Spec(A') \to S=\Spec(A)$ be faithfully flat and countably
		presented. A map $W \to S$ is bi-countably presented if and only if its
		base change $W_{S'} \to S'$ is bi-countably presented. The forward
		implication is base change. Conversely, write
		$S'_{\bullet}$ for the \v{C}ech nerve of $S' \to S$. Since
		$A \to A'$ is countably presented, each $S'_m \to S$ is countably
		presented by \Cref{lem:Affine_Countably_Presented_Calculus}. If
		$W_{S'}$ is a countable colimit of affines countably presented over
		$S'$, then each $W_{S'_m}$ is a countable colimit of affines countably
		presented over $S$. By fpqc-descent,
		\[
			W \simeq \abs{W_{S'_{\bullet}}},
		\]
		and the indexing category $\simp\catop$ is countable. Hence $W$ is
		again a countable colimit of affines countably presented over $S$.

		Second, let
		\[
			Z \simeq \colim_{i \in I} \Spec(B_i)
		\]
		be a countable affine presentation over $S$, and let
		$T=\Spec(C) \to S$ be affine. Any finite collection of sections
		$T \to Z_T:=Z\times_S T$ factors, after a faithfully flat countably
		presented base change $T' \to T$, through stages of the induced
		presentation
		\[
			Z_T \simeq \colim_{i \in I} \Spec(B_i\otimes_A C).
		\]
		Indeed, by the definition of sheafification, this factorization exists
		after some fpqc-cover of $T$. By
		\Cref{lem:Faithfully_Flat_Approximation}, and because the stages are
		countably presented over $T$, this cover can be refined by a
		faithfully flat countably presented affine cover. For finitely many
		sections we take a common refinement.

		We now prove the lemma. Let $T=\Spec(C) \to Z$ be an affine test
		object. We must show that
		\[
			T \times_Z S \to T
		\]
		is bi-countably presented. After base change along $T \to S$, the map
		$Z_T \to T$ has two sections: the section $t$ induced by the given map
		$T \to Z$, and the section $s_T$ obtained by pulling back $s$. By the
		fpqc-local reductions above, we may pass to a faithfully flat
		countably presented cover of $T$ and assume that both sections factor
		through stages of a countable presentation of $Z_T$.

		Let $\Aff_T^{\mathrm{cp}}$ be the full subcategory of affine
		$T$-schemes which are countably presented over $T$. It has finite
		limits by \Cref{lem:Affine_Countably_Presented_Calculus}. Let
		$\widetilde Z$ be the corresponding countable colimit of representables
		in $\PSh(\Aff_T^{\mathrm{cp}})$. The two sections $s_T,t\colon T \to
		Z_T$ now come from maps
		\[
			\ast \rightrightarrows \widetilde Z
		\]
		in this presheaf category, where $\ast$ denotes the terminal
		representable. Consider the presheaf fiber product
		\[
			P := \ast \times_{\widetilde Z} \ast.
		\]
		The objects $\ast$ and $\widetilde Z$ are $\aleph_1$-compact in
		$\PSh(\Aff_T^{\mathrm{cp}})$, and
		\Cref{prop:Kappa_Compact_Presheaves_Finite_Limits} shows that $P$ is
		again $\aleph_1$-compact. Since presheaf categories are generated by
		representables, $P$ is a countable colimit of representables from
		$\Aff_T^{\mathrm{cp}}$. Sheafification preserves finite limits, so the
		sheafification of $P$ is precisely
		\[
			T \times_{Z_T} T \simeq T \times_Z S.
		\]
		Thus, fpqc-locally on $T$, the map $T \times_Z S \to T$ is
		bi-countably presented. The first reduction descends this conclusion
		to $T$ itself.
	\end{proof}

	\begin{proposition}
		\label[proposition]{prop:Closure_Properties_Bicountably_Presented}
		The class $E$ is closed under composition, base change, sections, and formation of diagonals.
	\end{proposition}
	\begin{proof}		
		Closure under base change is immediate from the definition. For closure under composition, let $Z \to Y \to X$ be two maps in $E$ and let $\Spec(A) \to X$ be affine. Write
		\[
			Y_A := Y \times_X \Spec(A) \simeq \colim_{i \in I} \Spec(B_i),
		\]
		where $I$ is countable and each $A \to B_i$ is countably presented. By base change, $Z_A \to Y_A$ is bi-countably presented. Applying the definition of bi-countably presented to the affine test objects $\Spec(B_i) \to Y_A$, each
		\[
			Z_A \times_{Y_A} \Spec(B_i)
		\]
		is a countable colimit of affines $\Spec(C_{i,k})$ with $B_i \to C_{i,k}$ countably presented. Since countably presented maps are closed under composition, each $A \to C_{i,k}$ is countably presented. Universality of colimits in sheaves gives
		\[
			Z_A \simeq \colim_{i \in I} \left(Z_A \times_{Y_A} \Spec(B_i)\right)
				\simeq \colim_{(i,k)} \Spec(C_{i,k}),
		\]
		and the final indexing category is still countable. Hence $Z \to X$ lies in $E$.
		
		It remains to discuss sections and diagonals. The diagonal $\Delta_f\colon Y \to Y \times_X Y$ is a section of the first projection
		\[
			Y \times_X Y \to Y.
		\]
		This projection is a base change of $f$, hence lies in $E$. Since $E$ is already closed under base change, it is therefore enough to know that $E$ is closed under sections.
		
		Let $p\colon Y \to X$ lie in $E$, let $s\colon X \to Y$ be a section, and test $s$ against an affine map $a\colon \Spec(A) \to Y$. Base change $p$ along the composite $\Spec(A) \xrightarrow{a} Y \xrightarrow{p} X$. Since $p \in E$, the resulting map $Y_A := Y \times_X \Spec(A) \to \Spec(A)$ is bi-countably presented over $\Spec(A)$. The map $a$ gives one section of $Y_A \to \Spec(A)$, and $s \circ p \circ a$ gives another. Moreover
		\[
			X \times_Y \Spec(A) \simeq \Spec(A) \times_{Y_A} \Spec(A)
		\]
		with respect to these two sections. By \Cref{lem:Section_Of_Bicountably_Presented_Affine_Base}, the first section $\Spec(A) \to Y_A$ is bi-countably presented as a map to $Y_A$. Pulling it back along the second section shows that
		\[
			\Spec(A) \times_{Y_A} \Spec(A) \to \Spec(A)
		\]
		is bi-countably presented over $\Spec(A)$. Thus $s \in E$, so $E$ is closed under sections and therefore under diagonals.
	\end{proof}

	\section{A six-functor formalism for \texorpdfstring{$n\Pr$}{nPr} at \texorpdfstring{$n \geq 1$}{n >= 1}}

	We fix $n \geq 1$. Given a map $f\colon Y \to X$ of fpqc-stacks, we denote the resulting symmetric monoidal pullback functor by
	\[
		f_n^*\colon n\Pr_X \to n\Pr_Y,
	\]
	and its right adjoint by $f_{n,*}$. We denote by $\otimes_{n\Pr_X}$ the tensor product of $n\Pr_X$.

	\begin{remark}
		Throughout this talk, we will frequently abuse notation and regard $n\Pr_Y$ as an object of $(n+1)\Pr_X$ by setting
		\[
			n\Pr_Y := f_{n+1,*}(\unit_Y) \qin (n+1)\Pr_X.
		\]
		For $p_X\colon X \to \pt$, we get a `forgetful functor' $p_{X,n+1,*} \colon (n+1)\Pr_X \to (n+1)\Pr$, and we claim that the image of $f_{n+1,*}(\unit_Y)$ under this functor is $n\Pr_Y$. Indeed, if we write $Y \simeq \colim_{\Spec(A) \to Y} \Spec(A)$, then we may write $p_{X,n+1,*}f_{n+1,*}(\unit) = p_{Y,n+1,*}(\unit)$ as a limit over $p_{A,n+1,*}(\unit) = n\Pr_A$, which is $n\Pr_Y$ by definition.
	\end{remark}

	\begin{theorem}\label[theorem]{thm:Six_Functors_For_nPr}
		Let $f\colon Y \to X$ be a bi-countably presented map of fpqc-stacks, and let $n \geq 1$.
		\begin{enumerate}[(i)]
			\item The functor
			\[
				f_n^*\colon n\Pr_X \to n\Pr_Y
			\]
			admits a left adjoint $f_{n,\sharp}$ in $(n+1)\Pr_X$. In particular, $f_{n,\sharp}$ commutes with all small colimits and with tensoring.
			\item The object $n\Pr_Y \in (n+1)\Pr_X$ is $\aleph_1$-compact.
			\item The formation of $f_{n,\sharp}$ commutes with base change: for every cartesian square
			\[
				\begin{tikzcd}
					Y' \rar["v"] \dar["f'"'] & Y \dar["f"] \\
					X' \rar["u"'] & X
				\end{tikzcd}
			\]
			the Beck--Chevalley transformation
			\[
				u_n^* f_{n,\sharp} \to f'_{n,\sharp} v_n^*
			\]
			is an isomorphism.
			\item For $n \geq 2$, the object $f_{n,\sharp}(\unit) \in n\Pr_X$ is dualizable.
			\item For $n \geq 2$, the left adjoint $f_{n,\sharp}$ is also \emph{right} adjoint to $f_n^*$, i.e.\ $f_{n,\sharp} \simeq f_{n,*}$. In particular, $f_{n,*}$ also commutes with small colimits, tensoring, and base change.
		\end{enumerate}
	\end{theorem}

	In other words: at level $n \geq 1$, ``all maps are étale'', and at level $n \geq 2$, ``all maps are proper''.

	\begin{proof}
		\textbf{Step 1: The affine case.} We first assume that $X = \Spec(A)$ is affine and prove (i)--(iv). By assumption $Y = \colim_{i \in I} \Spec(B_i)$ is a countable colimit of affines $Y_i = \Spec(B_i)$, where each $B_i$ is countably presented over $A$. Consider
		\[
			f_n^*\colon n\Pr_X = n\Pr_A \to n\Pr_Y = \lim_{i \in I\catop} n\Pr_{B_i}.
		\]
		Since $n \geq 1$ and all transition maps are countably presented, the result from Talk 2 shows that all the pullback functors involved admit linear left adjoints. In particular, the limit defining $n\Pr_Y$ is along right adjoints, and hence agrees with the colimit along the corresponding left adjoints:
		\[
			n\Pr_Y \simeq \colim_{i \in I, \sharp} n\Pr_{B_i}.
		\]
		(We use countability of $I$ here so that this colimit takes place inside the $\aleph_1$-presentable setting $1\Pr$.) In this colimit description, the desired left adjoint is given as the colimit of the lower-$\sharp$ maps $n\Pr_{B_i} \to n\Pr_A$:
		\[
			f_{n,\sharp}\colon n\Pr_Y \simeq \colim_{i \in I, \sharp} n\Pr_{B_i} \to n\Pr_A.
		\]
		(See \citet[Corollary~1.3]{HorevYanovski2017Descent}.) This is a map in $(n+1)\Pr_A = \Mod_{n\Pr_A}(1\Pr)$, and in particular it commutes with all small colimits and with tensoring. This gives (i). We may explicitly write the resulting functor as
		\[
			f_{n,\sharp}\colon \lim_{i \in I\catop} n\Pr_{B_i} \to n\Pr_A, \qquad (Z_i) \mapsto \colim_{i \in I} (g_i)_{n,\sharp}(Z_i),
		\]
		with $g_i\colon \Spec(B_i) \to \Spec(A)$ the affine structure map.

		For (ii), since $\aleph_1$-compact objects in $(n+1)\Pr_A$ are closed under countable colimits, the $\aleph_1$-compactness of $n\Pr_Y$ follows from the $\aleph_1$-compactness of the $n\Pr_{B_i}$ established in Talk 2.
		
		For (iii), let $A \to A'$ be a base change, and write $B_i' := B_i \otimes_A A'$ and $Y' := Y \times_{\Spec(A)} \Spec(A')$. Since colimits in sheaves are universal,
		\[
			Y' \simeq \colim_{i \in I} \Spec(B_i').
		\]
		Hence the same construction applied over $A'$ gives
		\[
			f'_{n,\sharp}\colon n\Pr_{Y'} \to n\Pr_{A'}, \qquad (Z'_i) \mapsto \colim_{i \in I} (g'_i)_{n,\sharp}(Z'_i),
		\]
		where $g'_i\colon \Spec(B_i') \to \Spec(A')$ is the base change of $g_i\colon \Spec(B_i) \to \Spec(A)$. The base-change property is known from Talk 2:
		\[
			u_n^* (g_i)_{n,\sharp} \simeq (g_i')_{n,\sharp} (v_i)_n^*.
		\]
		Thus base change carries the diagram whose colimit defines $f_{n,\sharp}$ to the diagram whose colimit defines $f'_{n,\sharp}$. Since the indexing category $I$ is countable, this colimit is taken inside $1\Pr$, and we obtain the desired Beck--Chevalley isomorphism
		\[
			u_n^* f_{n,\sharp} \simeq f'_{n,\sharp} v_n^*.
		\]

		For (iv), we evaluate at the unit. In the affine case, we have $(g_i)_{n,\sharp} = (g_i)_{n,*}$ by ambidexterity from Talk 2, so
		\[
			f_{n,\sharp}(\unit) \simeq \colim_{i \in I} (g_i)_{n,*}(\unit) \simeq \colim_{i \in I,*} (n-1)\Pr_{B_i},
		\]
		where the colimit is along the lower-$\ast$ functors. For $n \geq 2$ (so that $n - 1 \geq 1$), the lower-$\ast$ functors $(n-1)\Pr_{B_i} \to (n-1)\Pr_{B_j}$ are themselves linear left adjoints. Hence the colimit above is a countable colimit of dualizable objects along internal left adjoint functors. Such a colimit is itself dualizable, by the following more general result:

		\textbf{Claim:} For $\Vv \in \CAlg(1\Pr)$, countable colimits of dualizable objects along internal left adjoints in $\Pr_{\Vv}$ are again dualizable.

		This is a result by Ramzi. In short, the reason is that the dualizable objects can be characterized as those $C$ that are retracts of $\Pp_{\Vv}(C^{\aleph_1})$ in a canonical way. Since this retract diagram is functorial in $C$, it shows closure under countable colimits.

		\textbf{Step 2: The general case.} For general $X$, we may write $X$ as a colimit of affine spectral schemes, providing an equivalence
		\[
			(n+1)\Pr_X \simeq \lim_{\Spec(A) \to X} (n+1)\Pr_A.
		\]
		This is a limit in $1\Pr = \PrL_{\aleph_1} \simeq \Rex_{\aleph_1}$, so on the level of underlying categories we have
		\[
			((n+1)\Pr_{X})^{\aleph_1} = \lim_{\Spec(A) \to X} ((n+1)\Pr_{A})^{\aleph_1}.
		\]
		Letting $Y_A := \Spec(A) \times_X Y$, we have already established that $n\Pr_{Y_A}$ is $\aleph_1$-compact in $(n+1)\Pr_A$, thus giving the $\aleph_1$-compactness of $n\Pr_Y$ in $(n+1)\Pr_X$. Moreover, to check that $f_n^*\colon n\Pr_X \to n\Pr_Y$ admits a left adjoint in $((n+1)\Pr_{X})^{\aleph_1}$, it will suffice to show that all maps $(f_A)_n^*\colon n\Pr_A \to n\Pr_{Y_A}$ admit left adjoints compatible with base change, which we established in Step 1. Similarly, the dualizability of $f_{n,\sharp}(\unit) \in n\Pr_X$ may be checked after pullback to affines.

		\textbf{Step 3: Ambidexterity.} For (v), we use the general ambidexterity result (\Cref{prop:General_Ambidexterity}) again: it suffices to check that the unit and counit of the adjunction $f_n^* \dashv f_{n,\sharp}$ admit right adjoints. For the counit $f_{n,\sharp} f_n^* \to \id$, we use that under the equivalence $\End_{(n+1)\Pr_X}(n\Pr_X,n\Pr_X) \simeq n\Pr_X$ this map corresponds to the map $f_{n,\sharp}(\unit) \to \unit$, so we need to show this admits a right adjoint in $n\Pr_X$. By (iv), the object $f_{n,\sharp}(\unit)$ is dualizable, so this is equivalent to its dual $\unit \to f_{n,\ast}(\unit)$ admitting a left adjoint in $n\Pr_X$. But this dual map is the map $(n-1)\Pr_X \to (n-1)\Pr_Y$, which has a left adjoint in $n\Pr_X$ by part (i) applied one level down (which uses $n - 1 \geq 1$).

		For the unit $\eta\colon \id \to f_n^* f_{n,\sharp}$ in $\End_{(n+1)\Pr_X}(n\Pr_Y, n\Pr_Y)$ we use a base-change argument. Consider the cartesian square
		\[
			\begin{tikzcd}
				Y \times_X Y \rar["p_2"] \dar["p_1"'] & Y \dar["f"] \\
				Y \rar["f"'] & X
			\end{tikzcd}
		\]
		together with the diagonal $\delta\colon Y \to Y \times_X Y$. Since the class $E$ is closed under formation of diagonals, $\delta$ is bi-countably presented, so the conclusions of parts (i)--(iv) are available for $\delta$ at level $n$. By base change (part (iii)), there is a canonical isomorphism
		\[
			f_n^* f_{n,\sharp} \simeq p_{1,n,\sharp}\, p_{2,n}^*.
		\]
		% Using $p_1 \circ \delta = p_2 \circ \delta = \id_Y$, we have $\delta_n^* p_{i,n}^* = \id$ and $p_{i,n,\sharp} \delta_{n,\sharp} = \id$ for $i = 1, 2$, and hence
		% \[
		% 	\id_{n\Pr_Y} = p_{1,n,\sharp}\, \delta_{n,\sharp}\, \delta_n^*\, p_{2,n}^*.
		% \]
		The level-$n$ counit $\epsilon_{\delta}\colon \delta_{n,\sharp} \delta_n^* \to \id_{n\Pr_{Y \times_X Y}}$ of $\delta_{n,\sharp} \dashv \delta_n^*$ then induces a $2$-morphism
		\[
			p_{1,n,\sharp}\, \epsilon_{\delta}\, p_{2,n}^*\colon \id_{n\Pr_Y} = p_{1,n,\sharp}\, \delta_{n,\sharp}\, \delta_n^*\, p_{2,n}^* \, \to p_{1,n,\sharp}\, p_{2,n}^* \simeq f_n^* f_{n,\sharp}.
		\]
		This $2$-morphism agrees with the unit $\eta_f$, which one verifies by checking that together with $\epsilon_f$ it satisfies the triangle identities (a consequence of base change). But since we already showed that the counit $\epsilon_{\delta}$ is a left adjoint and whiskering with the functors $p_{1,n,\sharp}$ and $p_{2,n}^*$ preserves adjunctions, this finishes the proof.
	\end{proof}

	\section{The Künneth formula}
	\label[section]{sec:Kunneth_Formula}

	As a consequence of the theorem, we obtain the following Künneth formula:

	\begin{corollary}\label[corollary]{cor:Kunneth_Higher_Pr}
		Let $f\colon Y \to X$ and $f'\colon Y' \to X$ be bi-countably presented maps, and let $n \geq 2$. Then the objects $(n-1)\Pr_Y$ and $(n-1)\Pr_{Y'}$ in $n\Pr_X$ are dualizable, and the \emph{Künneth formula}
		\[
			(n-1)\Pr_Y \otimes_{n\Pr_X} (n-1)\Pr_{Y'} \iso (n-1)\Pr_{Y \times_X Y'}
		\]
		holds, where the equivalence is implemented via the external tensor product.
	\end{corollary}

	\begin{proof}
		By definition, $(n-1)\Pr_Y$ as an object of $n\Pr_X$ is $f_{n,*}(\unit_{n\Pr_Y})$. Since $n \geq 2$, \Cref{thm:Six_Functors_For_nPr} gives ambidexterity $f_{n,*} \simeq f_{n,\sharp}$, and hence
		\[
			(n-1)\Pr_Y \simeq f_{n,\sharp}(\unit_{n\Pr_Y}).
		\]
		The same theorem then implies that this object is dualizable in $n\Pr_X$; similarly for $Y'$.

		For the Künneth formula, write $Z := Y \times_X Y'$, with projections
		\[
			\begin{tikzcd}
				Z \rar["q"] \dar["p"'] & Y' \dar["f'"] \\
				Y \rar["f"'] & X .
			\end{tikzcd}
		\]
		The maps $p$, $q$, and $f \circ p = f' \circ q$ are again bi-countably presented by \Cref{prop:Closure_Properties_Bicountably_Presented}, so the relevant $\sharp$-functors are available. Using the projection formula for $f_{n,\sharp}$, base change for the cartesian square above, and then the projection formula for $p_{n,\sharp}$, we then compute
		\[
		\begin{aligned}
			(n-1)\Pr_Y \otimes_{n\Pr_X} (n-1)\Pr_{Y'}
			&\simeq
			f_{n,\sharp}(\unit_{n\Pr_Y}) \otimes_{n\Pr_X} f'_{n,\sharp}(\unit_{n\Pr_{Y'}}) \\
			&\simeq
			f_{n,\sharp}\bigl(\unit_{n\Pr_Y} \otimes_{n\Pr_Y} f_n^* f'_{n,\sharp}(\unit_{n\Pr_{Y'}})\bigr) \\
			&\simeq
			f_{n,\sharp}\bigl(\unit_{n\Pr_Y} \otimes_{n\Pr_Y} p_{n,\sharp} q_n^*(\unit_{n\Pr_{Y'}})\bigr) \\
			&\simeq
			f_{n,\sharp} p_{n,\sharp}\bigl(p_n^*(\unit_{n\Pr_Y}) \otimes_{n\Pr_Z} q_n^*(\unit_{n\Pr_{Y'}})\bigr) \\
			&\simeq
			(f \circ p)_{n,\sharp}(\unit_{n\Pr_Z}) \\
			&\simeq
			(n-1)\Pr_Z.
		\end{aligned}
		\]
		In the penultimate line we used that pullback preserves monoidal units. Since $Z = Y \times_X Y'$, this is the desired equivalence.
	\end{proof}

	In particular, for $n \geq 2$, sheaves of presentable $(n-1)$-categories on a fiber product are determined by what they are on the factors and what sheaves of presentable $n$-categories are on the base. This is the first genuinely surprising statement about quasicoherent sheaves of higher categories: \emph{both} parts of \Cref{cor:Kunneth_Higher_Pr} fail at lower categorical levels:
	\begin{itemize}
		\item At level $n = 0$, it would say that $\Oo(Y/X)_0 \in \D(X)$ is dualizable (i.e.\ a perfect complex), or that $\Oo(Y/X)_0 \otimes_{\D(X)} \Oo(Y'/X)_0 \simeq \Oo((Y \times_X Y')/X)_0$, both of which are absurd without strong assumptions on $Y$ and $Y'$.
		\item At level $n = 1$, it is often the case that $\D(Y) \in 1\Pr_X$ is dualizable and the Künneth formula holds (cf.\ work by Gaitsgory), but one still needs some assumption such as ``$Y$ is a geometric stack over $X$''.
		\item For $n \geq 2$, these results are essentially \emph{always} true.
	\end{itemize}
	However, we cannot stay on a fixed categorical level (say $1\Pr$): the Künneth formula always uses the next categorical level on the base $X$.

	\section{Dropping the bi-countable-presentation hypothesis}

	We will now see that for some of the previous assertions we may drop the condition on the map being bi-countably presented if instead we require sufficient compactness conditions on $X$ and $Y$.

	\begin{proposition}\label[proposition]{prop:Kunneth_Aleph1_Compact}
		Let $f\colon Y \to X$ be a map between $\aleph_1$-compact objects of $\Shv_{\fpqc}(\Aff_{\S})$ (so both $X$ and $Y$ are countable colimits of affine spectral schemes by \Cref{lem:Affines_Aleph1_Compact_fpqc} below). For $n \geq 2$, the right adjoint
		\[
			f_{n,*}\colon n\Pr_Y \to n\Pr_X
		\]
		of $f_n^*$ admits a further right adjoint in $\Mod_{n\Pr_X}(\PrL)$, which is compatible with base change. Moreover, for any $f'\colon Y' \to X$ satisfying the same hypothesis, the Künneth formula
		\[
			(n-1)\Pr_Y \otimes_{n\Pr_X} (n-1)\Pr_{Y'} \simeq (n-1)\Pr_{Y \times_X Y'}
		\]
		holds.
	\end{proposition}

	\begin{remark}
		If $X$ is affine, the conclusion says that $f_{n,*}$ is a map in $\Mod_{n\Pr_A}(\PrL)$, but \emph{not} in general in $\Mod_{n\Pr_A}(\PrL_{\aleph_1}) = \Mod_{n\Pr_A}(1\Pr)$: it fails only to preserve $\aleph_1$-compact objects.
	\end{remark}

	The proof of the proposition requires various auxiliary lemmas, which are given afterwards.

	\begin{proof}
		We first reduce to the case where $X$ is affine. Since $X$ is $\aleph_1$-compact, by \Cref{lem:Affines_Aleph1_Compact_fpqc} it is a countable colimit of affines,
		\[
			X \simeq \colim_j \Spec(A_j).
		\]
		By \Cref{lem:Aleph1_Compact_fpqc_Finite_Limits},
		$\aleph_1$-compact objects of $\Shv_{\fpqc}(\Aff_{\S})$ are closed under
		finite limits; hence each base change
		$Y_j := Y \times_X \Spec(A_j)$ is again $\aleph_1$-compact.
		For such countable colimits, the descent formula
		\[
			n\Pr_X \simeq \lim_j n\Pr_{A_j}
		\]
		holds not only in $1\Pr = \PrL_{\aleph_1}$, but also in $\PrL$, hence on
		the level of underlying categories. Thus the desired
		assertions may be checked after pullback to the affine pieces
		$\Spec(A_j) \to X$. We may therefore assume that $X = \Spec(A)$.

		By \Cref{lem:Aleph1_Compact_Pro_Bicountably_Presented} below, we may
		choose a diagram $(Y_i)_{i \in I}$ over $X$, together with a compatible
		cone $Y \to Y_i$, such that each $f_i\colon Y_i \to X$ is
		bi-countably presented and
		\[
			n\Pr_Y \simeq \colim_{i \in I} n\Pr_{Y_i},
		\]
		where the transition maps are given by pullback. The pullback
		$f_n^*$ is a morphism in $1\Pr$, hence preserves
		$\aleph_1$-compact objects. For each $i$, the map $f_i$ is
		bi-countably presented. Hence, by
		\Cref{thm:Six_Functors_For_nPr}, for $n \geq 2$ the functor
		\[
			f_{i,n,*}\colon n\Pr_{Y_i} \to n\Pr_X
		\]
		preserves colimits, is $n\Pr_X$-linear, and satisfies base change. By
		\Cref{lem:Right_Adjoint_From_Filtered_Pro_Presentation} below, if
		$\alpha_i\colon n\Pr_{Y_i} \to n\Pr_Y$ denotes the colimit structure
		functor and $p_{ij}\colon Y_j \to Y_i$ denotes the transition map
		associated to a morphism $i \to j$ in $I$, then
		\[
			f_{n,*}(\alpha_i E)
			\simeq
			\colim_{j \in I_{i/}} f_{j,n,*}(p_{ij,n}^*E).
		\]
		Thus $f_{n,*}$ is assembled from the functors $f_{j,n,*}$ by taking
		$\aleph_1$-filtered colimits. Since filtered colimits of functors
		preserve colimits, $n\Pr_X$-linearity, and base change whenever the
		functors in the system do, it follows that $f_{n,*}$ also commutes with
		small colimits, tensors by $n\Pr_X$, and base change.

		With these properties established, the Künneth formula is proved formally just like in \Cref{cor:Kunneth_Higher_Pr}.
	\end{proof}

	\begin{lemma}
		\label[lemma]{lem:fpqc_Inclusion_Aleph1_Filtered_Colimits}
		\label[lemma]{lem:Affines_Aleph1_Compact_fpqc}
		\label[lemma]{lem:Aleph1_Compact_fpqc_Finite_Limits}
		The inclusion
		\[
			\Shv_{\fpqc}(\Aff_{\S}) \hookrightarrow \PSh(\Aff_{\S})
		\]
		preserves $\aleph_1$-filtered colimits. The same holds for the
		corresponding fpqc-sheaf categories on affines over any fixed affine
		base.

		Moreover, these sheaf categories are $\aleph_1$-compactly generated by
		affines, their $\aleph_1$-compact objects are closed under finite
		limits, and sheafification sends $\aleph_1$-compact presheaves to
		$\aleph_1$-compact sheaves. In particular, for every affine
		$X = \Spec(A)$, every affine object $\Spec(B)\to X$ is
		$\aleph_1$-compact in
		$\Shv_{\fpqc}(\Aff_{\S})_{/X}$, and every $\aleph_1$-compact object
		in this slice is a countable colimit of affine objects.
	\end{lemma}
	\begin{proof}
		By definition of the fpqc topology on affines, covering sieves are
		tested on finite covering families. Since finite sets are
		$\aleph_1$-small, \Cref{lem:Kappa_Compact_Covering_Sieves_Filtered_Colimits}
		shows that fpqc-sheaves are closed under $\aleph_1$-filtered colimits
		in presheaves. The same argument applies on the fpqc site of affines
		over a fixed affine base.

		Now apply \Cref{cor:Kappa_Compact_Sheaves_Finite_Limits} with
		$\kappa=\aleph_1$, both on $\Aff_{\S}$ and on the corresponding site
		over an affine base. This gives $\aleph_1$-compact generation, closure
		of $\aleph_1$-compact objects under finite limits, and preservation of
		$\aleph_1$-compact objects under sheafification. Since representable
		presheaves are compact, affine sheaves are $\aleph_1$-compact.

		The sheaf category over an affine $X$ is generated under colimits by
		affine objects over $X$. Hence its $\aleph_1$-compact objects are
		retracts of countable colimits of affine objects; absorbing the retract
		by the usual countable telescope gives the stated countable affine
		presentation.
	\end{proof}

	\begin{lemma}\label[lemma]{lem:Aleph1_Compact_Pro_Bicountably_Presented}
		Let $X = \Spec(A)$ be affine. Then every $\aleph_1$-compact object
		$Y \in \Shv_{\fpqc}(\Aff_{\S})_{/X}$ admits a compatible cone
		$Y \to Y_i$ over an $\aleph_1$-filtered category $I$, where each
		$Y_i \to X$ is bi-countably presented, such that, for every
		$n \geq 1$, the induced functor
		\[
			\colim_{i \in I} n\Pr_{Y_i} \to n\Pr_Y
		\]
		along pullback is an equivalence in $1\Pr = \PrL_{\aleph_1}$.
	\end{lemma}
	\begin{proof}
		By \Cref{lem:Affines_Aleph1_Compact_fpqc}, the affine objects over
		$X$ are $\aleph_1$-compact. Since affine objects
		generate $\Shv_{\fpqc}(\Aff_{\S})_{/X}$ under colimits, this sheaf
		category is $\aleph_1$-compactly generated by affines. Thus an
		$\aleph_1$-compact object over $X$ admits a presentation by countably
		many affines: we have
		\[
			Y \simeq \colim_{m \in M} \Spec(B_m)
		\]
		for some countable diagram $B \colon M\catop \to \CAlg_A(\Sp)$.
		Since $\CAlg_A(\Sp)$ is $\aleph_1$-compactly generated,
		\Cref{lem:Functor_Category_Presentable} shows that also
		$\Fun(M\catop,\CAlg_A(\Sp))$ is $\aleph_1$-compactly generated, with
		$\aleph_1$-compact objects given by the functors
		$M\catop \to \CAlg_A(\Sp)^{\aleph_1}$, i.e.\ the diagrams of
		countably presented $A$-algebras. Applying this to the diagram $B$, we
		obtain an $\aleph_1$-filtered category $I$ and diagrams
		$B_i\colon M\catop \to \CAlg_A(\Sp)$ such that each $B_i(m)$ is
		countably presented over $A$, and we have
		\[
			B \simeq \colim_{i \in I} B_i
		\]
		in $\Fun(M\catop,\CAlg_A(\Sp))$. Define
		\[
			Y_i := \colim_{m \in M} \Spec(B_i(m)).
		\]
		Since $M$ is countable and each $A \to B_i(m)$ is countably presented,
		each $Y_i \to X$ is bi-countably presented. The maps
		$B_i \to B$ induce a compatible cone $Y \to Y_i$ over $X$.

		It remains to check the asserted formula after applying $n\Pr_{(-)}$.
		First consider the affine case: if $B \simeq \colim_i B_i$ in
		$\CAlg_A(\Sp)$, then
		\[
			n\Pr_B
			\simeq
			\colim_i n\Pr_{B_i}.
		\]
		The base case $n = 0$ follows from
		\Cref{thm:Embeddings_CAlg(C)_Into_CAlg(Mod_C)} applied to
		$C = \Sp$: the functor
		$\CAlg_A(\Sp) \to \CAlg(1\Pr)_{0\Pr_A/}$, $B \mapsto 0\Pr_B$, is a
		left adjoint and hence preserves colimits. For higher $n$ the same
		argument applies inductively, applying the theorem with $C = n\Pr$
		and base algebra $(n-1)\Pr_A$.

		Now apply descent to the countable affine presentations of $Y$ and
		the $Y_i$. Since $B \simeq \colim_i B_i$ in the functor category, we
		have $B_m \simeq \colim_i B_i(m)$ for every $m \in M$. Hence
		\[
		\begin{aligned}
			n\Pr_Y
			&\simeq \lim_{m \in M\catop} n\Pr_{B_m} \\
			&\simeq \lim_{m \in M\catop}\colim_i n\Pr_{B_i(m)}.
		\end{aligned}
		\]
		Since $M$ is countable and $1\Pr$ is $\aleph_1$-compactly generated,
		$\aleph_1$-filtered colimits commute with this countable limit in
		$1\Pr$. Thus
		\[
		\begin{aligned}
			n\Pr_Y
			&\simeq \colim_i \lim_{m \in M\catop} n\Pr_{B_i(m)} \\
			&\simeq \colim_i n\Pr_{Y_i}.
		\end{aligned} \qedhere
		\]
	\end{proof}

	\begin{lemma}\label[lemma]{lem:Right_Adjoint_From_Filtered_Pro_Presentation}
		Let $X = \Spec(A)$, and let $(Y_i)_{i \in I}$ be a diagram over $X$,
		indexed by an $\aleph_1$-filtered category $I$, equipped with a
		compatible cone $Y \to Y_i$ over $X$. Assume that
		\[
			n\Pr_Y \simeq \colim_{i \in I} n\Pr_{Y_i}
		\]
		in $1\Pr$ along pullback. In particular, the structure functors
		$\alpha_i\colon n\Pr_{Y_i} \to n\Pr_Y$ preserve
		$\aleph_1$-compact objects, and compact generators of $n\Pr_Y$ are
		obtained from compact generators of the stages under countable
		colimits. Write $f_i\colon Y_i \to X$ and $f\colon Y \to X$ for the
		structure maps. For a morphism $i \to j$ in $I$, write
		$p_{ij}\colon Y_j \to Y_i$ for the induced transition map. Assume
		that $f_n^*$ preserves $\aleph_1$-compact objects and that each right
		adjoint
		\[
			f_{i,n,*}\colon n\Pr_{Y_i} \to n\Pr_X
		\]
		preserves small colimits. Then $f_{n,*}$ preserves small colimits,
		and, if $\alpha_i\colon n\Pr_{Y_i} \to n\Pr_Y$ denotes the colimit
		structure functor, then
		\[
			f_{n,*}(\alpha_i E)
			\simeq
			\colim_{j \in I_{i/}} f_{j,n,*}(p_{ij,n}^*E)
		\]
		for every $E \in n\Pr_{Y_i}$.
	\end{lemma}
	\begin{proof}
		For fixed $i$, define
		\[
			G_i(E) :=
			\colim_{j \in I_{i/}} f_{j,n,*}(p_{ij,n}^*E),
			\qquad E \in n\Pr_{Y_i}.
		\]
		Since each $f_{j,n,*}$ preserves colimits, so does $G_i$. If
		$i \to k$ is a morphism in $I$, then the inclusion
		$I_{k/} \to I_{i/}$ is cofinal, and hence
		\[
			G_k(p_{ik,n}^*E) \simeq G_i(E).
		\]
		Thus the restrictions of the $G_i$ to compact objects are compatible
		with the $1\Pr$-colimit presentation of $n\Pr_Y$. Since compact
		generators of $n\Pr_Y$ are generated under countable colimits by the
		objects $\alpha_iE$ with $E \in (n\Pr_{Y_i})^{\aleph_1}$, these
		compatible restrictions determine a functor on
		$(n\Pr_Y)^{\aleph_1}$ preserving countable colimits. Extending it
		cocontinuously gives a colimit-preserving functor
		\[
			G\colon n\Pr_Y \to n\Pr_X
		\]
		with $G(\alpha_iE) \simeq G_i(E)$ for compact $E$. Since both sides
		of this formula preserve colimits in $E$, the same formula holds for
		all $E \in n\Pr_{Y_i}$.

		We claim that $G$ is right adjoint to $f_n^*$. It suffices to test this
		on $\aleph_1$-compact generators $K \in (n\Pr_X)^{\aleph_1}$ and on
		$\aleph_1$-compact generators of $n\Pr_Y$ represented by some stage
		$n\Pr_{Y_i}$. Here we use the $1\Pr$-colimit presentation and the
		assumption that $f_n^*$ preserves compact objects. Using compactness
		of $K$ and the standard mapping-space formula for filtered colimits of
		categories, we compute
		\[
		\begin{aligned}
			\Hom_{n\Pr_X}(K,G(\alpha_iE))
			&\simeq
			\colim_{j \in I_{i/}}
			\Hom_{n\Pr_X}\bigl(K,f_{j,n,*}(p_{ij,n}^*E)\bigr) \\
			&\simeq
			\colim_{j \in I_{i/}}
			\Hom_{n\Pr_{Y_j}}\bigl(f_{j,n}^*K,p_{ij,n}^*E\bigr) \\
			&\simeq
			\colim_{j \in I_{i/}}
			\Hom_{n\Pr_{Y_j}}\bigl(p_{ij,n}^*f_{i,n}^*K,p_{ij,n}^*E\bigr) \\
			&\simeq
			\Hom_{n\Pr_Y}\bigl(f_n^*K,\alpha_iE\bigr).
		\end{aligned}
		\]
		Since $n\Pr_X$ is $\aleph_1$-compactly generated, this identifies
		$G$ with the right adjoint $f_{n,*}$.
	\end{proof}

	\section{The structure sheaf functor}

	We can summarize what we have shown in the language of Stefanich rings:

	\begin{corollary}\label[corollary]{cor:StRing_Functor_Properties}
		The functor
		\[
			\Shv_{\fpqc}(\Aff_{\S})^{\aleph_1} \to (\StRing_{\S/})\catop, \qquad X \mapsto \Oo(X)
		\]
		commutes with countable colimits and finite limits.
	\end{corollary} 

	If the target were itself an $\infty$-topos, then this functor would have all the properties of a morphism of $\infty$-topoi (except that we restrict to countable colimits rather than to all small colimits). In \Cref{chap:Gestalten}, we will make this idea precise using the concept of a \emph{countably toposic category}.

	\begin{proof}[Sketch of proof]
		Implicit in the statement is the claim that $\aleph_1$-compact objects are closed under finite limits in $\Shv_{\fpqc}(\Aff_{\S})$; see \Cref{cor:Kappa_Compact_Sheaves_Finite_Limits} for a proof of this fact.

		Commutation with countable colimits holds essentially by definition. To prove commutation with finite limits, it suffices to handle the terminal object (which works by definition of $\StRing_{\S/}$) and pullbacks. So consider a pullback diagram
		\[
			\begin{tikzcd}
				Y \times_X Y' \rar \dar & Y' \dar["{f'}"] \\
				Y \rar["f"'] & X
			\end{tikzcd}
		\]
		of $\aleph_1$-compact fpqc-stacks. A priori, to compute the pushout of Stefanich rings, we would start by forming the infinite tuple
		\[
			\big(\,\Gamma(Y) \otimes_{\Gamma(X)} \Gamma(Y'),\quad \D(Y) \otimes_{\D(X)} \D(Y'),\quad 1\Pr_Y \otimes_{1\Pr_X} 1\Pr_{Y'},\quad \ldots\,\big),
		\]
		consisting of all the levelwise pushouts. These come with natural maps from each entry to endomorphisms of the unit of the next, but these are typically \emph{not} isomorphisms; one needs to ``spectrify'', i.e.\ apply the left adjoint to the inclusion of Stefanich rings into the larger category that allows for non-invertible comparison maps. 
		
		In the next talk, we will explain a convenient way of computing such tensor products. What we will see there is that we can make a `first approximation' to the spectrification, given by:
		\[
			\big(\,\Oo(Y/X)_0 \otimes_{\D(X)} \Oo(Y'/X)_0,\quad \D(Y) \otimes_{1\Pr_X} \D(Y'),\quad 1\Pr_Y \otimes_{2\Pr_X} 1\Pr_{Y'},\quad \ldots\,\big).
		\]
		Here on each level the tensor product is taken at the next-higher categorical level. (For example, in the zeroth slot we form the tensor product not over $\Gamma(X)$, i.e.\ in $\Gamma(X)$-modules, but rather inside $\D(X)$; this is finer, since there is a canonical functor $\Mod_{\Gamma(X)}(\D(\S)) \to \D(X)$.) By \Cref{prop:Kunneth_Aleph1_Compact}, starting from the slot $1\Pr_Y \otimes_{2\Pr_X} 1\Pr_{Y'}$, this expression agrees with
		\[
			(\Omega^\infty\Gamma(Y \times_X Y'),\; \D(Y \times_X Y'),\; 1\Pr_{Y \times_X Y'},\; \ldots),
		\]
		which is already a Stefanich ring, namely the one corresponding to $Y \times_X Y'$. As the first few terms do not affect the Stefanich ring structure, this gives the desired result.
	\end{proof}

	\begin{remark}
		A general slogan in this area is that ``everything becomes affine under sufficient categorification''. The Künneth formulas above, which hold from level $1\Pr$ onwards, justify this slogan in some form. The slogan is a bit imprecise, however: stacks like $B^n \bbG_m$ are $n$-affine (in a sense to be made precise in the next talk) but not $k$-affine for $k < n$, so in particular $\bigsqcup_k B^k \bbG_m$ is not $n$-affine for any $n$. We will later see that even the infinite-dimensional affine space $\bbA^{\infty} = \colim_n \bbA^n$ is not $n$-affine for any $n$, although all its terms are $0$-affine and the transition maps are cohomologically smooth (since complete intersections are cohomologically smooth in the quasicoherent six-functor formalism).
	\end{remark}

	\chapter{Stefanich rings}
	\label[chapter]{chap:Stefanich_Rings}

	\emph{Talk given by Yiming.}

	Recall that we defined the $\aleph_1$-presentable category $\StRing$ of \emph{Stefanich rings} as
	\begin{align*}
		\StRing &:= \colim_{1\Pr}(\CAlg(\An) \xrightarrow{\Mod^{\bbP}_1} \CAlg(1\Pr) \xrightarrow{\Mod^{\bbP}_2} \CAlg(2\Pr) \xrightarrow{\Mod^{\bbP}_3} \dots) \\
		& \simeq \lim_{\CAT}(\CAlg(\An) \xleftarrow{\End^{\bbP}_1} \CAlg(1\Pr) \xleftarrow{\End^{\bbP}_2} \CAlg(2\Pr) \xleftarrow{\End^{\bbP}_3} \dots),
	\end{align*}
	where we write $\Mod^{\bbP}_n(R) := \Mod_R((n-1)\Pr)$ and $\End^{\bbP}_n(D) := \End_D(\unit_D)$.

	We denote the objects of this category as $A = (A_0,A_1,A_2, \dots)$, where $A_n \in \CAlg(n\Pr)$ and the structure isomorphisms $A_n \simeq \End^{\bbP}_{n+1}(A_{n+1})$ are left implicit in the notation. Similarly, morphisms from $A$ to $B$ are collections of maps $f_{n-1}^*\colon A_n \to B_n$ compatible with the structural isomorphisms.

	\begin{remark}
		Given an fpqc-stack $X$, the associated Stefanich ring has $A_n = (n-1)\Pr_X$. The indexing convention for $f_{n-1}^*\colon A_n \to B_n$ reflects the convention from the previous talk.
	\end{remark}

	We may equip $\StRing$ with the cocartesian monoidal structure; we will have a look at the resulting ``tensor product'' of Stefanich rings in more detail below. Since the functors $\End^{\bbP}_n$ preserve the monoidal unit, we see that the monoidal unit in $\StRing$ is the initial Stefanich ring $\bbP := (*, \An, 1\Pr, 2\Pr, \dots)$, with $n$-th component $\bbP_n := (n-1)\Pr$.
	
	Since $\bbP$ is initial, every Stefanich ring $A$ receives a unique map $a\colon\bbP \to A$. The individual maps $a_{n-1}^*\colon (n-1)\Pr \to A_n$ in $\CAlg(\Pr)$ are obtained using the identification $\CAlg(n\Pr) \simeq \CAlg(\Pr)_{(n-1)\Pr/}$. By passing to right adjoints, we obtain `forgetful functors'
	\[
		a_{n-1,*}\colon A_n \to (n-1)\Pr, \qquad B_{n-1} \mapsto \iHom_{A_n}(\unit_{A_n}, B_{n-1}).
	\]
	Given $B_{n-1} \in A_n$, we refer to $a_{n-1,*}(B_{n-1})$ as its \emph{underlying $(n-1)$-category.}

	\begin{notation}
		\label[notation]{nota:(A/A)_n}
		We will denote the monoidal unit of $A_n$ by
		\[
			(A/A)_{n-1} \; := \; \unit_{A_n} \; \in \; \CAlg(A_n).
		\]
		Observe that we have $a_{n-1,*}((A/A)_{n-1}) = \End_{A_n}(\unit_{A_n}) = A_{n-1}$.
	\end{notation}

	\section{Internal Stefanich rings}

	We will now provide a limit description of the slice category $\StRing_{A/}$ under any Stefanich ring $A$, exhibiting $A$-algebras in Stefanich rings as some form of internal Stefanich rings in $A$.

	\begin{construction}\label[construction]{con:Lift_To_Algebra}
		Let $f\colon A \to B$ be a morphism of Stefanich rings. For any $n \geq 1$, the symmetric monoidal functor $f_{n-1}^*\colon A_n \to B_n$ admits a lax symmetric monoidal right adjoint $f_{n-1,*}\colon B_n \to A_n$, which therefore induces a functor of the form
		\[
			f_{n-1,*}\colon \CAlg(B_n) \to \CAlg(A_n).
		\]
		Observe that the composite $\CAlg(B_n) \xrightarrow{f_{n-1,*}} \CAlg(A_n) \xrightarrow{\fgt} \CAlg((n-1)\Pr)$ is the forgetful functor $\CAlg(B_n) \to \CAlg((n-1)\Pr)$.

		We then define
		\[
			(B/A)_{n-1} \; := \; f_{n-1,*}(\unit_{B_n}) \qin \CAlg(A_n).
		\]
		Note that the underlying $(n-1)$-presentable category of $(B/A)_{n-1}$ is $B_{n-1}$. So we may think of $(B/A)_{n-1}$ as a lift of $B_{n-1}$ to a commutative algebra in $A_n$. For $f = \id_A$ this recovers \Cref{nota:(A/A)_n}.
	\end{construction}

	\begin{construction}
		\label[construction]{con:Internal_Module_Endomorphism_Adjunction}
		We now construct an internal form of the module-endomorphism adjunction: For a Stefanich ring $A$ and $m \geq 1$, we will produce an adjunction
		\[
			\begin{tikzcd}
				\CAlg(A_m) \rar[shift left]{\Mod^A_m} & \CAlg(A_{m+1}). \lar[shift left]{\End^A_m}
			\end{tikzcd}
		\]
		Recall from \Cref{nota:Module_Endomorphism_Adjunction} that, more generally, for $n \geq 1$ and $C \in \CAlg(n\Pr)$ we have an adjunction
		\[
			\begin{tikzcd}
				\CAlg(C) \rar[hook, shift left]{\Mod_{(-)}(C)} & \CAlg(\Mod_C(n\Pr)), \lar[shift left]{\End_{(-)}(\unit)}
			\end{tikzcd}
		\]
		where the left adjoint sends $R$ to $\Mod_R(C)$ and the right adjoint sends $D$ to $\End_D(\unit_D)$. Here we are using that $C$ is an algebra over $(n-1)\Pr$ in $1\Pr$, so that
		\[
			\Mod_C(1\Pr) \simeq \Mod_C(\Mod_{(n-1)\Pr}(1\Pr)) \simeq \Mod_C(n\Pr),
		\]
		giving the $(n-1)\Pr$-linear structure for free.

		Now let $A$ be a Stefanich ring. Applying this adjunction with $C = m\Pr$, we see that the structural equivalence
		\[
			A_m \simeq \End_{A_{m+1}}(\unit_{A_{m+1}})
		\]
		in $\CAlg(m\Pr)$ corresponds to a symmetric monoidal functor
		\[
			L^A_m\colon \Mod_{A_m}(m\Pr) \to A_{m+1}
		\]
		in $\CAlg((m+1)\Pr)$. It admits a right adjoint, which we denote by
		\[
			U^A_m\colon A_{m+1} \to \Mod_{A_m}(m\Pr), \qquad X \mapsto \iHom_{A_{m+1}}(\unit_{A_{m+1}},X).
		\]
		We use the same notation for the induced functors on commutative algebra objects. Finally, we write
		\[
			\Mod^A_m := L^A_m \circ \Mod_{(-)}(A_m)\colon \CAlg(A_m) \to \CAlg(A_{m+1})
		\]
		for the internal module functor, and
		\[
			\End^A_m := \End_{(-)}(\unit) \circ U^A_m\colon \CAlg(A_{m+1}) \to \CAlg(A_m)
		\]
		for its right adjoint. Explicitly, we have
		\[
			\Mod^A_m(R) = L^A_m(\Mod_R(A_m)) \qquad \text{and} \qquad \End^A_m(D) = \End_{U^A_m(D)}(\unit).
		\]
	\end{construction}

	\begin{observation}
		\label[observation]{obs:U^A_m_Preserves_Unit}
		The functor $U^A_m\colon \CAlg(A_{m+1}) \to \Mod_{A_m}(m\Pr)$ preserves the monoidal unit: the unit of the adjunction $L^A_m \dashv U^A_m$ takes the form
		\[
			A_m \to U^A_m(\unit_{A_{m+1}}) = \End_{A_{m+1}}(\unit_{A_{m+1}}),
		\]
		which is the structure isomorphism of the Stefanich ring $A$.
	\end{observation}

	\begin{proposition}\label[proposition]{prop:Mod_Transition_Fully_Faithful}
		Let $A \in \StRing$ and $m \geq 1$. The internal module functor
		\[
			\Mod^A_m\colon \CAlg(A_m) \to \CAlg(A_{m+1})
		\]
		is fully faithful.
	\end{proposition}
	\begin{proof}
		Since $\Mod^A_m$ is a morphism in $1\Pr$, by \Cref{lem:Fully_Faithful_On_Aleph1_Compacts} it suffices to check that the unit of the adjunction $\Mod^A_m \dashv \End^A_m$ is an equivalence on $\aleph_1$-compact objects, i.e.\ on countably presented objects of $\CAlg(A_m)$. Let $R \in \CAlg(A_m)$ be countably presented. Unwinding definitions, we get
		\[
			\End^A_m(\Mod^A_m(R)) \simeq \End_{U^A_mL^A_m(\Mod_R(A_m))}(\unit) = \End_{T(\Mod_R(A_m))}(\unit),
		\]
		where we use the abbreviated notation
		\[
			T := U^A_mL^A_m\colon \Mod_{A_m}(m\Pr) \to \Mod_{A_m}(m\Pr).
		\]
		To analyze this endomorphism category, note:
		\begin{itemize}
			\item The functors $U^A_m$ and $L^A_m$ are both maps in $(m+1)\Pr$, hence so is $T$. Since $m \geq 1$, this means that $T$ is (at least) a 2-functor, so it preserves adjunctions.
			\item By \Cref{obs:U^A_m_Preserves_Unit}, the unit map $A_m \to T(A_m)$ is an isomorphism.
		\end{itemize}
		We now consider the free $R$-module functor
		\[
			F_R\colon A_m \to \Mod_R(A_m), \qquad X \mapsto R \otimes X.
		\]
		It admits an $A_m$-linear right adjoint $G_R\colon \Mod_R(A_m) \to A_m$, given by forgetting the $R$-module structure. Since $R$ is countably presented, the functor $G_R$ preserves $\aleph_1$-compact objects, and so this adjunction lives inside $\Mod_{A_m}(m\Pr)$. So applying the $2$-functor $T$ to the adjunction $F_R \dashv G_R$, we obtain an adjunction
		\[
			T(F_R) \colon T(A_m) \rightleftarrows T(\Mod_R(A_m)) \noloc T(G_R)
		\]
		in $\Mod_{A_m}(m\Pr)$. Since $T$ is lax monoidal, it induces a functor on commutative algebras, and hence sends the symmetric monoidal functor $F_R$ to a symmetric monoidal functor $T(F_R)$; consequently the unit object of $T(\Mod_R(A_m))$ is $T(F_R)(\unit_{T(A_m)})$. Combining all of this, we may now compute:
		\begin{align*}
			\End^A_m(\Mod^A_m(R)) &= \End_{T(\Mod_R(A_m))}(\unit) \\
			& \simeq \End_{T(\Mod_R(A_m))}(T(F_R)(\unit_{T(A_m)})) \\
			& \simeq \iHom_{T(A_m)}(\unit_{T(A_m)}, T(G_R)T(F_R)(\unit_{T(A_m)})) \\
			& \simeq \iHom_{T(A_m)}(\unit_{T(A_m)}, T(G_R \circ F_R)(\unit_{T(A_m)})) \\
			& \simeq \iHom_{T(A_m)}(\unit_{T(A_m)}, T(R \otimes -)(\unit_{T(A_m)})) \\
			& \simeq \iHom_{A_m}(\unit_{A_m}, (R \otimes -)(\unit_{A_m})) \\
			& \simeq R,
		\end{align*}
		where in the second-to-last equivalence we use that under the equivalence $A_m \iso T(A_m)$, the endofunctor $T(R \otimes -)$ corresponds to $R \otimes -\colon A_m \to A_m$. Under this identification, the unit map $R \to \End^A_m(\Mod^A_m(R))$ is the map induced by the action of $R$ on the free rank-one $R$-module $F_R(\unit_{A_m})$. Under the adjunction $T(F_R) \dashv T(G_R)$, this action corresponds to the identity of $G_RF_R(\unit_{A_m}) \simeq R$. Thus the unit is an equivalence on countably presented $R$, and $\Mod^A_m$ is fully faithful.
	\end{proof}

	\begin{definition}\label[definition]{def:Internal_Stefanich_Rings}
		Let $A \in \StRing$. The category of \emph{internal Stefanich rings in $A$} is the colimit
		\[
			\StRing(A) := \colim_{1\Pr}(\CAlg(A_1) \xhookrightarrow{\Mod^A_1} \CAlg(A_2) \xhookrightarrow{\Mod^A_2} \dots)
		\]
		formed along the internal module functors. Equivalently, since each $\Mod^A_m$ is a left adjoint with right adjoint $\End^A_m$, it may be computed as the limit
		\[
			\StRing(A) \simeq \lim_{\CAT}(\CAlg(A_1) \xleftarrow{\End^A_1} \CAlg(A_2) \xleftarrow{\End^A_2} \dots).
		\]
	\end{definition}

	We now show that $A$-algebras in Stefanich rings are the same as Stefanich rings internal to $A$.

	\begin{proposition}\label[proposition]{prop:Slice_Description}
		Let $A \in \StRing$. The assignments $B \mapsto (B/A)_n$, for $n \geq 0$, determine an equivalence
		\[
			\StRing_{A/} \iso \StRing(A), \qquad B \mapsto ((B/A)_0, (B/A)_1, (B/A)_2, \dots).
		\]
	\end{proposition}
	\begin{proof}
		By passing to slices, we obtain
		\[
			\StRing_{A/} = \lim_{\CAT}(\dots \xrightarrow{\End^{\bbP}_3} \CAlg(2\Pr)_{A_2/} \xrightarrow{\End^{\bbP}_2} \CAlg(1\Pr)_{A_1/} \xrightarrow{\End^{\bbP}_1} \CAlg(\An)_{A_0/}).
		\]
		The entry $\CAlg(\An)_{A_0/}$ is redundant for the limit, and we will ignore it. Under the equivalences $\CAlg(m\Pr)_{A_m/} \simeq \CAlg(\Mod_{A_m}(m\Pr))$, the transition maps $\CAlg((m+1)\Pr)_{A_{m+1}/} \to \CAlg(m\Pr)_{A_m/}$ admit the following simple description. An object $D$ of the source is an $A_{m+1}$-linear monoidal category; the transition sends it to $\End_D(\unit_D)$, regarded over $\End_{A_{m+1}}(\unit_{A_{m+1}}) \simeq A_m$. Thus the transition factors as
		\[
			\CAlg(\Mod_{A_{m+1}}((m+1)\Pr)) \xrightarrow{\End_{(-)}(\unit)} \CAlg(A_{m+1}) \xrightarrow{U^A_m} \CAlg(\Mod_{A_m}(m\Pr)),
		\]
		where we use the notation from \Cref{con:Internal_Module_Endomorphism_Adjunction}. Hence we may alternatively compute $\StRing_{A/}$ as the limit of the chain
		\[
			\dots \to \CAlg(\Mod_{A_{m+1}}((m+1)\Pr)) \xrightarrow{\End_{(-)}(\unit)} \CAlg(A_{m+1}) \xrightarrow{U^A_m} \CAlg(\Mod_{A_m}(m\Pr)) \xrightarrow{\End_{(-)}(\unit)} \CAlg(A_m) \to \dots.
		\]
		Removing the intermediate $\CAlg(\Mod_{A_m}(m\Pr))$-vertices gives the limit over the composites
		\[
			\End^A_m = \End_{(-)}(\unit) \circ U^A_m\colon \CAlg(A_{m+1}) \to \CAlg(A_m),
		\]
		hence the identification
		\[
			\StRing_{A/} \simeq \lim_{\CAT}(\CAlg(A_1) \xleftarrow{\End^A_1} \CAlg(A_2) \xleftarrow{\End^A_2} \dots) = \StRing(A),
		\]
		where the last equality is the limit description of \Cref{def:Internal_Stefanich_Rings}. Under this equivalence, an object $f\colon A \to B$ is sent in degree $n \geq 0$ to
		\[
			f_{n,*}(\unit_{B_{n+1}}) = (B/A)_n \qin \CAlg(A_{n+1}),
		\]
		that is, to the tuple $((B/A)_0, (B/A)_1, (B/A)_2, \dots)$, as claimed.
	\end{proof}

	\begin{corollary}
		\label[corollary]{cor:Inclusion_n_Affine_A_Algebras}
		The functor $\StRing_{A/} \to \CAlg(A_{n+1}), B \mapsto (B/A)_n$ admits a fully faithful left adjoint
		\[
			\CAlg(A_{n+1}) \hookrightarrow \StRing_{A/}.
		\]
	\end{corollary}
	\begin{proof}
		By \Cref{prop:Mod_Transition_Fully_Faithful}, the internal module functors $\Mod^A_m$ are fully faithful left adjoints. Hence \Cref{lem:Colimit_Of_Fully_Faithful_Left_Adjoints} shows that, for every $n \geq 0$, the canonical functor
		\[
			\CAlg(A_{n+1}) \hookrightarrow \StRing(A)
		\]
		is fully faithful and left adjoint to the projection $\StRing(A) \to \CAlg(A_{n+1})$. The claim now follows from the equivalence $\StRing_{A/} \simeq \StRing(A)$ of \Cref{prop:Slice_Description}, under which this projection corresponds to the functor $B \mapsto (B/A)_n$.
	\end{proof}

	In \Cref{prop:Mod_Transition_Fully_Faithful} we showed that the composite
	\[
		\Mod^A_m\colon \CAlg(A_m) \xhookrightarrow{\Mod_{(-)}(A_m)} \CAlg(\Mod_{A_m}(m\Pr)) \xrightarrow{L^A_m} \CAlg(A_{m+1})
	\]
	is fully faithful. While the functor $\Mod_{(-)}(A_m)$ is also always fully faithful, the functor $L^A_m$ is usually not. Nevertheless, we still have the following statement:

	\begin{lemma}\label[lemma]{lem:Lm_Fully_Faithful_On_Dualizables}
		Let $A \in \StRing$ and $m \geq 1$. The functor
		\[
			L^A_m\colon \Mod_{A_m}(m\Pr) \to A_{m+1}
		\]
		is fully faithful on dualizable objects. (In particular, it is fully faithful on $\aleph_1$-filtered colimits of dualizable objects, cf.\ \Cref{lem:Fully_Faithful_On_Aleph1_Compacts}.)
	\end{lemma}
	\begin{proof}
		Set $C := \Mod_{A_m}(m\Pr)$ and $D := A_{m+1}$, and let $U^A_m$ denote the right adjoint of $L^A_m$. Let $X,Y \in C$ be dualizable. Since $L^A_m$ is symmetric monoidal, it preserves duals, and we have a commutative diagram
		\[
			\begin{tikzcd}
				\Hom_C(X \otimes Y^{\vee}, \unit_C) \rar["\sim"] \dar
					& \Hom_C(X,Y) \dar \\
				\Hom_D(L^A_m(X \otimes Y^{\vee}), \unit_D) \rar["\sim"]
					& \Hom_D(L^A_m(X),L^A_m(Y)).
			\end{tikzcd}
		\]
		The horizontal maps are the usual identifications coming from duality, given by tensoring with $Y$ and precomposing with the coevaluation $\unit_C \to Y^{\vee} \otimes Y$. Thus it suffices to show that the left vertical map is an equivalence. Composing this map with the adjunction equivalence
		\[
			\Hom_D(L^A_m(X \otimes Y^{\vee}),\unit_D) \simeq \Hom_C(X \otimes Y^{\vee},U^A_m(\unit_D))
		\]
		we see that it suffices to show that the unit map $\unit_C \to U^A_mL^A_m\unit_C = U^A_m\unit_D$ is an isomorphism. But this map is the canonical comparison map $A_m \iso \End_{A_{m+1}}(\unit_{A_{m+1}})$, which is an isomorphism as $A$ is a Stefanich ring.
	\end{proof}

	\section{Shifting}

	We will now define the \emph{shift functor} on $\StRing$.

	\begin{construction}
		The forgetful functors $\CAlg((n+1)\Pr) \to \CAlg(n\Pr)$ induce a functor
		\[
			\begin{tikzcd}
				\StRing \dar["\shift", swap] \rar{\simeq} & \lim( \dots \to \CAlg(3\Pr) \to \CAlg(2\Pr) \to \CAlg(1\Pr)) \dar \\
				\StRing \rar{=} & \lim( \dots \to \CAlg(2\Pr) \to \CAlg(1\Pr) \to \CAlg(\An))
			\end{tikzcd}
		\]
		given on objects by $A = (A_1, A_2, \dots) \mapsto (A_2, A_3, \dots)$.
	\end{construction}

	\begin{lemma}
		\label[lemma]{lem:Shift_Functor_Equivalence}
		The shift functor induces an equivalence
		\[
			\shift\colon \StRing \iso \StRing_{(1\Pr, 2\Pr,\dots)/}.
		\]
		In particular, for any $n \geq 0$, iterated shifting induces an equivalence
		\[
			\shift^n\colon \StRing \iso \StRing_{(n\Pr,(n+1)\Pr, \dots)/}.
		\]
	\end{lemma}
	\begin{proof}
		For each $n \geq 1$, we have the identification
		\[
			\CAlg((n+1)\Pr)
			= \CAlg(\Mod_{n\Pr}(1\Pr))
			= \CAlg(1\Pr)_{n\Pr/}
			= (\CAlg(1\Pr)_{(n-1)\Pr/})_{n\Pr/}
			= \CAlg(n\Pr)_{n\Pr/}.
		\]
		These identifications are compatible with the transition functors. Under them, the forgetful functor
		\[
			\CAlg((n+1)\Pr) \to \CAlg(n\Pr)
		\]
		is the projection $\CAlg(n\Pr)_{n\Pr/} \to \CAlg(n\Pr)$, while the same identification remembers the map from $n\Pr$.
		Since each omitted lower-degree term in the limit presentation of $\StRing$ is forced by the next one, we may start the limit at $\CAlg(2\Pr)$ and obtain
		\[
			\begin{aligned}
				\StRing
				&\simeq \lim_{\CAT}\bigl(\dots \to \CAlg(4\Pr) \to \CAlg(3\Pr) \to \CAlg(2\Pr)\bigr) \\
				&\simeq \lim_{\CAT}\bigl(\dots \to \CAlg(3\Pr)_{3\Pr/} \to \CAlg(2\Pr)_{2\Pr/} \to \CAlg(1\Pr)_{1\Pr/}\bigr) \\
				&\simeq \left(\lim_{\CAT}\bigl(\dots \to \CAlg(3\Pr) \to \CAlg(2\Pr) \to \CAlg(1\Pr)\bigr)\right)_{(1\Pr,2\Pr,\dots)/} \\
				&\simeq \StRing_{(1\Pr,2\Pr,\dots)/}.
			\end{aligned}
		\]
		Here we use that slices commute with limits in $\CAT$. The resulting functor is precisely $\shift$ by construction. The iterated statement follows by applying the same argument repeatedly.
	\end{proof}

	\section{Spectrification and colimits of Stefanich rings}
	\label[section]{sec:Spectrification}

	Throughout this section we fix a Stefanich ring $A$ and study colimits in the category $\StRing(A)$ of internal Stefanich rings in $A$. Limits and $\aleph_1$-filtered colimits in $\StRing(A)$ are computed pointwise, since they are preserved by the transition functors $\End^A_n$ in the limit description of \Cref{def:Internal_Stefanich_Rings}. Arbitrary colimits are harder to describe directly, since they are not (necessarily) preserved by the functors $\End^A_n$.

	\begin{definition}\label[definition]{def:PreStRing}
		Let
		\[
			\mathcal{P}_A := \prod_{n \geq 0} \CAlg(A_{n+1}),
		\]
		and define an endofunctor
		\[
			\End^A_\bullet\colon \mathcal{P}_A \to \mathcal{P}_A, \qquad (B_n)_{n \geq 0} \mapsto (\End^A_{n+1}(B_{n+1}))_{n \geq 0}.
		\]
		The category of \emph{internal preStefanich rings in $A$} is the lax equalizer
		\[
			\PreStRing(A)
			:=
			\mathcal{P}_A \times_{\mathcal{P}_A \times \mathcal{P}_A} \Ar(\mathcal{P}_A),
		\]
		where $\mathcal{P}_A \to \mathcal{P}_A \times \mathcal{P}_A$ is $(\id,\End^A_\bullet)$ and $\Ar(\mathcal{P}_A) \to \mathcal{P}_A \times \mathcal{P}_A$ is the source-target functor.
		Thus an internal preStefanich ring is a tuple $B = (B_0,B_1,B_2,\dots)$ with $B_n \in \CAlg(A_{n+1})$, together with structure maps
		\[
			\alpha_n\colon B_n \to \End^A_{n+1}(B_{n+1})
		\]
		that are not assumed to be equivalences. By \Cref{def:Internal_Stefanich_Rings}, the category $\StRing(A)$ identifies with the full subcategory of $\PreStRing(A)$ spanned by those objects for which all $\alpha_n$ are equivalences.
	\end{definition}

	\begin{lemma}\label[lemma]{lem:Colimits_In_PreStRing}
		The category $\PreStRing(A)$ admits small colimits, and the projection $\PreStRing(A) \to \mathcal{P}_A$ creates them. In particular, colimits in $\PreStRing(A)$ are computed pointwise.
	\end{lemma}
	\begin{proof}
		Let $I \to \PreStRing(A)$ be a diagram, written as $(B^i,\alpha^i)$ for $i \in I$. Let $B := \colim_i B^i$ in $\mathcal{P}_A$, so that $B_n \simeq \colim_i B^i_n$ in $\CAlg(A_{n+1})$. The maps $\alpha^i\colon B^i \to \End^A_\bullet(B^i)$ assemble to a morphism of diagrams $B^{\bullet} \to \End^A_\bullet(B^{\bullet})$. Hence we obtain a structure map
		\[
			B \simeq \colim_i B^i \to \colim_i \End^A_\bullet(B^i) \to \End^A_\bullet(\colim_i B^i) = \End^A_\bullet(B),
		\]
		where the second arrow is the canonical comparison map induced by the cocone $B^i \to B$. This defines an object $(B,\alpha)$ of $\PreStRing(A)$.

		For any $(Y,\beta) \in \PreStRing(A)$, the hom anima $\Hom_{\PreStRing(A)}((B,\alpha),(Y,\beta))$ is the homotopy equalizer of the two maps
		\[
			\Hom_{\mathcal{P}_A}(B,Y) \rightrightarrows \Hom_{\mathcal{P}_A}(B,\End^A_\bullet Y)
		\]
		defined by $\beta \circ (-)$ and $\End^A_\bullet(-)\circ \alpha$. Since $B$ is the colimit of the $B^i$ in $\mathcal{P}_A$, this homotopy equalizer identifies with
		\[
			\lim_i \Hom_{\PreStRing(A)}((B^i,\alpha^i),(Y,\beta)).
		\]
		Thus $(B,\alpha)$ has the universal property of the colimit.
	\end{proof}

	What we will do is to embed $\StRing(A)$ in $\PreStRing(A)$ and then reflect back.

	\begin{lemma}\label[lemma]{lem:Spectrification_Left_Adjoint}
		The inclusion $\StRing(A) \hookrightarrow \PreStRing(A)$ admits a left adjoint.
	\end{lemma}
	We refer to this left adjoint as \emph{spectrification}.
	\begin{proof}
		Define a functor
		\[
			T \colon \PreStRing(A) \to \PreStRing(A), \qquad (B_n)_{n \geq 0} \mapsto (B'_n)_{n \geq 0},
		\]
		by setting $B'_n = \End^A_{n+1}(B_{n+1})$, with structure maps obtained by applying the functors $\End^A_{n+1}$ to the structure maps of $B$. This admits a natural transformation $\eta\colon \id \to T$, given by the structure maps of an internal preStefanich ring. For $X \in \PreStRing(A)$, define $X^{(0)} := X$, $X^{(\alpha + 1)} := T(X^{(\alpha)})$, and
		\[
			X^{(\lambda)} := \colim_{\alpha < \lambda} X^{(\alpha)}
		\]
		for limit ordinals $\lambda$.

		We first show that $X^{(\omega_1)}$ lies in $\StRing(A)$. Each $\End^A_n$ is the right adjoint of the morphism $\Mod^A_n$ of $1\Pr$ from \Cref{con:Internal_Module_Endomorphism_Adjunction}, hence preserves $\aleph_1$-filtered colimits; therefore the functor $T$ preserves $\aleph_1$-filtered colimits. Hence
		\[
			T(X^{(\omega_1)})
			\simeq
			T(\colim_{\alpha < \omega_1} X^{(\alpha)})
			\simeq
			\colim_{\alpha < \omega_1} T(X^{(\alpha)})
			=
			\colim_{\alpha < \omega_1} X^{(\alpha+1)}
			\simeq
			X^{(\omega_1)}.
		\]
		The last isomorphism uses that the successor ordinals are cofinal in $\omega_1$. Under these identifications, the structure map $X^{(\omega_1)} \to T(X^{(\omega_1)})$ is the canonical map
		\[
			\colim_{\alpha < \omega_1} X^{(\alpha)} \to \colim_{\alpha < \omega_1} X^{(\alpha+1)},
		\]
		hence is an isomorphism. Thus $X^{(\omega_1)}$ is an internal Stefanich ring in $A$.

		It remains to prove the universal property. Let $Y \in \StRing(A)$. We claim by transfinite induction that the natural map
		\[
			\Hom_{\PreStRing(A)}(X^{(\alpha)},Y) \to \Hom_{\PreStRing(A)}(X,Y)
		\]
		is an equivalence for every ordinal $\alpha \leq \omega_1$. The claim is clear for $\alpha = 0$. For a successor ordinal, it suffices to show that
		\[
			\eta_{X^{(\alpha)}}^*\colon \Hom_{\PreStRing(A)}(T X^{(\alpha)},Y) \to \Hom_{\PreStRing(A)}(X^{(\alpha)},Y)
		\]
		is an equivalence. Since $Y$ lies in $\StRing(A)$, the map $\eta_Y\colon Y \to TY$ is an equivalence. An inverse to $\eta_{X^{(\alpha)}}^*$ is therefore given by applying $T$ to a map $X^{(\alpha)} \to Y$ and then using $\eta_Y^{-1}$ to identify $TY$ with $Y$.

		For a limit ordinal $\lambda$, the object $X^{(\lambda)}$ is the colimit of the $X^{(\alpha)}$ for $\alpha < \lambda$, so
		\[
			\Hom_{\PreStRing(A)}(X^{(\lambda)},Y)
			\simeq
			\lim_{\alpha < \lambda} \Hom_{\PreStRing(A)}(X^{(\alpha)},Y)
			\simeq
			\Hom_{\PreStRing(A)}(X,Y).
		\]
		Taking $\alpha = \omega_1$ gives
		\[
			\Hom_{\StRing(A)}(X^{(\omega_1)},Y)
			\simeq
			\Hom_{\PreStRing(A)}(X,Y),
		\]
		naturally in $X$ and $Y$. Thus $X \mapsto X^{(\omega_1)}$ is left adjoint to the inclusion.
	\end{proof}

	\begin{corollary}\label[corollary]{cor:Colimits_In_StRing}
		We may compute colimits in $\StRing(A)$ by computing them in $\PreStRing(A)$ (where they are computed pointwise) and then reflecting them back into $\StRing(A)$. \qed
	\end{corollary}

	\begin{remark}\label[remark]{rem:Absolute_Spectrification}
		The Stefanich ring $\bbP$ is initial in $\StRing$, so the slice $\StRing_{\bbP/}$ is all of $\StRing$. Hence \Cref{prop:Slice_Description} yields an equivalence $\StRing \simeq \StRing(\bbP)$, and we abbreviate $\PreStRing := \PreStRing(\bbP)$. Specializing the above to $A = \bbP$, colimits in $\StRing$ are computed pointwise in $\PreStRing$ and then spectrified.
	\end{remark}

	\begin{observation}
		\label[observation]{obs:Stefanich_Ring_In_High_Degrees}
		Let $B = (B_0, B_1, B_2, \dots)$ be an internal preStefanich ring in $A$. Assume that there is some $n \geq 0$ such that for each $m \geq n$ the canonical map $B_m \to \End_{m+1}^A(B_{m+1})$ is an isomorphism in $A_{m+1}$. Then the spectrification of $B$ agrees with $B$ in degrees $m \geq n$, and is in lower degrees given by repeatedly applying $\End^A_{\bullet}$.
	\end{observation}

	\section{Tensor products of Stefanich rings}
	\label[section]{sec:Tensor_Products_Stefanich_Rings}

	Consider a diagram of the form $B \leftarrow A \to C$. We wish to compute the pushout $B \otimes_A C$ in $\StRing$. By \Cref{rem:Absolute_Spectrification} (the case $A = \bbP$ of \Cref{cor:Colimits_In_StRing}), we can do this by computing the pushout in $\PreStRing$, where colimits are pointwise by \Cref{lem:Colimits_In_PreStRing}. This gives the preStefanich ring
	\[
		(B_0 \otimes_{A_0} C_0, B_1 \otimes_{A_1} C_1, \dots),
	\]
	and the actual tensor product is obtained by spectrifying this object.

	In certain situations, it is useful to do a relative version of the same construction which uses the slice description from \Cref{prop:Slice_Description}.

	\begin{construction}\label[construction]{con:Tensor_Product_In_Slice}
		Fix a Stefanich ring $A \in \StRing$. By \Cref{prop:Slice_Description} we have $\StRing_{A/} \simeq \StRing(A)$, and by \Cref{cor:Colimits_In_StRing} colimits in $\StRing(A)$ are computed by passing to $\PreStRing(A)$, where they are pointwise, and then spectrifying.

		Now let $B,C \in \StRing_{A/}$. Under the equivalence of \Cref{prop:Slice_Description}, they correspond to the objects of $\PreStRing(A)$ given by the sequences $((B/A)_m)_{m \geq 0}$ and $((C/A)_m)_{m \geq 0}$, where $(B/A)_m \in \CAlg(A_{m+1})$. Their coproduct in $\PreStRing(A)$ is therefore the pointwise tensor product
		\[
			\bigl((B/A)_0 \otimes (C/A)_0,\;
			(B/A)_1 \otimes (C/A)_1,\;
			\dots\bigr),
		\]
		where the $m$-th tensor product is formed in $\CAlg(A_{m+1})$. The coproduct in $\StRing_{A/}$ is obtained by spectrifying this internal preStefanich ring.
	\end{construction}

	Note that this method of computing the tensor product was already used at the end of the previous talk, in the proof of \Cref{cor:StRing_Functor_Properties}.

	\begin{construction}\label[construction]{con:Tensor_Product_By_Cobase_Change}
		In a yet more refined way, let $g\colon A \to C$ be a morphism of Stefanich rings. The cobase-change functor
		\[
			g^*\colon \StRing_{A/} \to \StRing_{C/}, \qquad B \mapsto B \otimes_A C
		\]
		is computed by passing from $\PreStRing(A)$ to $\PreStRing(C)$. Namely, first apply levelwise extension of scalars
		\[
			g_m^*\colon \CAlg(A_{m+1}) \to \CAlg(C_{m+1})
		\]
		to the sequence attached to $B$. This produces the internal preStefanich ring in $C$
		\[
			(g_0^*(B/A)_0,\; g_1^*(B/A)_1,\; g_2^*(B/A)_2,\; \dots).
		\]
		The structure maps are obtained by applying $g_m^*$ to the structure maps of $B$ and then using the natural comparison
		\[
			g_m^*(\End^A_{m+1}(R)) \to \End^C_{m+1}(g_{m+1}^*(R)),
			\qquad R \in \CAlg(A_{m+2}).
		\]
		In other words, for $R_m = (B/A)_m$, the structure map is the composite
		\[
			g_m^*R_m \to g_m^*\End^A_{m+1}(R_{m+1}) \to \End^C_{m+1}(g_{m+1}^*R_{m+1}).
		\]
		Spectrifying this pre-object in $\PreStRing(C)$ gives $B \otimes_A C$.
	\end{construction}

	\emph{During the seminar, this talk also covered the theory of affine maps of Stefanich rings. For the sake of exposition, that material has been collected in \Cref{sec:Affine_Maps} of the following chapter.}

	\chapter{Affine, proper and étale maps of Stefanich rings}
	\label[chapter]{chap:Properties_Of_Maps}

	\emph{Talk given by Tashi.}

	In this talk, we will introduce a variety of properties of maps of Stefanich rings: affine maps, prim maps, proper maps, suave maps, and étale maps.

	\section{Recollections and notations}
	\label[section]{sec:Recollections_And_Notations}

	Before we start, let us recall the notation for the transition maps attached to a morphism of Stefanich rings. First recall the external construction. Let
	\[
		\bbP = (\ast, \An,1\Pr,2\Pr,\dots)
	\]
	denote the initial Stefanich ring, with $\bbP_n = (n-1)\Pr$. The absolute transition adjunctions are the adjunctions
	\[
		\Mod^{\bbP}_n\colon \CAlg((n-1)\Pr)
		\rightleftarrows
		\CAlg(n\Pr) \noloc
		\End^{\bbP}_n.
	\]
	Thus a Stefanich ring $A$ consists of objects $A_n \in \CAlg(n\Pr)$ together with structure equivalences
	\[
		A_n \iso \End^{\bbP}_{n+1}(A_{n+1}) = \End_{A_{n+1}}(\unit_{A_{n+1}}).
	\]
	By adjunction, these are equivalently encoded by symmetric monoidal left adjoints
	\[
		L^A_n\colon \Mod_{A_n}(n\Pr) \to A_{n+1}.
	\]

	In the previous talk, we repeated this construction internally to a fixed Stefanich ring $A$. Namely, for $m \geq 1$ there are internal module-endomorphism adjunctions
	\[
		\Mod^A_m\colon \CAlg(A_m)
		\rightleftarrows
		\CAlg(A_{m+1}) \noloc
		\End^A_m.
	\]
	These adjunctions give the category $\StRing(A)$ of Stefanich rings internal to $A$, and \Cref{prop:Slice_Description} identifies it with the slice category $\StRing_{A/}$. Under this equivalence, a map $f\colon A \to B$ is sent to the internal Stefanich ring $B/A$ whose degree $n$ term is
	\[
		(B/A)_n = f_{n,*}(\unit_{B_{n+1}}) \qin \CAlg(A_{n+1}),
	\]
	as in \Cref{con:Lift_To_Algebra}. In particular, $B/A$ comes equipped with internal structure equivalences
	\[
		(B/A)_n \iso \End^A_{n+1}((B/A)_{n+1}).
	\]
	Adjointing these equivalences gives maps in $\CAlg(A_{n+2})$
	\[
		L^{B/A}_n\colon
		\Mod^A_{n+1}((B/A)_n)
		=
		L^A_{n+1}\bigl(\Mod_{(B/A)_n}(A_{n+1})\bigr)
		\to
		(B/A)_{n+1}.
	\]
	For $B=A$, this gives the corresponding maps $L^{A/A}_n$ for the unit internal Stefanich ring $A/A$.

	Now regard $f$ as a morphism $A/A \to B/A$ in $\StRing(A)$. In degree $n$, this gives a map
	\[
		(f/A)^*_{n-1}\colon (A/A)_n \to (B/A)_n
	\]
	in $\CAlg(A_{n+1})$; the shifted subscript follows the convention used for pullback functors. The fact that these maps define a morphism of internal Stefanich rings means that they are compatible with the internal structure equivalences. After adjunction, this compatibility says that the square
	\[
		\begin{tikzcd}[column sep = huge]
			\Mod^A_{n+1}((A/A)_n) \rar["\Mod^A_{n+1}((f/A)^*_{n-1})"] \dar["L^{A/A}_n"'] & \Mod^A_{n+1}((B/A)_n) \dar["L^{B/A}_n"] \\
			(A/A)_{n+1} \rar["(f/A)^*_n"'] & (B/A)_{n+1}
		\end{tikzcd}
	\]
	commutes. Since $\Mod^A_{n+1}$ preserves monoidal units, $L^{A/A}_n$ is an equivalence, and thus we see that the map $(f/A)^*_n$ factors as follows:
	\[
		(A/A)_{n+1} \to \Mod^A_{n+1}((B/A)_n) \xrightarrow{L^{B/A}_n} (B/A)_{n+1}.
	\]

	For later use, we also fix notation for internal Hom objects relative to $A$. Let $m \geq 1$ and let $D \in \CAlg(A_{m+1})$. For $X,Y \in D$, we write
	\[
		\iHom^A_D(X,Y) := \iHom_{U^A_m(D)}(X,Y) \qin A_m,
	\]
	where $U^A_m\colon A_{m+1} \to \Mod_{A_m}(m\Pr)$ is the right adjoint from \Cref{con:Internal_Module_Endomorphism_Adjunction}. In particular,
	\[
		\End^A_m(D) = \iHom^A_D(\unit_D,\unit_D).
	\]

	\section{Affine maps}
	\label[section]{sec:Affine_Maps}

	\emph{During the seminar, this material was presented as part of the previous talk; it has been collected here for the sake of exposition.}

	\begin{definition}\label[definition]{def:N_Affine}
		Consider $n \geq 0$, and let $A \in \StRing$ be a Stefanich ring. Then $B \in \StRing_{A/}$ is called \emph{$n$-affine} if it lies in the essential image of the functor $\CAlg(A_{n+1}) \hookrightarrow \StRing_{A/}$.
	\end{definition}

	Thus an $n$-affine object over $A$ is obtained from a single algebra
	\[
		R \in \CAlg(A_{n+1})
	\]
	by inserting it as the relative algebra in degree $n$ and then applying the left adjoint transition functors in all higher degrees. In other words, if $B$ is the $n$-affine object determined by $R$, then
	\[
		(B/A)_n \simeq R,
		\qquad
		(B/A)_{m+1} \simeq \Mod^A_{m+1}((B/A)_m)
		\quad (m \geq n),
	\]
	while the lower relative algebras are forced by the right adjoint transitions $\End^A_m$. Equivalently, for an arbitrary $B \in \StRing_{A/}$, the object $B$ is $n$-affine precisely when the relative structure maps
	\[
		L^{B/A}_m\colon \Mod^A_{m+1}((B/A)_m) \to (B/A)_{m+1}
	\]
	are equivalences for all $m \geq n$.

	The fully faithful embedding $\CAlg(A_{n+1}) \hookrightarrow \StRing_{A/}$ already says that the recursive construction of an $n$-affine object does not lose information. We will also need a slightly different recovery statement: in high degrees, the relative algebra $(B/A)_{m+1}$ is obtained by applying $L^A_{m+1}$ to the $A_{m+1}$-linear category $\Mod_{(B/A)_m}(A_{m+1})$, and we want to recover this linear category from its image in $A_{m+2}$. This is where \Cref{lem:Lm_Fully_Faithful_On_Dualizables} enters. If $R \in \CAlg(A_{m+1})$ is countably presented, then $\Mod_R(A_{m+1})$ is dualizable in $\Mod_{A_{m+1}}((m+1)\Pr)$. More generally, after writing $R$ as an $\aleph_1$-filtered colimit of countably presented algebras, and using that $R \mapsto \Mod_R(A_{m+1})$ preserves such colimits, $\Mod_R(A_{m+1})$ is an $\aleph_1$-filtered colimit of dualizable objects. On this class, the functor
	\[
		L^A_{m+1}\colon \Mod_{A_{m+1}}((m+1)\Pr) \to A_{m+2}
	\]
	is fully faithful.

	\begin{remark}
		We have a factorization $\CAlg(A_{n+1}) \hookrightarrow \CAlg(A_{n+2}) \hookrightarrow \StRing_{A/}$. In particular, any $n$-affine Stefanich ring under $A$ is also $(n+1)$-affine.
	\end{remark}

	Here are some basic properties of affineness:

	\begin{proposition}\label[proposition]{prop:Affineness_Properties}
		Let $n \geq 0$ and $A \in \StRing$.
		\begin{enumerate}
			\item The class of $n$-affines over $A$ is closed under small colimits.
			\item If $B \in \StRing_{A/}$ is $n$-affine, then the comparison map $\Mod_{(B/A)_m}(A_{m+1}) \to B_{m+1}$ is an equivalence for $m \geq n$.
			\item Given a map $g\colon A \to C$ in $\StRing$ and if $B \in \StRing_{A/}$ is $n$-affine, then also $g^*B = B \otimes_A C \in \StRing_{C/}$ is $n$-affine. Moreover, we have
			\[
				g_m^*(B/A)_m \simeq (B \otimes_A C/C)_m
			\]
			for $m \geq n$.
			\item Given maps $f\colon A \to B$ and $h\colon B \to C$ in $\StRing$ with composite $g\colon A \to C$, if $f$ is $n$-affine, then $g$ is $n$-affine if and only if $h$ is.
			\item If $B \in \StRing_{A/}$ is $n$-affine, then the codiagonal $\nabla\colon B \otimes_A B \to B$ is $n$-affine; consequently all iterated codiagonals of $A \to B$ are $n$-affine.
		\end{enumerate}
	\end{proposition}
	\begin{proof}
		(1) is clear, since $\CAlg(A_{n+1}) \hookrightarrow \StRing_{A/}$ is a fully faithful left adjoint.

		(2) Let $m \geq n$, and let $U^A_{m+1}$ denote the right adjoint of
		\[
			L^A_{m+1}\colon \Mod_{A_{m+1}}((m+1)\Pr) \to A_{m+2}.
		\]
		The comparison map in question is the unit map
		\[
			\Mod_{(B/A)_m}(A_{m+1})
			\to
			U^A_{m+1}L^A_{m+1}\bigl(\Mod_{(B/A)_m}(A_{m+1})\bigr),
		\]
		followed by the identification of the target with $B_{m+1}$. Indeed, by the recursive description above, since $B$ is $n$-affine over $A$, we have
		\[
			\Mod^A_{m+1}((B/A)_m)
			\simeq
			(B/A)_{m+1}.
		\]
		Moreover, by the discussion preceding the proposition, after writing $(B/A)_m$ as an $\aleph_1$-filtered colimit of countably presented algebras, the object $\Mod_{(B/A)_m}(A_{m+1})$ is an $\aleph_1$-filtered colimit of dualizable objects. Hence \Cref{lem:Lm_Fully_Faithful_On_Dualizables} implies that the above unit map is an equivalence. Applying the displayed identification, the target becomes
		\[
			U^A_{m+1}((B/A)_{m+1})
			=
			\fgt((B/A)_{m+1})
			\simeq
			B_{m+1},
		\]
		as desired.

		(3) We use the description of cobase change from \Cref{con:Tensor_Product_By_Cobase_Change}. The pre-object whose spectrification computes $g^*B$ is represented in $\PreStRing(C)$ by the sequence
		\[
			(g_0^*(B/A)_0,\; g_1^*(B/A)_1,\; g_2^*(B/A)_2,\; \dots).
		\]
		For $m \geq n$, the recursive description of $n$-affineness gives
		\[
			(B/A)_{m+1}
			\simeq
			\Mod^A_{m+1}((B/A)_m).
		\]
		Extension of scalars is compatible with these transition functors, in the sense that for $R \in \CAlg(A_{m+1})$ we have
		\[
			g_{m+1}^*\Mod^A_{m+1}(R)
			\simeq
			\Mod^C_{m+1}(g_m^*R).
		\]
		Thus the transition maps of the above preStefanich ring are already equivalences in all degrees $m \geq n$. Spectrification therefore leaves those degrees unchanged, and the result is the $n$-affine Stefanich ring over $C$ determined by $g_n^*(B/A)_n$. In particular,
		\[
			g_m^*(B/A)_m \simeq (B \otimes_A C/C)_m
		\]
		for $m \geq n$.

		(4) Assume $B$ is $n$-affine over $A$, so that it is the image of $(B/A)_n \in \CAlg(A_{n+1})$. Slicing the fully faithful functor
		\[
			\CAlg(A_{n+1}) \hookrightarrow \StRing_{A/}
		\]
		at this object gives a fully faithful functor
		\[
			\CAlg(A_{n+1})_{(B/A)_n/}
			\hookrightarrow
			\StRing_{B/},
		\]
		whose essential image consists precisely of those objects $C \in \StRing_{B/}$ for which the composite $A \to C$ is $n$-affine.
		On the other hand, by part (2),
		\[
			B_{n+1}
			\simeq
			\Mod_{(B/A)_n}(A_{n+1}),
		\]
		and therefore
		\[
			\CAlg(B_{n+1})
			\simeq
			\CAlg(A_{n+1})_{(B/A)_n/}.
		\]
		Under this equivalence, the functor $\CAlg(B_{n+1}) \hookrightarrow \StRing_{B/}$ has the same essential image as the sliced functor above. Thus $C$ is $n$-affine over $B$ if and only if it is $n$-affine over $A$.

		(5) Let $\iota\colon B \to B \otimes_A B$ be a coprojection, i.e.\ the cobase change of the structure map $f\colon A \to B$ along itself; it is $n$-affine by (3). Hence (4), applied to $A \xrightarrow{f} B \xrightarrow{\iota} B \otimes_A B$, shows that the structure map $\eta\colon A \to B \otimes_A B$ is $n$-affine. As $\nabla \circ \eta = f$, a second application of (4), now to $A \xrightarrow{\eta} B \otimes_A B \xrightarrow{\nabla} B$, shows that $\nabla$ is $n$-affine. The statement for iterated codiagonals follows by induction, since each is the codiagonal of the $n$-affine map obtained at the previous step.
	\end{proof}

	\begin{remark}
		The converse to part (2) is \emph{not} necessarily true. Let $k$ be a field, let
		\[
			f\colon * \to B^2 \bbG_m
		\]
		be the base point, and let $A \to B$ be the corresponding map of Stefanich rings. One can show that $f$ is $1$-affine, so the comparison in part (2) is an equivalence for $m \geq 1$. It is also an equivalence for $m = 0$: here $A_1 = B_1 = \D(k)$ and $(B/A)_0 = k$, so the comparison is just
		\[
			\Mod_k(\D(k)) \simeq \D(k).
		\]
		Nevertheless $f$ is not $0$-affine. Indeed, if it were $0$-affine, then $B$ would be determined over $A$ by the algebra $(B/A)_0 \in \CAlg(A_1)$. But $(B/A)_0 = k$ is the unit algebra, and the $0$-affine object corresponding to the unit is just the identity object $A \in \StRing_{A/}$. This would force $A \simeq B$, i.e.\ $* \simeq B^2 \bbG_m$, which is false.
	\end{remark}

	\section{Prim maps}
	\label[section]{sec:Prim_Maps}

	\begin{definition}\label[definition]{def:N_Prim}
		The map $f\colon A \to B$ is called \emph{$n$-prim} if for all $m \geq n$ the map
		\[
			(f/A)^*_m\colon (A/A)_{m+1} \to (B/A)_{m+1}
		\]
		admits a right adjoint $(f/A)_{m,*}\colon (B/A)_{m+1}\to (A/A)_{m+1}$ in $A_{m+2}$.
	\end{definition}

	\begin{observation}
		\label[observation]{obs:Endomorphism_Object_Via_Right_Adjoint}
		If $f$ is $n$-prim, the composite $(f/A)_{m,*}(f/A)^*_m\colon (A/A)_{m+1} \to (A/A)_{m+1}$ defines an endomorphism of the unit in $A_{m+2}$. We claim that it is sent to $(B/A)_m$ under the equivalence $\End_{A_{m+2}}(\unit) \iso A_{m+1}$ that evaluates an endomorphism at the $\unit \in A_{m+1}$. Indeed, since the forgetful 2-functor $A_{m+2} \to (m+2)\Pr$ preserves adjunctions, the internal adjunction $(f/A)^*_m \dashv (f/A)_{m,*}$ in $A_{m+2}$ refines the external adjunction $f^*_m \dashv f_{m,*}$, and we compute
		\[
			(f/A)_{m,*}(f/A)^*_m(\unit_{A_{m+1}}) \simeq f_{m,*}f^*_m(\unit_{A_{m+1}}) = (B/A)_m.
		\]
	\end{observation}

	\begin{proposition}\label[proposition]{prop:Prim_Properties}
		Let $A \in \StRing$.
		\begin{enumerate}
			\item If $B/A$ is $n$-affine and countably presented, then $B/A$ is $n$-prim.
			\item The $n$-prim Stefanich $A$-algebras are closed under $\aleph_0$-filtered colimits.
			\item Consider a pushout square of Stefanich rings
			\[
				\begin{tikzcd}
					A \rar["f"] \dar["g"'] \drar[pushout] & B \dar["g'"] \\
					C \rar["f'"'] & C \otimes_A B.
				\end{tikzcd}
			\]
			If $f$ is $n$-prim, then so is $f'$, and for every $m \geq n$ the canonical map
			\[
				g^*_m(B/A)_m \to (B \otimes_A C / C)_m
			\]
			is an isomorphism (i.e.\ we do not need to ``spectrify''). Moreover, the internal Beck--Chevalley map
			\[
				(g/A)^*_m (f/A)_{m,*} \to (f'/A)_{m,*}(g'/A)^*_m
			\]
			is an isomorphism of morphisms $(B/A)_{m+1} \to (C/A)_{m+1}$ in $A_{m+2}$.
			\item The $n$-prim maps are closed under composition.
		\end{enumerate}
	\end{proposition}
	\begin{proof}
		(1) The map $(f/A)^*_m\colon (A/A)_{m+1} \to (B/A)_{m+1}$ factors as
		\[
			(A/A)_{m+1} \to \Mod^A_{m+1}((B/A)_m) \xrightarrow{L^{B/A}_m} (B/A)_{m+1}.
		\]
		If $B/A$ is $n$-affine, then the second map is an equivalence for $m \geq n$, and so this composite has a right adjoint in $A_{m+2}$ if and only if the first map admits a right adjoint. By definition, the first map is the image under the 2-functor $L^A_{m+1}\colon \Mod_{A_{m+1}}((m+2)\Pr) \to A_{m+2}$ of the free module functor
		\[
			A_{m+1} \to \Mod_{(B/A)_m}(A_{m+1}).
		\]
		If $(B/A)_m \in \CAlg(A_{m+1})$ is countably presented, then this free module functor admits a right adjoint in $\Mod_{A_{m+1}}((m+2)\Pr)$ by part \eqref{it:Restriction_Scalars_Internal_Right_Adjoint} of \Cref{rmk:Basic_Properties_Extension_Scalars}.
		
		(4) Let $g\colon B \to C$ be a further $n$-prim map, with composite $h = g \circ f\colon A \to C$. Fix $m \geq n$. Since $g$ is $n$-prim, the relative map $(g/B)^*_m\colon \unit \to (C/B)_{m+1}$ admits a right adjoint in $B_{m+2}$. The forgetful functor $B_{m+2} \to A_{m+2}$ is a $2$-functor of the underlying $2$-categories, and any $2$-functor preserves adjunctions; applying it to the previous adjunction exhibits a right adjoint of
		\[
			(B/A)_{m+1} \xrightarrow{(g/A)^*_m} (C/A)_{m+1}
		\]
		in $A_{m+2}$. Since $f$ is $n$-prim, the relative map $(f/A)^*_m\colon \unit \to (B/A)_{m+1}$ also admits a right adjoint in $A_{m+2}$. Composing the two right adjoints exhibits a right adjoint of
		\[
			(h/A)^*_m\colon \unit \xrightarrow{(f/A)^*_m} (B/A)_{m+1} \xrightarrow{(g/A)^*_m} (C/A)_{m+1}
		\]
		in $A_{m+2}$. As $m \geq n$ was arbitrary, $h$ is $n$-prim.

		(3) Recall from \Cref{con:Tensor_Product_By_Cobase_Change} that the relative tensor product $C \otimes_A B \in \StRing_{C/}$ may be computed as the spectrification of the internal preStefanich ring
		\[
			(g_0^*(B/A)_0, g_1^*(B/A)_1, g_2^*(B/A)_2, \dots)
		\]
		in $C$. We show that in degrees $m \geq n$ the levels $g^*_m(B/A)_m$ already assemble into a Stefanich ring, i.e.\ that the canonical maps $g^*_m(B/A)_m \to \End^C_{m+1}(g^*_{m+1}(B/A)_{m+1})$ are isomorphisms; it then follows from \Cref{obs:Stefanich_Ring_In_High_Degrees} that $g^*(B/A)_\bullet$ already agrees with $(C \otimes_A B / C)_\bullet$ in these degrees.

		Fix $m \geq n$. Since $f$ is $n$-prim, we have the internal adjunction $(f/A)^*_m \dashv (f/A)_{m,*}$ in $A_{m+2}$, and by the observation preceding the proposition the algebra $(B/A)_m \in \CAlg(A_{m+1})$ is the value of the composite
		\[
			(A/A)_{m+1} \xrightarrow{(f/A)^*_m} (B/A)_{m+1} \xrightarrow{(f/A)_{m,*}} (A/A)_{m+1}
		\]
		under the equivalence $\End_{A_{m+2}}(\unit) \iso A_{m+1}$. Applying the $2$-functor $g^*_{m+1}\colon A_{m+2} \to C_{m+2}$ to this adjunction yields an internal adjunction
		\[
			(C/C)_{m+1} \xrightarrow{g^*_{m+1}((f/A)^*_m)} g^*_{m+1}((B/A)_{m+1}) \xrightarrow{g^*_{m+1}((f/A)_{m,*})} (C/C)_{m+1}
		\]
		in $C_{m+2}$. The functor $g^*_{m+1}((f/A)^*_m)$ is symmetric monoidal and so preserves the monoidal unit. We may now compute the $C_{m+1}$-linear endomorphisms of the unit in $g^*_{m+1}((B/A)_{m+1}) \in \CAlg(C_{m+2})$ as
		\begin{align*}
			\End^C_{m+1}(g^*_{m+1}((B/A)_{m+1})) &= \underline{\Hom}_{g^*_{m+1}((B/A)_{m+1})}(\unit, \unit) \\
			&= \underline{\Hom}_{g^*_{m+1}((B/A)_{m+1})}(g^*_{m+1}((f/A)^*_m)\unit, g^*_{m+1}((f/A)^*_m)\unit) \\
			&\simeq \underline{\Hom}_{(C/C)_{m+1}}(\unit, g^*_{m+1}\big((f/A)_{m,*}(f/A)^*_m\big)\unit) \\
			&\simeq g^*_{m+1}\big((f/A)_{m,*}(f/A)^*_m\big)\unit \\
			&\overset{(1)}{\simeq} g_m^*\big((f/A)_{m,*}(f/A)^*_m\unit\big) \\
			&\overset{(2)}{\simeq} g_m^*((B/A)_m),
		\end{align*}
		where the equivalence (1) uses the fact that the map of Stefanich rings $g\colon A \to C$ comes with a commutative diagram of the form
		\[
		\begin{tikzcd}
			\End_{A_{m+2}}(\unit) \rar{g^*_{m+1}} \dar{\ev_{\unit}}[swap]{\simeq} & \End_{C_{m+2}}(\unit) \dar{\ev_{\unit}}[swap]{\simeq} \\
			A_{m+1} \rar{g^*_m} & C_{m+1},
		\end{tikzcd}
		\]
		while (2) holds by \Cref{obs:Endomorphism_Object_Via_Right_Adjoint}. 
		We conclude that $g^*(B/A)_m = (C \otimes_A B / C)_m$ for $m \geq n$. Moreover the map $(f'/C)_m^*\colon (C/C)_{m+1} \to g^*(B/A)_{m+1} = (C \otimes_A B / C)_{m+1}$ admits the right adjoint $g^*_{m+1}((f/A)_{m,*})$ in $C_{m+2}$, so $f'$ is $n$-prim.

		It remains to prove the internal Beck--Chevalley statement: we need to show that the square
		\[
		\begin{tikzcd}
			(A/A)_{m+1} \dar[swap]{(f/A)^*_m} \rar{(g/A)^*_m} & (C/A)_{m+1} \dar{(f'/A)^*_m} \\
			(B/A)_{m+1} \rar{(g'/A)^*_m} & (B \otimes_A C/A)_{m+1}
		\end{tikzcd}
		\]
		in $A_{m+2}$ is vertically right adjointable. By the preceding proof, this square is isomorphic to the naturality square
		\[
		\begin{tikzcd}
			(A/A)_{m+1} \dar[swap]{(f/A)^*_m} \rar{\eta} & g_{m+1,*}g^*_{m+1}(A/A)_{m+1} \dar{g_{m+1,*}g^*_{m+1}(f/A)^*_m} \\
			(B/A)_{m+1} \rar{\eta} & g_{m+1,*}g^*_{m+1}(B/A)_{m+1}
		\end{tikzcd}
		\]
		of the unit $\eta\colon \id \to g_{m+1,*}g^*_{m+1}$ of the 2-categorical adjunction $g_{m+1}^*\colon A_{m+2} \rightleftarrows C_{m+2}: g_{m+1,*}$. The vertical adjointability of this square is formal.

		(2) This is a consequence of the following lemma.
	\end{proof}

	\begin{lemma}\label[lemma]{lem:Filtered_Colimit_Of_Adjoints}
		Consider a functor $X\colon I \to \bbC \in 2\Pr$ where $I$ is filtered and $0 \in I$ is an initial object. Assume that each $F_m\colon X_0 \to X_m$ has a right adjoint $F_m^R$. Then the functor $F\colon X_0 \to \colim X$ admits a right adjoint $F^R\colon \colim_m X_m \to X_0$, and we have
		\[
			F^R F = \colim_{m \in I} F_m^R F_m \qin \bbC(X_0, X_0).
		\]
	\end{lemma}
	\begin{proof}
		We argue by the (enriched) Yoneda lemma. For each object $Y \in \bbC$, the represented functor
		\[
			\bbC(-, Y)\colon \bbC^{\catop} \to 1\Pr
		\]
		sends a $1$-morphism that is a left adjoint in $\bbC$ to a right adjoint functor of categories. Conversely, a $1$-morphism $G$ of $\bbC$ is a left adjoint as soon as $\bbC(-, Y)(G)$ admits a right adjoint for every $Y$, and these right adjoints are compatible with the maps $\bbC(-, Y) \to \bbC(-, Y')$ induced by morphisms $Y \to Y'$ (the Beck--Chevalley condition); as usual, it suffices to test this on $Y = X_0$ and $Y = \colim X$. We may therefore assume $\bbC = 1\Pr$.

		In $1\Pr$ this is the standard comparison between the colimit presentation $\colim_{m \in I} X_m$ in $\PrL$ and the limit presentation $\lim_{m \in I^{\catop}} X_m$ formed along the right adjoints. Writing $\iota_m\colon X_m \to \colim X$ for the structure maps, each $\iota_m$ admits a right adjoint $\iota_m^R$; since $0 \in I$ is initial we have $F = \iota_0$, and hence a right adjoint $F^R := \iota_0^R$. Under the limit presentation $F^R$ is computed by the partial right adjoints, and the Beck--Chevalley condition required above holds by construction. Tracing through this description, the composite $F^R F$ is computed as the colimit
		\[
			F^R F = \colim_{m \in I} F_m^R F_m \qin \bbC(X_0, X_0),
		\]
		the colimit being filtered, so that it is again computed in $\bbC(X_0, X_0)$ compatibly with the structure maps; the unit $\id \to F^R F$ is the inclusion of the initial index $m = 0$, for which $F_0 = \id_{X_0}$. Alternatively, one writes down the unit and counit of the adjunction explicitly and checks the triangle identities directly.
	\end{proof}

	\begin{lemma}[Prim projection formula]
		\label[lemma]{lem:Prim_Projection_Formula}
		Let $f\colon A \to B$ be an $n$-prim map of Stefanich rings. Then for all $m \geq n$, $X \in A_{m+1}$ and $Y \in B_{m+1}$, the canonical map
		\[
			X \otimes_{A_{m+1}} f_{m,*}(Y) \to f_{m,*}\bigl(f_m^*(X) \otimes_{B_{m+1}} Y\bigr)
		\]
		is an isomorphism.
	\end{lemma}
	\begin{proof}
		Fix $m \geq n$. By the definition of $n$-primness, the relative pullback
		\[
			(f/A)^*_m\colon (A/A)_{m+1} \to (B/A)_{m+1}
		\]
		admits a right adjoint $(f/A)_{m,*}$ as a $1$-morphism in $A_{m+2}$. Apply the $2$-functor underlying
		\[
			U^A_{m+1}\colon A_{m+2} \to \Mod_{A_{m+1}}((m+1)\Pr).
		\]
		Since $2$-functors preserve adjunctions, this gives an adjunction of $A_{m+1}$-module objects whose underlying adjunction is
		\[
			f_m^*\colon A_{m+1} \rightleftarrows B_{m+1} \noloc f_{m,*}.
		\]
		In particular, the right adjoint $f_{m,*}$ is $A_{m+1}$-linear, where the action of $A_{m+1}$ on $B_{m+1}$ is induced by the symmetric monoidal functor $f_m^*$. Its linearity isomorphism is therefore
		\[
			X \otimes_{A_{m+1}} f_{m,*}(Y)
			\iso
			f_{m,*}\bigl(f_m^*(X) \otimes_{B_{m+1}} Y\bigr).
		\]
		This isomorphism is adjoint to the composite obtained from the counit $f_m^*f_{m,*}(Y) \to Y$; hence it is precisely the canonical projection map in the statement.
	\end{proof}

	The relative objects of $n$-prim maps satisfy a Künneth formula for tensor products over $A$:

	\begin{lemma}[Künneth for prim maps]\label[lemma]{lem:Kunneth_Prim}
		Let $f\colon A \to B$ and $g\colon A \to C$ be $n$-prim maps of Stefanich rings. Then
		\[
			(B \otimes_A C/A)_n \cong (B/A)_n \otimes_{A_{n+1}} (C/A)_n.
		\]
	\end{lemma}
	\begin{proof}
		Let $\pr\colon C \to B \otimes_A C$ be the cobase change of $f$ along $g$, so that $\pr \circ g\colon A \to B \otimes_A C$ is the structure map. By functoriality of the relative $*$-pushforward,
		\[
			(B \otimes_A C/A)_n = (\pr \circ g)_{n,*}(\unit) = g_{n,*}\bigl(\pr_{n,*}(\unit)\bigr) = g_{n,*}\bigl((B \otimes_A C/C)_n\bigr).
		\]
		Since $\pr$ is the cobase change of the $n$-prim map $f$ along $g$, \Cref{prop:Prim_Properties}(3) gives $(B \otimes_A C/C)_n \simeq g_n^*\bigl((B/A)_n\bigr)$, and therefore
		\[
			(B \otimes_A C/A)_n \simeq g_{n,*} g_n^*\bigl((B/A)_n\bigr).
		\]
		As $g$ is $n$-prim, the right adjoint $g_{n,*}$ satisfies the projection formula, so
		\[
			(B \otimes_A C/A)_n \simeq g_{n,*} g_n^*\bigl((B/A)_n\bigr) \simeq (B/A)_n \otimes_{A_{n+1}} g_{n,*}(\unit) = (B/A)_n \otimes_{A_{n+1}} (C/A)_n. \qedhere
		\]
	\end{proof}

	\section{Proper maps}
	\label[section]{sec:Proper_Maps}

	Note that one missing feature of the previous proposition is that $n$-prim maps are not necessarily closed under diagonals.

	\begin{example}\label[example]{ex:DeRham_Stack}
		Let $k$ be a field of characteristic $0$ (for simplicity), and let $X$ be a variety over $k$. The map of Stefanich rings corresponding to the de Rham stack $X_{\mathrm{dR}} \to \Spec(k)$ is $0$-prim over $k$, but its diagonal is not. Thus $n$-primness genuinely fails to be stable under passage to diagonals.
	\end{example}

	This motivates the following definition.

	\begin{definition}\label[definition]{def:N_Proper}
		A map $f\colon A \to B$ of Stefanich rings is called \emph{$n$-proper} if $f$ and all its iterated codiagonals
		\[
			\nabla_f\colon B \otimes_A B \to B, \qquad \nabla_{\nabla_f}\colon B \otimes_{B \otimes_A B} B \to B, \qquad \dots
		\]
		are all $n$-prim.
	\end{definition}

	\begin{proposition}\label[proposition]{prop:Proper_Properties}
		Let $A \in \StRing$.
		\begin{enumerate}
			\item If $B \in \StRing_{A/}$ is $n$-affine and countably presented, then $B$ is $n$-proper.
			\item The full subcategory of $\StRing_{A/}$ spanned by the $n$-proper $A$-algebras is closed under countable colimits.
			\item The subcategory in (2) is also closed under cobase change.
			\item Let $f\colon A \to B$ be $n$-proper, and let $g\colon B \to C$ be a map. Then $g$ is $n$-prim if and only if $gf$ is $n$-prim. Moreover, $g$ is $n$-proper if and only if $gf$ is $n$-proper.
		\end{enumerate}
	\end{proposition}
	\begin{proof}
		All four statements follow from \Cref{prop:Prim_Properties}, together with the stability of $n$-affine maps under passing to iterated codiagonals (\Cref{prop:Affineness_Properties}(5)).

		(1) If $B$ is $n$-affine and countably presented over $A$, then so is every iterated codiagonal of $f\colon A \to B$, by the stability just mentioned and the fact that countable presentability is preserved under cobase change. Hence each iterated codiagonal is $n$-prim by \Cref{prop:Prim_Properties}(1), so $f$ is $n$-proper.

		(2) It suffices to treat sifted colimits and finite coproducts separately. For a sifted colimit, every iterated codiagonal of the colimit map is itself a sifted colimit of cobase changes of the corresponding iterated codiagonals of the diagram, so \Cref{prop:Prim_Properties}(2) and (3) apply degreewise. Finite coproducts reduce to stability under cobase change and composition, which are parts (3) and (4) below.

		(3) The formation of iterated codiagonals commutes with cobase change. Hence, if all codiagonals of $f$ are $n$-prim, then \Cref{prop:Prim_Properties}(3) shows that all codiagonals of the cobase change $f'$ are $n$-prim as well, i.e.\ $f'$ is $n$-proper.

		(4) An iterated codiagonal of a composite can be written as a composite of cobase changes of the iterated codiagonals of the two individual maps. Together with \Cref{prop:Prim_Properties}(3) and (4), this shows that $n$-prim and $n$-proper maps are closed under composition; in particular, if $g$ is $n$-prim (resp.\ $n$-proper), then so is $gf$. For the converse, assume $f$ is $n$-proper and $gf$ is $n$-prim (resp.\ $n$-proper). Factor $g$ as
		\[
			B \to B \otimes_A C \to C,
		\]
		where the first map is the cobase change of $gf\colon A \to C$ along $f$, hence $n$-prim (resp.\ $n$-proper), and the second map is a cobase change of the codiagonal $\nabla_f\colon B \otimes_A B \to B$, hence $n$-proper, and in particular $n$-prim. By the composition statement just proved, $g$ is $n$-prim (resp.\ $n$-proper).
	\end{proof}

	Being $0$-proper is closely related to being \emph{rigid} in the sense of Gaitsgory--Rozenblyum \citep{GaitsgoryRozenbluym2017StudyDAG}.

	\begin{proposition}\label[proposition]{prop:Proper_Rigid}
		Let $A \in \StRing$, and let $B$ be a countably presented $1$-affine Stefanich $A$-algebra. Then $B$ is $0$-proper over $A$ if and only if $(B/A)_1 \in \CAlg(A_2)$ is \emph{rigid} in $A_2$, i.e.\ the following two conditions hold.
		\begin{enumerate}
			\item The unit map $\unit \to (B/A)_1$ admits a right adjoint in $A_2$.
			\item The multiplication map $(B/A)_1 \otimes_{A_2} (B/A)_1 \to (B/A)_1$ admits a right adjoint in $\Mod_{(B/A)_1 \otimes_{A_2} (B/A)_1}(A_2)$.
		\end{enumerate}
	\end{proposition}
	\begin{proof}
		Since $B$ is $1$-affine and countably presented, it is automatically $1$-proper by \Cref{prop:Proper_Properties}(1). Thus being $0$-proper amounts to having, in addition, the relevant right adjoints at the lowest level.

		If $B$ is $0$-proper, then (1) and (2) hold by applying the $0$-prim condition to $f\colon A \to B$ and to its codiagonal $\nabla_f\colon B \otimes_A B \to B$, using that by \Cref{prop:Affineness_Properties}(2) we have $(B \otimes_A B)_2 \simeq \Mod_{(B/A)_1 \otimes_{A_2} (B/A)_1}(A_2)$.

		Conversely, it suffices to show that (1) and (2) imply that the codiagonal $\nabla_f\colon B \otimes_A B \to B$ is $0$-affine. Indeed, the codiagonal is then certainly $0$-proper, and together with $A \to B$ being $0$-prim this shows that $f$ is $0$-proper. Now the codiagonal is $1$-affine and $0$-prim, and it admits a section, the coprojection $B \to B \otimes_A B$, which is $1$-affine as a cobase change of $f$. Hence it is $0$-affine by \Cref{lem:Reduce_Affineness}.
	\end{proof}

	The following lemma is often useful to reduce the level of affineness.

	\begin{lemma}\label[lemma]{lem:Reduce_Affineness}
		Let $f\colon A \to B$ be a map of Stefanich rings. Assume that $f$ is $(n+1)$-affine and $n$-prim, and that there is an $(n+1)$-affine map $g\colon B \to A$ with $fg = \id_B$. Then $f$ is $n$-affine.
	\end{lemma}
	\begin{proof}
		By shifting we may assume $n = 0$. Replacing $A$ by the image of $(B/A)_0 \in \CAlg(A_1)$ in $\StRing_{A/}$, we may assume that $(B/A)_0 = \unit \in \CAlg(A_1)$; we must then show that $A \to B$ is an equivalence. By $1$-affineness it suffices to show that $\unit \to (B/A)_1$ is an equivalence.

		Since $f$ is $0$-prim, the map $(f/A)^*_0\colon \unit \to (B/A)_1$ admits a right adjoint $(f/A)_{0,*}$ in $A_2$. Moreover, the unit $\id \to (f/A)_{0,*} (f/A)^*_0$ in $A_1 \simeq \End_{A_2}(\unit)$ is an equivalence, since it is the map $\unit \to (B/A)_0 = \unit$. We would like to conclude that $\unit \to (B/A)_1$ is an equivalence because $g$ splits it; this is not quite true, but the following weaker statement is. After applying $g_{1,*}$, the map $\unit \to (B/A)_1$ becomes the map $(A/B)_1 \to \unit$, which is split by $g_0^*\colon \unit \to (A/B)_1$. Since $g_{1,*}$ preserves the fully faithful adjunction, $\unit \to (B/A)_1$ becomes an equivalence after applying $g_{1,*}$. But $g$ is $1$-affine, so $g_{1,*}\colon A_2 \simeq \Mod_{(A/B)_1}(B_2) \to B_2$ is conservative, which finishes the proof.
	\end{proof}

	\section{Suave maps}
	\label[section]{sec:Suave_Maps}

	\begin{definition}\label[definition]{def:N_Suave}
		A map $f\colon A \to B$ of Stefanich rings is called \emph{$n$-suave} if the following two conditions are satisfied.
		\begin{enumerate}
			\item For each $m \geq n$, the map $(f/A)^*_m\colon \unit = (A/A)_{m+1} \to (B/A)_{m+1}$ admits a left adjoint
			\[
				(f/A)_{m,\sharp}\colon (B/A)_{m+1} \to (A/A)_{m+1}
			\]
			in $A_{m+2}$.
		\end{enumerate}
		This in particular results in a left adjoint $f_{m,\sharp}\colon B_{m+1} \to A_{m+1}$ to $f^*_m$. We define the \emph{$m$-suave dualizing object of $B$ over $A$} as $(B/A)^!_m := f_{m,\sharp}(\unit_{B_{m+1}}) \in A_{m+1}$. The counit of the adjunction defines a map
		\[
			(f/A)_{m-1,!}\colon (B/A)^!_m \to (A/A)_m
		\]
		in $A_{m+1}$.
		\begin{enumerate}
			\item[(2)] For each $m \geq n+1$, the map $(f/A)_{m-1,!}$ admits a right adjoint
			\[
				(f/A)^!_{m-1}\colon (A/A)_m \to (B/A)^!_m
			\]
			in $A_{m+1}$.
		\end{enumerate}
	\end{definition}

	Since the functor $f_{m,\sharp}$ is oplax symmetric monoidal, the object $(B/A)^!_m$ is canonically a commutative coalgebra in $A_{m+1}$. Moreover, if we write $(-)^{\vee} := \iHom_{A_{m+1}}(-, \unit) \colon A_{m+1}\catop \to A_{m+1}$ for the weak duality functor, then we get by adjunction that
	\[
		\bigl((B/A)^!_m\bigr)^{\vee} = \iHom_{A_{m+1}}(f_{m,\sharp}(\unit_{B_{m+1}}), \unit_{A_{m+1}}) \simeq \iHom_{A_{m+1}}(\unit_{A_{m+1}},f_{m,*}(\unit_{B_{m+1}})) \simeq (B/A)_m.
	\]
	Under this equivalence, the counit map $(f/A)_{m-1,!}\colon (B/A)^!_m \to (A/A)_m$ dualizes to the unit map $(f/A)^*_{m-1}\colon (A/A)_m \to (B/A)_m$. Since the duality functor $(-)^{\vee}\colon (A/A)_{m+1} \to (A/A)_{m+1}$ preserves adjunctions, we obtain an isomorphism
	\[
		(f/A)_{m-1,\sharp} \cong ((f/A)^!_{m-1})^{\vee}
	\]
	of maps $(B/A)_m \to (A/A)_m$ in $A_{m+1}$. In particular, condition (2) is a strengthening of the existence of the left adjoint for $(A/A)_m \to (B/A)_m$. One should think of the existence of the right adjoint $(B/A)^!_m \to \unit$ as the ``correct'' condition to put, but its mere formulation requires the existence of the left adjoint one level higher.

	\begin{lemma}
		\label[lemma]{lem:Suave_Dualizing_Object_Via_Adjunction}
		Let $f\colon A \to B$ be $n$-suave. For each $m \geq n$, there is a preferred equivalence
		\[
			f_{m,!}f_m^!\unit_{A_{m+1}} \simeq (B/A)_m^!.
		\]
	\end{lemma}
	\begin{proof}
		Note that the duality functor $(-)^{\vee}\colon A_{m+2}\catop \to A_{m+2}$ restricts to the identity on $\End_{A_{m+2}}(\unit)$. In particular, the composite $(f/A)_{m,!}(f/A)_m^!\colon (A/A)_{m+1} \to (A/A)_{m+1}$ is identified with the composite $((f/A)_{m,!})^{\vee}((f/A)_m^!)^{\vee}$. By the previous observation, this is identified with the composite
		\[
			(A/A)_{m+1} \xrightarrow{(f/A)^*_{m}} (B/A)_{m+1} \xrightarrow{(f/A)_{m,\sharp}} (A/A)_{m+1}.
		\]
		All in all, we conclude that
		\[
			f_{m,!}f_m^!\unit_{A_{m+1}} \simeq f_{m,\sharp}f_m^*\unit_{A_{m+1}} = (B/A)^!_m.\qedhere
		\]
	\end{proof}

	\begin{proposition}\label[proposition]{prop:Suave_Properties}
		\begin{enumerate}
			\item The $n$-suave maps are closed under cobase change. More precisely, consider a pushout square of Stefanich rings
			\[
				\begin{tikzcd}
					A \rar["f"] \dar["g"'] \drar[pushout] & B \dar["g'"] \\
					C \rar["f'"'] & B \otimes_A C.
				\end{tikzcd}
			\]
			If $f$ is $n$-suave, then so is $f'$, and for $m \geq n$ the canonical map
			\[
				g_m^*(B/A)^!_m \to (B \otimes_A C/C)^!_m
			\]
			is an isomorphism. Moreover, the internal Beck--Chevalley map
			\[
				(f'/A)_{m,\sharp}(g'/A)^*_m \to (g/A)^*_m(f/A)_{m,\sharp}
			\]
			is an isomorphism of morphisms $(B/A)_{m+1} \to (C/A)_{m+1}$ in $A_{m+2}$.
			\item The $n$-suave maps are closed under composition.
		\end{enumerate}
	\end{proposition}
	\begin{proof}
		(1) By \Cref{con:Tensor_Product_By_Cobase_Change}, the relative tensor product $B \otimes_A C \in \StRing_{C/}$ is the spectrification of the internal preStefanich ring $(g^*_m(B/A)_m)_{m \geq 0}$ in $C$. Following \Cref{lem:Spectrification_Left_Adjoint}, the first spectrification step $T$ replaces the $m$-th level by $\End^C_{m+1}(g^*_{m+1}(B/A)_{m+1})$. We compute this endomorphism object, show that the result already forms a Stefanich $C$-algebra in degrees $m \geq n$, and read off both the suave structure of $f'$ and the base change of the dualizing object. Throughout we write $(-)^\vee := \iHom_{C_{m+1}}(-,\unit)$ for the weak duality functor.

		\emph{The endomorphism object.} Fix $m \geq n$. The unit map of the algebra $g^*_{m+1}(B/A)_{m+1} \in \CAlg(C_{m+2})$ is the symmetric monoidal functor
		\[
			u := g^*_{m+1}\bigl((f/A)^*_m\bigr)\colon \unit \to g^*_{m+1}(B/A)_{m+1}.
		\]
		Since $f$ is $n$-suave, $(f/A)^*_m$ admits a left adjoint $(f/A)_{m,\sharp}$ in $A_{m+2}$. As the $2$-functor $g^*_{m+1}\colon A_{m+2} \to C_{m+2}$ preserves adjunctions, it follows that $u$ admits the left adjoint $u_\sharp := g^*_{m+1}\bigl((f/A)_{m,\sharp}\bigr)$. The $C_{m+1}$-linear adjunction $u_\sharp \dashv u$ then yields
		\[
			\End^C_{m+1}\bigl(g^*_{m+1}(B/A)_{m+1}\bigr) = \iHom_{g^*_{m+1}(B/A)_{m+1}}(u\unit, u\unit) \simeq \iHom_{(C/C)_{m+1}}(u_\sharp u\unit, \unit) = (u_\sharp u\unit)^\vee.
		\]
		This is dual to the computation in the prim case \Cref{prop:Prim_Properties}(3): there one moves the right adjoint $(f/A)_{m,*}$ into the second variable of $\iHom$, here one moves the left adjoint $(f/A)_{m,\sharp}$ out of the first. Exactly as in \Cref{obs:Endomorphism_Object_Via_Right_Adjoint}, the endomorphism $(f/A)_{m,\sharp}(f/A)^*_m$ of $\unit$ in $A_{m+2}$ evaluates at the unit to $f_{m,\sharp}f^*_m\unit = (B/A)^!_m$; applying $g^*$ and using the compatibility of $g^*$ with evaluation at the unit, we obtain $u_\sharp u\unit \simeq g^*_m\bigl((B/A)^!_m\bigr)$ and therefore
		\[
			\End^C_{m+1}\bigl(g^*_{m+1}(B/A)_{m+1}\bigr) \simeq \bigl(g^*_m(B/A)^!_m\bigr)^\vee \qquad (m \geq n).
		\]

		\emph{This already forms a Stefanich $C$-algebra.} Set $P_m := \bigl(g^*_m(B/A)^!_m\bigr)^\vee$; we verify the Stefanich condition $\End_{P_m}(\unit) \simeq P_{m-1}$ for $m \geq n+1$. The second suave condition for $f$ gives an adjunction
		\[
			\unit \xrightarrow{(f/A)^!_{m-1}} (B/A)^!_m \xrightarrow{(f/A)_{m-1,!}} \unit \qin A_{m+1},
		\]
		with composite $(B/A)^!_{m-1} \in A_m = \End_{A_{m+1}}(\unit)$. Applying $g^*_m\colon A_{m+1} \to C_{m+1}$ yields
		\[
			\unit \xrightarrow{g^*_m(f/A)^!_{m-1}} g^*_m(B/A)^!_m \xrightarrow{g^*_m(f/A)_{m-1,!}} \unit \qin C_{m+1},
		\]
		with composite $g^*_{m-1}(B/A)^!_{m-1}$, and dualizing via $(-)^\vee$ turns this into an adjunction
		\[
			\unit \to P_m \to \unit \qin C_{m+1},
		\]
		again with composite $g^*_{m-1}(B/A)^!_{m-1}$, in which $P_m$ is a commutative algebra with the first map as unit and the second as its left adjoint. By the predual computation above, the composite is a predual of $\End_{P_m}(\unit)$, so
		\[
			\End_{P_m}(\unit) = \bigl(g^*_{m-1}(B/A)^!_{m-1}\bigr)^\vee = P_{m-1},
		\]
		as required. We conclude that a single application of $\End_{(-)}(\unit)$ already produces a Stefanich ring in degrees $m \geq n$, so by \Cref{lem:Spectrification_Left_Adjoint} the spectrification agrees with it there:
		\[
			(B \otimes_A C/C)_m \simeq \bigl(g^*_m(B/A)^!_m\bigr)^\vee \qquad (m \geq n).
		\]
		
		\emph{Warning:} Unlike the prim case, the right-hand side cannot in general be replaced by $g^*_m(B/A)_m$: the comparison map $g^*_m(B/A)_m \to \bigl(g^*_m(B/A)^!_m\bigr)^\vee$ is an isomorphism only when $g^*_m$ commutes with the weak dual of $(B/A)^!_m$. By \Cref{lem:Linear_Adjoint_Preserves_Dualizable} below this holds for example when $f$ is in addition $n$-prim, so that $(B/A)^!_m$ is dualizable. But in general, the relative algebra genuinely requires spectrifying and it is only the dualizing object that base-changes, as we now record.

		\emph{Suaveness of $f'$.} The previous computation shows that for $m \geq n$, the map $(f'/C)^*_{m}\colon \unit \to (B \otimes_A C)_{m+1}$ is the weak dual of the map
		\[
			g_{m+1}^*((f/A)_{m,!})\colon g_{m+1}^*(B/A)^!_{m+1} \to g_{m+1}^*(A/A) = \unit.
		\]
		Since $(f/A)_{m,!}$ admits a right adjoint $(f/A)^!_m$, so does $g_{m+1}^*((f/A)_{m,!})$, and hence $(f'/C)^*_{m}$ admits a left adjoint $(f'/C)_{m,\sharp}$. We may then compute
		\begin{align*}
			(B \otimes_A C/C)^!_m &= (f'/C)_{m,\sharp}(f'/C)^*_{m}\unit \\
			&= g_{m+1}^*((f/A)_{m,!}(f/A)^!_{m})\unit \\
			&\overset{(1)}{\simeq} g_{m}^*((f/A)_{m,!}(f/A)^!_{m}\unit) \\
			&\overset{(2)}{\simeq} g_m^*((B/A)^!_m),
		\end{align*}
		where (1) holds since $g_{m+1}^*$ restricts to $g_m^*$ on endomorphisms of the unit and (2) is \Cref{lem:Suave_Dualizing_Object_Via_Adjunction}. Under this equivalence, the counit map $(f'/C)_{m-1,!}\colon (B \otimes_A C/C)^!_m \to \unit$ is the image under $g_m^*$ of the counit map $(f/A)_{m-1,!}\colon (B/A)^!_m \to \unit$, and hence it admits a right adjoint. This shows that $f'$ is $n$-suave, and that we have
		\[
			g^*_m(B/A)^!_m \iso (B \otimes_A C/C)^!_m \qquad (m \geq n).
		\]

		\emph{Internal Beck--Chevalley.} We now deduce the internal Beck--Chevalley isomorphism by dualizing the argument given in the prim case in \Cref{prop:Prim_Properties}(3). We need to show that the commutative square
		\[
		\begin{tikzcd}
			(A/A)_{m+1} \dar[swap]{(f/A)^*_{m}} \rar{(g/A)^*_m} & (C/A)_{m+1} \dar{(f'/A)^*_m} \\
			(B/A)_{m+1} \rar{(g'/A)^*_m} & (B \otimes_A C/C)_{m+1}
		\end{tikzcd}
		\]
		in $A_{m+2}$ is vertically left adjointable. By the computations we have done before, this square may be identified with the naturality square
		\[
		\begin{tikzcd}
			(A/A)_{m+1} \dar[swap]{((f/A)^!_m)^{\vee}} \rar{\eta} & g_{m+1,*}g_{m+1}^*(A/A)_{m+1} \dar{g_{m+1,*}g_{m+1}^*(((f/A)^!_m)^{\vee})} \\
			((B/A)_{m+1}^!)^{\vee} \rar{\eta} & g_{m+1,*}g_{m+1}^*((B/A)_{m+1}^!)^{\vee}
		\end{tikzcd}
		\]
		of the unit $\eta\colon \id \to g_{m+1,*}g_{m+1}^*$ of the adjunction $g^*_{m+1}\colon A_{m+2} \rightleftarrows C_{m+2} \noloc g_{m+1,*}$. In other words, the square in question is the image of the functor $((f/A)^!_m)^{\vee}$ under the 2-functor $A_{m+2} \to \Fun([1],C_{m+2})$ corresponding to $\eta$. Since this 2-functor preserves adjunctions, it follows that the given commutative square is vertically left adjointable.

		(2) Let $f\colon A \to B$ and $g\colon B \to C$ be $n$-suave, with composite $h\colon A \to C$. For $m \geq n$ the functor $f_{m+2,*}\colon B_{m+2} \to A_{m+2}$ preserves adjunctions, so the left adjoint of $(g/B)^*_m$ in $B_{m+2}$ results in a left adjoint of the induced map $(g/A)^*_m\colon (B/A)_{m+1} \to (C/A)_{m+1}$ in $A_{m+2}$. Since adjoints compose, this gives a left adjoint for the composite
		\[
			(h/A)^*_m\colon (A/A)_{m+1} \xrightarrow{(f/A)^*_m} (B/A)_{m+1} \xrightarrow{(g/A)^*_m} (C/A)_{m+1},
		\]
		which is the first suave condition for $h$. For the second condition, the resulting maps
		\[
			(C/A)^!_m \to (B/A)^!_m \to (A/A)_m
		\]
		each admit a right adjoint for $m \geq n+1$ (the first since $g$ is $n$-suave, the second since $f$ is), hence so does their composite. Thus $h$ is $n$-suave.
	\end{proof}

	\begin{lemma}[Suave projection formula]
		\label[lemma]{lem:Suave_Projection_Formula}
		Let $f\colon A \to B$ be an $n$-suave map of Stefanich rings. Then for all $m \geq n$, $X \in A_{m+1}$ and $Y \in B_{m+1}$, the canonical map
		\[
			f_{m,\sharp}\bigl(Y \otimes_{B_{m+1}} f_m^*(X)\bigr) \to f_{m,\sharp}(Y) \otimes_{A_{m+1}} X
		\]
		is an isomorphism.
	\end{lemma}
	\begin{proof}
		Fix $m \geq n$. By the definition of $n$-suaveness, the relative pullback
		\[
			(f/A)^*_m\colon (A/A)_{m+1} \to (B/A)_{m+1}
		\]
		admits a left adjoint $(f/A)_{m,\sharp}$ as a $1$-morphism in $A_{m+2}$. Apply the $2$-functor underlying
		\[
			U^A_{m+1}\colon A_{m+2} \to \Mod_{A_{m+1}}((m+1)\Pr).
		\]
		Since $2$-functors preserve adjunctions, this gives an adjunction of $A_{m+1}$-module objects whose underlying adjunction is
		\[
			f_{m,\sharp}\colon B_{m+1} \rightleftarrows A_{m+1} \noloc f_m^*.
		\]
		In particular, the left adjoint $f_{m,\sharp}$ is $A_{m+1}$-linear, where the action of $A_{m+1}$ on $B_{m+1}$ is induced by the symmetric monoidal functor $f_m^*$. Its linearity isomorphism is therefore
		\[
			f_{m,\sharp}\bigl(Y \otimes_{B_{m+1}} f_m^*(X)\bigr)
			\iso
			f_{m,\sharp}(Y) \otimes_{A_{m+1}} X.
		\]
		This isomorphism is adjoint to the composite obtained from the unit $Y \to f_m^*f_{m,\sharp}(Y)$; hence it is precisely the canonical projection map in the statement.
	\end{proof}

	The dualizing objects of $n$-suave maps satisfy a Künneth formula, dual to \Cref{lem:Kunneth_Prim}.

	\begin{lemma}[Künneth for suave maps]\label[lemma]{lem:Kunneth_Suave}
		Let $f\colon A \to B$ and $g\colon A \to C$ be $n$-suave maps of Stefanich rings. Then
		\[
			(B \otimes_A C/A)^!_n \cong (B/A)^!_n \otimes_{A_{n+1}} (C/A)^!_n.
		\]
	\end{lemma}
	\begin{proof}
		Let $\pr\colon C \to B \otimes_A C$ be the cobase change of $f$ along $g$, so that $\pr \circ g\colon A \to B \otimes_A C$ is the structure map. By functoriality of the relative $\sharp$-pushforward,
		\[
			(B \otimes_A C/A)^!_n = (\pr \circ g)_{n,\sharp}(\unit) \simeq g_{n,\sharp}\bigl(\pr_{n,\sharp}(\unit)\bigr) = g_{n,\sharp}\bigl((B \otimes_A C/C)^!_n\bigr).
		\]
		Since $\pr$ is the cobase change of the $n$-suave map $f$ along $g$, \Cref{prop:Suave_Properties}(1) gives $(B \otimes_A C/C)^!_n \simeq g_n^*\bigl((B/A)^!_n\bigr)$, and therefore
		\[
			(B \otimes_A C/A)^!_n \simeq g_{n,\sharp} g_n^*\bigl((B/A)^!_n\bigr) \simeq (B/A)^!_n \otimes_{A_{n+1}} g_{n,\sharp}(\unit) = (B/A)^!_n \otimes_{A_{n+1}} (C/A)^!_n,
		\]
		using the projection formula for the left adjoint $g_{n,\sharp}$.
	\end{proof}

	\section{Étale maps}

	\begin{definition}\label[definition]{def:N_Etale}
		We say that $f\colon A \to B$ is \emph{$n$-\'etale} if $f$ and all its codiagonals are $n$-suave.
	\end{definition}

	\begin{proposition}\label[proposition]{prop:Etale_Properties}
		Let $A \in \StRing$ and $n \geq 0$.
		\begin{enumerate}
			\item If $f\colon A \to B$ is $n$-\'etale and $g\colon A \to C$ is any map of Stefanich rings, then the cobase change $f'\colon C \to B \otimes_A C$ is $n$-\'etale.
			\item Let $f\colon A \to B$ be $n$-\'etale, and let $g\colon B \to C$ be a map. Then $g$ is $n$-suave if and only if $gf$ is $n$-suave, and $g$ is $n$-\'etale if and only if $gf$ is $n$-\'etale.
		\end{enumerate}
	\end{proposition}
	\begin{proof}
		Since $f$ is $n$-\'etale precisely when $f$ and all of its iterated codiagonals are $n$-suave, both statements follow from \Cref{prop:Suave_Properties} by exactly the argument used to deduce \Cref{prop:Proper_Properties} from \Cref{prop:Prim_Properties}. Concretely, the iterated codiagonals of a cobase change (resp.\ of a composite) are cobase changes (resp.\ composites of cobase changes) of the iterated codiagonals of the original map(s), and $n$-suave maps are closed under cobase change and composition by \Cref{prop:Suave_Properties}.
	\end{proof}

	\section{Ambidexterity}
	\label[section]{sec:Ambidexterity_Stefanich_Rings}

	In the previous talks we saw that, for $n \geq 1$, all bi-countably presented fpqc stacks give rise to ambidextrous adjunctions, so that the corresponding maps of Stefanich rings are $1$-\'etale and $2$-proper. Abstracting that argument, with the general ambidexterity result \Cref{prop:General_Ambidexterity} in place of the geometric input, yields the following.

	\begin{proposition}\label[proposition]{prop:Ambidexterity_Stefanich_Rings}
		Let $f\colon A \to B$ be a map of Stefanich rings.
		\begin{enumerate}
			\item If $f$ and $\nabla_f\colon B \otimes_A B \to B$ are both $n$-prim, then $f$ is $(n+1)$-suave.
			\item If $f$ is $n$-proper, then $f$ is $(n+1)$-\'etale.
			\item If $f$ and $\nabla_f$ are $n$-suave, then $f$ is $(n+1)$-prim.
			\item If $f$ is $n$-\'etale, then $f$ is $(n+1)$-proper.
		\end{enumerate}
		Moreover, in any of these cases, we have an identification $f_{m,\sharp} \cong f_{m,*}$ for each $m \geq n+1$.
	\end{proposition}
	\begin{proof}
		By shifting we may assume $n = 0$.

		(1) Suppose $f$ and $\nabla_f$ are $0$-prim. Fix $m \geq 1$ and consider the map
		\[
			(f/A)^*_m\colon \unit = (A/A)_{m+1} \to (B/A)_{m+1}
		\]
		in the $(\infty,3)$-category $A_{m+2}$ (which lies in $(m+2)\Pr$, hence in $3\Pr$, since $m \geq 1$). As $f$ is $0$-prim, $(f/A)^*_m$ admits a right adjoint $(f/A)_{m,*}$ in $A_{m+2}$; we verify that the unit $\eta$ and counit $\epsilon$ of this adjunction admit right adjoints, so that \Cref{prop:General_Ambidexterity} applies.

		For $\eta$, evaluation at the unit gives an equivalence $\End_{A_{m+2}}(\unit) \iso A_{m+1}$, under which, as in the observation following \Cref{def:N_Prim}, the unit $\eta\colon \id \to (f/A)_{m,*}(f/A)^*_m$ corresponds to the map $(f/A)^*_{m-1}\colon \unit \to (B/A)_m$. Since $f$ is $0$-prim and $m - 1 \geq 0$, this map admits a right adjoint in $A_{m+1}$.

		For $\epsilon$, the counit lies in $\End_{A_{m+2}}((B/A)_{m+1})$. As in the proof of the ambidexterity property of the functors $f_{n,*}$ in \Cref{sec:Ambidexterity}, there is an equivalence
		\[
			\End_{A_{m+2}}\bigl((B/A)_{m+1}\bigr) \simeq \Mod_{(B \otimes_A B/A)_m}(A_{m+2}),
		\]
		under which $\epsilon$ corresponds to the relative map of the codiagonal $\nabla_f\colon B \otimes_A B \to B$ at level $m$. Since $\nabla_f$ is $0$-prim, this map admits a right adjoint.

		Hence \Cref{prop:General_Ambidexterity}, applied in $A_{m+2}$, shows that $(f/A)_{m,*}$ is also a left adjoint $(f/A)_{m,\sharp}$ of $(f/A)^*_m$, and provides a canonical identification $(f/A)_{m,\sharp} \simeq (f/A)_{m,*}$ for $m \geq 1$. By \Cref{lem:Linear_Adjoint_Preserves_Dualizable} this common functor preserves dualizable objects; in particular $(B/A)^!_m = (f/A)_{m,\sharp}(\unit)$ is dualizable for $m \geq 1$. As recorded after \Cref{def:N_Suave}, the second suave condition is automatic once $(B/A)^!_m$ is dualizable. Thus $f$ is $1$-suave.

		(2) If $f$ is $0$-proper, then every iterated codiagonal $\xi$ of $f$ has the property that $\xi$ and $\nabla_{\xi}$ are again $0$-prim, since the iterated codiagonals of $\xi$ are among those of $f$. By part (1), each such $\xi$ is $1$-suave, so $f$ is $1$-\'etale.

		(3) Suppose $f$ and $\nabla_f$ are $0$-suave. For $m \geq 1$, the map $(f/A)^*_m$ admits a left adjoint $(f/A)_{m,\sharp}$ in $A_{m+2}$, as $f$ is $0$-suave. We apply the dual version of \Cref{prop:General_Ambidexterity} (cf.\ the remark following it), and so must check that the unit and counit of $(f/A)_{m,\sharp} \dashv (f/A)^*_m$ admit right adjoints. The counit corresponds, via the equivalence $\End_{A_{m+2}}(\unit) \iso A_{m+1}$, to the map $(f/A)_{m-1,!}$, which admits a right adjoint by the second condition in the definition of $0$-suave for $f$. The unit corresponds, under the relation-object identification used in part (1), to the analogous datum for the codiagonal $\nabla_f$, which admits a right adjoint since $\nabla_f$ is $0$-suave. Hence the dual ambidexterity result shows that $(f/A)_{m,\sharp}$ is also a right adjoint of $(f/A)^*_m$ for $m \geq 1$, so $f$ is $1$-prim.

		(4) If $f$ is $0$-\'etale, then every iterated codiagonal $\xi$ of $f$ has $\xi$ and $\nabla_{\xi}$ both $0$-suave, so $\xi$ is $1$-prim by part (3); hence $f$ is $1$-proper.

		Finally, the argument above produces the ambidextrous identification for $(f/A)^*_m$ with $m \geq 1$ in the shifted ($n = 0$) picture; unshifting, the canonical identification of \Cref{prop:General_Ambidexterity} reads $(f/A)_{m,\sharp} \simeq (f/A)_{m,*}$ for $m \geq n+1$ in cases (1) and (3), and the same identification holds in cases (2) and (4), applied to $f$ itself. By \Cref{lem:Linear_Adjoint_Preserves_Dualizable} this common functor preserves dualizable objects.
	\end{proof}

	The following lemma was used above; in particular, being both $n$-prim and $n$-suave forces $f_{m,*}$ and $f_{m,\sharp}$ to preserve dualizable objects.

	\begin{lemma}\label[lemma]{lem:Linear_Adjoint_Preserves_Dualizable}
		Let $\varphi\colon \Cc \to \Dd$ be a morphism in $\CAlg(\PrL)$, and suppose $\varphi$ admits both a $\Cc$-linear left adjoint $\varphi_{\sharp}$ and a $\Cc$-linear right adjoint $\varphi_{*}$. Then both $\varphi_{\sharp}$ and $\varphi_{*}$ preserve dualizable objects.
	\end{lemma}
	\begin{proof}
		Let $N \in \Dd$ be dualizable, with dual $N^{\vee}$. Since $\varphi$ is symmetric monoidal, it preserves dualizable objects, and $\Cc$-linearity of the adjoints yields the projection formulas
		\[
			\varphi_{\sharp}\bigl(Y \otimes \varphi(M)\bigr) \simeq \varphi_{\sharp}(Y) \otimes M, \qquad \varphi_{*}\bigl(Y \otimes \varphi(M)\bigr) \simeq \varphi_{*}(Y) \otimes M
		\]
		for $Y \in \Dd$ and $M \in \Cc$. Using the adjunction $\varphi_{\sharp} \dashv \varphi \dashv \varphi_{*}$ together with these formulas, we compute, naturally in $M \in \Cc$,
		\[
			\Hom_{\Cc}(\varphi_{\sharp} N, M)
			\simeq \varphi_{*}\Hom_{\Dd}(N, \varphi(M))
			\simeq \varphi_{*}\bigl(N^{\vee} \otimes \varphi(M)\bigr)
			\simeq \varphi_{*}(N^{\vee}) \otimes M.
		\]
		Setting $M = \unit$ gives $\Hom_{\Cc}(\varphi_{\sharp} N, \unit) \simeq \varphi_{*}(N^{\vee})$, so the displayed natural equivalence becomes $\Hom_{\Cc}(\varphi_{\sharp} N, M) \simeq \Hom_{\Cc}(\varphi_{\sharp} N, \unit) \otimes M$. Hence $\varphi_{\sharp} N$ is dualizable, with dual $\varphi_{*}(N^{\vee})$. The symmetric computation, exchanging the roles of the two adjoints, shows that $\varphi_{*}$ likewise preserves dualizable objects.
	\end{proof}

	\begin{corollary}
		Let $n \geq 0$ and assume that $f\colon A \to B$ is a map of Stefanich rings which is both $n$-suave and $n$-prim. Then the suave dualizing object $(B/A)^!_m$ is dualizable in $A_{m+1}$ for all $m \geq n$, with dual $(B/A)_m$.
	\end{corollary}
	\begin{proof}
		We have $(B/A)^!_m = f_{m,\sharp}(\unit_{A_{m+1}})$, which is dualizable by \Cref{lem:Linear_Adjoint_Preserves_Dualizable}. Its dual is $f_{m,*}(\unit_{A_{m+1}}) = (B/A)_m$.
	\end{proof}

	\chapter{Gestalten}
	\label[chapter]{chap:Gestalten}

	\emph{Talk given by Giovanni.}

	In the previous talk we introduced suave and étale maps of Stefanich rings. The goal of this talk is to assemble the $n$-étale maps over a fixed Stefanich ring $A$ into a category and to show that its opposite is a \emph{countably toposic category}, in the sense made precise below. Passing to opposites throughout, this exhibits the category of \emph{Gestalten}, the opposite of $\StRing$, as a setting in which the étale maps play the role of the open immersions in a topos.

	Throughout we freely use the relative objects $(B/A)_m$ and $(B/A)^!_m$ attached to a map $f\colon A \to B$ of Stefanich rings, together with the relative maps $(f/A)^*_m$, $(f/A)_{m,\sharp}$, $(f/A)_{m,!}$, etc., from \Cref{chap:Properties_Of_Maps}.

	\section{Countable limits of étale algebras}
	\label[section]{sec:Countable_Limits_Etale_Algebras}

	Recall from \Cref{def:N_Suave,def:N_Etale} that, for $n \geq 0$, a map $f\colon A \to B$ is \emph{$n$-suave} if for each $m \geq n$ the relative map $(f/A)^*_m\colon \unit = (A/A)_{m+1} \to (B/A)_{m+1}$ admits a left adjoint $(f/A)_{m,\sharp}$ in $A_{m+2}$, whose value $(B/A)^!_m := (f/A)_{m,\sharp}(\unit)$ we call the relative \emph{dualizing object}, and if moreover, for each $m \geq n+1$, the resulting counit $(f/A)_{m-1,!}\colon (B/A)^!_m \to (A/A)_m$ admits a right adjoint; and that $f$ is \emph{$n$-\'etale} if $f$ and all its iterated codiagonals are $n$-suave. By \Cref{prop:Suave_Properties,prop:Etale_Properties}, the $n$-suave and $n$-\'etale maps are stable under cobase change and composition, the dualizing objects satisfy the base-change formula
	\[
		g_m^*(B/A)^!_m \simeq (B \otimes_A C/C)^!_m \qquad (m \geq n)
	\]
	for any $g\colon A \to C$, and $n$-\'etale maps enjoy right cancellation: if $f\colon A \to B$ is $n$-\'etale and $g\colon B \to C$ is arbitrary, then $g$ is $n$-suave (resp.\ $n$-\'etale) if and only if $gf$ is.

	We now turn to the behaviour of suave algebras under \emph{limits}. The classical theory offers a cautionary example.

	\begin{example}
		Static étale algebras are not closed under countable limits. For instance, if $k$ is a field, then every finite product $k^m$ is étale over $k$, but the limit $\lim(\cdots \to k^3 \to k^2 \to k) \simeq \prod_{\N} k$ is not étale over $k$, since it is not of finite presentation (it is not even noetherian).
	\end{example}

	Étale Stefanich rings do not suffer from this defect, essentially because of ambidexterity. The key adjointability input is the following.

	\begin{lemma}\label[lemma]{lem:Adjoints_Countable_Limits}
		Let $\bbC \in 2\Pr$ be a $2$-presentable category and let $I$ be a countable category.
		\begin{enumerate}
			\item Let $X_\bullet\colon I \to \bbC$ be a diagram and let $F\colon Y \to \lim_{i \in I} X_i$ be a morphism. Write $F_i := \pr_i \circ F\colon Y \to X_i$ for its composite with the projection to $X_i$. If each $F_i$ admits a left adjoint $F_{i,\sharp}\colon X_i \to Y$, then $F$ admits a left adjoint $F_\sharp$, and there is a canonical equivalence
			\[
				F_\sharp F \simeq \colim_{i \in I} F_{i,\sharp} F_i \qin \End_{\bbC}(Y).
			\]
			\item Let $X_\bullet\colon I \to \bbC$ be a diagram and let $G\colon \colim_{i \in I} X_i \to Y$ be a morphism. Write $G_i := G \circ \iota_i\colon X_i \to Y$ for its composite with the colimit structure map. If each $G_i$ admits a right adjoint $G_{i,*}\colon Y \to X_i$, then $G$ admits a right adjoint $G_*$, and there is a canonical equivalence
			\[
				G_* G \simeq \lim_{i \in I} G_{i,*} G_i \qin \End_{\bbC}(Y).
			\]
		\end{enumerate}
	\end{lemma}
	\begin{proof}
		The two statements are dual, so we treat (1). As in the proof of \Cref{lem:Filtered_Colimit_Of_Adjoints}, the represented functors $\bbC(-,Z)\colon \bbC\catop \to 1\Pr$ jointly detect and reflect adjunctions of $1$-morphisms (and the displayed formula may be checked after applying them), so we may assume $\bbC = 1\Pr = \PrL_{\aleph_1}$. There a countable limit $\lim_{i \in I} X_i$ is computed as the corresponding limit in $\Cat$, equivalently as the limit in $\PrR$ formed along the right adjoints of the projections. A morphism $F$ into such a limit admits a left adjoint as soon as each component $F_i$ does, by assembling the partial left adjoints; this is where countability of $I$ enters, ensuring the relevant Beck--Chevalley colimits exist and are computed in $\PrL$. Tracing through the limit presentation, the composite $F_\sharp F$ is computed as the colimit $\colim_{i \in I} F_{i,\sharp} F_i$, the transition maps being the mates of the identifications $X_\alpha F_i \simeq F_j$ for $\alpha\colon i \to j$ in $I$.
	\end{proof}

	\begin{proposition}\label[proposition]{prop:Suave_Countable_Limits}
		Let $A \in \StRing$ and $n \geq 0$.
		\begin{enumerate}
			\item The $n$-suave $A$-algebras are closed under countable limits in $\StRing_{A/}$. Moreover, for $m \geq n$ the formation of the dualizing object sends countable limits of $n$-suave algebras to colimits: if $B = \lim_{i \in I} B_i$ with each $B_i$ $n$-suave over $A$, then
			\[
				(B/A)^!_m \simeq \colim_{i \in I} (B_i/A)^!_m.
			\]
			\item For any map $g\colon A \to C$, the base-change functor $- \otimes_A C\colon \StRing_{A/} \to \StRing_{C/}$ preserves countable limits of $n$-suave algebras.
		\end{enumerate}
	\end{proposition}
	\begin{proof}
		By shifting (\Cref{lem:Shift_Functor_Equivalence}) we may assume $n = 0$. Let $I$ be countable and let $B_\bullet\colon I \to \StRing_{A/}$ be a diagram of $0$-suave algebras, with limit $B := \lim_{i \in I} B_i$, which exists since $\StRing_{A/}$ is presentable. Since limits of internal Stefanich rings are computed levelwise (\Cref{prop:Slice_Description}), we have $(B/A)_m \simeq \lim_{i \in I} (B_i/A)_m$ in $\CAlg(A_{m+1})$ for every $m$.

		\emph{Part (1).} Fix $m \geq 0$. To produce the dualizing object we must show that the map
		\[
			\unit = (A/A)_{m+1} \to (B/A)_{m+1} \simeq \lim_{i \in I} (B_i/A)_{m+1}
		\]
		admits a left adjoint in $A_{m+2}$. For each $i$ the map $\unit \to (B_i/A)_{m+1}$ admits a left adjoint, since $B_i$ is $0$-suave over $A$. Applying \Cref{lem:Adjoints_Countable_Limits}(1) in $A_{m+2}$ produces the desired left adjoint, together with the identification
		\[
			(B/A)^!_m \simeq \colim_{i \in I} (B_i/A)^!_m
		\]
		in $A_{m+1}$. For the second suave condition, under this identification the counit $(B/A)^!_m \to \unit$ is the colimit over $I$ of the counits $(B_i/A)^!_m \to \unit$, each of which admits a right adjoint because $B_i$ is $0$-suave. As $(B/A)^!_m$ is the countable colimit of the $(B_i/A)^!_m$, \Cref{lem:Adjoints_Countable_Limits}(2) shows that the counit $(B/A)^!_m \to \unit$ admits a right adjoint. Hence $B$ is $0$-suave, proving (1).

		\emph{Part (2).} Fix $g\colon A \to C$. As $0$-suave maps are stable under base change (\Cref{prop:Suave_Properties}(1)), the $C$-algebras $B \otimes_A C$ and $B_i \otimes_A C$ are $0$-suave, and we must show that the canonical map
		\[
			\Bigl(\lim_{i \in I} B_i\Bigr) \otimes_A C \to \lim_{i \in I} (B_i \otimes_A C)
		\]
		is an equivalence. By the base-change formula of \Cref{prop:Suave_Properties}(1) and the fact that $g_{m+1}^*$ preserves colimits, we have for every $m \geq 0$
		\[
			(B \otimes_A C/C)^!_m
			\simeq g_{m+1}^*(B/A)^!_m
			\simeq g_{m+1}^* \colim_{i \in I} (B_i/A)^!_m
			\simeq \colim_{i \in I} (B_i \otimes_A C/C)^!_m.
		\]
		On the other hand, applying part (1) to the diagram $i \mapsto B_i \otimes_A C$ of $0$-suave $C$-algebras gives
		\[
			\Bigl(\bigl(\lim_{i \in I} (B_i \otimes_A C)\bigr)/C\Bigr)^!_m \simeq \colim_{i \in I} (B_i \otimes_A C/C)^!_m.
		\]
		Thus both $C$-algebras have the same dualizing objects in every degree. Since for suave algebras $\bigl((B/A)^!_m\bigr)^{\vee} \simeq (B/A)_m$ (see the discussion after \Cref{def:N_Suave}), they also have the same relative levels, showing that the canonical map is an equivalence.
	\end{proof}

	\section{Countably toposic categories}

	Recall the definition of a topos:

	\begin{definition}\label[definition]{def:Topos}
		A \emph{topos} is an accessible left exact localization of a presheaf category $\PSh(C)$ on a small category $C$; that is, a category $\Xx$ for which there exist a small category $C$ and an adjunction $\PSh(C) \rightleftarrows \Xx$ whose left adjoint preserves finite limits and whose right adjoint is fully faithful and accessible.
	\end{definition}

	In particular, every topos is presentable. In ordinary category theory, Giraud's theorem characterizes topoi by intrinsic exactness properties; the same holds in the $\infty$-categorical setting, the third axiom below (effectivity of groupoids) being the one that distinguishes $\infty$-topoi from $1$-topoi.

	\begin{definition}\label[definition]{def:Giraud_Axioms}
		A presentable category $\Xx$ \emph{satisfies the Giraud axioms} if:
		\begin{enumerate}
			\item[(G1)] \emph{Colimits are universal:} for every map $f\colon X \to Y$ in $\Xx$, the pullback functor $f^*\colon \Xx_{/Y} \to \Xx_{/X}$ preserves all colimits;
			\item[(G2)] \emph{Coproducts are disjoint:} for all $X, Y \in \Xx$, the pullback of the span $X \to X \sqcup Y \leftarrow Y$ is initial;
			\item[(G3)] \emph{Groupoid objects are effective:} for every groupoid object $X_\bullet\colon \simp\catop \to \Xx$ with geometric realization $\abs{X_\bullet}$, the canonical map from $X_\bullet$ to the Čech nerve of $X_0 \to \abs{X_\bullet}$ is an equivalence; equivalently, $X_n \simeq X_0 \times_{\abs{X_\bullet}} \cdots \times_{\abs{X_\bullet}} X_0$ for all $n \geq 0$.
		\end{enumerate}
	\end{definition}

	\begin{theorem}[\protect{\citet[Theorem~6.1.0.6]{lurie2009HTT}}]\label[theorem]{thm:Topos_Characterization}
		Let $\Xx$ be a presentable category. The following are equivalent:
		\begin{enumerate}
			\item $\Xx$ is a topos.
			\item For some (equivalently, any) regular cardinal $\kappa$ such that $\Xx$ is $\kappa$-compactly generated and $\Xx^{\kappa}$ is closed under finite limits, the colimit-preserving extension $\PSh(\Xx^{\kappa}) \to \Xx$ of the inclusion preserves finite limits.
			\item $\Xx$ satisfies the Giraud axioms.
		\end{enumerate}
	\end{theorem}

	We are interested in the following `countable analogue' of the definition of a topos:

	\begin{definition}\label[definition]{def:Countably_Toposic}
		A small category $\Yy$ with finite limits and countable colimits is a \emph{countably toposic category} if the presentable category $\Ind_{\aleph_1}(\Yy)$ is a topos.
	\end{definition}

	\begin{remark}
		Scholze uses ``countably topos'', and Giovanni used ``countable topos'' during his talk; we avoid both, the former being grammatically awkward and the latter misleading, since the category in question is not countable and not a topos.
	\end{remark}

	Since $\Yy$ has countable colimits, the unit $\Yy \to \Ind_{\aleph_1}(\Yy)$ identifies $\Yy$ with the full subcategory $\Ind_{\aleph_1}(\Yy)^{\aleph_1}$ of $\aleph_1$-compact objects. We deduce the following characterization.

	\begin{proposition}\label[proposition]{prop:Countably_Toposic_Characterization}
		Let $\Yy$ be a small category with finite limits and countable colimits. The following are equivalent.
		\begin{enumerate}
			\item $\Yy$ is a countably toposic category.
			\item The countable-colimit-preserving extension $\PSh(\Yy)^{\aleph_1} \to \Yy$ of the identity preserves finite limits.
			\item There exist a small category $C$ with finite limits and countable colimits and a finite-limit-preserving, countable-colimit-preserving localization $\PSh(C)^{\aleph_1} \to \Yy$.
			\item $\Yy$ satisfies the \emph{countable Giraud axioms}: countable colimits are universal, coproducts are disjoint, and groupoid objects are effective.
		\end{enumerate}
	\end{proposition}
	\begin{proof}
		$(1) \Rightarrow (2)$: If $\Yy$ is countably toposic then $\Ind_{\aleph_1}(\Yy)$ is a topos. Using $\Yy \simeq \Ind_{\aleph_1}(\Yy)^{\aleph_1}$ together with the finite limits of $\Yy$, we may apply \Cref{thm:Topos_Characterization}(2) with $\kappa = \aleph_1$: the colimit-preserving functor $\PSh(\Yy) \to \Ind_{\aleph_1}(\Yy)$ preserves finite limits. Restricting to $\aleph_1$-compact objects, which are closed under finite limits in both source and target by \Cref{prop:Kappa_Compact_Presheaves_Finite_Limits} and \Cref{thm:Topos_Characterization}(2), yields (2).

		$(2) \Rightarrow (3)$: Take $C = \Yy$. The functor $\PSh(\Yy)^{\aleph_1} \to \Yy$ preserves countable colimits and finite limits, and it is a localization: its fully faithful right adjoint is induced by the Yoneda embedding $\Yy \hookrightarrow \PSh(\Yy)^{\aleph_1}$.

		$(3) \Rightarrow (4)$: Let $L\colon \PSh(C)^{\aleph_1} \to \Yy$ be a left exact localization as in (3). First, $\PSh(C)^{\aleph_1}$ satisfies the countable Giraud axioms: indeed $\PSh(C) \simeq \Ind_{\aleph_1}(\PSh(C)^{\aleph_1})$ is a presheaf topos, hence satisfies the Giraud axioms, and these restrict to the $\aleph_1$-compact objects, which are closed under finite limits by \Cref{prop:Kappa_Compact_Presheaves_Finite_Limits}. It then suffices to know that the countable Giraud axioms are stable under left exact localizations, which is \Cref{lem:Countable_Giraud_Localization} below.

		$(4) \Rightarrow (1)$: Since $\Ind_{\aleph_1}(\Yy)$ is presentable, by \Cref{thm:Topos_Characterization} it is enough to verify the Giraud axioms for $\Ind_{\aleph_1}(\Yy)$. Now $\Ind_{\aleph_1}(\Yy)$ is generated under $\aleph_1$-filtered colimits by $\Yy$, and finite limits commute with $\aleph_1$-filtered colimits; hence universality of colimits and disjointness of coproducts propagate from $\Yy$ to $\Ind_{\aleph_1}(\Yy)$. For effectivity of groupoid objects one uses in addition that $\simp\catop$ is countable, so that the canonical functor
		\[
			\Fun(\simp\catop, \Ind_{\aleph_1}(\Yy)) \to \Ind_{\aleph_1}\bigl(\Fun(\simp\catop, \Yy)\bigr)
		\]
		is an equivalence, while both the groupoid condition and the comparison with the Čech nerve are finite-limit conditions after passing to geometric realizations.
	\end{proof}

	\begin{lemma}\label[lemma]{lem:Countable_Giraud_Localization}
		The countable Giraud axioms are stable under left exact localizations.
	\end{lemma}
	\begin{proof}
		Let $L\colon \Xx \rightleftarrows \Yy \noloc \iota$ be a left exact localization, with $\Xx$ satisfying the countable Giraud axioms. Finite limits in $\Yy$ are computed in $\Xx$ through the fully faithful right adjoint $\iota$, while colimits in $\Yy$ are obtained by applying $L$ to the corresponding colimits in $\Xx$. Universality of countable colimits and disjointness of coproducts then transport from $\Xx$ to $\Yy$ exactly as in the localization step of the proof of \Cref{thm:Topos_Characterization}, using that $L$ is left exact and colimit-preserving and that $\iota$ is left exact and fully faithful. Effectivity of groupoid objects is preserved because $L$ carries groupoid objects to groupoid objects and commutes with the formation of Čech nerves and geometric realizations, while $\iota$ reflects equivalences.
	\end{proof}

	\section{The category of étale algebras as a topos}
	\label[section]{sec:Etale_Algebras_Form_Topos}

	We now prove that $(\StRing^{n\text{-}\et}_{A/})\catop$ is a countably toposic category, where $\StRing^{n\text{-}\et}_{A/} \subseteq \StRing_{A/}$ denotes the full subcategory on the $n$-\'etale algebras. We first record that this is a legitimate small category with finite limits.

	\begin{remark}\label[remark]{rmk:Etale_Small_Finite_Colimits}
		Let $A \in \StRing$ and $n \geq 0$.
		\begin{enumerate}
			\item $\StRing^{n\text{-}\et}_{A/}$ has finite colimits: the identity $A = A$ is $n$-\'etale, giving an initial object, and $n$-\'etale maps are stable under cobase change (\Cref{prop:Etale_Properties}(1)), giving pushouts.
			\item $\StRing^{n\text{-}\et}_{A/}$ is essentially small. Indeed, if $A \to B$ is $n$-\'etale then it is $(n+1)$-proper (\Cref{prop:Ambidexterity_Stefanich_Rings}), and for $m \geq n+1$ ambidexterity identifies $(B/A)_m \simeq (B/A)^!_m$ as a dualizable object of $A_{m+1}$ (\Cref{prop:Ambidexterity_Stefanich_Rings} and \Cref{lem:Linear_Adjoint_Preserves_Dualizable}). Since the unit of $A_{m+1}$ is $\aleph_1$-compact, the dualizable objects are $\aleph_1$-compact and hence form an essentially small subcategory (cf.\ \Cref{chap:Compact_Objects}); as a Stefanich algebra under $A$ is determined by the sequence $((B/A)_m)_{m \geq n+1}$ via the identification $\StRing_{A/} \simeq \lim(\CAlg(A_1) \to \CAlg(A_2) \to \cdots)$, the $n$-\'etale $A$-algebras form an essentially small category.
		\end{enumerate}
		Consequently $C := (\StRing^{n\text{-}\et}_{A/})\catop$ is a small category with finite limits.
	\end{remark}

	The key technical input is the following. Let $C$ be a small category with finite limits, and suppose given an $\aleph_1$-accessible left exact localization $L\colon \PSh(C) \rightleftarrows \Shv(C) \noloc \iota$. (The notation is suggestive, but we do \emph{not} assume that $\Shv(C)$ is the category of sheaves for a Grothendieck topology.) Assume that $C$ is \emph{subcanonical}, meaning that the composite $C \xrightarrow{y} \PSh(C) \xrightarrow{L} \Shv(C)$ is fully faithful. Fix $A \in \StRing$ and let
	\[
		\Oo\colon C\catop \to \StRing^{n\text{-}\et}_{A/}
	\]
	be a finite-colimit-preserving functor which is a \emph{sheaf}, in the sense that its right Kan extension to $\PSh(C)$ inverts $L$-equivalences. Equivalently, the right Kan extension of $\Oo$ along $y\catop$ factors through a limit-preserving functor
	\[
		\Oo\colon \Shv(C)\catop \to \StRing_{A/}, \qquad \Oo(F) \simeq \lim_{(y(X) \to F) \in C_{/F}} \Oo(X),
	\]
	which exists since $\StRing_{A/}$ is presentable.

	\begin{theorem}\label[theorem]{thm:Sheaf_Valued_In_Etale}
		In the situation above, the functor $\Oo\colon (\Shv(C)^{\aleph_1})\catop \to \StRing_{A/}$ factors through $\StRing^{n\text{-}\et}_{A/}$, and it preserves finite colimits.
	\end{theorem}

	The proof occupies the rest of this section; we first deduce the main result.

	\begin{theorem}\label[theorem]{thm:Etale_Countably_Toposic}
		Let $A \in \StRing$ and $n \geq 0$. Then $(\StRing^{n\text{-}\et}_{A/})\catop$ is a countably toposic category.
	\end{theorem}
	\begin{proof}
		By \Cref{rmk:Etale_Small_Finite_Colimits}, $C := (\StRing^{n\text{-}\et}_{A/})\catop$ is a small category with finite limits. Equip it with the trivial topology, so that $\PSh(C) = \Shv(C)$ and $C$ is trivially subcanonical, and apply \Cref{thm:Sheaf_Valued_In_Etale} to the identity sheaf $\Oo = \id\colon C\catop = \StRing^{n\text{-}\et}_{A/} \to \StRing^{n\text{-}\et}_{A/}$, which preserves finite colimits. We obtain a finite-colimit-preserving functor
		\[
			(\PSh(C)^{\aleph_1})\catop \to \StRing^{n\text{-}\et}_{A/}.
		\]
		Passing to opposite categories gives a functor
		\[
			\PSh(C)^{\aleph_1} \to (\StRing^{n\text{-}\et}_{A/})\catop = C
		\]
		which preserves finite limits and countable colimits and is a localization with fully faithful right adjoint the Yoneda embedding. By the equivalence $(2) \Leftrightarrow (1)$ of \Cref{prop:Countably_Toposic_Characterization}, $C$ is a countably toposic category.
	\end{proof}

	\begin{corollary}\label[corollary]{cor:Etale_Countable_Limits}
		The full subcategory $\StRing^{n\text{-}\et}_{A/} \subseteq \StRing_{A/}$ is closed under countable limits.
	\end{corollary}
	\begin{proof}
		By \Cref{thm:Etale_Countably_Toposic} the category $(\StRing^{n\text{-}\et}_{A/})\catop$ is countably toposic, hence in particular has countable colimits; equivalently, $\StRing^{n\text{-}\et}_{A/}$ has countable limits. That these agree with the limits computed in $\StRing_{A/}$ follows from \Cref{prop:Suave_Countable_Limits}(1): a countable limit of $n$-suave algebras formed in $\StRing_{A/}$ is again $n$-suave. The iterated codiagonals of such a limit are computed as countable limits of the corresponding codiagonals, using that base change preserves countable limits of suave algebras (\Cref{prop:Suave_Countable_Limits}(2)), and these are again $n$-suave, so the limit is $n$-\'etale.
	\end{proof}

	\begin{proof}[Proof of \Cref{thm:Sheaf_Valued_In_Etale}]
		The ``preserves finite colimits'' assertion is immediate once the factorization is established: by construction $\Oo$ sends finite limits in $\Shv(C)^{\aleph_1}$ to finite colimits in $\StRing_{A/}$. So we must show that for every $\aleph_1$-compact $X \in \Shv(C)^{\aleph_1}$ the algebra $\Oo(X)$ is $n$-\'etale over $A$. Our strategy is to show:
		\begin{itemize}
			\item for every map $f\colon Y \to X$ in $\Shv(C)^{\aleph_1}$ the map $\Oo(f)\colon \Oo(X) \to \Oo(Y)$ is $n$-suave (Claim 1);
			\item $\Oo$ satisfies a Künneth formula (Claim 2).
		\end{itemize}
		Granting both, the iterated codiagonals of $\Oo(f)$ are again of the form $\Oo(f')$: by Künneth they correspond to the diagonals $Y \to Y \times_X Y$, $Y \to Y \times_{(Y \times_X Y)} Y$, \dots, hence are $n$-suave by Claim 1. Thus $\Oo(f)$ is $n$-\'etale; taking $X$ to be the terminal object of $\Shv(C)^{\aleph_1}$, for which $\Oo(X) = \Oo(\ast) = A$, shows that $\Oo(Y)$ is $n$-\'etale over $A$.

		\textbf{Claim 1.} \emph{For every map $f\colon Y \to X$ in $\Shv(C)^{\aleph_1}$, the map $\Oo(f)\colon \Oo(X) \to \Oo(Y)$ is $n$-suave.}

		Suppose first that $X \in C$ is representable. Since $Y$ is $\aleph_1$-compact, in the slice $\Shv(C)^{\aleph_1}_{/X}$ it may be written as a countable colimit $Y \simeq \colim_{i \in I} Y_i$ of representables $Y_i \in C_{/X}$. As $\Oo$ sends countable colimits in $\Shv(C)$ to countable limits in $\StRing_{A/}$,
		\[
			\Oo(Y) \simeq \lim_{i \in I} \Oo(Y_i) \qin \StRing_{\Oo(X)/}.
		\]
		Each map $\Oo(X) \to \Oo(Y_i)$ arises from the map $Y_i \to X$ in $C$, hence is $n$-\'etale and in particular $n$-suave. By \Cref{prop:Suave_Countable_Limits}(1) the countable limit $\Oo(X) \to \Oo(Y)$ is again $n$-suave, with
		\[
			(\Oo(Y)/\Oo(X))^!_m \simeq \colim_{i \in I} (\Oo(Y_i)/\Oo(X))^!_m \qquad (m \geq n).
		\]
		Moreover, \Cref{prop:Suave_Countable_Limits}(2) shows that for a map $g\colon X' \to X$ of representables base change is compatible:
		\[
			\Oo(X') \otimes_{\Oo(X)} \Oo(Y) \simeq \Oo(X' \times_X Y), \qquad
			\Oo(X') \otimes_{\Oo(X)} (\Oo(Y)/\Oo(X))^!_m \simeq (\Oo(X' \times_X Y)/\Oo(X'))^!_m.
		\]

		Now let $X \in \Shv(C)^{\aleph_1}$ be arbitrary. Here we cannot apply \Cref{prop:Suave_Countable_Limits} directly, since the base need not be representable; we argue by descent. For each map $g\colon X' \to X$ with $X' \in C$ representable, form the pullback $X' \times_X Y$, which is again $\aleph_1$-compact, since $\aleph_1$-compact sheaves are closed under finite limits (\Cref{cor:Kappa_Compact_Sheaves_Finite_Limits}). By the representable case the map $\Oo(X') \to \Oo(X' \times_X Y)$ is $n$-suave, with dualizing objects $(\Oo(X' \times_X Y)/\Oo(X'))^!_m \in \Oo(X')_{m+1}$ for $m \geq n$. For a further representable $X'' \to X' \to X$, the base-change compatibility above gives
		\[
			\Oo(X'') \otimes_{\Oo(X')} (\Oo(X' \times_X Y)/\Oo(X'))^!_m \simeq (\Oo(X'' \times_X Y)/\Oo(X''))^!_m,
		\]
		so the assignment $X' \mapsto (\Oo(X' \times_X Y)/\Oo(X'))^!_m$ is a descent datum for the sheaf $\Oo(-)_{m+1}$. As $\Oo$ is a sheaf, it descends to a cocommutative coalgebra object $B^!_m \in \Oo(X)_{m+1}$ with $g^*B^!_m \simeq (\Oo(X' \times_X Y)/\Oo(X'))^!_m$ for every representable $g\colon X' \to X$. The trick we use is to relate properties of $B^!_m$ to properties of the map $\Oo(X) \to \Oo(Y)$.

		For $m \geq n+1$, the counit $(\Oo(X' \times_X Y)/\Oo(X'))^!_m \to \unit$ admits a right adjoint, since $\Oo(X') \to \Oo(X' \times_X Y)$ is $n$-suave. These right adjoints are compatible with base change, so the counit $h_{m-1,!}\colon B^!_m \to \unit$ admits a right adjoint in $\Oo(X)_{m+1}$. It follows that the induced map $\unit \to (B^!_m)^{\vee}$ on duals admits a \emph{left} adjoint. This dual may be computed as follows:
		\begin{align*}
			(B^!_m)^{\vee}
			&\simeq \lim_{g\colon X' \to X} g_{m,*}\bigl((\Oo(X' \times_X Y)/\Oo(X'))^!_m\bigr)^{\vee} \\
			& \simeq \lim_{g\colon X' \to X} g_{m,*}(\Oo(X' \times_X Y)/\Oo(X'))_m \\
			& \simeq (\Oo(Y)/\Oo(X))_m,
		\end{align*}
		where we used $\bigl((B/A)^!_m\bigr)^{\vee} \simeq (B/A)_m$ for suave algebras and, in the last step, that $\Oo$ is a sheaf (so that $\Oo(Y) \simeq \lim_{g\colon X' \to X} \Oo(X' \times_X Y)$). This in particular shows that the map $(f/\Oo(X))^*_{m-1}\colon \unit \to (B^!_m)^{\vee} = (\Oo(Y)/\Oo(X))_m$ admits a left adjoint
		\[
			(f/\Oo(X))_{m-1,\sharp}\colon (\Oo(Y)/\Oo(X))_m \to (\Oo(X)/\Oo(X))_m.
		\]
		We claim that the value at the unit is $B^!_{m-1}$, showing that $B^!_m \simeq (\Oo(Y)/\Oo(X))^!_m$. This concludes the proof, as we already know that $B^!_m \to \unit$ admits a right adjoint.

		To see this, recall from \Cref{lem:Suave_Dualizing_Object_Via_Adjunction} that we may alternatively compute it as the image of the unit under the composite
		\[
			\unit \xrightarrow{h^!_{m-1}} B^!_m \xrightarrow{h_{m-1,!}} \unit.
		\]
		By descent, this is
		\[
			h^!_{m-1}h_{m-1,!}\unit \simeq \lim_{g\colon X' \to X} g_{m-1}^!g_{m-1,!}\unit \simeq B^!_{m-1}.
		\]
		We conclude that $\Oo(X) \to \Oo(Y)$ is $n$-suave.

		\textbf{Claim 2 (Künneth formula).} \emph{For every pullback square}
		\[
			\begin{tikzcd}
				Y' \rar \dar & Y \dar \\
				X' \rar & X
			\end{tikzcd}
		\]
		\emph{in $\Shv(C)^{\aleph_1}$, the natural map $\Oo(X') \otimes_{\Oo(X)} \Oo(Y) \to \Oo(Y')$ is an equivalence in $\StRing_{A/}$.}

		Set $Y' = X' \times_X Y$. Suppose first that $X' = S \in C$ is representable, with structure map $g\colon S \to X$. By Claim 1 and the construction of $B^!_m$, the dualizing objects of $\Oo(X) \to \Oo(Y)$ pull back along $g$ to $\Oo(S) \otimes_{\Oo(X)} (\Oo(Y)/\Oo(X))^!_m \simeq (\Oo(S \times_X Y)/\Oo(S))^!_m$. Since base change of an $n$-suave map is computed on dualizing objects (\Cref{prop:Suave_Properties}(1)), it follows that $\Oo(S) \otimes_{\Oo(X)} \Oo(Y) \simeq \Oo(S \times_X Y)$.

		For general $X' \to X$, consider $\varphi\colon \Oo(X') \otimes_{\Oo(X)} \Oo(Y) \to \Oo(Y')$. We check that $\varphi$ is an equivalence after base change along every representable $S \to X'$. The source becomes
		\[
			\Oo(S) \otimes_{\Oo(X')}\bigl(\Oo(X') \otimes_{\Oo(X)} \Oo(Y)\bigr)
			\simeq \Oo(S) \otimes_{\Oo(X)} \Oo(Y)
			\simeq \Oo(S \times_X Y),
		\]
		while the target becomes
		\[
			\Oo(S) \otimes_{\Oo(X')} \Oo(Y')
			\simeq \Oo(S \times_{X'} Y')
			\simeq \Oo(S \times_X Y),
		\]
		both by the representable case (the second applied to $S \to X'$ and $Y' \to X'$); under these identifications $\varphi$ becomes the identity of $\Oo(S \times_X Y)$. Since $\Oo$ is a sheaf, equivalences of $\Oo(X')$-algebras may be detected after base change to representables over $X'$, so $\varphi$ is an equivalence. This proves the Künneth formula, and with it the theorem.
	\end{proof}

	\section{Gestalten}

	\begin{definition}\label[definition]{def:Gestalten}
		The category of \emph{Gestalten} is the opposite category $\Gest := \StRing\catop$. For $A \in \StRing$ we write $\Gest(A) \in \Gest$ for the corresponding Gestalt.
	\end{definition}

	Dualizing the properties of maps of Stefanich rings introduced in \Cref{chap:Properties_Of_Maps}, we obtain notions of \emph{$n$-affine}, \emph{$n$-prim}, \emph{$n$-proper}, \emph{$n$-suave}, and \emph{$n$-\'etale} map of Gestalten. In this language, \Cref{thm:Etale_Countably_Toposic} asserts that for every $A \in \StRing$ and $n \geq 0$ the category of $n$-\'etale Gestalten over $\Gest(A)$ is countably toposic.

	\begin{corollary}\label[corollary]{cor:Etale_Descent_Properties}
		Let $f\colon X \to Y$ be a map in $\Gest$, let $Y' \to Y$ be an effective epimorphism, and let $f'\colon X' = X \times_Y Y' \to Y'$ be the base change. Then $f$ has property $\Pp_n$ if and only if $f'$ does, where $\Pp_n$ is any of $n$-affine, $n$-prim, $n$-proper, $n$-suave, $n$-\'etale.
	\end{corollary}
	\begin{proof}[Proof sketch]
		The ``only if'' direction is stability of each property under base change, established in \Cref{chap:Properties_Of_Maps}. The ``if'' direction is descent along the effective epimorphism $Y' \to Y$: the defining adjointability conditions for each $\Pp_n$ may be checked after pulling back along $Y' \to Y$, since $Y' \to Y$ is an effective epimorphism and the relevant relative objects satisfy descent, as in the descent argument of \Cref{thm:Sheaf_Valued_In_Etale}.
	\end{proof}

	\chapter{Injections and surjections of Gestalten}
	\label[chapter]{chap:Injections_Surjections_Gestalten}

	\emph{Talk given by Shai.}

	The goal of this chapter is to understand injective and surjective maps of Gestalten, and to relate surjectivity to Mathew's notion of \emph{descendability} \citep{Mathew_Galois,Mathew2018DescentNilpotence}.

	Throughout, for $X \in \Gest$ we write $\Oo(X)_n \in \CAlg(n\Pr)$ for the $n$-th level of the corresponding Stefanich ring, and for a map $f\colon Y \to X$ of Gestalten we write $(\Oo(Y)/\Oo(X))_m$ and $(\Oo(Y)/\Oo(X))^!_m$, with transition maps $f^*_m, f_{m,*}, f_{m,\sharp}, f_{m,!}$, for the relative objects and maps of \Cref{chap:Properties_Of_Maps}, dualized along $\Gest = \StRing\catop$. Dualizing the fully faithful embedding $\CAlg(A_{n+1}) \hookrightarrow \StRing_{A/}$ of \Cref{def:N_Affine}, whose essential image is the $n$-affine objects, we write $A \mapsto \Gest_X(A)$ for the resulting equivalence
	\[
		\CAlg(\Oo(X)_{n+1}) \simeq \bigl(\Gest_X^{n\text{-affine}}\bigr)\catop,
	\]
	so that $\Gest_X(A) \to X$ is the $n$-affine Gestalt over $X$ corresponding to $A \in \CAlg(\Oo(X)_{n+1})$.

	\section{Injections of Gestalten}

	\begin{definition}\label[definition]{def:Injective_Gestalten}
		A map $f\colon Y \to X$ of Gestalten is \emph{injective} if its diagonal $\Delta_f\colon Y \to Y \times_X Y$ is an isomorphism. Equivalently, the Čech nerve of $f$ is constant.
	\end{definition}

	\subsection{Closed immersions and idempotent algebras}

	We first relate closed immersions to idempotent \emph{algebras}, via the relative algebras $(B/A)_m$ of an $n$-prim map (\Cref{def:N_Prim,def:N_Proper}).

	\begin{proposition}\label[proposition]{prop:Injective_Prim_Idempotent_Algebras}
		Let $f\colon Y \to X$ be an injective map of Gestalten. Then $f$ is $n$-prim if and only if it is $n$-proper. In this case $f$ is $n$-affine and the algebra $(\Oo(Y)/\Oo(X))_n$ is idempotent in $\Oo(X)_{n+1}$. This assignment defines an equivalence of posets
		\[
			\bigl\{\text{$n$-proper injections $Z \hookrightarrow X$}\bigr\} \;\longleftrightarrow\; \bigl\{\text{idempotent algebras in $\Oo(X)_{n+1}$}\bigr\}.
		\]
	\end{proposition}
	\begin{proof}
		As the diagonal of $f$ is an isomorphism, so are all its iterated diagonals, and isomorphisms are automatically $n$-prim; hence $n$-prim is equivalent to $n$-proper for injective maps. By \Cref{lem:Kunneth_Prim} and injectivity, the diagonal of $f$ induces an isomorphism
		\[
			(\Oo(Y)/\Oo(X))_n \otimes_{\Oo(X)_{n+1}} (\Oo(Y)/\Oo(X))_n \simeq (\Oo(Y \times_X Y)/\Oo(X))_n \iso (\Oo(Y)/\Oo(X))_n,
		\]
		exhibiting $(\Oo(Y)/\Oo(X))_n$ as an idempotent algebra. In particular, we see that the functor
		\[
			(\Oo(-)/\Oo(X))_n\colon \Gest_X\catop \to \CAlg(\Oo(X)_{n+1})
		\]
		carries $n$-proper injections to idempotent algebras. In the other direction, recall from \Cref{def:N_Affine} that this functor admits a fully faithful, symmetric monoidal left adjoint, which we will denote by
		\[
			\Gest_X(-)\colon \CAlg(\Oo(X)_{n+1}) \rightleftarrows \Gest_X\catop.
		\]
		Given an idempotent algebra $A$ in $\Oo(X)_{n+1}$, we show that the map $\Gest_X(A) \to X$ is an $n$-proper injection. Since $\Gest_X(-)$ is symmetric monoidal, it sends the codiagonal $A \otimes_{\Oo(X)_{n+1}} A \iso A$ of $A$ to the diagonal $\Gest_X(A) \to \Gest_X(A) \times_X \Gest_X(A)$, and so since $A$ is idempotent this is an isomorphism. This shows that $\Gest_X(A) \to X$ is an injection. Writing $A_m := (\Oo(\Gest_X(A))/\Oo(X))_m$ for $m \geq n$, each $A_m$ is idempotent (as it is the image of $A$ under the inclusion $\CAlg(\Oo(X)_{n+1}) \hookrightarrow \CAlg(\Oo(X)_{m+1})$) and, by affineness, $\Oo(\Gest_X(A))_{m+1} = \Mod_{A_m}(\Oo(X)_{m+1})$. This gives an internal adjunction in $\Mod_{\Oo(X)_{m+1}}(1\Pr)$
		\[
			A_m \otimes -\colon \Oo(X)_{m+1} \to \Mod_{A_m}(\Oo(X)_{m+1}),
		\]
		whose right adjoint is fully faithful since $A_m$ is idempotent. Applying $\Mod_{\Oo(X)_{m+1}}(1\Pr) \to \Oo(X)_{m+2}$, we obtain an internal adjunction in $\Oo(X)_{m+2}$ of the form
		\[
			\unit \rightleftarrows (\Oo(\Gest_X(A))/\Oo(X))_{m+1}.
		\]
		As $m \geq n$ was arbitrary, $\Gest_X(A) \to X$ is an $n$-prim injection.
		
		We conclude that there is an adjunction of posets
		\[
			(\Oo(-)/\Oo(X))_n\colon \bigl\{\text{$n$-proper injections $Z \hookrightarrow X$}\bigr\} \;\leftrightarrows\; \bigl\{\text{idempotent algebras in $\Oo(X)_{n+1}$}\bigr\} \noloc \; \Gest_X(-).
		\]
		Since $\Gest_X(-)$ is fully faithful, the unit $A \iso (\Oo(\Gest_X(A))/\Oo(X))_n$ is an isomorphism. For the counit, let $Z \hookrightarrow X$ be an $n$-proper injection; we must show its $n$-affinization map $Z \to \Gest_X\bigl((\Oo(Z)/\Oo(X))_n\bigr)$ is an isomorphism. By cancellation properties, this map is again an $n$-proper injection, now with trivial relative level-$n$ algebra. It thus suffices to show that any $n$-proper injection $Z \to X$ satisfying $\unit \iso (\Oo(Z)/\Oo(X))_n \simeq \unit$ is an isomorphism.

		So let $i\colon Z \hookrightarrow X$ be an $n$-proper injection with $\unit \iso (\Oo(Z)/\Oo(X))_n$. We prove by induction on $m \geq n$ that $\unit \iso (\Oo(Z)/\Oo(X))_m$; the case $m = n$ is the hypothesis. Assume the claim for $m - 1$. By assumption on $i$, the map $i_m^*\colon \unit_X \to (\Oo(Z)/\Oo(X))_m$ has an internal right adjoint $i_{m,*}$ in $\Oo(X)_{m+1}$, and by the inductive hypothesis the unit of this adjunction, $\unit \to (\Oo(Z)/\Oo(X))_{m-1}$, is an isomorphism in $\End_{\Oo(X)_{m+1}}(\unit) = \Oo(X)_m$. Hence $i_m^*$ is a colocalization with $i_{m,*} i_m^* \simeq \id$, so $\unit$ is a retract of $(\Oo(Z)/\Oo(X))_m$; consequently the map $\unit \to (\Oo(Z)/\Oo(X))_m$ is a retract of the isomorphism
		\[
			(\Oo(Z)/\Oo(X))_m \iso (\Oo(Z)/\Oo(X))_m \otimes (\Oo(Z)/\Oo(X))_m,
		\]
		hence is itself an isomorphism. This completes the induction and the proof.
	\end{proof}

	\begin{definition}\label[definition]{def:Closed_Immersion}
		A map of Gestalten is a \emph{closed immersion} if it is a $0$-prim injection. By \Cref{prop:Injective_Prim_Idempotent_Algebras}, closed immersions $Z \hookrightarrow X$ are classified by idempotent algebras in $\Oo(X)_1$.
	\end{definition}

	\begin{example}\label[example]{ex:Closed_Immersion_Zp}
		The idempotent algebra $\Z[\tfrac{1}{p}] \in \D(\Z)$ corresponds to a \emph{closed} immersion $\Gest(\Z[\tfrac{1}{p}]) \hookrightarrow \Gest(\Z)$.
	\end{example}

	\subsection{Open immersions and idempotent coalgebras}

	We now dualize the previous subsection, relating open immersions to idempotent \emph{coalgebras}. Recall from \Cref{def:N_Suave,def:N_Etale} that an $n$-suave map $f\colon A \to B$ comes with dualizing objects $(B/A)^!_m := f_{m,\sharp}(\unit)$, built from the left adjoints $f_{m,\sharp}$ of its structure maps, and that $f$ is \emph{$n$-étale} if all its codiagonals are $n$-suave. Since $f_{m,\sharp}$ is left adjoint to a symmetric monoidal functor $f^*_m$, it is automatically oplax symmetric monoidal. Consequently, $(B/A)^!_m$ is a cocommutative coalgebra in $A_{m+1}$, whose (weak) dual is the algebra $(B/A)_m$; the counit $f_{m-1,!}\colon (B/A)^!_m \to \unit$ similarly dualizes to the unit $f^*_{m-1}$ of $(B/A)_m$.

	\begin{lemma}\label[lemma]{lem:Smashing_Idempotent_Coalgebras}
		Let $\Vv \in \CAlg(1\Pr)$, and let $D$ be an idempotent cocommutative coalgebra in $\Vv$. Then the inclusion
		\[
			\iota_D\colon \coMod_D(\Vv) \hookrightarrow \Vv
		\]
		admits the symmetric monoidal right adjoint $-\otimes D$, and exhibits $\coMod_D(\Vv)$ as a smashing colocalization of $\Vv$. Moreover, for idempotent coalgebras $D,D'$ in $\Vv$ there is a canonical equivalence
		\[
			\coMod_D(\Vv) \otimes_\Vv \coMod_{D'}(\Vv) \simeq \coMod_{D \otimes_\Vv D'}(\Vv).
		\]
		Finally, regarded as an object of $\Mod_\Vv(1\Pr)$, the category $W_D := \coMod_D(\Vv)$ is simultaneously an idempotent algebra and an idempotent coalgebra, with multiplication and comultiplication mutually inverse. Consequently, if $\Ww \in \CAlg(1\Pr)$ and $F\colon \Mod_\Vv(1\Pr) \to \Ww$ is any symmetric monoidal functor, then
		\[
			\Mod_{F(W_D)}(\Ww) \simeq \coMod_{F(W_D)}(\Ww),
		\]
		both sides identifying with the full subcategory of $F(W_D)$-local objects.
	\end{lemma}
	\begin{proof}
		Since $D \to D \otimes_\Vv D$ is an equivalence, the comonad $T_D := -\otimes D$ on $\Vv$ is idempotent. Thus its comodules are precisely the $T_D$-colocal objects, and the inclusion of this full subcategory admits $T_D$ as right adjoint. The localized tensor product has unit $D$, and the idempotence of $D$ makes $T_D$ symmetric monoidal as a functor $\Vv \to \coMod_D(\Vv)$.

		The tensor formula follows from the same description: both sides are smashing colocalizations of $\Vv$, and the right adjoints of their inclusions into $\Vv$ are canonically identified with $-\otimes D \otimes_\Vv D'$.

		Under this equivalence with $D'=D$, the multiplication of $W_D$ is induced by the inverse of the comultiplication $D \iso D \otimes_\Vv D$. The colocalization adjunction also regards $W_D$ as an idempotent coalgebra in $\Mod_\Vv(1\Pr)$, with counit $\iota_D$. These algebra and coalgebra structures have inverse multiplication and comultiplication. This property is preserved by any symmetric monoidal functor $F$. For an object whose multiplication and comultiplication are inverse equivalences, module structures and comodule structures are the same datum: both say that the canonical map to tensoring with that object is an equivalence, with the action and coaction inverse to each other.
	\end{proof}

	\begin{proposition}\label[proposition]{prop:Suave_Etale_Idempotent_Coalgebras}
		Let $f\colon Y \to X$ be an injective map of Gestalten. Then $f$ is $n$-suave if and only if it is $n$-étale. In this case $(\Oo(Y)/\Oo(X))^!_n$ is an idempotent coalgebra in $\Oo(X)_{n+1}$, and this assignment defines an equivalence of posets
		\[
			\bigl\{\text{$n$-étale injections $U \hookrightarrow X$}\bigr\} \;\longleftrightarrow\; \bigl\{\text{idempotent coalgebras in $\Oo(X)_{n+1}$}\bigr\},
		\]
		whose inverse sends $C$ to $\Gest_X\bigl(\coMod_C(\Oo(X)_{n+1})\bigr)$.
	\end{proposition}
	\begin{proof}
		As $\Delta_f$ is an isomorphism, all iterated diagonals of $f$ are isomorphisms, and isomorphisms are $n$-suave; so $n$-suave is equivalent to $n$-étale for injective maps.

		Assume $f$ is $n$-suave. By injectivity, we have $Y \iso Y \times_X Y$, and so by the Künneth formula for dualizing objects from \Cref{lem:Kunneth_Suave} we get an isomorphism
		\[
			(\Oo(Y)/\Oo(X))^!_n \otimes_{\Oo(X)_{n+1}} (\Oo(Y)/\Oo(X))^!_n \simeq (\Oo(Y \times_X Y)/\Oo(X))^!_n \iso (\Oo(Y)/\Oo(X))^!_n,
		\]
		exhibiting $(\Oo(Y)/\Oo(X))^!_n$ as an idempotent coalgebra.

		Conversely, let $C$ be an idempotent coalgebra in $\Oo(X)_{n+1}$. Its comodules form a full subcategory
		\[
			\iota\colon \coMod_C(\Oo(X)_{n+1}) \hookrightarrow \Oo(X)_{n+1},
		\]
		on which the cofree functor $- \otimes C$ is a symmetric monoidal right adjoint to $\iota$, exhibiting $\coMod_C(\Oo(X)_{n+1})$ as a smashing colocalization of $\Oo(X)_{n+1}$. Through the structure map $-\otimes C$ this is a commutative algebra in $\Oo(X)_{n+2}$, and so determines an $(n+1)$-affine Gestalt
		\[
			U := \Gest_X\bigl(\coMod_C(\Oo(X)_{n+1})\bigr).
		\]
		We claim $U \to X$ is the desired $n$-étale injection, with $(\Oo(U)/\Oo(X))^!_n = C$.

		First, $U \to X$ is an injection. Under the symmetric monoidal equivalence $\CAlg(\Oo(X)_{n+2}) \simeq \bigl(\Gest_X^{(n+1)\text{-affine}}\bigr)\catop$, its diagonal corresponds to the codiagonal of the algebra $\coMod_C(\Oo(X)_{n+1})$ in $\Oo(X)_{n+2}$. Since
		\[
			\coMod_C(\Oo(X)_{n+1}) \otimes_{\Oo(X)_{n+1}} \coMod_C(\Oo(X)_{n+1}) \simeq \coMod_{C \otimes_{\Oo(X)_{n+1}} C}(\Oo(X)_{n+1}),
		\]
		this codiagonal is the functor induced by the inverse of the comultiplication $C \iso C \otimes_{\Oo(X)_{n+1}} C$; hence the codiagonal is an isomorphism and $U \to X$ is an injection.

		It remains to show that $U \to X$ is $n$-suave; being an injection, it is then $n$-étale. The idea is to realize each level of $U$ as comodules over a canonical idempotent coalgebra, bootstrapping up from $C$.

		We shall use the smashing-idempotent package from \Cref{lem:Smashing_Idempotent_Coalgebras}. In particular, the conclusions are preserved after applying the symmetric monoidal transition functors $\Mod_{\Oo(X)_{m+1}}(1\Pr) \to \Oo(X)_{m+2}$.

		Now define idempotent coalgebras recursively by $C_n := C$ and
		\[
			C_{m+1} := \coMod_{C_m}(\Oo(X)_{m+1}) \qin \Oo(X)_{m+2},
		\]
		each one an idempotent coalgebra by \Cref{lem:Smashing_Idempotent_Coalgebras}, taken with $\Vv = \Oo(X)_{m+1}$ and $D = C_m$. We show by induction that
		\[
			\Oo(U)_{m+1} = \coMod_{C_m}(\Oo(X)_{m+1}) \qquad (m \geq n).
		\]
		The base case $m = n$ is the construction of $U$. For the induction step, $(n+1)$-affineness gives $\Oo(U)_{m+2} = \Mod_{(\Oo(U)/\Oo(X))_{m+1}}(\Oo(X)_{m+2})$, and by the inductive hypothesis the relative algebra $(\Oo(U)/\Oo(X))_{m+1}$ is $\coMod_{C_m}(\Oo(X)_{m+1}) = C_{m+1}$; \Cref{lem:Smashing_Idempotent_Coalgebras} identifies modules and comodules over $C_{m+1}$ in $\Oo(X)_{m+2}$, so
		\[
			\Oo(U)_{m+2} = \Mod_{C_{m+1}}(\Oo(X)_{m+2}) = \coMod_{C_{m+1}}(\Oo(X)_{m+2}),
		\]
		completing the induction.

		At each level $m \geq n$ the structure map $f^*_m\colon \unit \to (\Oo(U)/\Oo(X))_{m+1} = \coMod_{C_m}(\Oo(X)_{m+1})$ is the cofree functor $- \otimes C_m$, whose left adjoint in $\Oo(X)_{m+2}$ is the inclusion $\iota_m$; this is the required left adjoint $f_{m,\sharp}$, and $(\Oo(U)/\Oo(X))^!_m = \iota_m(\unit) = C_m$. The same smashing-idempotent structure provides the right adjoint to each counit $f_{m-1,!}$, so both conditions of \Cref{def:N_Suave} hold and $U \to X$ is $n$-suave. In particular $(\Oo(U)/\Oo(X))^!_n = C_n = C$, as claimed.

		It remains to recover an $n$-étale injection $U \hookrightarrow X$ from its dualizing coalgebra $C := (\Oo(U)/\Oo(X))^!_n$. The injection $U \to X$ factors through its $(n+1)$-affinization $U \to \Gest_X\bigl(\coMod_C(\Oo(X)_{n+1})\bigr)$, again an $n$-étale injection, now with trivial dualizing coalgebra; replacing $X$ by this Gestalt, it suffices to show that an $n$-étale injection $U \to X$ with $(\Oo(U)/\Oo(X))^!_n \simeq \unit_X$ is an isomorphism. We show by induction that $(\Oo(U)/\Oo(X))^!_m \simeq \unit_X$ for all $m \geq n$. Assume the claim for $m$. As $j\colon U \to X$ is $n$-suave, the map $j_{m,!}\colon (\Oo(U)/\Oo(X))^!_{m+1} \to \unit_X$ admits a right adjoint $j^!_m$ in $\Oo(X)_{m+2}$, and the composite
		\[
			\unit \xrightarrow{j^!_m} (\Oo(U)/\Oo(X))^!_{m+1} \xrightarrow{j_{m,!}} \unit_X
		\]
		is, as an object of $\End_{\Oo(X)_{m+2}}(\unit) = \Oo(X)_{m+1}$, given by $(\Oo(U)/\Oo(X))^!_m$, while the counit of the adjunction is the isomorphism $j_{m-1,!}\colon (\Oo(U)/\Oo(X))^!_m \iso \unit_X$ furnished by the inductive hypothesis. So $j_{m,!}$ is a localization with $j_{m,!} \circ j^!_m \simeq \id$, exhibiting $\unit$ as a retract of the idempotent $(\Oo(U)/\Oo(X))^!_{m+1}$. Hence the counit $j_{m,!}\colon (\Oo(U)/\Oo(X))^!_{m+1} \to \unit_X$ is a retract of the multiplication map $(\Oo(U)/\Oo(X))^!_{m+1} \otimes_{\Oo(X)_{m+2}} (\Oo(U)/\Oo(X))^!_{m+1} \to (\Oo(U)/\Oo(X))^!_{m+1}$, which is an isomorphism since $(\Oo(U)/\Oo(X))^!_{m+1}$ is idempotent. It follows that $j_{m,!}$ is itself an isomorphism, completing the induction.
	\end{proof}

	\begin{definition}\label[definition]{def:Open_Immersion}
		A map of Gestalten is an \emph{open immersion} if it is a $0$-étale injection. By \Cref{prop:Suave_Etale_Idempotent_Coalgebras}, open immersions $U \hookrightarrow X$ are classified by idempotent coalgebras in $\Oo(X)_1$.
	\end{definition}

	\section{Complements}

	\begin{definition}\label[definition]{def:Complement_Gestalten}
		Let $f\colon Y \to X$ be a map of Gestalten. Its \emph{complement} is the subpresheaf $X \setminus Y$ of $\Hom_{\Gest}(-,X)$ given by
		\[
			(X \setminus Y)(S) = \{S \to X \mid S \times_X Y = \emptyset\}.
		\]
		If this presheaf is representable by some $Z \to X$, we call $Z$ the complement of $Y$ in $X$.
	\end{definition}

	\begin{remark}
		In general $X \setminus Y$ need not be representable. If $f$ lies in $\Gest^{n\text{-ét}}$ for some $n$, then it might follow from abstract topos-theoretic arguments that $X \setminus Y$ is representable.
	\end{remark}

	\begin{proposition}\label[proposition]{prop:Complement_Construction}
		Let $j\colon U \hookrightarrow X$ be an $n$-étale injection, corresponding to the idempotent coalgebra $C := (\Oo(U)/\Oo(X))^!_n$ in $\Oo(X)_{n+1}$, and let
		\[
			A := \cofib(C \to \unit) \qin \Oo(X)_{n+1}.
		\]
		Then $A$ is an idempotent algebra, and $\Gest_X(A) \hookrightarrow X$ is the complement of $U$.
	\end{proposition}
	\begin{proof}
		Since the tensor product preserves cofibers in each variable, idempotency of $C$ gives
		\[
			A \otimes C = \cofib(C \otimes C \to C) = 0, \qquad\text{whence}\qquad A \otimes A = \cofib(A \otimes C \to A) = \cofib(0 \to A) = A,
		\]
		so $A$ is an idempotent algebra, complementary to $C$: the sequence $C \to \unit \to A$ is a cofiber sequence with $A \otimes C = 0$.

		Let $g\colon S \to X$ be a map. On one side, $S \times_X U \hookrightarrow S$ is the $n$-étale injection over $S$ with dualizing coalgebra $(\Oo(S \times_X U)/\Oo(S))^!_n = g_n^*(C)$, so $S \times_X U$ is empty if and only if $g_n^*(C) = 0$. On the other, $S \times_X \Gest_X(A) \hookrightarrow S$ is the closed immersion over $S$ with idempotent algebra $(\Oo(S \times_X \Gest_X(A))/\Oo(S))_n = g_n^*(A)$, so $g$ factors through $\Gest_X(A) \hookrightarrow X$ if and only if $g_n^*(A) \simeq \unit$.

		Applying $g_n^*$ to the cofiber sequence $C \to \unit \to A$ yields a complementary pair $g_n^*(C), g_n^*(A)$ in $\Oo(S)_{n+1}$, again with $g_n^*(A) = \cofib(g_n^*(C) \to \unit)$ and $g_n^*(C) \otimes g_n^*(A) = 0$. For such a pair one member vanishes precisely when the other is the unit: if $g_n^*(C) = 0$ then $g_n^*(A) = \cofib(0 \to \unit) = \unit$, and if $g_n^*(A) = \unit$ then $g_n^*(C) = g_n^*(C) \otimes g_n^*(A) = g_n^*(C \otimes A) = 0$. Hence $g$ factors through $\Gest_X(A)$ if and only if $S \times_X U = \emptyset$; that is, $\Gest_X(A)$ represents the complement $X \setminus U$.
	\end{proof}

	\begin{remark}
		If instead one is given an $n$-proper injection $i\colon Z \hookrightarrow X$, one regards $i$ as $(n+1)$-étale by ambidexterity (\Cref{prop:Ambidexterity_Stefanich_Rings}) and applies \Cref{prop:Complement_Construction} one level higher to obtain its complement.
	\end{remark}

	\begin{warning}
		The passage $C \mapsto A = \cofib(C \to \unit)$ is not reversible when $\Oo(X)_{n+1}$ is not stable\footnote{Note that stability can only happen for $n = 0$.}: it may happen that $A = 0$ even though $C \to \unit$ is not an isomorphism. So an open immersion can have empty closed complement without being everything.
	\end{warning}

	\begin{corollary}\label[corollary]{cor:Complement_Open_Closed}
		The complement of an open immersion always yields a closed immersion. The converse holds in Gestalten over $\Spec(\S)$. \qed
	\end{corollary}

	\begin{corollary}
		Consider a 1-presentable stable category $E \in 1\Pr_{\st}$. Then there are (order-preserving or order-reversing) bijections between the following four posets:
		\begin{enumerate}
			\item Closed immersions $Z \hookrightarrow \Gest_{\S}(E)$;
			\item Idempotent algebras $A$ in $E$;
			\item Open immersions $U \hookrightarrow \Gest_{\S}(E)$;
			\item Idempotent coalgebras $C$ in $E$. \qednow
		\end{enumerate} 
	\end{corollary}

	\begin{example}\label[example]{ex:Complement_Zp}
		The map $\Gest(\Z[\tfrac{1}{p}]) \to \Gest(\Z)$ of \Cref{ex:Closed_Immersion_Zp} is a closed immersion, given by the idempotent algebra $\Z[\tfrac{1}{p}] \in \D(\Z)$. As $\D(\Z)$ is stable, its complement is the open immersion given by the complementary idempotent coalgebra $C := \fib(\Z \to \Z[\tfrac{1}{p}]) \simeq (\Z[\tfrac{1}{p}]/\Z)[-1]$, the Gestalt corresponding to
		\[
			\coMod_C(\D(\Z)) = \{M \mid M \otimes \Z[\tfrac{1}{p}] = 0\} \xrightarrow[\sim]{M \mapsto M^{\wedge}_p} \D(\Z)^{\wedge}_p.
		\]
		The two immersions assemble into the arithmetic fracture recollement of $\D(\Z)$: 
		\[
			\begin{tikzcd}[column sep=5em]
				\D(\Z[\tfrac{1}{p}]) \arrow[r, hook] \arrow[r, bend left, leftarrow, "{(-)[1/p]}", shift left] & \D(\Z) \arrow[l, bend left, shift left] \arrow[r, "C \otimes -"] \arrow[r, hookleftarrow,bend left, shift left] \arrow[r, hookleftarrow,bend right, shift right, "M \mapsto M^{\wedge}_p"] & \coMod_C(\D(\Z)).
			\end{tikzcd}
		\]
		Tensoring the cofiber sequence $C \to \unit \to \Z[\tfrac{1}{p}]$ with $M$ recovers the fracture fiber sequence
		\[
			C \otimes M \to M \to M[\tfrac{1}{p}].
		\]
	\end{example}

	\section{Surjections of Gestalten}
	\label[section]{sec:Surjections_Of_Gestalten}

	We will now introduce the dual notion of a \emph{surjection} of Gestalten. Even though the category $\Gest$ as a whole is not a topos, the suave surjections of Gestalten will admit a behavior similar to that of effective epimorphisms in a topos.

	\begin{definition}\label[definition]{def:Surjective_Gestalten}
		A map $g\colon Y \to X$ of Gestalten is \emph{surjective}, or a \emph{cover}, if the map
		\[
			\im(g) := \colim_{[m] \in \simp\catop} Y^{\times_X(m+1)} \to X
		\]
		is an isomorphism.
	\end{definition}

	\begin{remark}
		Assume that $g$ is $n$-étale, so that it is a morphism in the countably toposic category $\Gest^{n\text{-ét}}_{/X}$ of \Cref{thm:Etale_Countably_Toposic}. Then $g$ is surjective precisely when it is a cover in the sense of topos theory.
	\end{remark}

	\begin{lemma}
		\label[lemma]{lem:Pullback_Along_Surjection_Conservative}
		Let $f\colon Y \to X$ be an $n$-suave surjection of Gestalten. Then the pullback functor $f^*\colon \Gest_{/X} \to \Gest_{/Y}$ is conservative.
	\end{lemma}
	\begin{proof}
		Let $\varphi\colon A \to B$ be a morphism over $X$ such that $f^*\varphi\colon A \times_X Y \to B \times_X Y$ is an isomorphism over $Y$; we must show $\varphi$ is an isomorphism. Since $f$ is surjective, $\colim_{[m] \in \simp\catop} Y^{\times_X(m+1)} \simeq X$. Each map $Y^{\times_X(m+1)} \to X$ is $n$-suave by iterated base change, so \Cref{prop:Suave_Countable_Limits}(2) gives, for any $T \to X$,
		\[
			T \times_X \colim_{[m]} Y^{\times_X(m+1)} \simeq \colim_{[m]} (T \times_X Y)^{\times_Y(m+1)}.
		\]
		Applying this with $T = A$ and using surjectivity of $f$:
		\[
			A \simeq A \times_X \colim_{[m]} Y^{\times_X(m+1)} \simeq \colim_{[m]} (A \times_X Y)^{\times_Y(m+1)},
		\]
		and likewise $B \simeq \colim_{[m]} (B \times_X Y)^{\times_Y(m+1)}$. Since $A \times_X Y \simeq B \times_X Y$ over $Y$, taking iterated fiber products over $Y$ gives $(A \times_X Y)^{\times_Y(m+1)} \simeq (B \times_X Y)^{\times_Y(m+1)}$ for all $m \geq 0$. Taking colimits, $\varphi$ induces an equivalence $A \simeq B$ over $X$.
	\end{proof}

	\begin{lemma}\label[lemma]{lem:Image_Suave_Injection}
		Let $f\colon Y \to X$ be an $n$-suave map of Gestalten. Then the image $U := \colim_{[m] \in \simp\catop} Y^{\times_X(m+1)}$ is an $n$-suave injection $U \hookrightarrow X$.
	\end{lemma}
	\begin{proof}
		By base change and stability under colimits, each $Y^{\times_X(m+1)} \to X$ is $n$-suave, hence so is $U \to X$. It remains to show $U \to X$ is a monomorphism, i.e.\ that the diagonal $U \to U \times_X U$ is an isomorphism, or equivalently that the first projection $U \times_X U \to U$ is an isomorphism.

		It will suffice to show that the projection $Y \times_X U \to Y$ is an isomorphism. By \Cref{prop:Suave_Countable_Limits}(2), base change along $Y \to X$ commutes with the colimit of the $n$-suave diagram $[m] \mapsto Y^{\times_X(m+1)}$, giving
		\[
			Y \times_X U \simeq \colim_{[m]} Y^{\times_X(m+2)}.
		\]
		The simplicial object $[m] \mapsto Y^{\times_X(m+2)}$ carries extra degeneracies $s_{-1}\colon Y^{\times_X(m+2)} \to Y^{\times_X(m+3)}$ given by inserting a copy of $Y$ at position $0$ via the diagonal $\Delta\colon Y \to Y \times_X Y$. These make it a split augmented simplicial object with augmentation $Y$ (via the first projection), so $Y \times_X U \simeq \colim_{[m]} Y^{\times_X(m+2)} \iso Y$.
	\end{proof}

	\begin{proposition}\label[proposition]{prop:Epi_Mono_Factorization_Suave}
		Every $n$-suave map $f\colon Y \to X$ of Gestalten admits a (necessarily unique) factorization
		\[
			Y \twoheadrightarrow U \hookrightarrow X
		\]
		into an $n$-suave surjection followed by an $n$-suave injection.
	\end{proposition}
	\begin{proof}
		The factorization is provided by \Cref{lem:Image_Suave_Injection}: set $U := \colim_{[m] \in \simp\catop} Y^{\times_X(m+1)}$, so that $U \hookrightarrow X$ is an $n$-suave injection by that lemma, and $Y \to U$ is an $n$-suave surjection by definition. For uniqueness, any factorization $Y \twoheadrightarrow V \hookrightarrow X$ with $V \to X$ a monomorphism gives $Y^{\times_V(m+1)} \simeq Y^{\times_X(m+1)}$ for all $m$, hence $V \simeq \colim_{[m]} Y^{\times_V(m+1)} \simeq U$.
	\end{proof}

	\begin{corollary}
		\label[corollary]{cor:Surjective_And_Injective_Implies_Iso}
		An $n$-suave map of Gestalten is an isomorphism if and only if it is both surjective and injective.\qed
	\end{corollary}

	\begin{proposition}\label[proposition]{prop:Surjective_Criteria}
		Let $f\colon Y \to X$ be an $n$-suave map of Gestalten. The following are equivalent.
		\begin{enumerate}
			\item $f$ is surjective.
			\item The map $\colim_{[m]} (\Oo(Y^{\times_X(m+1)})/\Oo(X))^!_n \to \unit$ is an isomorphism.
			\item The map $\colim_{[m]} \bigl((\Oo(Y)/\Oo(X))^!_n\bigr)^{\otimes(m+1)} \to \unit$ is an isomorphism.
			\item The pullback functor $f_n^*\colon \Oo(X)_{n+1} \to \Oo(Y)_{n+1}$ is conservative.
		\end{enumerate}
	\end{proposition}
	\begin{proof}
		(1) $\Leftrightarrow$ (2): Consider the epi-mono factorization $Y \twoheadrightarrow U \hookrightarrow X$ of $f$ from \Cref{prop:Epi_Mono_Factorization_Suave}. Then $f$ is surjective if and only if $U \to X$ is an isomorphism. Since this map is an $n$-suave injection, hence $n$-affine by \Cref{prop:Suave_Etale_Idempotent_Coalgebras}, it is an isomorphism if and only if it induces an isomorphism on $(\Oo(-)/\Oo(X))_n$, which is condition (2).

		(2) $\Leftrightarrow$ (3) is \Cref{lem:Kunneth_Suave}.

		(1) $\Rightarrow$ (4): By \Cref{lem:Pullback_Along_Surjection_Conservative}, the pullback functor $f^*\colon \Gest_{/X} \to \Gest_{/Y}$ is conservative. Since $n$-affine objects of $\Gest_{/X}$ correspond, via \Cref{prop:Suave_Etale_Idempotent_Coalgebras}, to objects of $\Oo(X)_{n+1}$, and $f^*$ restricts to $f_n^*$ on these, condition~(4) follows.

		(4) $\Rightarrow$ (3): it suffices to check the map of (3) becomes an isomorphism after base change to $\Oo(Y)_{n+1}$. But then the projection $Y \times_X Y \to Y$ is split by the diagonal, so the colimit system acquires extra degeneracies and is therefore already a colimit diagram.
	\end{proof}

	For $n = 0$ these conditions are easy to check directly. For higher $n$ one instead wants a criterion phrased in terms of \emph{generation}: being a cover is closely related to the functor
	\[
		f_{n-1,!}\colon \Hom_{\Oo(X)_{n+1}}\bigl(\unit,(\Oo(Y)/\Oo(X))^!_n\bigr) \to \Oo(X)_n
	\]
	generating the target quickly. If $f$ is in addition $(n-1)$-proper, the source identifies with $\Oo(Y)_n$ and the displayed functor with $f_{n-1,*}\colon \Oo(Y)_n \to \Oo(X)_n$.

	We record the following $2$-categorical lemma, a special case of the condensation formalism of \citet{GaiottoJohnsonFreyd2019Condensations}.

	\begin{lemma}[Gaiotto--Johnson-Freyd]\label[lemma]{lem:Gaiotto_Johnson_Freyd}
		Let $\bbC \in 2\Pr$, let $X, Y \in \bbC$, and suppose there are maps $F\colon Y \to X$ and $G\colon X \to Y$ such that $\id_Y$ is a retract of $GF$ in $\End_{\bbC}(Y)$. Then $Y$ lies in the $1\Pr$-linear subcategory of $\bbC$ generated by $X$.
	\end{lemma}
	\begin{proof}
		We may replace $\bbC$ by the $1\Pr$-linear subcategory atomically generated by $X$ and $Y$, so that $X$ is atomic in $\bbC$. Let $i\colon \bbC_X \subset \bbC$ be the subcategory generated by $X$; it admits a linear right adjoint $i^R\colon \bbC \to \bbC_X$, and we set $Y' := i^R(Y) \in \bbC_X$, regarded also as an object of $\bbC$ via $i$. There is a comparison map $H\colon Y' \to Y$, which we must show is an isomorphism. Applying $i^R$ to $F$ and $G$, the object $Y'$ inherits maps $F'\colon Y' \to X$ and $G'\colon X \to Y'$ such that $\id_{Y'}$ is a retract of $G'F'$.

		By the Yoneda embedding we may assume $\bbC = 1\Pr$, so that $X, Y, Y'$ are presentable categories and $H\colon Y' \to Y$ is a functor we must show is an equivalence. Every object $y \in Y$ is a retract
		\[
			y \to GF(y) = H(G'F(y)) \to y,
		\]
		so every object of $Y$ is a retract of an object in the image of $H$, and it suffices to show $H$ is fully faithful. We first check this on the image of $G'$. For $x_1, x_2 \in X$, the composite
		\[
			\Map_Y(Gx_1,Gx_2) \xrightarrow{F} \Map_X(FGx_1,FGx_2) \xrightarrow{G} \Map_Y(GFGx_1,GFGx_2) \to \Map_Y(Gx_1,Gx_2),
		\]
		where the last map uses the two structure maps of the retraction $\id_Y \Rightarrow GF \Rightarrow \id_Y$, is the identity, so $\Map_Y(Gx_1,Gx_2)$ is exhibited as a retract of $\Map_X(FGx_1,FGx_2)$. Likewise $\Map_{Y'}(G'x_1,G'x_2)$ is a retract of $\Map_X(F'G'x_1,F'G'x_2)$. Since $F'G' = FG$, these are retracts of the same anima $\Map_X(FGx_1,FGx_2)$, compatibly with $H$ (which intertwines $G$ with $G'$ and $F$ with $F'$), so $H$ induces an isomorphism $\Map_{Y'}(G'x_1,G'x_2) \iso \Map_Y(Gx_1,Gx_2)$. As every object of $Y$ is a retract of one in the image of $G$, this shows $H$ is fully faithful, hence an isomorphism.
	\end{proof}

	\begin{proposition}\label[proposition]{prop:Surjective_Retract_Criterion}
		Let $f\colon Y \to X$ be $n$-suave, $n \geq 1$. Consider the functor
		\[
			\colim_{[m] \in \simp\catop} \Hom_{\Oo(X)_{n+1}}\bigl(\unit,(\Oo(Y^{\times_X(m+1)})/\Oo(X))^!_n\bigr) \to \Oo(X)_n \qin 1\Pr.
		\]
		If $\unit$ is a retract of an object in the image of this functor, then $f$ is surjective.
	\end{proposition}
	\begin{proof}
		Let $C := \colim_{[m] \in \simp\catop} (\Oo(Y^{\times_X(m+1)})/\Oo(X))^!_n$, the idempotent coalgebra corresponding to the image of $f$ (\Cref{prop:Surjective_Criteria}), with counit $F\colon C \to \unit$. The hypothesis exhibits $\unit$ as a retract of an object in the image of
		\[
			F\colon \Hom_{\Oo(X)_{n+1}}(\unit,C) \to \Hom_{\Oo(X)_{n+1}}(\unit,\unit) = \Oo(X)_n,
		\]
		so there is $G\colon \unit \to C$ with $\id$ a retract of $FG$. By \Cref{lem:Gaiotto_Johnson_Freyd} applied in $\Oo(X)_{n+1}$ with $Y = \unit$, $X = C$, the unit $\unit$ lies in the $1\Pr$-linear subcategory of $\Oo(X)_{n+1}$ generated by $C$. The category of $C$-comodules is a $1\Pr$-linear subcategory containing $C$, hence contains this generated subcategory; so $\unit$ is a $C$-comodule, $\coMod_C(\Oo(X)_{n+1}) = \Oo(X)_{n+1}$, and $C = \unit$. By \Cref{prop:Surjective_Criteria}(2), $f$ is surjective.
	\end{proof}

	We also record a converse, assuming $X$ is $n$-prim.

	\begin{proposition}\label[proposition]{prop:Surjective_Generation}
		Let $n \geq 1$ and let $X$ be an $n$-prim Gestalt (or, more generally, assume the functor $n\Pr \to \Oo(X)_{n+1}$ has a colimit-preserving right adjoint). Let $f\colon Y \to X$ be an $n$-suave surjection. Then
		\[
			\colim_{[m] \in \simp\catop} \Hom_{\Oo(X)_{n+1}}\bigl(\unit,(\Oo(Y^{\times_X(m+1)})/\Oo(X))^!_n\bigr) \iso \Oo(X)_n,
		\]
		and the image of $f_{n-1,!}$ generates $\Oo(X)_n$ under simplicial colimits.
	\end{proposition}
	\begin{proof}
		The colimit-preserving right adjoint is $\Hom_{\Oo(X)_{n+1}}(\unit,-)$, so the displayed isomorphism follows from \Cref{prop:Surjective_Criteria}(2). Any $M \in \Oo(X)_n$ is the geometric realization of the simplicial object obtained by composing
		\[
			\Oo(X)_n \to \Hom_{\Oo(X)_{n+1}}\bigl(\unit,(\Oo(Y^{\times_X(\bullet+1)})/\Oo(X))^!_n\bigr) \to \Oo(X)_n,
		\]
		where the first map is the right adjoint, on underlying categories, of the second. This can be checked by conservativity after base change along $f$, where this bar resolution is split. The colimit takes values in the image of $f_{n-1,!}$.
	\end{proof}

	\begin{definition}
		A map of Gestalten is called \emph{eventually suave} if it is $m$-suave for some $m \geq 0$.
	\end{definition}

	\begin{corollary}\label[corollary]{cor:Surjective_Affine_Truncation}
		Let $f\colon Y \to X$ be a countably presented, $n$-prim, eventually suave map of Gestalten. Then the map
		\[
			Y \to \Gest_X\bigl((\Oo(Y)/\Oo(X))_n\bigr)
		\]
		is surjective. In particular, $f$ is surjective if and only if $\Gest_X\bigl((\Oo(Y)/\Oo(X))_n\bigr) \to X$ is surjective.
	\end{corollary}
	\begin{proof}
		Replacing $X$ by $\Gest_X\bigl((\Oo(Y)/\Oo(X))_n\bigr)$, we may assume $(\Oo(Y)/\Oo(X))_n = \unit$. Assume $f$ is $m$-suave, $m \geq n+2$ without loss of generality; we induct on $m$.

		\emph{Base case $m = n+2$.} Shifting, we may assume $n = 0$, so $f_{0,*}(\unit) = \unit$. As $f$ is both $2$-prim (being $0$-prim) and $2$-suave, $(\Oo(Y)/\Oo(X))^!_2 \in \Oo(X)_3$ is dualizable with dual $(\Oo(Y)/\Oo(X))_2$. As $f$ is $1$-prim, the unit map $f_1^*\colon \unit \to (\Oo(Y)/\Oo(X))_2$ has a right adjoint $f_{1,*}$ in $\Oo(X)_3$, so dually the counit $f_{1,!}\colon (\Oo(Y)/\Oo(X))^!_2 \to \unit$ has a left adjoint $f^!_1$ in $\Oo(X)_3$. Dualizing identifies the composite
		\[
			\unit \xrightarrow{f^!_1} (\Oo(Y)/\Oo(X))^!_2 \xrightarrow{f_{1,!}} \unit \qquad\text{with}\qquad \unit \xrightarrow{f_1^*} (\Oo(Y)/\Oo(X))_2 \xrightarrow{f_{1,*}} \unit,
		\]
		which, as an object of $\Oo(X)_2 = \End_{\Oo(X)_3}(\unit)$, is given by tensoring with $f_{1,*}(\unit)$. Now $f_{1,*}(\unit) \in \Oo(X)_2$ has $\unit$ as a retract: the unit $\unit \to f_{1,*}(\unit)$ has a right adjoint (as $f$ is $0$-prim), and the composite $\unit \to f_{1,*}(\unit) \to \unit$ is tensoring with $f_{0,*}(\unit) = \unit$. So $\unit \in \Oo(X)_2$ is a retract of an object in the image of
		\[
			f_{1,!}\colon \Hom_{\Oo(X)_3}\bigl(\unit,(\Oo(Y)/\Oo(X))^!_2\bigr) \to \Oo(X)_2,
		\]
		and \Cref{prop:Surjective_Retract_Criterion} (with $n=2$) shows that $f$ is surjective.

		\emph{Induction step.} Assume the claim for $(m-1)$-suave maps, and let $f$ be $n$-prim and $m$-suave. Let $Y' := \Gest_X\bigl((\Oo(Y)/\Oo(X))_{m-2}\bigr)$ be the $(m-2)$-affine truncation of $Y$. By the base case, $Y \to Y'$ is surjective. The map $g\colon Y' \to X$ is $(m-2)$-affine and still $n$-prim: for $k \leq m-2$ the unit $\unit \to g_{k,*}(\unit)$ agrees with $\unit \to f_{k,*}(\unit)$, while for $k > m-2$ it has a right adjoint as $g$ is $(m-2)$-affine and countably presented. So $g$ is $n$-prim and, being countably presented and $(m-2)$-affine, $(m-2)$-proper (\Cref{prop:Proper_Properties}(1)); hence $g$ is $(m-1)$-suave by ambidexterity (\Cref{prop:Ambidexterity_Stefanich_Rings}), and by induction $g$ is surjective. As $Y \to Y'$ is surjective as well, so is the composite $Y \to X$.
	\end{proof}

	\begin{remark}
		Scholze does not assume countable presentability in \Cref{cor:Surjective_Affine_Truncation}; we do not know how to prove the statement without this assumption.
	\end{remark}

	\section{Descendability}

	We now relate surjectivity to Mathew's notion of descendability \citep[§§3--4]{Mathew_Galois}, see also \citet{Mathew2018DescentNilpotence}.

	\begin{definition}\label[definition]{def:Descendable}
		Let $\bbC$ be a stably symmetric monoidal category and $R \in \Alg(\bbC)$.
		\begin{enumerate}
			\item An object $M \in \bbC$ is \emph{$R$-nilpotent} if it lies in the thick $\otimes$-ideal $\Thick^{\otimes}(R)$ generated by $R$.
			\item $R$ is \emph{descendable} if the unit $\unit \in \bbC$ is $R$-nilpotent.
		\end{enumerate}
	\end{definition}

	\begin{proposition}[{\citet[Proposition 2.32]{Mathew2018DescentNilpotence}}]\label[proposition]{prop:Descendable_Criteria}
		Let $R \in \Alg(\bbC)$ with $I := \fib(\unit \to R)$. The following are equivalent.
		\begin{enumerate}
			\item $R$ is descendable.
			\item The map $I \to \unit$ is $\otimes$-nilpotent: $I^{\otimes N} \to \unit$ is null for some $N$.
			\item For some finite $m$, the unit map $\unit \to \lim_{\simp_{\leq m}} R^{\otimes \bullet+1}$, into the $m$-th partial totalization of the Amitsur (cobar) complex of $R$, admits a retract.
		\end{enumerate}
	\end{proposition}

	\begin{corollary}\label[corollary]{cor:Descendable_Implies_Surjective}
		Let $f\colon Y \to X$ be a countably presented, $0$-prim, eventually suave map of Gestalten over $\Spec(\S)$. If $f_{0,*}(\unit) \in \Oo(X)_1$ is descendable, then $f$ is surjective.
	\end{corollary}
	\begin{proof}
		By \Cref{cor:Surjective_Affine_Truncation} we may assume $Y$ is $0$-affine over $X$. Then $f$ is $0$-proper, so by ambidexterity $(\Oo(Y)/\Oo(X))^!_1 = (\Oo(Y)/\Oo(X))_1$, and likewise for the finite fiber powers of $Y$ over $X$. For each finite $m$, the functor
		\[
			\colim_{\simp\catop_{\leq m}} \Hom_{\Oo(X)_2}\bigl(\unit,(\Oo(Y^{\times_X(\bullet+1)})/\Oo(X))^!_1\bigr) = \colim_{\simp\catop_{\leq m}} (\Oo(Y^{\times_X(\bullet+1)})/\Oo(X))_1 \to \Oo(X)_1
		\]
		has a left adjoint, whose unit is the finite limit of the units at the individual levels: the indexing category is finite, and finite limits agree with finite colimits in the stable category $\Oo(X)_1$, so this finite limit of left adjoints is again computed levelwise. Evaluated at $\unit \in \Oo(X)_1$, this unit is the map
		\[
			\unit \to \lim_{\simp_{\leq m}} f_{0,*}(\unit)^{\otimes \bullet+1}
		\]
		into the $m$-th partial totalization of the Amitsur complex of $f_{0,*}(\unit)$. By descendability of $f_{0,*}(\unit)$ and \Cref{prop:Descendable_Criteria}(3), some finite $m$ makes this map admit a retract. Then $\unit$ is a retract of an object in the image of the displayed functor, and \Cref{prop:Surjective_Retract_Criterion} shows that $f$ is surjective.
	\end{proof}

	\begin{corollary}\label[corollary]{cor:Surjective_Iff_Descendable}
		Let $X$ be a $0$-prim Gestalt over $\Spec(\S)$, and let $f\colon Y \to X$ be a countably presented, $0$-prim, eventually suave map. Then $f$ is surjective if and only if $f_{0,*}(\unit)$ is descendable in $\Oo(X)_1$.
	\end{corollary}
	\begin{proof}
		The implication $\Leftarrow$ is \Cref{cor:Descendable_Implies_Surjective}. Conversely, suppose $f$ is surjective; by \Cref{cor:Surjective_Affine_Truncation} we may assume $f$ is $0$-affine. The image of $f_{0,*}\colon \Oo(Y)_1 \to \Oo(X)_1$ lies in $f_{0,*}(\unit)$-modules, so by \Cref{prop:Surjective_Generation} every object of $\Oo(X)_1$, in particular $\unit$, is a geometric realization of $f_{0,*}(\unit)$-modules. As $X$ is $0$-prim, the unit map $\unit \to \Oo(X)_1$ admits the internal right adjoint $\Hom_{\Oo(X)_1}(\unit,-)$, so $\unit \in \Oo(X)_1$ is compact; it is therefore already a retract of a finite stage of this realization, i.e.\ a retract of an object of the thick subcategory generated by $f_{0,*}(\unit)$-modules. Thus $\unit \in \Thick^{\otimes}(f_{0,*}(\unit))$, which is descendability.
	\end{proof}

	% Talks 9a and 9b are the two halves of the single ninth talk. The chapter heading,
	% the table of contents, and chapter cross-references display "9a"/"9b" (via
	% \thechapter during the \chapter call), but all internal numbering (sections,
	% theorems) runs continuously as "9.X" across both halves: \thechapter is reset to
	% plain "9" for the body, and the section/theorem counters are carried from 9a into
	% 9b below. The chapter counter still advances normally, so subsequent talks are
	% shifted down by one (Six-functor = Talk 10).
	\newcounter{savedsectionctr}
	\newcounter{savedtheoremctr}
	\renewcommand{\thechapter}{9a}
	\chapter{Affine and coaffine maps of Gestalten}
	\label[chapter]{chap:Affine_Coaffine_Maps_Gestalten}
	\renewcommand{\thechapter}{9}

	\emph{During the seminar, the contents of Chapters \ref{chap:Affine_Coaffine_Maps_Gestalten} and \ref{chap:Examples} were discussed in the same talk, which came before Talk \ref{chap:Injections_Surjections_Gestalten}. The material has been expanded and reorganized into two separate chapters for the sake of clarity of exposition.}

	The goal of this chapter is to isolate the affine-coaffine factorization system and use it to study how affineness interacts with diagonals.

	Throughout the chapter, we keep using the notations $\Gest_X(-)\colon (\StRing_{\Oo(X)/})\catop \to \Gest_{/X}$ and $\Oo(-)/\Oo(X)\colon (\Gest_{/X})\catop \to \StRing_{\Oo(X)/}$ for the inverse equivalence between Gestalten over a fixed Gestalt $X$ and Stefanich rings over the global sections $\Oo(X)$ of $X$.

	\section{The affine-coaffine factorization system}

	\begin{definition}
		Recall from \Cref{cor:Inclusion_n_Affine_A_Algebras} that for a Stefanich ring $A$, the functor $\StRing_{A/} \to \CAlg(A_{n+1}), B \mapsto (B/A)_n$ admits a fully faithful left adjoint $\CAlg(A_{n+1}) \to \StRing_{A/}$. This defines a functor
		\[
			(-)^{n\text{-affine}}\colon \StRing_{A/} \to \StRing_{A/}, B \mapsto (B/A)^{n\text{-affine}}
		\]
		equipped with a natural transformation $(B/A)^{n\text{-affine}} \to B$ of Stefanich $A$-algebras. We refer to this as the \emph{$n$-affinization of $B$ over $A$}.

		By dualizing, we similarly obtain for every map $Y \to X$ of Gestalten a factorization
		\[
			Y \to (Y/X)^{n\text{-affine}} \to X,
		\]
		with $(Y/X)^{n\text{-affine}}$ the \emph{$n$-affinization of $Y$ over $X$}.
	\end{definition}

	\begin{definition}\label[definition]{def:N_Coaffine}
		A map $f\colon Y \to X$ of Gestalten is called \emph{$n$-coaffine} if the unit map
		\[
			\unit_{\Oo(X)_{n+1}} \to (\Oo(Y)/\Oo(X))_n
		\]
		is an isomorphism in $\CAlg(\Oo(X)_{n+1})$.
	\end{definition}

	\begin{observation}\label[observation]{obs:Affine_Coaffine_Monotonicity}
		If $n < m$, then every $n$-affine map is $m$-affine, and every $m$-coaffine map is $n$-coaffine.
	\end{observation}
	\begin{proof}
		The assertion for affine maps follows, with $A=\Oo(X)$, from the factorization $\CAlg(A_{n+1}) \hookrightarrow \CAlg(A_{m+1}) \hookrightarrow \StRing_{A/}$. For coaffine maps, if $(\Oo(Y)/\Oo(X))_m$ is the unit, then applying the endomorphism transition functors $\End^{\Oo(X)}_m,\dots,\End^{\Oo(X)}_{n+1}$ identifies each lower relative algebra with the corresponding unit.
	\end{proof}

	\begin{corollary}\label[corollary]{cor:Affine_Coaffine_Implies_Iso}
		If $f\colon Y \to X$ is both $n$-affine and $n$-coaffine, then $f$ is an isomorphism.
	\end{corollary}
	\begin{proof}
		An $n$-affine map is classified by its relative algebra $(\Oo(Y)/\Oo(X))_n \in \CAlg(\Oo(X)_{n+1})$. If $f$ is also $n$-coaffine, this algebra is the unit. The $n$-affine Gestalt over $X$ classified by the unit algebra is the identity of $X$.
	\end{proof}

	\begin{lemma}\label[lemma]{lem:Affinization_Preserves_Affineness}
		Let $X$ be a Gestalt, let $m, n \geq 0$, and let $\phi\colon Z \to W$ be an $m$-affine map between Gestalten over $X$. Then its $n$-affinization over $X$,
		\[
			\phi^{n\text{-affine}}\colon (Z/X)^{n\text{-affine}} \to (W/X)^{n\text{-affine}},
		\]
		is again $m$-affine.
	\end{lemma}
	\begin{proof}
		Write $Z' := (Z/X)^{n\text{-affine}}$ and $W' := (W/X)^{n\text{-affine}}$, both $n$-affine over $X$ by construction. Thus $\phi^{n\text{-affine}}\colon Z' \to W'$ is a map between $n$-affine Gestalten over $X$, and is therefore itself $n$-affine by the cancellation property of \Cref{prop:Affineness_Properties}(4).

		If $m \geq n$, this already gives the claim, since an $n$-affine map is $m$-affine by \Cref{obs:Affine_Coaffine_Monotonicity}.

		If $m \leq n$, it remains to upgrade $n$-affineness to $m$-affineness. By the recursive criterion for affineness following \Cref{def:N_Affine}, a map that is $n$-affine is $m$-affine as soon as its relative structure maps
		\[
			L^{\Oo(Z')/\Oo(W')}_i\colon \Mod^{\Oo(W')}_{i+1}\bigl((\Oo(Z')/\Oo(W'))_i\bigr) \to (\Oo(Z')/\Oo(W'))_{i+1}
		\]
		are equivalences in the remaining degrees $m \leq i < n$. In these degrees the $n$-affinization changes neither the categories $\Oo(-)_{i+1}$ nor the relative algebra $(\Oo(Z')/\Oo(W'))_i$ of the map: they coincide with the corresponding data for $\phi$. Hence each such structure map agrees with the one for $\phi$, which is an equivalence because $\phi$ is $m$-affine. This verifies the criterion, so $\phi^{n\text{-affine}}$ is $m$-affine.
	\end{proof}

	We record the following observation.

	\begin{proposition}\label[proposition]{prop:Affine_Coaffine_Factorization_System}
		The $n$-coaffine maps and the $n$-affine maps form a factorization system on $\Gest$. The factorization of a map $Y \to X$ is
		\[
			Y \to (Y/X)^{n\text{-affine}} \to X.
		\]
	\end{proposition}
	\begin{proof}
		Write $Y' := (Y/X)^{n\text{-affine}}$. By construction, $Y' \to X$ is $n$-affine and the map $Y \to Y'$ identifies the relative algebra of $Y'$ over $X$ in degree $n$ with that of $Y$ over $X$. After replacing the base $X$ by $Y'$, this says precisely that
		\[
			\unit_{\Oo(Y')_{n+1}} \to (\Oo(Y)/\Oo(Y'))_n
		\]
		is an isomorphism, i.e.\ that $Y \to Y'$ is $n$-coaffine.

		It remains to check orthogonality. Let $i\colon Y \to Z$ be $n$-coaffine, let $p\colon W \to X$ be $n$-affine, and consider a commutative square
		\[
		\begin{tikzcd}
			Y \rar \dar[swap]{i} & W \dar{p} \\
			Z \rar & X.
		\end{tikzcd}
		\]
		We regard $Y$ and $Z$ as objects of $\Gest_{/X}$ through the bottom map. Since $p$ is $n$-affine, with classifying algebra $R := (\Oo(W)/\Oo(X))_n \in \CAlg(\Oo(X)_{n+1})$, maps into $W$ over $X$ are computed by maps out of $R$:
		\[
			\Map_{\Gest_{/X}}(T,W)
			\simeq
			\Map_{\CAlg(\Oo(X)_{n+1})}\bigl(R,(\Oo(T)/\Oo(X))_n\bigr)
		\]
		naturally in $T \in \Gest_{/X}$. Applying this to $T=Z$ and $T=Y$, the map
		\[
			\Map_{\Gest_{/X}}(Z,W) \to \Map_{\Gest_{/X}}(Y,W)
		\]
		is an equivalence, because $i$ being $n$-coaffine implies that
		\[
			(\Oo(Z)/\Oo(X))_n \to (\Oo(Y)/\Oo(X))_n
		\]
		is an isomorphism: if $A=\Oo(X)$, $C=\Oo(Z)$, and $B=\Oo(Y)$, then the relative algebra for $A \to B$ is obtained by applying the right adjoint for $A \to C$ to the relative algebra for $C \to B$, which is the unit by $n$-coaffineness. This shows that the square we started with has a contractible anima of fillers. 
		
		The usual retract argument then identifies the two classes with the corresponding left and right orthogonal classes.
	\end{proof}

	\section{Affineness and diagonals}

	The main result of this section, \Cref{prop:Affineness_And_Diagonals} below, shows that, under suitable hypotheses, passing to the diagonal drops the level of affineness by precisely one. Part (2) of this proposition generalizes results of Gaitsgory on 1-affineness \citep{Gaitsgory_1Affineness}, and related results for the general case appeared in Stefanich's thesis \citep{Stefanich_HigherQCoh}.

	\begin{lemma}\label[lemma]{lem:Affine_Equivalence_On_Classifying_Level}
		Let $f\colon Y \to X$ be a map of Gestalten. If
		\[
			f_n^*\colon \Oo(X)_{n+1} \to \Oo(Y)_{n+1}
		\]
		is an equivalence, then $f$ is $n$-coaffine.
	\end{lemma}
	\begin{proof}
		Let $f_{n,*}$ denote the right adjoint of $f_n^*$. Since $f_n^*$ is an equivalence, the adjunction unit identifies the unit of $\Oo(X)_{n+1}$ with
		\[
			f_{n,*}\unit_{\Oo(Y)_{n+1}} = (\Oo(Y)/\Oo(X))_n.
		\]
		This is exactly the condition that $f$ is $n$-coaffine.
	\end{proof}

	\begin{lemma}\label[lemma]{lem:Prim_Coaffine_Diagonal_Coaffine}
		Let $h\colon Y \to X$ be an $(n+1)$-prim and $(n+1)$-coaffine map of Gestalten. Then the diagonal $\Delta_h\colon Y \to Y \times_X Y$ is $n$-coaffine.
		In particular, if $\Delta_h$ is $n$-affine, $h$ is injective.
	\end{lemma}
	\begin{proof}
		By $(n+1)$-coaffineness, the unit map
		\[
			\unit_{\Oo(X)_{n+2}} \to (\Oo(Y)/\Oo(X))_{n+1}
		\]
		is an isomorphism. The Künneth formula for prim maps (\Cref{lem:Kunneth_Prim}), applied to the two copies of $h$, therefore gives
		\[
			(\Oo(Y \times_X Y)/\Oo(X))_{n+1}
			\simeq
			(\Oo(Y)/\Oo(X))_{n+1}
			\otimes_{\Oo(X)_{n+2}}
			(\Oo(Y)/\Oo(X))_{n+1}
			\simeq
			\unit_{\Oo(X)_{n+2}}.
		\]
		Under this identification, the map induced by $\Delta_h$ on relative algebras over $X$ is the multiplication map $\unit \otimes \unit \to \unit$, hence is an isomorphism. Passing to underlying categories gives an equivalence
		\[
			(\Delta_h)_n^*\colon \Oo(Y \times_X Y)_{n+1} \to \Oo(Y)_{n+1}.
		\]
		Thus \Cref{lem:Affine_Equivalence_On_Classifying_Level} shows that $\Delta_h$ is $n$-coaffine.

		If $\Delta_h$ is also $n$-affine, then \Cref{cor:Affine_Coaffine_Implies_Iso} shows that $\Delta_h$ is an isomorphism. Thus $h$ is injective.
	\end{proof}

	\begin{proposition}
		\label[proposition]{prop:Affineness_And_Diagonals}
		Let $n \geq 0$, and let $f\colon Y \to X$ be an eventually suave map of Gestalten.
		\begin{enumerate}
			\item If $f$ is $(n+1)$-affine and $\Delta_f$ is $n$-prim, then $\Delta_f$ is $n$-affine.
		\end{enumerate}
		Assume additionally that both $f$ and its $(n+1)$-affinization $(Y/X)^{(n+1)\text{-affine}} \to X$ are countably presented.\footnote{Scholze does not assume these countability hypotheses; at present we do not know how to prove this implication without them.}
		\begin{enumerate}
			\item[(2)] If $\Delta_f$ is $n$-affine and $f$ is $(n+1)$-prim, then $f$ is $(n+1)$-affine.
		\end{enumerate}
		In particular, if $f$ is $n$-proper, then $f$ is $(n+1)$-affine if and only if $\Delta_f$ is $n$-affine.
	\end{proposition}
	\begin{proof}
		Part (1) is an instance of \Cref{lem:Reduce_Affineness}, applied to the codiagonal
		\[
			\Delta_f^*\colon A = \Oo(Y \times_X Y) \to B = \Oo(Y),
		\]
		which is what the diagonal $\Delta_f$ becomes under the equivalence $\Gest_{/X} \simeq (\StRing_{\Oo(X)/})\catop$. We verify the three hypotheses of that lemma. First, since $n$-affine maps are stable under passing to codiagonals (\Cref{prop:Affineness_Properties}(5)), the codiagonal $\Delta_f^*$ of the $(n+1)$-affine map $f$ is again $(n+1)$-affine. Second, $\Delta_f^*$ is $n$-prim by assumption. Third, a projection $\pr\colon Y \times_X Y \to Y$ yields a section $\pr^*\colon \Oo(Y) \to \Oo(Y \times_X Y)$ of $\Delta_f^*$, and $\pr^*$ is $(n+1)$-affine because $\pr$ is a base change of $f$ (\Cref{prop:Affineness_Properties}(3)). Thus \Cref{lem:Reduce_Affineness} applies and shows that $\Delta_f^*$, and therefore $\Delta_f$, is $n$-affine.

		For part (2), consider the unique factorization of $f$ as an $(n+1)$-coaffine map followed by an $(n+1)$-affine map
		\[
			Y \xrightarrow{h} Y' := (Y/X)^{(n+1)\text{-affine}} \xrightarrow{g} X.
		\]
		We need to show that $h$ is an isomorphism. Since $f$ is countably presented, \Cref{cor:Surjective_Affine_Truncation}, applied to the $(n+1)$-prim, eventually suave map $f$, shows that $h$ is surjective. Moreover $h$ is eventually suave: the map $g$ is $(n+1)$-affine and countably presented, hence $(n+1)$-proper (\Cref{prop:Proper_Properties}(1)) and therefore $(n+2)$-\'etale by ambidexterity (\Cref{prop:Ambidexterity_Stefanich_Rings}(2)); since $g$ is \'etale, right cancellation for \'etale maps (\Cref{prop:Etale_Properties}) applied to $f = g \circ h$, together with the eventual suaveness of $f$, shows that $h$ is eventually suave. It therefore suffices by \Cref{cor:Surjective_And_Injective_Implies_Iso} to show that $h$ is a monomorphism, i.e.\ that $\Delta_h$ is an isomorphism.

		By construction, $h$ is $(n+1)$-coaffine. We now check that $h$ is also $(n+1)$-prim. Recall from above that $g$ is $(n+1)$-proper. Applying the cancellation statement \Cref{prop:Proper_Properties}(4) to the factorization $f=g \circ h$, the $(n+1)$-primness of $f$ implies the $(n+1)$-primness of $h$. It then follows from \Cref{lem:Prim_Coaffine_Diagonal_Coaffine} that $\Delta_h$ is $n$-coaffine, and so it remains to show that $\Delta_h$ is also $n$-affine.

		Since $\Delta_f$ is assumed $n$-affine and $n$-affine maps are closed under left cancellation and base change, it suffices to show that $\Delta_g$ is $n$-affine:
		\[
		\begin{tikzcd}
			Y \dar[swap]{\Delta_h} \drar{\Delta_f} \\
			Y \times_{Y'} Y \rar \drar[pullback] \dar & Y \times_X Y \dar \\
			Y' \rar{\Delta_g} & Y' \times_X Y'.
		\end{tikzcd}
		\]
		By assumption, the map $\Delta_f\colon Y \to Y \times_X Y$ is $n$-affine. Since $(n+1)$-affinization preserves $n$-affine maps by \Cref{lem:Affinization_Preserves_Affineness}, it thus remains to show that $\Delta_g$ is the image of $\Delta_f$ under the functor
		\[
			(-)^{(n+1)\text{-affine}}\colon \Gest_{/X} \to \Gest_{/X}^{(n+1)\text{-affine}}.
		\]
		But indeed, since the map $f\colon Y \to X$ is $(n+1)$-prim, the Künneth formula for prim maps (\Cref{lem:Kunneth_Prim}) implies that we have
		\[
			(\Oo(Y \times_X Y)/\Oo(X))_{n+1} \iso (\Oo(Y)/\Oo(X))_{n+1} \otimes_{\Oo(X)_{n+2}} (\Oo(Y)/\Oo(X))_{n+1}.
		\]
		Under the equivalence
		\[
			\CAlg(\Oo(X)_{n+2}) \simeq \bigl(\Gest_X^{(n+1)\text{-affine}}\bigr)\catop,
		\]
		the object $Y'$ corresponds to $(\Oo(Y)/\Oo(X))_{n+1}$, and fiber products of $(n+1)$-affine Gestalten over $X$ correspond to tensor products of the corresponding algebras in $\CAlg(\Oo(X)_{n+2})$. This identifies the $(n+1)$-affinization of $Y \times_X Y$ over $X$ with $Y' \times_X Y'$, and identifies the $(n+1)$-affinization of $\Delta_f$ with $\Delta_g$, as claimed. This concludes the proof.
	\end{proof}

	% Carry the section and theorem counters over from talk 9a so that 9b continues the
	% same "9.X" sequence (e.g. if 9a ends at 9.9, 9b starts at 9.10).
	\setcounter{savedsectionctr}{\value{section}}
	\setcounter{savedtheoremctr}{\value{theorem}}
	\renewcommand{\thechapter}{9b}
	\chapter{Examples of Gestalten}
	\label[chapter]{chap:Examples}
	\renewcommand{\thechapter}{9}
	\setcounter{section}{\value{savedsectionctr}}
	\setcounter{theorem}{\value{savedtheoremctr}}

	\emph{Talk given by Sil.} 
	
	The goal of this chapter is to discuss examples of Gestalten.

	\begin{convention}
		Throughout this chapter, we will fix $n$ large enough, and work inside $\Gest_{/X}^{n\text{-}\et}$ for some Gestalt $X$, mostly $X = *$ or $X = \Gest(\D(R))$ for some $R \in \CAlg(\Sp)$. Since the inclusions $\Gest_{/X}^{n\text{-}\et} \hookrightarrow \Gest_{/X}^{(n+1)\text{-}\et}$ are fully faithful, it doesn't matter which precise choice of $n$ we pick as long as all objects that are involved are $n$-étale.
	\end{convention}

	\section{Basic examples}

	We start with some elementary examples.

	\begin{example}[The point]
		The \emph{terminal Gestalt} $* \in \Gest_{/X}$ is the identity map $\id_X\colon X \to X$. This is $0$-affine, $0$-étale, and $0$-proper.
	\end{example}

	\begin{example}[The empty set]
		The \emph{initial Gestalt} $\emptyset \in \Gest_{/X}$ is the Gestalt associated to the zero ring:
		\[
			\emptyset = \Gest(*, *, *, \dots).
		\]
		We claim that this is always 1-affine. For each $n \geq 0$, we have $(\Oo(\emptyset)/\Oo(X))_n = *$. But for $n \geq 2$, the category $\Oo(X)_{n}$ is a module over the pointed category $1\Pr$, hence $\Oo(X)_n$ itself is pointed, so that $\Oo(\emptyset)_n = \Mod_{(\Oo(\emptyset)/\Oo(X))_{n-1}}(\Oo(X)_n) = \Mod_*(\Oo(X)_n) = *$ as desired.

		Since $\emptyset$ is 1-affine, it is in particular 1-proper. Moreover, it is $0$-étale, as it is a colimit (an empty colimit) of 0-étale objects over $X$. It is $0$-affine if and only if $\Oo(X)_1$ is pointed, or equivalently whenever $X$ is a Gestalt over $\Gest(\An_*)$.
	\end{example}

	\begin{example}[Finite sets]
		Let $I$ be a finite set. Then we have the corresponding Gestalt
		\[
			I := \bigsqcup_{x \in I} * \in \Gest_{/X}.
		\]
		If $X = \Gest(A)$, we may express this as $\Gest(\prod_{I} A_0, \prod_I A_1, \prod_I A_2, \dots)$, since colimits of Gestalten correspond to limits of Stefanich rings.

		Note that $I$ is 0-étale over $X$, as a colimit of 0-étale Gestalten. We claim it is also 1-affine: we have $(\Oo(I)/\Oo(X))_{n-1} = \prod_I \unit \in \CAlg(A_n)$, and so
		\[
			\Oo(I)_n = \Mod_{(\Oo(I)/\Oo(X))_{n-1}}(\Oo(X)_n) = \Mod_{\prod_I \unit}(A_n) \simeq \prod_I \Mod_{\unit}(A_n) \simeq \prod_{I} A_n
		\]
		for all $n \geq 2$; here we use that $A_n$ semiadditive, since it is a module over $1\Pr$, which is semiadditive. If $A_1$ is semiadditive, then $I$ is even 0-affine.
	\end{example}

	\begin{example}[Countable sets]
		Similarly as in the previous example, for any countable set $I$ we may define
		\[
			I = \coprod_I * \in \Gest_{/X}.
		\]
		Once again this is $0$-étale and 1-affine; this uses the fact that $1\Pr$ is not just semiadditive but even countably semiadditive, in that countable products agree with countable coproducts.
	\end{example}

	\begin{example}[Countable animae]
		Generalizing the previous two examples, there exists a unique countable-colimit-preserving functor
		\[
			\An^{\mathrm{ctbl}} \to \Gest_{/X}^{0\text{-}\et}
		\]
		given on the point by $* \mapsto *$. Every countable anima is $0$-étale, hence 1-proper. However, it may not be 1-affine.

		As a counterexample, we will show by induction on $n$ that the Gestalt $B^{n+1}(\Z/2)$ is not $n$-affine over $\Gest(\D(\Q))$. 
		
		We first show that it is $0$-proper. It is $0$-étale, hence $1$-proper, and its diagonal has fibers $B^n(\Z/2)$, so by induction on $n$ it remains to show that $B^{n+1}(\Z/2)$ is $0$-prim. In other words, we need to show that the functor
		\[
			\D(\Q) \to \D(\Q)^{B^{n+1}(\Z/2)} \simeq \Mod_{\Q[\Omega(B^{n+1}(\Z/2))]}(\D(\Q)) \simeq \Mod_{\Q[B^n(\Z/2)]}(\D(\Q))
		\]
		has a $\D(\Q)$-linear right adjoint. For $n \geq 1$, we have $H_*(B^n(\Z/2);\Q) \cong \Q$, so that $\Q[B^n(\Z/2)]$ is the monoidal unit in $\D(\Q)$ and this functor is even an isomorphism. In the remaining case $n = 0$, the functor is $\D(\Q) \to \Fun(B(\Z/2),\D(\Q))$, whose right adjoint is given by taking $\Z/2$-fixed points; we must show that this right adjoint preserves colimits. But since $\D(\Q)$ is rational, the norm map $(-)_{h\Z/2} \to (-)^{h\Z/2}$ is an isomorphism, so the right adjoint is also a left adjoint and hence preserves colimits.

		We may now use \Cref{prop:Affineness_And_Diagonals}(1) to inductively show that $B^{n+1}(\Z/2)$ is not $n$-affine for any $n$. To see this, assume for contradiction that $B^{n+1}(\Z/2)$ is $n$-affine. Then its diagonal is $(n-1)$-affine, hence by the pullback square
		\[
		\begin{tikzcd}
			B^{n}(\Z/2) \rar \dar \drar[pullback] & B^{n+1}(\Z/2) \dar{\Delta} \\
			* \rar & B^{n+1}(\Z/2) \times B^{n+1}(\Z/2)
		\end{tikzcd}
		\]
		we would get that $B^{n}(\Z/2)$ is $(n-1)$-affine. It thus suffices to show that $B(\Z/2)$ is not $0$-affine. Indeed, if it were, then the category $\D(\Q)^{B(\Z/2)}$ of $\Z/2$-representations of $\Q$-modules would be equivalent to the module category $\Mod_{\Q^{B(\Z/2)}}(\D(\Q))$, which is equivalent to $\D(\Q)$ since again $B(\Z/2)$ has trivial rational cohomology. But this is a contradiction as there exist non-trivial $\Q$-representations.
	\end{example}

	\section{Examples from algebraic geometry}

	Recall that for a Gestalt of the form $X = \Gest(\D(R))$ for some ring spectrum $R \in \CAlg(\Sp)$, we defined a functor
	\[
		\Gest(-)\colon \Shv_{\fpqc}(\Aff_R)^{\aleph_1} \to \Gest_{/X}, \qquad Y \mapsto \Gest(\Omega^{\infty}\Gamma(Y), \D(Y), 1\Pr_Y, 2\Pr_Y, \dots).
	\]
	In fact, we can define things more generally.

	\begin{example}[Affine line]
		\label[example]{ex:Affine_Line}
		There is a functor
		\[
			\Oo(-)_0\colon (\Gest_{/X})\catop \to \An, \qquad Y \mapsto \Oo(Y)_0.
		\]
		Writing $A := \Oo(X) \in \StRing$, this is represented by the \emph{affine line}
		\[
			\bbA^1_X := \Gest_X(A_0\{x\}),
		\]
		where $A_0\{x\}$ is the free $E_{\infty}$-algebra on one generator $x$ over $A_0$:
		\[
			\Hom_{\Gest_{/X}}(Y,\bbA^1_X) \simeq \Hom_{\CAlg(A_1)}(A_0\{x\}, (\Oo(Y)/\Oo(X))_0) \simeq \Oo(Y)_0,
		\]
		where the first equivalence is the adjunction $\Gest_X(-) \dashv (\Oo(-)/\Oo(X))_0$ of \Cref{def:N_Affine}, using $\bbA^1_X = \Gest_X(A_0\{x\})$. Note that $\bbA^1_X$ is $0$-affine, hence $0$-proper and $1$-étale.

		Over the absolute base, this is $\Gest(*\{x\})$, where
		\[
			*\{x\} = \bigsqcup_{n \geq 0} */\Sigma_n
		\]
		is the free $E_{\infty}$-monoid (in animae) on one generator.
	\end{example}

	\begin{example}[Flat affine line]
		Over $\S$, or even over $\Z$, the free $E_{\infty}$-algebra $\Z\{x\} = \bigoplus_{n \geq 0} \Z[*/\Sigma_n]$  is \emph{not} isomorphic to the polynomial algebra $\Z[x]$; rather, it has contributions coming from the group homology of symmetric groups. After base change to $\Q$, the two versions of the polynomial algebra agree.

		Over a general base $X = \Gest(A)$, we will denote this other version of the affine line by
		\[
			\bbG_a := \Gest_X(A_0[x]),
		\]
		where $A_0[x] = \bigsqcup_{n \geq 0} A_0 \cdot x^n$ is the polynomial algebra over $A_0$. This is 0-affine, and hence $0$-proper and $1$-étale.

		Over the absolute base, we have $\bbG_a = \Gest(\N)$, the natural numbers.
	\end{example}

	\begin{example}[Multiplicative group]
		Consider the subfunctor
		\[
			\GL_1 \subseteq \bbA^1_X
		\]
		sending a Gestalt $Y \in \Gest_{/X}$ to the invertible elements $\Oo(Y)^{\times}_0$ of the commutative ring $\Oo(Y)_0$. This is representable by the Gestalt
		\[
			\GL_1 := \Gest_X(A_0\{x^{\pm}\}),
		\]
		the free $E_{\infty}$-algebra over $A_0$ on an invertible element $x$.

		There is again another version:
		\[
			\bbG_m := \GL_1 \times_{\bbA^1} \bbG_a = \Gest_X(A_0[x^{\pm1}]).
		\]
		Over the absolute base, this is $\Gest(\Z)$, where we regard $\Z$ as a ``multiplicative group'': it represents the functor taking $Y$ to $\Hom_{\CAlg(\An)}(\Z,\Oo(Y)_0)$.

		Over $\D(\Q)$ (or just $\D_{\geq 0}(\Q)$), $\GL_1$ and $\bbG_m$ are isomorphic.
	\end{example}

	\begin{example}[Projective line]
		We define the \emph{projective line} as the pushout
		\[
			\bbP^1_X := \bbA^1_X \sqcup_{\GL_1} \bbA^1_X
		\]
		in $\Gest_{/X}$. Here the two maps $\GL_1 \to \bbA^1$ are the maps represented by the algebra maps $A_0\{x\} \to A_0\{x^{\pm}\}$ given by $x \mapsto x$ and $x \mapsto x^{-1}$.

		As a pushout of 1-étale Gestalten, this is 1-étale. It is not 0-affine, but its diagonal is 0-affine: The product of two copies of $\bbA^1_X \sqcup_{\GL_1} \bbA^1_X$ is covered by four copies of $\bbA^1 \times \bbA^1$ and the preimage of the diagonal is either $\bbA^1$ embedded diagonally, or $\GL_1$ embedded via $x \mapsto (x,x^{-1})$, and both maps are $0$-affine as maps between $0$-affines.
	\end{example}

	While we may already define $\bbP^1_X$ over the absolute base, it turns out that before passing to the stable setting of $\D(\S)$, this is not well-behaved, and is different from the object one might expect. For concreteness, we explain the situation over a field $k$.

	\begin{proposition}
		\label[proposition]{prop:Counterexample_Zariski_Cover}
		Let $k$ be a field. Then the map
		\[
			\Gest_X(k[T^{\pm 1}]) \sqcup \Gest_X(k[T,\frac{1}{1-T}]) \to \Gest_X(k[T])
		\]
		over $X = \Gest(\D_{\geq 0}(k))$ is not surjective. After base change to $\Gest(\D(k))$, it becomes surjective.
	\end{proposition}
	\begin{proof}
		The fact that it is surjective after base change to $\D(k)$ follows from fpqc descent for the affine functor $\Spec(A) \mapsto \Oo(\Spec(A))$, assembled from \Cref{thm:fpqc_Hyperdescent_nPr}, together with \Cref{thm:Sheaf_Valued_In_Etale}. Indeed, after base change the displayed map is the image of the usual Zariski cover $D(T) \cup D(1-T) = \Spec(k[T])$, whose \v{C}ech nerve realizes to $\Spec(k[T])$ in $\Shv_{\fpqc}(\Aff_k)^{\aleph_1}$.

		We now show that the original map is not surjective. If the displayed map were surjective, then \Cref{prop:Surjective_Generation}, applied with $n=1$, would imply that the image of
		\[
			\D_{\geq 0}(k[T^{\pm 1}]) \times \D_{\geq 0}(k[T,\frac{1}{1-T}]) \to \D_{\geq 0}(k[T])
		\]
		generates $\D_{\geq 0}(k[T])$ under small colimits.

		However, $k[T] \in \D_{\geq 0}(k[T])$ receives no nonzero maps from any object in this image. Indeed, if $M \in \D_{\geq 0}(k[T^{\pm 1}])$, then every $k[T]$-linear map $M \to k[T]$ is detected on $\pi_0(M)$, and its image in $k[T]$ is divisible by all powers of $T$, hence lies in $\bigcap_{r \geq 0} T^r k[T] = 0$. Similarly, maps from objects of $\D_{\geq 0}(k[T,\frac{1}{1-T}])$ have image in $\bigcap_{r \geq 0} (1-T)^r k[T] = 0$. The full subcategory of objects mapping trivially to $k[T]$ is closed under colimits, so $k[T]$ cannot be generated by the displayed image.
	\end{proof}

	\begin{remark}
		After base change to $\Gest(\D(\S))$, the descent argument in the previous proof, together with \Cref{thm:fpqc_Hyperdescent_nPr,thm:Sheaf_Valued_In_Etale}, shows that
		\[
			\bbA^1 \sqcup_{\GL_1} \bbA^1
		\]
		comes from the obvious fpqc stack, namely the derived scheme obtained by gluing two affine lines along the invertible locus. In particular, it is $0$-proper. If one instead forms the version
		\[
			\bbG_a \sqcup_{\bbG_m} \bbG_a
		\]
		using ordinary polynomial algebras rather than free $E_{\infty}$-algebras, then the resulting Gestalt is not only $0$-proper but even $0$-suave, which is apparently not the case for the free $E_{\infty}$-version above. Over $\D(\Q)$, the two choices of polynomial algebra agree, so the distinction between these two projective lines disappears.
	\end{remark}

	For our last example we work over $\D(\Q)$.

	\begin{example}
		Consider the functor
		\[
			\pic(-)\colon (\Gest_{/\D(\Q)})\catop \to \An, X \mapsto \Oo(X)_1^{\times}.
		\]
		This is represented by a 2-étale Gestalt
		\[
			\Pic \in \Gest_{/\D(\Q)}.
		\]
	\end{example}

	\begin{construction}
		Over the base $X_0 := \Gest(\D(\Q))$, let
		\[
			\mu_2 := \Gest_{X_0}(\Q[x]/(x^2-1))
		\]
		denote the group of second roots of unity. Since $2$ is invertible in $\Q$, this is isomorphic to the constant sheaf $\underline{\Z/2}$.

		We construct a map
		\[
			\Pic \to \mu_2
		\]
		of sheaves of spectra on $\Gest_{/\D(\Q)}$, where we identify $\CGrp(\An)$-valued sheaves with $\Sp_{\geq 0}$-valued sheaves. Given an invertible object $L \in \Oo(X)^{\times}_1 = \pic(X)$, the swap map $c\colon L \otimes L \cong L \otimes L$ provides a map
		\[
			\unit \cong L^{-\otimes 2} \otimes (L \otimes L) \iso L^{-\otimes 2} \otimes (L \otimes L) \cong \unit,
		\]
		which provides an element of $\mu_2$.
	\end{construction}

	\begin{proposition}
		\label[proposition]{prop:Relation_Pic_BG_m}
		There is a fiber sequence of sheaves of spectra on $\Gest_{/\D(\Q)}$ of the form
		\[
			*/\bbG_m \to \Pic \to \mu_2.
		\]
	\end{proposition}
	\begin{proof}
		The map $\Pic \to \mu_2$ is surjective, since the $\Q$-module $\Q[1]$ has $-1$ as a swap. In particular, we get a fiber sequence
		\[
			\fib(\Pic \to \mu_2) \to \Pic \to \mu_2.
		\]
		We will now show that the map $* \to \fib(\Pic \to \mu_2)$ is surjective, i.e.\ that any $L \in \Oo(X)^{\times}_1$ such that the swap map $c\colon L \otimes L \to L \otimes L$ is trivial is already locally trivial. Write $A = \Oo(X)$. The action of $\Sigma_n$ on $L^{\otimes n}$ is uniquely trivial for all $n \geq 1$. Indeed, this action is given by a map
		\[
			\Sigma_n \to \Aut(L^{\otimes n}) \simeq \Aut(\unit) = A_0^{\times}.
		\]
		Since $X$ lies over $\D(\Q)$, the commutative algebra $A_0$ refines to an object of $\CAlg(\D_{\geq 0}(\Q))$. Thus $A_0^{\times}$ is a connective $E_{\infty}$-group whose higher homotopy groups are $\Q$-vector spaces. The induced map $\Sigma_n \to \pi_0(A_0^{\times})$ is trivial because transpositions act trivially by assumption, and transpositions generate $\Sigma_n$. Since $\Sigma_n$ is finite and the higher homotopy groups are rational vector spaces, there are no higher obstructions or choices, so the map $\Sigma_n \to A_0^{\times}$ is uniquely nullhomotopic.

		We may therefore form the commutative algebra
		\[
			B := \bigoplus_{n \in \Z} L^{\otimes n} \qin \CAlg(A_1),
		\]
		with multiplication induced by the identifications $L^{\otimes m} \otimes L^{\otimes n} \simeq L^{\otimes(m+n)}$. Tensoring with $L$ shifts the grading, so
		\[
			L \otimes B \simeq B
		\]
		in $\Mod_B(A_1)$. Hence $L$ becomes trivial after base change along the $0$-affine map $\Gest_X(B) \to X$. Moreover, the inclusion $\unit \to B$ of the degree-zero summand splits in $A_1$ by projection onto degree zero, so $B \otimes -$ is conservative. By the conservative criterion of \Cref{prop:Surjective_Criteria}, the map $\Gest_X(B) \to X$ is a cover. This proves that $* \to \fib(\Pic \to \mu_2)$ is surjective.

		Finally, $* \times_{\Pic} *$ classifies automorphisms of the trivial line bundle, which is precisely $\bbG_m$. Since $* \to \fib(\Pic \to \mu_2)$ is surjective, its \v{C}ech nerve presents the fiber as the quotient $*/\bbG_m$, giving the desired identification $\fib(\Pic \to \mu_2) \simeq */\bbG_m$.
	\end{proof}

	\begin{remark}
		We have $\pi_0\Pic(\D(\Q)) = \Z$, given by the shifts $\Q[n]$ for $n \in \Z$. However, the proof shows that the classes with $n$ even are locally trivial in Gestalten. To trivialize them, one has to pass to the periodic $E_{\infty}$-rings
		\[
			\bigoplus_{m \in \Z} \Q[nm].
		\]
		Only for even $n$ does this admit a commutative multiplication.

		We will later study higher versions of this story, comparing $B^n\bbG_m$ with the moduli space of invertible objects of $A_n$. In the terminology used in the literature, this studies ``invertible topological quantum field theories''; these are expected to be classified by a spectrum essentially given by an Anderson dual of the sphere. From this viewpoint, the $\mu_2$ above should be seen as dual to the first stable homotopy group of spheres, $\pi_1 \S = \Z/2$. We will in fact see that the moduli space of invertible objects in the $A_n$ defines a sheaf of spectra on Gestalten which is an internal Anderson dual of the sphere, proving a precise version of this expectation.
	\end{remark}

	\section{Examples from condensed and analytic mathematics}

	So far our examples have started from commutative algebras in simple symmetric monoidal base categories, such as $\D(\S)$ or $\D(k)$, producing $0$-affine Gestalten, from which more general ones were obtained by gluing, i.e.\ by taking colimits. We now turn to examples of a more analytic flavour, coming from condensed and analytic mathematics.

	\begin{construction}[Light condensed animae]
		\label[construction]{con:Condensed_Gestalten}
		Generalizing the finite and countable sets considered above, we work over $X = \Gest(\D(\S))$ and send a light profinite set $S = \lim_i S_i$ to the Gestalt
		\[
			\Gest_{\Gest(\D(\S))}\bigl(\Cont(S,\S)\bigr) = \lim_i \bigl(S_i \times \Gest(\D(\S))\bigr),
		\]
		where $\Cont(S,\S)$ is the Stefanich ring of continuous $\S$-valued functions on $S$, and $S_i \times \Gest(\D(\S))$ denotes the finite coproduct of copies of the base indexed by $S_i$. This could already be defined over the absolute base $\Gest(\An)$; for \Cref{prop:Internal_Real_Numbers} below it is important to work stably. By hyperdescent (\Cref{thm:fpqc_Hyperdescent_nPr}), this assignment takes hypercovers to colimits in Gestalten, so the functor
		\[
			\ProFin^{\mathrm{light}} \to \Gest^{0\text{-affine}}_{/\Gest(\D(\S))}
		\]
		extends to a countable-colimit- and finite-limit-preserving functor
		\[
			\CondAn^{\mathrm{light}} \to \Gest^{1\text{-}\et}_{/\Gest(\D(\S))}.
		\]
	\end{construction}

	As an important example, we can identify the internal real number object.

	\begin{proposition}[{\citet[Prop.~9.1]{Scholze_Gestalten}}]
		\label[proposition]{prop:Internal_Real_Numbers}
		The internal real number object of Cauchy reals in $\Gest_{/\Gest(\D(\S))}$ is given by the image of the real numbers $\R \in \CondAn^{\mathrm{light}}$ under \Cref{con:Condensed_Gestalten}, and is $1$-affine, equal to
		\[
			\Gest\bigl(\Shv(\R, \D(\S))\bigr).
		\]
	\end{proposition}
	\begin{proof}
		The internal construction of light profinite sets yields the functor $\ProFin^{\mathrm{light}} \to \Gest$ used in \Cref{con:Condensed_Gestalten}. Under its extension to light condensed animae, the real numbers $\R \in \CondAn^{\mathrm{light}}$ are sent to $\Gest(\Shv(\R, \D(\S)))$, using descendability of $f_{\S}$ for a locally finite cover $f\colon C \to \R$ of $\R$ by a disjoint union of Cantor sets.\footnote{This can be checked most cleanly using a descendability criterion for sheaves due to Clausen.}

		Now let $\mathrm{Cauchy}$ denote the internal object of Cauchy sequences of rational numbers with fixed convergence, i.e.\ sequences $(x_1, x_2, \dots)$ of rational numbers with $|x_i - x_j| \leq 1/\min(i,j)$ for all $i$ and $j$. This is closed inside $\prod_{\N} \Q$, and the latter is $1$-affine as a countable limit of $1$-affine Gestalten; in fact, $\mathrm{Cauchy}$ is precisely the Gestalt corresponding to the category of $\D(\S)$-valued sheaves on the topological space of Cauchy sequences. Using this, one constructs a map
		\[
			\mathrm{Cauchy} \to \Gest\bigl(\Shv(\R, \D(\S))\bigr).
		\]
		This is surjective, as the cover $f\colon C \to \R$ above factors over $\mathrm{Cauchy}$. Moreover, the preimage of $0 \in \R$ is precisely the internal object $\mathrm{Null} \subseteq \mathrm{Cauchy}$ of nullsequences (automatically with fixed convergence, $|x_i| \leq 1/i$). Thus the induced map
		\[
			\mathrm{Cauchy}/\mathrm{Null} \to \Gest\bigl(\Shv(\R, \D(\S))\bigr)
		\]
		is an isomorphism, as desired.
	\end{proof}

	\begin{remark}
		The internal real numbers object can also be defined over the absolute base. The proof shows that it is then given by the image in Gestalten of
		\[
			\Gest\bigl(\Shv(\mathrm{Cauchy}, \An)\bigr) \to \Gest\bigl(\Shv(\R, \An)\bigr),
		\]
		where $\mathrm{Cauchy}$ now denotes the topological space of Cauchy sequences with fixed convergence. It is not known whether this map is surjective before passing to $\D(\S)$. In the absence of surjectivity, the image is a priori only a $2$-étale injection into $\Gest(\Shv(\R, \An))$, and hence might only be $3$-affine.
	\end{remark}

	One can build more general examples by allowing more general presentable symmetric monoidal categories as coefficients. We mention two such families.

	\begin{example}[Almost rings and Tamarkin's category]
		Almost rings yield examples of Gestalten, thereby incorporating ``almost mathematics''. A related example is Tamarkin's category $W$ of completely and continuously $\R$-filtered $\S$-modules. The map
		\[
			\Gest(W) \to \Gest(\D(\S))
		\]
		is an $\R_{>0}$-torsor, where $\R_{>0} = \Gest(\Shv(\R_{>0}, \D(\S)))$ denotes the internal positive real number object, acting on $W$ by multiplicatively rescaling the filtration. In other words, working over $\Gest(\D(\S))$ there is a nontrivial map
		\[
			* \to */\R_{>0},
		\]
		exhibiting highly nontrivial internal structure on the base point.
	\end{example}

	\begin{example}[Analytic rings and analytic stacks]
		Analytic rings yield examples of Gestalten. Any analytic ring $A$, consisting of a light condensed ring together with a chosen full subcategory $\D(A)$ of ``complete'' modules, defines a symmetric monoidal category $\D(A)$ and hence a Stefanich ring; this assembles into a functor
		\[
			\AnRing \to \StRing, \qquad A \mapsto (\Omega^{\infty}\End_{\D(A)}(\unit), \D(A), 1\Pr_{\D(A)}, 2\Pr_{\D(A)}, \dots).
		\]
		Analytic stacks are defined precisely so that this extends to a functor
		\[
			\AnStck \to \Gest.
		\]
		In general, the resulting Gestalten are only colimits of $1$-affine Gestalten, hence only $2$-étale. In practice, after fixing a base analytic ring $A$, the stacks under consideration are $1$-étale over $\Gest(\D(A))$.
	\end{example}

	All of the examples above are colimits of $1$-affine Gestalten. A natural source of interesting Gestalten that go \emph{beyond} such colimits are six-functor formalisms, which we turn to in the next chapter.

	% From here on the chapter counter runs one ahead of the displayed talk number,
	% because the ninth talk occupies two chapters (9a, 9b) above. Display = counter - 1.
	\renewcommand{\thechapter}{\the\numexpr\value{chapter}-1\relax}
	\chapter{Six-functor formalisms and Gestalten}
	\label[chapter]{chap:Transmutations}

	\emph{Talk given by Ou. The material from Sections \ref{sec:Translation}, \ref{sec:Extension_To_Stacks} and \ref{sec:6FF_On_Gest} were not treated during the seminar, and has been added afterwards for completeness.}

	Today we discuss how to assign a Stefanich ring $A_{\D}$ to a six-functor formalism $\D$, resulting in a functor $[-]_{\D}\colon C \to \Gest$ called the \emph{transmutation} associated to $\D$.

	Recall that a \emph{six-functor formalism} on a category $C$ is a lax symmetric monoidal functor
	\[
		\D\colon \Span(C, E) \to 1\Pr,
	\]
	assigning to each $X \in C$ a presentably symmetric monoidal category $\D(X)$ of coefficients, to each map $f$ an adjoint pair $f^* \dashv f_*$, and to each map $f$ in a wide subcategory $E$ of $C$ closed under base change an adjunction $f_! \dashv f^!$, compatible via base change and the projection formula (see \Cref{chap:Six_Functor_Formalisms} for some recollections). As was observed in \citep[Lecture~X]{Scholze_Six_Functors}, it is an empirical fact that many six-functor formalisms on the category $C = \Sch$ of schemes (separated of finite type, say) can be factored as a composite
	\[
		\Span(\Sch) \xrightarrow{X \mapsto X^{?}} \Span(\AnStck) \xrightarrow{\D_{\qc}} 1\Pr,
	\]
	where $\AnStck$ is the category of analytic stacks and $\D_{\qc}$ is a six-functor formalism on $\AnStck$ given by quasi-coherent sheaves. The functor $X \mapsto X^{?}$ from schemes to analytic stacks is called a \emph{transmutation}, following \citep[Remark~2.3.8]{Bhatt2022FGauges}. Here are some examples:
	
	\begin{enumerate}
		\item For $\D(\Z)$-valued sheaves we have $\D(X,\Z) \simeq \D_{\qc}(X_{\mathrm{Betti}}),$ where $X_{\mathrm{Betti}}$ is the Betti stack of $X$.
		\item For algebraic D-modules, we have $D\mathrm{mod}(X) \simeq \D_{\qc}(X_{\mathrm{dR}}),$ where $X_{\mathrm{dR}}$ is the de Rham stack of $X$.
		\item There is an assignment $X \mapsto X^{\Delta_A}$ encoding prismatic cohomology.
		\item There is an assignment $X \mapsto X^{\square}_{\mathrm{dR}}$ encoding the solid de Rham D-module formalism.
	\end{enumerate}
	
	The transmutation perspective has the virtue of black-boxing all the category theory to the theory of analytic stacks, and reducing the construction of six-functor formalisms to a geometric construction $X \mapsto X^{?}$. (In fact, the situation is even better: the functor $X \mapsto X^{?}$ often commutes with finite limits and is compatible with Zariski gluing, which means that it is completely determined by the ring stack $(\bbA^1)^{?}$. See \Cref{sec:Ring_Gestalten} for details.)

	The goal of the talk is to develop an analogous story of transmutations that works for six-functor formalisms over an arbitrary geometric category $C$ in place of schemes, where we replace the category of analytic stacks with the category $\Gest$ of Gestalten. More precisely, given a presentable six-functor formalism $\D\colon \Span(C,E) \to 1\Pr$, we will define a finite-limit-preserving transmutation functor $[-]_{\D}\colon C \to \Gest_{/[*]_{\D}}$ fitting in a triangle
	\[
	\begin{tikzcd}[column sep=huge]
		C\catop \rar[dashed,"{[-]_{\D}}"] \arrow[dr,swap,"{\D^{*,\otimes}}"] & \Gest\catop \dar["{\Oo(-)_1}"] \\
		& \CAlg(1\Pr),
	\end{tikzcd}
	\]
	where $\D^{*,\otimes}$ is the $(f^*,\otimes)$-part of $\D$. Moreover, we will show that the Gestalt $[X]_{\D}$ is $1$-affine over $[*]_{\D}$ for every $X \in C$ if and only if $\D$ satisfies the categorical Künneth formula, in the sense that the canonical functor
	\[
		\D(X) \otimes_{\D(*)} \D(Y) \xrightarrow{\boxtimes} \D(X \times Y)
	\]
	is an equivalence for all $X,Y \in C$. We will also construct a six-functor formalism on $\Gest$ and discuss when also the shriek-functoriality of $\D$ can be recovered from $[-]_{\D}$.

	\section{The category of kernels \texorpdfstring{$\bbK_{\D}$}{K\_D}}
	\label[section]{sec:Category_Of_Kernels}

	Fix a category $C$ with finite limits and a presentable six-functor formalism $\D\colon \Span(C) \to 1\Pr$. In this section, we introduce and study a presentable 2-category $\bbK_{\D}$, which is the main ingredient for the construction of the transmutation functor $[-]_{\D}$.

	\begin{definition}[{\citet[Def.~2.4]{Aoki_PhD_Thesis}}]\label[definition]{def:Category_Of_Kernels}
		Let $C$ be a category with finite limits. Regarding the six-functor formalism $\D$ as a commutative algebra in $\Fun(\Span(C),1\Pr)$ with the Day convolution symmetric monoidal structure, the \emph{$2$-category of kernels of $\D$} is the module category
		\[
			\bbK_{\D} := \Mod_{\D}(\Fun(\Span(C),1\Pr)).
		\]
	\end{definition}

	This comes with a symmetric monoidal functor from $\Span(C)\catop$: the Yoneda embedding followed by the free $\D$-module functor:
	\[
		\Span(C)\catop \to \Fun(\Span(C),1\Pr) \to \bbK_{\D}, \qquad X \mapsto [X]_{\bbK_{\D}}.
	\]
	Since $[X]_{\bbK_{\D}}$ is the free $\D$-module on the representable at $X$, the free-forgetful adjunction and the Yoneda lemma give for each $\D$-module $M\in\bbK_{\D}$ an equivalence
	\[
		\Hom_{\bbK_{\D}}([X]_{\bbK_{\D}}, M) \iso M(X).
	\]
	In particular, as the unit of $\bbK_{\D}$ is $\D$ itself, we see that
	\[
		\Hom_{\bbK_{\D}}([X]_{\bbK_{\D}}, \unit) \iso \D(X).
	\]
	Every object of $\Span(C)$ is canonically self-dual \citep[2.3.9]{Heyer_Mann}, and the embedding is symmetric monoidal, so $[Y]_{\bbK_{\D}}\simeq[Y]_{\bbK_{\D}}^{\vee}$ and $[X]_{\bbK_{\D}}\otimes[Y]_{\bbK_{\D}}\simeq[X\times Y]_{\bbK_{\D}}$; hence
	\[
		\Hom_{\bbK_{\D}}([X]_{\bbK_{\D}}, [Y]_{\bbK_{\D}}) \simeq \Hom_{\bbK_{\D}}([X\times Y]_{\bbK_{\D}}, \unit) = \D(X\times Y).
	\]
	With more care, one may show that the composition functor in $\bbK_{\D}$ between representables is given by the formula
	\[
		\D(Y \times Z) \times \D(X \times Y) \to \D(X \times Z), \qquad (B, A) \mapsto \pr_{XZ,!}(\pr_{XY}^*A \otimes \pr_{YZ}^*B).
	\]

	Note that this definition of the 2-category of kernels is different from the one that usually appears in the literature, e.g.\ in \citep{Heyer_Mann}: Usually, one considers the canonical self-enrichment of $\Span(C)$ via its closed monoidal structure, and then transfers the enrichment via the lax monoidal functor $\D\colon \Span(C) \to \Cat$. (See \Cref{def:Kernel_2_Category}.) The following result shows that $\bbK_{\D}$ is the $1\Pr$-enriched presheaves on this 2-category.

	\begin{theorem}[{\citet[Thm.~2.6]{Aoki_PhD_Thesis}}]\label[theorem]{thm:Aoki_Enriched_Presheaves}
		Let $S$ be a small closed symmetric monoidal category and $\Vv \in \CAlg(\PrL)$, and let $\D\colon S \to \Vv$ be a lax symmetric monoidal functor. Then there is a canonical equivalence
		\[
			\PSh^{\Vv}(\widetilde{S}\catop) \simeq \Mod_{\D}(\Fun(S,\Vv)),
		\]
		where $\widetilde{S}$ is the $\Vv$-enriched category obtained by base changing the self-enrichment of $S$ along $\D$.
	\end{theorem}
	\begin{proof}
		The argument proceeds by monadicity.

		Write $U\colon\An\to\Vv$ for the unit functor. The \emph{trivial} $\Vv$-enrichment of $S$ is the enriched category $S^{\mathrm{triv}}$ with the same objects as $S$ and
		\[
			\uHom_{S^{\mathrm{triv}}}(a,b) := U(\Hom_S(a,b)).
		\]
		In this case, $\Vv$-enriched presheaves are given by
		\[
			\PSh^{\Vv}((S^{\mathrm{triv}})\catop) \simeq \Fun(S,\Vv).
		\]
		On the other hand, the closed symmetric monoidal structure on $S$ makes it enriched over itself, with enriched hom $\uHom_S(a,b)\in S$, and change of enrichment along $\D$ gives the $\Vv$-enrichment
		\[
			\uHom_{\widetilde S}(a,b) := \D(\uHom_S(a,b)).
		\]

		Equip $\Fun(S,\Vv)$ with Day convolution. Its monoidal unit is the functor
		\[
			\unit_{\mathrm{Day}}(s) = U(\Hom_S(\unit_S,s)).
		\]
		The lax symmetric monoidal structure on $\D$ is equivalently a commutative algebra structure on $\D\in\Fun(S,\Vv)$, and its unit $\unit_{\mathrm{Day}}\to\D$ induces an identity-on-objects $\Vv$-enriched functor
		\[
			i\colon S^{\mathrm{triv}}\to\widetilde S,
		\]
		which on enriched homs is given by the map
		\[
			U(\Hom_S(a,b))
			\simeq U(\Hom_S(\unit_S,\uHom_S(a,b)))
			\longrightarrow \D(\uHom_S(a,b)).
		\]
		Put
		\[
			\mathcal C:=\Fun(S,\Vv), \qquad \mathcal P:=\PSh^{\Vv}(\widetilde S\catop).
		\]
		Restriction and enriched left Kan extension along $i$ give an adjunction
		\[
			i_!\colon \mathcal C \rightleftarrows \mathcal P\colon i^*.
		\]
		Since $i$ is the identity on objects, $i^*$ is conservative. Colimits of enriched functors are computed objectwise, so $i^*$ also preserves all colimits. The Barr--Beck theorem therefore identifies
		\[
			\mathcal P \simeq \Mod_T(\mathcal C), \qquad T:=i^*i_!.
		\]

		The category $\mathcal P$ carries the enriched Day convolution symmetric monoidal structure. Its unit is the enriched functor represented by $\unit_S$, and hence
		\[
			i^*(\unit_{\mathcal P})(s)
			= \uHom_{\widetilde S}(\unit_S,s)
			= \D(\uHom_S(\unit_S,s))
			\simeq \D(s).
		\]
		In fact, this is an identification $i^*(\unit_{\mathcal P})\simeq\D$ of commutative algebras in $\mathcal C$: the algebra structure on the left is induced by the lax symmetric monoidal structure on $i^*$, which by construction is the lax symmetric monoidal structure on $\D$.

		It follows formally that $i^*$ lifts to a functor
		\[
			\widetilde{i^*}\colon \mathcal P\to\Mod_{\D}(\mathcal C),
		\]
		where for $Y\in\mathcal P$ the action map on $i^*Y$ is given by the lax monoidal structure on $i^*$:
		\[
			\D\otimes i^*Y
			\simeq i^*(\unit_{\mathcal P})\otimes i^*Y
			\longrightarrow i^*(\unit_{\mathcal P}\otimes Y)
			\simeq i^*Y.
		\]
		Since $\widetilde{i^*}$ lies over $i^*$, it determines a canonical morphism of monads
		\[
			\alpha\colon (-)\otimes\D\longrightarrow T.
		\]
		Concretely, its component at $P\in\mathcal C$ is the projection map
		\[
			\alpha_P\colon
			P\otimes\D
			\simeq P\otimes i^*(\unit_{\mathcal P})
			\longrightarrow i^*(i_!P\otimes\unit_{\mathcal P})
			\simeq i^*i_!P.
		\]
		The fact that $\alpha$ is a morphism of monads, including all of its higher coherences, is built into its construction from the lift $\widetilde{i^*}$.

		It remains only to show that $\alpha$ is an equivalence. Both $T=i^*i_!$ and $(-)\otimes\D$ preserve colimits and are $\Vv$-linear. Moreover, $\mathcal C=\PSh^{\Vv}((S^{\mathrm{triv}})\catop)$ is generated under colimits and the $\Vv$-action by the enriched representables
		\[
			h_a:=\uHom_{S^{\mathrm{triv}}}(a,-)=U(\Hom_S(a,-)), \qquad a\in S.
		\]
		It therefore suffices to prove that $\alpha_{h_a}$ is an equivalence for every $a\in S$. Enriched left Kan extension takes the representable $h_a$ to the corresponding representable of $\widetilde S$, and hence
		\[
			T(h_a)(s)=i^*i_!(h_a)(s)\simeq\uHom_{\widetilde S}(a,s)=\D(\uHom_S(a,s)).
		\]
		On the other hand, the Day-convolution Yoneda formula, together with the adjunction $a\otimes-\dashv\uHom_S(a,-)$, gives a canonical equivalence
		\[
			(h_a\otimes\D)(s)\simeq\D(\uHom_S(a,s)).
		\]
		Under these identifications, $\alpha_{h_a}$ is the identity of the enriched representable $\uHom_{\widetilde S}(a,-)$, and is therefore an equivalence. We conclude that $\alpha$ is an equivalence of monads, and hence
		\[
			\PSh^{\Vv}(\widetilde S\catop)
			=\mathcal P
			\simeq \Mod_T(\mathcal C)
			\simeq \Mod_{\D}(\mathcal C)
			=\Mod_{\D}(\Fun(S,\Vv)),
		\]
		as desired.
	\end{proof}

	\begin{corollary}
		The category of kernels $\bbK_{\D}$ from \Cref{def:Category_Of_Kernels} is the $1\Pr$-enriched presheaves on the 2-category of kernels $K_{\D}$ from \Cref{def:Kernel_2_Category}.
	\end{corollary}
	\begin{proof}
		Take $\Vv = 1\Pr$ and $S = \Span(C)$.
	\end{proof}

	We will now show that the category of kernels is not very interesting if $\D$ satisfies the categorical Künneth formula: $\bbK_{\D}$ is simply the module category $\Mod_{\D(*)}(1\Pr)$ in that case. We start by constructing a comparison functor.

	\begin{construction}\label[construction]{con:Restricted_Objects_Functor}
		The monoidal unit of $\bbK_{\D}$ is $\unit=[*]_{\bbK_{\D}}=\D$. Taking $X=*$ in the formula
		\[
			\Hom_{\bbK_{\D}}([X]_{\bbK_{\D}},M)=M(X)
		\]
		gives
		\[
			\End_{\bbK_{\D}}(\unit)=\D(*).
		\]
		The resulting action of $\D(*)$ on $\unit$ defines an induction functor
		\[
			L_{\D}\colon\Mod_{\D(*)}(1\Pr)\to\bbK_{\D},
			\qquad N\longmapsto N\otimes_{\D(*)}\unit.
		\]
		It is left adjoint to evaluation at the point,
		\[
			G_{\D}\colon\bbK_{\D}\to\Mod_{\D(*)}(1\Pr),
			\qquad M\longmapsto M(*)=\Hom_{\bbK_{\D}}(\unit,M),
		\]
		where the $\D(*)$-action on $M(*)$ is given by precomposition. Indeed, the adjunction $L_{\D}\dashv G_{\D}$ is the tensor-hom adjunction for the $\D(*)$-module object $\unit\in\bbK_{\D}$.
	\end{construction}

	\begin{proposition}
		\label[proposition]{prop:Restricted_Objects}
		The induction functor $L_{\D}\colon\Mod_{\D(*)}(1\Pr)\to\bbK_{\D}$ of \Cref{con:Restricted_Objects_Functor} is fully faithful, with essential image the \emph{restricted} objects: $M\in\bbK_{\D}$ lies in it if and only if $\D(X)\otimes_{\D(*)}M(*)\to M(X)$ is an equivalence for all $X$.
	\end{proposition}
	\begin{proof}
		Work with $\bbK_{\D} = \Mod_{\D}(\Fun(\Span(C),1\Pr))$ and the formula $\Hom_{\bbK_{\D}}([X]_{\bbK_{\D}}, M) = M(X)$ of \Cref{def:Category_Of_Kernels}.

		\emph{Evaluation is atomic.} Recall that an object $A$ of a $1\Pr$-module category $\mathcal C$ is \emph{$1\Pr$-atomic} if the $1\Pr$-valued functor $\Hom_{\mathcal C}(A,-)$ preserves colimits and $1\Pr$-tensors. The formula above identifies the corepresentable of $[X]_{\bbK_{\D}}$ with the evaluation $\mathrm{ev}_X\colon M\mapsto M(X)$. Colimits in $\bbK_{\D}$ are created by the forgetful functor to $\Fun(\Span(C),1\Pr)$, where they are pointwise, so $\mathrm{ev}_X$ preserves colimits and is $1\Pr$-linear; thus each $[X]_{\bbK_{\D}}$ is $1\Pr$-atomic. In particular $\unit = [*]_{\bbK_{\D}}$ is atomic, with $\End_{\bbK_{\D}}(\unit) = \mathrm{ev}_*(\unit) = \D(*)$.

		\emph{Full faithfulness.} By \Cref{con:Restricted_Objects_Functor}, we have an adjunction $L_{\D}\dashv G_{\D}$. The functor $L_{\D}$ is fully faithful if and only if the unit $N\to G_{\D}L_{\D}(N)$ is an equivalence. As evaluation at $*$ is colimit-preserving and $1\Pr$-linear, it commutes with $N\otimes_{\D(*)}(-)$, so
		\[
			G_{\D}L_{\D}(N) = (N\otimes_{\D(*)}\unit)(*) \simeq N\otimes_{\D(*)}\unit(*) = N\otimes_{\D(*)}\D(*) = N,
		\]
		using $\unit(*) = \D(*)$. Hence $G_{\D}L_{\D}\simeq\id$.

		\emph{Essential image.} As $L_{\D}$ is fully faithful with right adjoint $G_{\D}$, an object $M$ lies in its image if and only if the counit $L_{\D}G_{\D}(M)\to M$ is an equivalence. Evaluating $L_{\D}G_{\D}(M) = M(*)\otimes_{\D(*)}\unit$ at $X$ (again using that $\mathrm{ev}_X$ commutes with $M(*)\otimes_{\D(*)}(-)$, and $\unit(X) = \D(X)$) gives $\D(X)\otimes_{\D(*)}M(*)$, and the counit at $X$ is the canonical map $\D(X)\otimes_{\D(*)}M(*)\to M(X)$. So $M$ is restricted if and only if this is an equivalence for all $X$.
	\end{proof}

	\begin{corollary}
		\label[corollary]{cor:Categorical_Kunneth_Iff_Kernel_Category_Is_Modules}
		The inclusion $L_{\D}\colon \Mod_{\D(*)}(1\Pr) \hookrightarrow \bbK_{\D}$ is an equivalence if and only if $\D$ satisfies the categorical Künneth formula: for all objects $X,Y \in C$, we have
		\[
			\D(X)\otimes_{\D(*)}\D(Y)\xrightarrow{\ \sim\ }\D(X\times Y).
		\]
	\end{corollary}
	\begin{proof}
		The objects in the image of $L_{\D}$ are closed under colimits and $1\Pr$-tensoring, and $\bbK_{\D}$ is $1\Pr$-atomically generated by the $[Y]_{\bbK_{\D}}$, so $L_{\D}$ is an equivalence iff every $[Y]_{\bbK_{\D}}$ is restricted. By the $\Hom_{\bbK_{\D}}$-formula, $[Y]_{\bbK_{\D}}(*) = \D(Y)$ and $[Y]_{\bbK_{\D}}(X) = \D(X\times Y)$, so this holds iff $\D(X)\otimes_{\D(*)}\D(Y)\to\D(X\times Y)$ is an equivalence for all $X,Y$, which is the categorical Künneth formula.
	\end{proof}

	\section{Transmutation}
	\label[section]{sec:Transmutation}

	Our next goal is to use the category of kernels construction to define the transmutation functor associated to $\D$. For each object $X \in C$, we may form the category of kernels associated to the restricted six-functor formalism
	\[
		\D_X \colon \Span(C_{/X}) \to \Span(C) \xrightarrow{\D} 1\Pr,
	\]
	giving rise to 
	\[
		\bbK_{\D,X} := \Mod_{\D_X}(\Fun(\Span(C_{/X}),1\Pr)) \in \CAlg(2\Pr).
	\]
	For each map $f\colon Y \to X$ in $C$, left Kan extension along the symmetric monoidal pullback functor $\Span(f^*) \colon \Span(C_{/X}) \to \Span(C_{/Y})$ induces a symmetric monoidal restriction functor
	\[
		\Span(f^*)_!\colon \Fun(\Span(C_{/X}),1\Pr) \to \Fun(\Span(C_{/Y}), 1\Pr),
	\]
	which by construction sends $\D_X$ to $\D_Y$ and hence induces a symmetric monoidal functor
	\[
		f_1^*\colon \bbK_{\D,X}\to\bbK_{\D,Y}.
	\]
	All these constructions are functorial in $f$, resulting in a functor
	\[
		C\catop\to\CAlg(2\Pr).
	\]

	\begin{proposition}[{\citet[Prop.~9.4]{Scholze_Gestalten}}]\label[proposition]{prop:Adjoints_Category_Of_Kernels}
		For every morphism $f\colon Y\to X$, the functor $f_1^*\colon \bbK_{\D,X}\to\bbK_{\D,Y}$ admits canonically equivalent left and right adjoints, $f_{1,\sharp}\simeq f_{1,*}$. They satisfy the projection formula and their formation is compatible with base change.
	\end{proposition}
	\begin{proof}
		For the duration of the proof, write (and use the same notation for all other maps)
		\[
			f_!\colon C_{/Y}\to C_{/X}
		\]
		for postcomposition with $f$. It is left adjoint to the base-change functor $f^*\colon C_{/X}\to C_{/Y}$. The naturality squares of the unit and counit of $f_!\dashv f^*$ are cartesian: those for the counit are the usual base-change squares for the projections $W\times_XY\to W$, while those for the unit are the corresponding squares for the graph maps $Z\to Z\times_XY$. Applying both parts of \citep[Corollary~C.21]{BachmannHoyois2021Norms} therefore gives a pair of adjunctions
		\[
			\Span(f_!)\dashv\Span(f^*)\dashv\Span(f_!).
		\]
		Passing to $1\Pr$-enriched presheaves by precomposition reverses both adjunctions, giving
		\[
			\Span(f^*)^*
			\dashv
			\Span(f_!)^*
			\dashv
			\Span(f^*)^*.
		\]
		Here the middle functor goes from $\Fun(\Span(C_{/X}),1\Pr)$ to $\Fun(\Span(C_{/Y}),1\Pr)$. It is symmetric monoidal: the base-change functor $f^*$ preserves finite products, so $\Span(f^*)$ is symmetric monoidal, and the adjunction $\Span(f^*)\dashv\Span(f_!)$ identifies $\Span(f_!)^*$ with left Kan extension along $\Span(f^*)$. To see that this adjunction passes further to module categories over $\D_X$ and $\D_Y$, we need to verify the projection formula
		\[
			\Span(f^*)^*N\otimes M
			\xrightarrow{\ \sim\ }
			\Span(f^*)^*\bigl(N\otimes\Span(f_!)^*M\bigr).
		\]
		It is enough to check this when $M$ and $N$ are representable, since both sides preserve colimits and $1\Pr$-tensors in both variables. For objects $A\to X$ and $B\to Y$, the assertion reduces to the canonical isomorphism
		\[
			f_!\bigl(B\times_Y f^*A\bigr)
			\simeq
			f_!B\times_XA,
		\]
		which follows directly from the defining pullback square for $f^*A$. All in all, we obtain the desired pair of adjunctions
		\[
			f_{1,\sharp} \dashv f_1^* \dashv f_{1,*},
		\]
		where by construction we have an identification $f_{1,\sharp} = f_{1,*}$.
		The projection formula exhibits the adjoint $\Span(f_!)^*$ as a $\D_X$-linear functor, so the two adjunctions pass to the corresponding module categories and give the projection formula for $f_{1,\sharp}$ and $f_{1,*}$.
		
		It remains to show that these adjoints are compatible with base change. Consider a cartesian square
		\[
			\begin{tikzcd}
				Y' \rar["f'"] \dar[swap,"g'"] \drar[pullback] & X' \dar["g"] \\
				Y \rar["f"'] & X.
			\end{tikzcd}
		\]
		There is a canonical equivalence $f^*g_!\simeq g'_!f'^*$ between functors on slice categories. Passing successively to span categories, enriched presheaves, and modules gives
		\[
			g_1^*f_{1,*}\simeq f'_{1,*}g_1'^*.
		\]
		Under the canonical equivalences $f_{1,\sharp}\simeq f_{1,*}$ and $f'_{1,\sharp}\simeq f'_{1,*}$, this is also the base-change equivalence for the left adjoints.
	\end{proof}

	The category of kernels already carries nontrivial geometry. The assignment
	\[
		C \to \Gest, \qquad X \mapsto \Gest(\bbK_{\D,X}),
	\]
	takes values in the $2$-affine Gestalten, and \Cref{prop:Adjoints_Category_Of_Kernels} shows that for each map $f\colon X \to Y$ the induced map $\Gest(\bbK_{\D,X}) \to \Gest(\bbK_{\D,Y})$ is $2$-proper and $1$-prim. It is not $1$-suave in general, though: the missing left adjoint $f_{2,\sharp}$ exists precisely when the assignment $X \mapsto \Gest(\bbK_{\D,X})$ satisfies the Künneth formula, which it need not. To force the Künneth formula to hold (and thereby obtain maps that are $1$-proper and $1$-étale), we iterate the category of kernels construction.

	\begin{construction}[Iterated categories of kernels]
		\label[construction]{con:Iterated_Categories_Of_Kernels}
		Set $\D^{(1)}:=\D$. By \Cref{prop:Adjoints_Category_Of_Kernels}, the functors associated with a map $f\colon Y\to X$ have adjoints
		\[
			f_{1,\sharp}\simeq f_{1,*}\colon \bbK_{\D,Y}\longrightarrow\bbK_{\D,X}
		\]
		that are compatible with base change and projection formulas. Together with the lax symmetric monoidal structure of the kernel construction, these adjunctions assemble the assignment $X\mapsto\bbK_{\D,X}$ into a presentable six-functor formalism
		\[
			\D^{(2)}\colon\Span(C)\longrightarrow 2\Pr,
			\qquad
			\D^{(2)}(X):=\bbK_{\D,X}.
		\]

		More generally, suppose that $\D^{(n)}\colon\Span(C)\to n\Pr$ has been constructed. Restricting it to $C_{/X}$ and applying the kernel construction in $n\Pr$ defines
		\[
			\D^{(n+1)}(X)
			\; :=\; \bbK_{\D^{(n)},X}
			\; := \; \Mod_{\D^{(n)}_X}\bigl(\Fun(\Span(C_{/X}),n\Pr)\bigr)
			\qin\CAlg((n+1)\Pr).
		\]
		The higher-categorical form of \Cref{prop:Adjoints_Category_Of_Kernels} again provides base-change-compatible left and right adjoints, and hence makes $\D^{(n+1)}$ a presentable six-functor formalism.

		For $X\in C$ and $n\geq1$, put $A_{\D,X,n}:=\D^{(n)}(X)$, and set
		\[
			A_{\D,X,0}:=\End_{\D(X)}(\unit).
		\]
		By the enriched-presheaf description of the kernel construction in \Cref{thm:Aoki_Enriched_Presheaves}, there are canonical equivalences
		\[
			A_{\D,X,n}\simeq\End_{A_{\D,X,n+1}}(\unit)
			\qquad(n\geq0).
		\]
		Consequently the sequence
		\[
			A_{\D,X}
			:=\bigl(\End_{\D(X)}(\unit),\ \D(X),\ \bbK_{\D,X},\ \bbK_{\D^{(2)},X},\ \ldots\bigr)
		\]
		is a Stefanich ring, functorially in $X$.
	\end{construction}

	\begin{definition}[Transmutation]\label[definition]{def:Transmutation}
		The structure map $X\to *$ induces a map of Stefanich rings $A_{\D,*}\to A_{\D,X}$. The \emph{transmutation} of $\D$ is the functor
		\[
			[-]_{\D}\colon C\to\Gest_{/\Gest(A_{\D,*})}, \qquad X\mapsto [X]_{\D} := \Gest(A_{\D,X}).
		\]
		Thus every $[X]_{\D}$ is canonically a Gestalt over $[*]_{\D}:=\Gest(A_{\D,*})$.
	\end{definition}

	\begin{theorem}
		\label[theorem]{thm:Transmutation}
		For every presentable six-functor formalism $\D$, the transmutation $[-]_{\D}$ satisfies:
		\begin{enumerate}
			\item It commutes with finite limits;
			\item For each $f\colon Y\to X$ in $C$, the map $[f]_{\D}\colon [Y]_{\D}\to[X]_{\D}$ is $1$-étale and $1$-proper, and satisfies $f_{1,\sharp}\simeq f_{1,*}$;
			\item The composite $\Oo(-)_1\circ[-]_{\D}\colon C\catop \to \CAlg(1\Pr)$ recovers the $(f^*,\otimes)$-functoriality of $\D$.
		\end{enumerate}
	\end{theorem}
	\begin{proof}
		\emph{Adjointability.} By construction the $m$-th level of $A_{\D,X}$ has pullback the kernel pullback for $\D^{(m-1)}$, so \Cref{prop:Adjoints_Category_Of_Kernels} applied to $\D^{(m-1)}$ gives, for every $m\geq 2$, base-change-compatible left and right adjoints to $f_{m-1}^*\colon \D^{(m)}(X)\to\D^{(m)}(Y)$ with $f_{m-1,\sharp}\simeq f_{m-1,*}$. Unwinding the definitions of primness and suaveness for Stefanich rings, this is exactly the internal right-adjoint datum for $1$-primness and the internal left-adjoint datum for $1$-suaveness. The extra condition in $n$-suaveness (\Cref{def:N_Suave}) also holds: by the levelwise ambidexterity and \Cref{lem:Linear_Adjoint_Preserves_Dualizable}, the adjoints preserve dualizable objects, so the dualizing object $(B/A)^!_m = f_{m,\sharp}(\unit)$ is dualizable, and for such objects the extra right-adjoint condition follows from the available left-adjoint one. Hence every $[f]_{\D}$ is $1$-prim and $1$-suave.

		\emph{Finite limits.} The terminal object goes to $[*]_{\D} = \Gest(A_{\D,*})$. For pullbacks, after passing to a slice it suffices to preserve binary products. Given $p_X\colon X\to S$ and $p_Y\colon Y\to S$, put $P = X\times_S Y$, with projections $q_X,q_Y$ and structure map $p\colon P\to S$. The two projections induce a map of Stefanich rings over $A_{\D,S}$ of the form
		\[
			\Phi\colon A_{\D,X}\otimes_{A_{\D,S}}A_{\D,Y} \to A_{\D,P}.
		\]
		We must show that this map is an isomorphism. Recall from \Cref{con:Tensor_Product_In_Slice} the explicit description of the relative tensor product over the base $A_{\D,S}$: we first form the levelwise tensor product in $\PreStRing(A_{\D,S})$, and then we spectrify. It will thus suffice to show that the levelwise tensor product in $\PreStRing(A_{\D,S})$ is isomorphic to $A_{\D,P}$ (so that spectrifying is unnecessary). In level $n$, the comparison map in $\PreStRing(A_{\D,S})$ is given by the map
		\[
			p_{X,n*}\unit_X\otimes p_{Y,n*}\unit_Y\to p_{n*}\unit_P
		\]
		in $\CAlg((A_{\D,S})_{n+1})$. But this follows from the projection formulas and base change established in \Cref{prop:Adjoints_Category_Of_Kernels}:
		\[
			p_{X,n*}\unit_X\otimes p_{Y,n*}\unit_Y \simeq p_{X,n*}\bigl(\unit_X\otimes p_X^*p_{Y,n*}\unit_Y\bigr) \simeq p_{X,n*}q_{X,n*}\unit_P = p_{n*}\unit_P,
		\]
		using $p_X^*p_{Y,n*}\unit_Y\simeq q_{X,n*}q_Y^*\unit_Y = q_{X,n*}\unit_P$. We conclude that $A_{\D,X}\otimes_{A_{\D,S}}A_{\D,Y}\iso A_{\D,P}$, and so $[P]_{\D}\iso[X]_{\D}\times_{[S]_{\D}}[Y]_{\D}$.

		\emph{Properness/étaleness.} In the Gestalt formalism $1$-properness (resp.\ $1$-étaleness) means $1$-primness (resp.\ $1$-suaveness) for the map and all its iterated diagonals. As $[-]_{\D}$ preserves finite limits it preserves diagonals, $[\Delta_f]_{\D}\simeq\Delta_{[f]_{\D}}$, so applying the first part to $f$ and to its iterated diagonals shows $[f]_{\D}$ is $1$-proper and $1$-étale.

		\emph{Recovering $\D^{*,\otimes}$.} By construction, the first level of $A_{\D,X}$ is $\Oo([X]_{\D})_1 = \D(X)$, with pullback and tensor functoriality the original $(f^*,\otimes)$ of $\D$; hence $\Oo(-)_1\circ[-]_{\D}$ recovers $\D^{*,\otimes}$.
	\end{proof}

	\begin{remark}
		\label[remark]{rem:Why_Iterate}
		The proof of finite-limit preservation makes clear why in \Cref{con:Iterated_Categories_Of_Kernels} we need to iterate the kernel construction indefinitely: to compare the degree-$n$ relative algebras, we work in the degree-$(n+1)$ algebra of the base object.
	\end{remark}

	\begin{corollary}[Künneth implies affineness]
		\label[corollary]{cor:Kunneth_Transmutation_Affine}
		Let $n\geq1$, and suppose that $\D^{(n)}$ satisfies the relative categorical Künneth formula: for every $S\in C$ and all $X,Y\in C_{/S}$, the canonical map
		\[
			\D^{(n)}(X)\otimes_{\D^{(n)}(S)}\D^{(n)}(Y)
			\longrightarrow
			\D^{(n)}(X\times_SY)
		\]
		is an equivalence. Then the Gestalt $[*]_{\D}$ is $n$-affine.
	\end{corollary}
	\begin{proof}
		We show more generally that, for every $m\geq n$, the coefficient system $\D^{(m)}$ satisfies the same relative Künneth formula, and that
		\[
			\D^{(m+1)}(S)\simeq \Mod_{\D^{(m)}(S)}(m\Pr)
		\]
		for every $S\in C$.

		Assume this has been proved in degree $m$. For a fixed $S\in C$, the restriction of $\D^{(m)}$ to the slice $C_{/S}$ satisfies the categorical Künneth formula in the sense of \Cref{cor:Categorical_Kunneth_Iff_Kernel_Category_Is_Modules}. Applying the same corollary with $1\Pr$ replaced by $m\Pr$ identifies the category of kernels of this restricted formalism with
		\[
			\Mod_{\D^{(m)}(S)}(m\Pr).
		\]
		By construction this category of kernels is exactly $\D^{(m+1)}(S)$.

		Now let $X,Y\in C_{/S}$. Using the identifications just obtained and the tensor product formula for module categories, we get
		\[
		\begin{aligned}
			\D^{(m+1)}(X)\otimes_{\D^{(m+1)}(S)}\D^{(m+1)}(Y)
			&\simeq
			\Mod_{\D^{(m)}(X)}(m\Pr)
			\otimes_{\Mod_{\D^{(m)}(S)}(m\Pr)}
			\Mod_{\D^{(m)}(Y)}(m\Pr)\\
			&\simeq
			\Mod_{\D^{(m)}(X)\otimes_{\D^{(m)}(S)}\D^{(m)}(Y)}(m\Pr)\\
			&\simeq
			\Mod_{\D^{(m)}(X\times_SY)}(m\Pr)\\
			&\simeq
			\D^{(m+1)}(X\times_SY).
		\end{aligned}
		\]
		Thus $\D^{(m+1)}$ again satisfies relative Künneth, completing the induction.

		Taking $S=*$, the Stefanich ring of $[*]_{\D}$ has $m$-th level $\D^{(m)}(*)$ for $m\geq1$. The equivalences
		\[
			\D^{(m+1)}(*)\simeq \Mod_{\D^{(m)}(*)}(m\Pr)
			\qquad(m\geq n)
		\]
		are precisely the affine comparison maps in degrees $m\geq n$. Hence $[*]_{\D}$ is $n$-affine.
	\end{proof}

	\section{Dictionary: six-functor notions vs.\ Gestalten}
	\label[section]{sec:Translation}

	Consider a presentable six-functor formalism $\D\colon \Span(C) \to 1\Pr$, and consider a map $f\colon Y \to X$ in $C$. Using the transmutation functor from the previous section, we now have two competing notions of cohomological properties (suaveness, primness, étaleness, properness):
	\begin{itemize}
		\item We may consider the map $[f]_{\D}$ of Gestalten, and use the definitions from the framework of Gestalten established in \Cref{chap:Properties_Of_Maps};
		\item We also have the notions from the six-functor calculus in the 2-category of kernels, recalled in \Cref{sec:Geometric_Properties}.
	\end{itemize}
	We now show that these two notions agree with each other.

	The comparison is concentrated entirely at the bottom layer. By \Cref{thm:Transmutation} every $[f]_{\D}$ is automatically $1$-\'etale and $1$-proper with $f_{1,\sharp}\simeq f_{1,*}$, so the layers above the bottom are forced by ambidexterity and separate no maps. Writing $A := \Oo([X]_{\D})$ and $B := \Oo([Y]_{\D})$, the morphism $[f]_{\D}$ is attached to a map $A\to B$ of Stefanich rings, whose lowest two layers are the coefficient category and the module $2$-category of kernels,
	\[
		A_1 = \D(X), \qquad A_2 = \bbK_{\D,X}.
	\]
	By \Cref{thm:Aoki_Enriched_Presheaves} the module $2$-category $\bbK_{\D,X}$ is the $\D_X$-enriched presheaf category on the kernel $2$-category $K_{\D,X}$ (\Cref{def:Kernel_2_Category}), so the fully faithful $2$-Yoneda embedding $K_{\D,X}\hookrightarrow\bbK_{\D,X}$ preserves and reflects adjunctions. By the construction of \Cref{sec:Transmutation}, the level-$0$ relative datum $(B/A)_1 \in \CAlg(A_2)$ has unit map
	\[
		(f/A)^*_0\colon \unit \longrightarrow (B/A)_1,
	\]
	and, under the embedding above, it is the unit kernel $\unit_Y$ of $f$ regarded as a $1$-morphism $[X]\to[Y]$ in $K_{\D,X}$. At every higher level $m\geq 1$ the relative pullback $(f/A)^*_m$ already admits adjoints on both sides (\Cref{thm:Transmutation}). Thus the only conditions of \Cref{def:N_Suave,def:N_Prim} not automatic for $[f]_{\D}$ sit at level $m=0$, where they concern precisely the unit kernel $\unit_Y$.

	We will now spell this out in a bit more detail.

	\begin{lemma}[Suave and prim maps]\label[lemma]{lem:Dictionary_Suave_Prim}
		Let $f\colon Y\to X$ be a morphism in $C$.
		\begin{enumerate}
			\item $f$ is $\D$-suave (\Cref{def:Suave_Prim_Morphisms}) if and only if $[f]_{\D}$ is $0$-suave (\Cref{def:N_Suave}).
			\item $f$ is $\D$-prim (\Cref{def:Suave_Prim_Morphisms}) if and only if $[f]_{\D}$ is $0$-prim (\Cref{def:N_Prim}).
		\end{enumerate}
	\end{lemma}
	\begin{proof}
		(2) By \Cref{def:N_Prim} and the reduction above, $[f]_{\D}$ is $0$-prim if and only if its one non-automatic condition holds: the level-$0$ datum $\unit\to(B/A)_1$ admits a right adjoint in $A_2 = \bbK_{\D,X}$. Since $2$-Yoneda reflects adjunctions, this says that the unit kernel $\unit_Y$, regarded as the $1$-morphism $[X]\to[Y]$ in $K_{\D,X}$, is a \emph{left} adjoint, i.e.\ $\unit_Y\in\Prim_f(Y)$ (\Cref{def:Suave_Prim_Objects}), i.e.\ $f$ is $\D$-prim (\Cref{def:Suave_Prim_Morphisms}).

		(1) Dually, $0$-suaveness of $[f]_{\D}$ asks that each $(f/A)^*_m$ admit a \emph{left} adjoint, subject to the secondary dualizability conditions of \Cref{def:N_Suave}. By the reduction above and \Cref{thm:Transmutation} every requirement is automatic except the left adjoint at $m=0$ and the secondary condition at $m=1$. The latter holds because the level-$1$ dualizing object $(B/A)^!_1 = f_{1,\sharp}(\unit)$ is dualizable (the ambidextrous level-$1$ adjoints preserve dualizable objects by \Cref{lem:Linear_Adjoint_Preserves_Dualizable}) and so the required right adjoint is the dual of an already-available left adjoint (\Cref{obs:Invertible_Vs_Dualizable}). As in~(2), the left adjoint at $m=0$ says that $\unit_Y$ is a \emph{right} adjoint $1$-morphism $[X]\to[Y]$ in $K_{\D,X}$. Via the self-duality $K_{\D,X}\simeq K_{\D,X}\catop$, this is equivalent to saying that $\unit_Y$ is a left adjoint $[Y]\to[X]$, i.e.\ that $\unit_Y\in\Suave_f(Y)$. We conclude that $[f]_{\D}$ is $0$-suave if and only if $f$ is $\D$-suave.
	\end{proof}

	\begin{corollary}[\'Etale and proper maps]\label[corollary]{lem:Dictionary_Etale_Proper}
		Let $f\colon Y\to X$ be a \emph{truncated} morphism in $C$.
		\begin{enumerate}
			\item $f$ is $\D$-\'etale (\Cref{def:Cohomologically_Etale_Proper}) if and only if $[f]_{\D}$ is $0$-\'etale (\Cref{def:N_Etale}).
			\item $f$ is $\D$-proper (\Cref{def:Cohomologically_Etale_Proper}) if and only if $[f]_{\D}$ is $0$-proper (\Cref{def:N_Proper}).
		\end{enumerate}
	\end{corollary}
	\begin{proof}
		Since $[-]_{\D}$ preserves finite limits, it carries the tower of iterated diagonals of $f$ to the tower of iterated codiagonals of $[f]_{\D}$. Since both $\D$-\'etaleness (resp.\ $\D$-properness) of $f$ and $0$-\'etaleness (resp.\ $0$-properness) of $[f]_{\D}$ ask for suaveness (resp.\ primness) for all iterated diagonals, the claim follows from \Cref{lem:Dictionary_Suave_Prim} applied at each stage. (Truncatedness of $f$ ensures the tower of iterated diagonals terminates at isomorphisms.)
	\end{proof}

	\begin{center}
		\begin{tabular}{ll}
			\toprule
			$f$ in $C$ & $[f]_{\D}$ in $\Gest$ \\
			\midrule
			$\D$-suave & $0$-suave \\
			$\D$-prim & $0$-prim \\
			$\D$-étale & truncated and $0$-étale \\
			$\D$-proper & truncated and $0$-proper \\
			\bottomrule
		\end{tabular}
	\end{center}

	\section{Extension to stacks}
	\label[section]{sec:Extension_To_Stacks}

	The transmutation extends from $C$ to stacks for the Grothendieck topology intrinsically attached to $\D$. Following \citet[Definition~5.3]{Scholze_Six_Functors}, call a family $\{f_i\colon Y_i\to X\}_{i\in I}$ \emph{universally $\D$-$!$-descending} if, for every $T\to X$, the canonical functor
	\[
		\D(T) \longrightarrow
		\lim_{[p]\in\simp}^{!}
		\prod_{(i_0,\ldots,i_p)\in I^{p+1}}
		\D\bigl(T\times_XY_{i_0}\times_X\cdots\times_XY_{i_p}\bigr)
	\]
	is an equivalence. Here the limit is formed using exceptional pullbacks. Let $J_{\D}$ be the Grothendieck topology generated by these families, and write $\Shv_{\D}(C)$ for the resulting category of sheaves.

	\begin{proposition}
		\label[proposition]{prop:Transmutation_Universal_Shriek_Descent}
		Let $\{f_i\colon Y_i\to X\}_{i\in I}$ be universally $\D$-$!$-descending. Then
		\[
			\coprod_{i\in I}[Y_i]_{\D}\longrightarrow[X]_{\D}
		\]
		is a cover of Gestalten.
	\end{proposition}
	\begin{proof}
		Both the hypothesis and the conclusion are stable under base change. We may therefore replace $C$ by $C_{/X}$ and assume that $X=\ast$.

		First suppose that $I$ is countable. By \Cref{thm:Transmutation}, each $[Y_i]_{\D}\to[\ast]_{\D}$ is $1$-\'etale, and so is their coproduct. Under the kernel description, the relative degree-$1$ dualizing object associated to $[Z]_{\D}\to[\ast]_{\D}$ is the corepresentable $[Z]_{\bbK_{\D}}$: for $p\colon Z\to\ast$, the identification $p_{1,\sharp}\simeq p_{1,*}$ gives
		\[
			p_{1,\sharp}(\unit_Z)\simeq p_{1,*}(\unit_Z)\simeq[Z]_{\bbK_{\D}}.
		\]
		Thus the cover criterion of \Cref{prop:Surjective_Criteria} reduces the claim to showing that the augmentation
		\[
			\colim_{[p]\in\simp\catop}
			\coprod_{(i_0,\ldots,i_p)\in I^{p+1}}
			\bigl[Y_{i_0}\times\cdots\times Y_{i_p}\bigr]_{\bbK_{\D}}
			\longrightarrow
			\unit
			\qin\bbK_{\D}=\Oo([\ast]_{\D})_2
		\]
		is an equivalence. The category $\bbK_{\D}$ is $1\Pr$-atomically generated by the corepresentables $[T]_{\bbK_{\D}}$ for $T\in C$. Applying $\Hom_{\bbK_{\D}}([T]_{\bbK_{\D}},-)$ to the displayed map gives
		\[
			\colim_{[p]\in\simp\catop}
			\coprod_{(i_0,\ldots,i_p)\in I^{p+1}}
			\D\bigl(T\times Y_{i_0}\times\cdots\times Y_{i_p}\bigr)
			\longrightarrow \D(T).
		\]
		This is the left adjoint of the $!$-descent functor for the base change $T\to\ast$, and is therefore an equivalence by hypothesis. Hence the augmentation in $\bbK_{\D}$ is an equivalence, as required.

		For a general indexing set $I$, let $C_I$ denote the source of the analogous augmentation. The same atomic-generation argument, now using universal $!$-descent for the entire family, shows that
		\[
			C_I\xrightarrow{\ \sim\ }\unit
		\]
		is an equivalence. Moreover, $C_I$ is the $\aleph_1$-filtered colimit of the analogous objects $C_{I_0}$ as $I_0$ ranges over the countable subsets of $I$. Since $\unit=[\ast]_{\bbK_{\D}}$ is $\aleph_1$-compact, the inverse $\unit\to C_I$ to the displayed equivalence factors through some $C_{I_0}$ with $I_0$ countable. Consequently, the identity of $\unit$ lies in the image of
		\[
			\Hom_{\bbK_{\D}}(\unit,C_{I_0})
			\longrightarrow
			\End_{\bbK_{\D}}(\unit).
		\]
		Using the identification of $C_{I_0}$ with the idempotent coalgebra of the countable family, \Cref{prop:Surjective_Retract_Criterion} shows that $\coprod_{i\in I_0}[Y_i]_{\D}\to[\ast]_{\D}$ is a cover. Hence the same is true for the original family.
	\end{proof}

	\begin{corollary}[{\citet[Corollary~9.8]{Scholze_Gestalten}}]
		\label[corollary]{cor:Transmutation_Extends_To_Stacks}
		The transmutation extends uniquely to a functor
		\[
			[-]_{\D}\colon
			\Shv_{\D}(C)^{\aleph_1}
			\longrightarrow
			\Gest^{1\text{-}\et}_{/[*]_{\D}}
			\colon X\longmapsto[X]_{\D}
		\]
		which preserves finite limits and countable colimits.
	\end{corollary}
	\begin{proof}
		The target is countably toposic by \Cref{thm:Etale_Countably_Toposic}. By \Cref{prop:Transmutation_Universal_Shriek_Descent}, the finite-limit-preserving functor $[-]_{\D}\colon C\to\Gest^{1\text{-}\et}_{/[*]_{\D}}$ carries the generating $J_{\D}$-covers to effective epimorphisms. The universal property of the $\aleph_1$-compact $J_{\D}$-sheaf category therefore gives the asserted extension. Its construction as the left exact localization of the countable presheaf category shows that it preserves finite limits and countable colimits.
	\end{proof}

	Thus the $\D$-topology is detected by the canonical topology on Gestalten: universal $!$-descent becomes effective descent after transmutation. The result is deliberately only a statement about the $\aleph_1$-compact part of the $\D$-topos; this is the level on which the countable colimits used in the Gestalt formalism are available.

	\section{A six-functor formalism on Gestalten}
	\label[section]{sec:6FF_On_Gest}

	Given a presentable six-functor formalism $\D\colon \Span(C) \to 1\Pr$, we constructed a pullback-preserving functor $[-]_{\D}\colon C \to \Gest$ fitting in a commutative triangle as follows:
	\[
	\begin{tikzcd}[column sep=huge]
		C\catop \rar[dashed,"{[-]_{\D}}"] \arrow[dr,swap,"{\D^{*,\otimes}}"] & \Gest\catop \dar["{\Oo(-)_1}"] \\
		& \CAlg(1\Pr).
	\end{tikzcd}
	\]
	A natural follow-up question arises: Can we also recover the shriek-functoriality of $\D$? Or, even better: does $\Gest$ admit a six-functor formalism whose restriction along $[-]_{\D}$ agrees with $\D$?
	
	The goal of this section is to give some partial answers to this question. Let's start by informally describing $!$-functors for a map of Gestalten. In the previous section we saw that for any map $f$ in $C$, the induced map $[f]_{\D}$ of Gestalten is $1$-étale and $1$-proper, and moreover comes with an identification $f_{1,\sharp} \cong f_{1,*}$. These conditions already yield $!$-functors: We get the counit map
	\[
		f_{1,*}f_1^* \cong f_{1,\sharp}f_1^* \to \id
	\]
	which, when evaluated at the unit, yields a map
	\[
		f_{0,!}\colon (\Oo(Y)/\Oo(X))_1 \to \unit
	\]
	in $\Oo(X)_2$. Said differently, the identification $f_{1,\sharp} \cong f_{1,*}$ yields an identification $(\Oo(Y)/\Oo(X))_1^! \cong (\Oo(Y)/\Oo(X))_1$, along which one can transfer $f_{0,!}$. Taking $\Hom_{\Oo(X)_2}(\unit,-)$, this in particular produces a map
	\[
		f_{0,!}\colon \Oo(Y)_1 \to \Oo(X)_1.
	\]
	
	This suggests that there exists a six-functor formalism on the category $\Gest$ with the following class of $!$-able maps:

	\begin{definition}
		\label[definition]{def:Quasi_Shrieksable}
		A map $f\colon Y \to X$ of Gestalten is \emph{quasi-$!$-able} if $f$ is $1$-étale and $1$-proper. The map $f$ is \emph{$!$-able} if it is quasi-$!$-able and some iterated diagonal of $f$ is $0$-proper. The map $f$ is \emph{co-$!$-able} if it is quasi-$!$-able and some iterated diagonal of $f$ is $0$-étale.
	\end{definition}

	Recall from \Cref{chap:Properties_Of_Maps} that these $!$-able maps automatically satisfy base change:

	\begin{proposition}[Prim and suave base change]\label[proposition]{prop:Prim_Suave_Base_Change}
		Consider a cartesian diagram
		\[\begin{tikzcd}
			Y' \rar["f'"] \dar["g'"'] & X' \dar["g"] \\
			Y \rar["f"'] & X
		\end{tikzcd}\]
		of Gestalten, and let $n \geq 0$.
		\begin{enumerate}
			\item If $g$ is $n$-prim, then for every $m \geq n$ the Beck--Chevalley transformation
			\[
				f_m^* g_{m,*} \to g'_{m,*} f'^*_m
			\]
			of functors $\Oo(X')_{m+1} \to \Oo(Y)_{m+1}$ is an isomorphism.
			\item If $f$ is $n$-suave, then for every $m \geq n$ the Beck--Chevalley transformation
			\[
				f'_{m,\sharp} g'^*_m \to g_m^* f_{m,\sharp}
			\]
			of functors $\Oo(Y)_{m+1} \to \Oo(X')_{m+1}$ is an isomorphism.
		\end{enumerate}
		In either case the Beck--Chevalley transformation refines to an isomorphism of $1$-morphisms in the $2$-category $\Oo(X)_{m+2}$.
	\end{proposition}
	\begin{proof}
		For (1), apply \Cref{prop:Prim_Properties}(3) to the pushout square of Stefanich rings obtained by applying $\Oo(-)$, with the prim map $g^*\colon \Oo(X) \to \Oo(X')$ and base-change map $f^*\colon \Oo(X) \to \Oo(Y)$. For (2), apply \Cref{prop:Suave_Properties}(1) to the same pushout square, now with the suave map $f^*\colon \Oo(X) \to \Oo(Y)$ and base-change map $g^*\colon \Oo(X) \to \Oo(X')$.
	\end{proof}

	Also recall that if functors admit both left and right adjoints, it is important to control their interaction in terms of a \emph{double} Beck--Chevalley transformation.

	\begin{proposition}
		\label[proposition]{prop:Double_BC_Gest}
		Consider a cartesian diagram
		\[\begin{tikzcd}
			Y' \rar["f'"] \dar["g'"'] & X' \dar["g"] \\
			Y \rar["f"'] & X
		\end{tikzcd}\]
		of Gestalten, and assume that $f$ is $n$-suave while $g$ is $n$-prim. Then the double Beck--Chevalley map
		\[
			f_{n,\sharp}g'_{n,*} \to g_{n,*}f'_{n,\sharp}
		\]
		of maps from $\Oo(Y')_{n+1}$ to $\Oo(X)_{n+1}$ is an isomorphism. In fact, it is even an isomorphism of maps from $(\Oo(Y')/\Oo(X))_{n+1}$ to $\unit$ in $\Oo(X)_{n+2}$.
	\end{proposition}
	\begin{proof}
		First, we claim that there is a natural transformation
		\[
			f_{n+1,*}f_{n+1}^* \xrightarrow{f_{n,\sharp}} \id
		\]
		in $\End(\Oo(X)_{n+2})$. If $f$ is $n$-étale, so that $f_{n+1,*}$ is linear, this is just obtained from $f_{n+1,*}\unit \xrightarrow{f_{n,\sharp}} \unit$ in $\Oo(X)_{n+2}$ by letting $\Oo(X)_{n+2}$ act on itself. In general, one first has a natural transformation
		\[
			\id \xrightarrow{f_{0,!}} f_{n+1,\sharp}f_{n+1}^*
		\]
		in $\End(\Oo(X)_{n+2})$, by the similar argument, and then passes to right adjoints.

		Write $h := g f' = f g'\colon Y' \to X$ for the long diagonal of the square. Now consider the natural morphism $\unit \to g_{n+1,*}\unit$ in $\Oo(X)_{n+2}$. Applying $f_{n+1,*}f_{n+1}^*$ to it, this becomes $f_{n+1,*}\unit \to h_{n+1,*}\unit$, and the existence of $f_{n,\sharp}$ amounts to the commutative square
		\[\begin{tikzcd}
			f_{n+1,*}\unit \rar \dar["{f_{n,\sharp}}"'] & h_{n+1,*}\unit \dar["{f'_{n,\sharp}}"] \\
			\unit \rar & g_{n+1,*}\unit
		\end{tikzcd}\]
		expressing the base change compatibility of $f_{n,\sharp}$.

		On the other hand, we can also apply this reasoning to the right adjoint $g_{n,*}\colon g_{n+1,*}\unit \to \unit$ of $\unit \to g_{n+1,*}\unit$. Applying $f_{n+1,*}f_{n+1}^*$ to it, this becomes $g'_{n,*}\colon h_{n+1,*}\unit \to f_{n+1,*}\unit$, and then $f_{n,\sharp}\colon f_{n+1,*}f_{n+1}^* \to \id$ gives the desired commutative diagram.
	\end{proof}

	In particular, we get the following corollary. (In principle one could state this with $1$ replaced by any $n \geq 0$.)

	\begin{corollary}\label[corollary]{cor:Sharp_Star_Twist}
		Let $f\colon Y \to X$ be a map of Gestalten that is $1$-étale and $1$-proper. Then $f_{1,\sharp}, f_{1,*}\colon \Oo(Y)_2 \to \Oo(X)_2$ agree up to twist. More precisely, there are natural isomorphisms
		\[
			f_{1,\sharp}(p_{1,*}\Delta_{f,1,\sharp}\unit \otimes -) \cong f_{1,*}, \qquad f_{1,*}(p_{1,\sharp}\Delta_{f,1,*}\unit \otimes -) \cong f_{1,\sharp}
		\]
		and
		\[
			p_{1,*}\Delta_{f,1,\sharp}\unit \otimes p_{1,\sharp}\Delta_{f,1,*}\unit \cong \unit \in \Oo(Y)_2,
		\]
		so both tensor factors are invertible. Here $Y \xrightarrow{\Delta_f} Y \times_X Y \xrightarrow{p} Y$ are the diagonal and first projection.
	\end{corollary}
	\begin{proof}
		Indeed, using linearity of all intervening functors, as well as the double Beck--Chevalley isomorphism,
		\begin{align*}
			f_{1,\sharp}(p_{1,*}\Delta_{f,1,\sharp}\unit \otimes M)
			&\cong f_{1,\sharp}p_{1,*}\Delta_{f,1,\sharp}\Delta_{f,1}^*p_1^*M \\
			&\cong f_{1,\sharp}p_{1,*}\Delta_{f,1,\sharp}M \\
			&\cong f_{1,*}q_{1,\sharp}\Delta_{f,1,\sharp}M \\
			&\cong f_{1,*}M,
		\end{align*}
		where $q\colon Y \times_X Y \to Y$ is the second projection. The other isomorphism is analogous. Moreover,
		\begin{align*}
			p_{1,*}\Delta_{f,1,\sharp}\unit \otimes p_{1,\sharp}\Delta_{f,1,*}\unit
			&\cong p_{1,*}\Delta_{f,1,\sharp}\Delta_{f,1}^*p_1^*p_{1,\sharp}\Delta_{f,1,*}\unit \\
			&\cong p_{1,*}\Delta_{f,1,\sharp}p_{1,\sharp}\Delta_{f,1,*}\unit \\
			&\cong p_{1,*}\Delta_{f,1,*}\unit \cong \unit. \qedhere
		\end{align*}
	\end{proof}

	In particular, for any $1$-étale and $1$-proper $f\colon Y \to X$, we get a canonically associated invertible $\mathbb{D}_f \in \Oo(Y)_2$. If $f$ is $0$-étale or $0$-proper, then one gets a trivialization $\mathbb{D}_f \cong \unit$ by ambidexterity. If one relates $f$ and $\Delta_f\colon Y \to Y \times_X Y$, one finds that $\mathbb{D}_f$ and $\mathbb{D}_{\Delta_f}$ are inverse. Thus, if $f$ is $!$-able or co-$!$-able, one gets a trivialization $\mathbb{D}_f \cong \unit$ by passing to some iterated diagonal.

	Once one has such a trivialization, we get the counit map
	\[
		f_{0,!}\colon f_{1,*}\unit \cong f_{1,\sharp}\unit \to \unit
	\]
	yielding a map $f_{0,!}\colon \Oo(Y)_1 \to \Oo(X)_1$ compatible with base change and composition. It is reasonable to expect that these define a six-functor formalism on $!$-able maps (or co-$!$-able maps), but at present it is not yet clear how to resolve all coherence issues.

	Restricting to truncated maps, we can however get the desired six-functor formalism. Let
	\[
		E_{\mathrm{tr}}\subseteq\Gest
	\]
	denote the wide subcategory of truncated, $1$-\'etale and $1$-proper maps.

	\begin{proposition}
		\label[proposition]{prop:6FF_Truncated_Gest}
		The assignment $X\mapsto\Oo(X)_1$ extends to a presentable six-functor formalism
		\[
			\Oo(-)_1\colon\Span(\Gest,E_{\mathrm{tr}})\longrightarrow1\Pr.
		\]
		If $f\colon Y\to X$ lies in $E_{\mathrm{tr}}$, the exceptional pushforward of this formalism is the map $f_{0,!}\colon\Oo(Y)_1\to\Oo(X)_1$ constructed above.
	\end{proposition}
	\begin{proof}
		\textbf{Step 1: the slicewise construction.}
		Fix a Gestalt $X$, and let
		\[
			C_X:=(E_{\mathrm{tr}})_{/X}
		\]
		be the category of truncated, $1$-\'etale and $1$-proper Gestalten over $X$. This category has finite limits, and every morphism in $C_X$ again lies in $E_{\mathrm{tr}}$: indeed, the relevant properties are stable under base change and composition, and the cancellation properties for the $1$-\'etale and $1$-proper conditions are part of \Cref{prop:Etale_Properties,prop:Proper_Properties}. Thus we may apply \Cref{thm:Span_2_Universal} with all three classes equal to $C_X$.

		Set $A:=\Oo(X)$. The relative algebras define a symmetric monoidal functor to $\CAlg(A_3)$; after forgetting the algebra structures, we regard this as a lax symmetric monoidal functor
		\[
			\Phi_X\colon C_X\catop\longrightarrow A_3,
			\qquad
			(Y\to X)\longmapsto (\Oo(Y)/A)_2.
		\]
		The symmetric monoidal structure is supplied by the Künneth formula \Cref{lem:Kunneth_Prim}, applied in degree $2$. If $f\colon Y\to Z$ lies in $C_X$, then the induced map $(f/A)^*_1$ admits both internal adjoints, since $f$ is $1$-\'etale and $1$-proper. The Beck--Chevalley, double Beck--Chevalley, and projection-formula requirements are exactly \Cref{prop:Prim_Suave_Base_Change,prop:Double_BC_Gest,lem:Suave_Projection_Formula,lem:Prim_Projection_Formula}. Hence \Cref{thm:Span_2_Universal} produces a lax symmetric monoidal $2$-functor
		\[
			\widetilde\Phi_X\colon
			\SpanTwoP{C_X}{C_X}{C_X}{C_X}\longrightarrow A_3
		\]
		extending $\Phi_X$.

		The terminal object of $C_X$ is $\id_X$, and $\Phi_X(\id_X)=(A/A)_2\simeq\unit_{A_3}$. Passing to endomorphisms of the unit as in \Cref{con:Endo_Unit} gives a lax symmetric monoidal functor
		\[
			\Span(C_X)
			\simeq
			\Hom_{\SpanTwoP{C_X}{C_X}{C_X}{C_X}}(\id_X,\id_X)
			\longrightarrow
			\End_{A_3}(\unit)
			\simeq A_2.
		\]
		An object $Y\to X$ of $C_X$ corresponds to the endomorphism span
		\[
			\id_X\xleftarrow{p_Y}Y\xrightarrow{p_Y}\id_X.
		\]
		Its image is the endomorphism $(p_Y/A)_{1,*}(p_Y/A)^*_1$ of the unit in $A_3$. Under the equivalence $\End_{A_3}(\unit)\simeq A_2$, evaluation at the unit identifies this object with
		\[
			(p_Y/A)_{1,*}(p_Y/A)^*_1(\unit)
			\simeq(\Oo(Y)/A)_1.
		\]
		For a morphism $f\colon Y\to Z$ in $C_X$, the forward span $Y\xleftarrow{\id}Y\xrightarrow{f}Z$ acts through the counit of $(f/A)_{1,\sharp}\dashv(f/A)^*_1$. Since $f$ belongs to both the left and the right class, \Cref{thm:Span_2_Universal} identifies $(f/A)_{1,\sharp}$ with $(f/A)_{1,*}$; after evaluation at the unit, this is precisely the map $f_{0,!}\colon\Oo(Y)_1\to\Oo(Z)_1$ constructed above. Forgetting along $A_2=\Oo(X)_2\to1\Pr$, we therefore obtain a lax symmetric monoidal functor
		\[
			D_X\colon\Span((E_{\mathrm{tr}})_{/X})\longrightarrow1\Pr,
			\qquad
			(Y\to X)\longmapsto\Oo(Y)_1.
		\]

		\textbf{Step 2: assembling the slices.}
		We now apply \Cref{thm:Slicewise_Reconstruction}. For $Y\to X$ in $(E_{\mathrm{tr}})_{/X}$ and $T\to X$ in $\Gest_{/X}$, define
		\[
			\D_{/X}(Y)(T):=D_T(T\times_XY)=\Oo(T\times_XY)_1.
		\]
		A span $Y\xleftarrow{}W\xrightarrow{}Z$ in $(E_{\mathrm{tr}})_{/X}$ pulls back along $T\to X$ to a span in $(E_{\mathrm{tr}})_{/T}$, and applying $D_T$ defines the corresponding functor
		\[
			\D_{/X}(Y)(T)\longrightarrow\D_{/X}(Z)(T).
		\]
		Thus we obtain functors
		\[
			\D_{/X}\colon\Span((E_{\mathrm{tr}})_{/X})\longrightarrow\Fun(\Gest_{/X}\catop,1\Pr).
		\]
		Their contravariant compatibility in $X$ is the evident base-change identification
		\[
			T'\times_{X'}(X'\times_XY)\simeq T'\times_XY.
		\]
		If $e\colon X'\to X$ lies in $E_{\mathrm{tr}}$, the covariant compatibility required in \Cref{thm:Slicewise_Reconstruction} is checked objectwise by the same pullback identity:
		\[
			\D_{/X}(e_{\sharp}Y)(T)
			=\Oo(T\times_XY)_1
			\simeq
			\Oo((T\times_XX')\times_{X'}Y)_1
			=(e_*\D_{/X'}(Y))(T).
		\]
		Therefore \Cref{thm:Slicewise_Reconstruction}, using the expected monoidal refinement discussed in \Cref{rmk:Slicewise_Reconstruction_Monoidal}, assembles the $D_X$ into a lax symmetric monoidal functor
		\[
			\Oo(-)_1\colon\Span(\Gest,E_{\mathrm{tr}})\longrightarrow1\Pr.
		\]
		Its value at $X$ is $\D_{/X}(\id_X)(\id_X)=\Oo(X)_1$, and the description of the exceptional pushforward in \Cref{thm:Slicewise_Reconstruction} identifies the map attached to $f\colon Y\to X$ with the forward-span map in $D_X$, hence with $f_{0,!}$. All values are presentable, so the remaining right adjoints exist by \Cref{rmk:Presentable_Targets}.
	\end{proof}

	\begin{question}
		Can the six-functor formalism on $\Gest$ constructed in \Cref{prop:6FF_Truncated_Gest} be extended to a functor
		\[
			\Oo(-)_1\colon \Span(\Gest, \Gest^{!}) \to 1\Pr,
		\]
		defined on the span category of \emph{all} $!$-able maps of Gestalten (\Cref{def:Quasi_Shrieksable})? Here $\Gest^{!} \subseteq \Gest$ denotes the wide subcategory of $!$-able maps.
	\end{question}

	\begin{question}
		If so, can the original six-functor formalism $\D$ be recovered from the six-functor formalism on $\Gest$ by precomposing with the transmutation, i.e.\ as the composite
		\[
			\Span(C,E) \xrightarrow{\Span([-]_{\D})} \Span(\Gest, \Gest^{!}) \xrightarrow{\Oo(-)_1} 1\Pr?
		\]
	\end{question}

	\chapter{Universality of the motivic Gestalt}
	\label[chapter]{chap:SH}

	\emph{Talk given by Bastiaan.}

	\emph{Essentially all material in this chapter is taken from \citet{Aoki_PhD_Thesis}.}

	In classical algebraic geometry, one studies smooth projective varieties over a field $k$ by means of cohomology theories, like Betti cohomology, de Rham cohomology, and $\ell$-adic cohomology. In order to understand the relation between these theories, Grothendieck envisioned the concept a universal abelian category $\Mot(k)$ of motives, together with a ``universal cohomology theory'' of the form $\Sch \to \Mot(k)$; every other cohomology theory ought to factor through $\Mot(k)$ by means of a \emph{realization functor} from $\Mot(k)$ to the category of vector spaces.

	While Grothendieck's original vision of $\Mot(k)$ has remained conjectural, it was later substantially realized by Voevodsky \citep{Voevodsky_Triangulated} through his triangulated category of motives $\mathrm{DM}(k)$. Closely related is the stable motivic homotopy theory $\SH(k)$ of \citet{MorelVoevodsky1999Motivic}, obtained from the category of smooth $k$-schemes by asking for presentability, Nisnevich descent, $\bbA^1$-invariance and $\bbP^1$-stability. For the purposes of this chapter, the categories $\SH(k)$ are the main objects of interest.

	The motivic stable category $\SH(X)$ can be defined for any scheme $X$, and the assignment $X \mapsto \SH(X)$ can be upgraded to a six-functor formalism
	\[
		\SH\colon \Span(\Sch_S) \to 1\Pr,
	\]
	where $\Sch_S$ is the category of separated finite type schemes over some base scheme $S$. By a result of Drew--Gallauer \citep{DrewGallauer2022Universal}, formulated in the closely related language of coefficient systems, this is the universal six-functor formalism on $\Sch_S$ satisfying four axioms: smooth base change and the projection formula, localization, $\bbA^1$-invariance, and $\bbP^1$-stability.

	Applying the category-of-kernels construction provides a presentably symmetric monoidal 2-category
	\[
		\SH^{(2)}(S):=\bbK_{\SH,S}\in\CAlg(2\Pr).
	\]
	Its corepresentable objects give a symmetric monoidal functor
	\[
		\Sch_S\catop\longrightarrow\SH^{(2)}(S), \qquad X\longmapsto[X],
	\]
	and its defining property is the formula
	\[
		\Hom_{\SH^{(2)}(S)}\bigl([X],[Y]\bigr)
		\simeq
		\SH(X\times_SY).
	\]
	The goal of today's talk is to go through the following results by \citet{Aoki_PhD_Thesis}:
	\begin{itemize}
		\item The functor $\Sch_S\catop \to \SH^{(2)}(S), X \mapsto [X]$ is universal among symmetric monoidal functors from schemes into stable presentably symmetric monoidal $2$-categories which satisfy the corresponding $2$-categorical versions of the same four axioms.
		\item The Gestalt $[\Spec(\Z)]_{\SH}$ is the free Gestalt $S$ equipped with a ring object $R \in \CRing(\Gest_{/S})$ which is 1-affine, 0-suave, homologically trivial and `sutured' (roughly: $0 \hookrightarrow R$ and $R^{\times} \hookrightarrow R$ form a recollement).
	\end{itemize}

	\section{Motivic spectra and the six operations}
	\label[section]{sec:Motivic_Spectra}

	We first recall the motivic six-functor formalism that will be used throughout this chapter. We let $\Sch = \Sch^{\sep,\mathrm{f.t.}}_{\Z}$ denote the category of separated $\Z$-schemes of finite type. For $X\in\Sch$, let $\Sm_X$ be the category of smooth $X$-schemes of finite presentation. The category of motivic spectra over $X$ is the symmetric monoidal localization
	\[
		\SH(X):= \Shv_{\Nis,\bbA^1}(\Sm_X;\Sp)
		\bigl[(\Sigma^\infty\bbP^1_X)^{\otimes-1}\bigr].
	\]
	Thus we first impose Nisnevich descent and $\bbA^1$-invariance, and then invert the pointed projective line. This is a stable presentably symmetric monoidal category.

	Base change of smooth schemes induces a contravariant functor
	\[
		\SH^{*,\otimes}\colon
		\Sch\catop\longrightarrow\CAlg(1\Pr),
		\qquad X\longmapsto\SH(X).
	\]
	We generally suppress the superscript and write this functor simply as $\SH$. Its basic properties follow directly from the construction. Every $f^*$ is exact, colimit-preserving, and symmetric monoidal, hence has a right adjoint $f_*$. If $f$ is smooth, then $f^*$ also has a left adjoint $f_\sharp$; it satisfies smooth base change and the projection formula. The localization theorem gives a recollement
	\[
		\SH(U)\xleftarrow{\ j^*\ }\SH(X)\xrightarrow{\ i^*\ }\SH(Z)
	\]
	for every closed immersion $i\colon Z\hookrightarrow X$ with open complement $j\colon U\hookrightarrow X$. Finally, for the projection $\pi\colon\bbA^1_X\to X$ and its zero section $s$, $\pi^*$ is fully faithful and the Tate functor
	\[
		\pi_\sharp s_*\colon\SH(X)\longrightarrow\SH(X)
	\]
	is an equivalence. The first assertion is $\bbA^1$-invariance, while the second is exactly the effect of inverting $\bbP^1$.

	The non-formal input is the following theorem of Ayoub, in the extended form due to Cisinski--D\'eglise. We first isolate the axioms to which it applies.

	\begin{definition}
		\label[definition]{def:Motivic_Coefficient_System}
		Let $S$ be noetherian and finite-dimensional. A \emph{motivic coefficient system} over $S$ is a lax symmetric monoidal functor
		\[
			D\colon\Sch_S\catop\longrightarrow1\Pr_{\st}
		\]
		satisfying the following four conditions.
		\begin{enumerate}
			\item Smooth maps admit left adjoints $f_\sharp$, and these left adjoints satisfy smooth base change and the projection formula.
			\item For every closed immersion $i\colon Z\hookrightarrow X$ with open complement $j\colon U\hookrightarrow X$, the diagram
			\[
				D(U)\xleftarrow{\ j^*\ }D(X)\xrightarrow{\ i^*\ }D(Z)
			\]
			is a recollement.
			\item For $\pi\colon\bbA^1_X\to X$, the functor $\pi^*$ is fully faithful.
			\item If $s\colon X\to\bbA^1_X$ is the zero section, then $\pi_\sharp s_*$ is an equivalence.
		\end{enumerate}
		A morphism of motivic coefficient systems is a lax symmetric monoidal natural transformation which is compatible with the left adjoints for smooth maps.
		We write
		\[
			\Fun^{\otimes\textup{-}\lax}_{\mathrm{mot}}\bigl(\Sch_S\catop,1\Pr_{\st}\bigr)
		\]
		for the category of motivic coefficient systems and their morphisms. Thus the objects are the functors satisfying the four conditions above, and the morphisms are the smooth-compatible lax symmetric monoidal natural transformations just described.
	\end{definition}

	The discussion before the definition says exactly that $X\mapsto\SH(X)$ is a motivic coefficient system. It satisfies the following universal property:

	\begin{theorem}[Drew--Gallauer {\citep{DrewGallauer2022Universal}}]
		\label[theorem]{thm:Drew_Gallauer_Universal_SH}
		The motivic coefficient system
		\[
			\SH\colon\Sch_S\catop\longrightarrow1\Pr_{\st}
		\]
		is an initial object of $\Fun^{\otimes\textup{-}\lax}_{\mathrm{mot}}\bigl(\Sch_S\catop,1\Pr_{\st}\bigr)$.
	\end{theorem}

	\begin{theorem}[Ayoub--Cisinski--D\'eglise {\citep{Ayoub2007SixOperations,CisinskiDeglise2019Triangulated}}]
		\label[theorem]{thm:Ayoub_Cisinski_Deglise_Axiomatic_Proper_Base_Change}
		Let $D\colon\Sch_S\catop\to1\Pr_{\st}$ be a motivic coefficient system. Then for every proper map $p\colon Y\to X$, the functor $p^*$ admits a right adjoint $p_*$. These right adjoints satisfy proper base change and the projection formula. Moreover, if
		\[
		\begin{tikzcd}
			Y' \rar["g'"] \dar[swap,"p'"] \drar[pullback] & Y \dar["p"] \\
			X' \rar[swap,"g"] & X
		\end{tikzcd}
		\]
		is cartesian with $p$ proper and $g$ smooth, then the smooth--proper Beck--Chevalley map
		\[
			g_\sharp p'_*\longrightarrow p_*g'_\sharp
		\]
		is an equivalence.

		Finally, if $\alpha\colon D\to D'$ is a morphism of motivic coefficient systems, then $\alpha$ is also compatible with the right adjoints for proper maps.
	\end{theorem}

	\begin{remark}
		\label[remark]{rmk:Nonmonoidal_ACD_Proper}
		Below I will also use a non-monoidal version of this theorem: after forgetting tensor products and projection formulas, the same hypotheses should imply the existence of proper right adjoints and the proper and smooth--proper base-change equivalences. I do not actually know for sure at the moment whether this version is true in this generality, and I will soon ask Denis-Charles Cisinski about it. Aoki only states the corresponding non-monoidal result for projective morphisms in the quasiprojective setting, and perhaps one should restrict to that case until this point is clarified.
	\end{remark}

	\begin{remark}
		Ayoub proves the theorem in the quasiprojective setting, where proper morphisms are projective. Cisinski--D\'eglise remove this quasiprojectivity restriction and construct the exceptional functoriality for all separated finite type morphisms; in particular $p_!\simeq p_*$ for proper $p$, with the proper base-change and projection-formula compatibilities above.
	\end{remark}

	\begin{corollary}
		\label[corollary]{cor:SH_Six_Functor_Formalism}
		The functor $\SH^{*,\otimes}$ extends to a presentable six-functor formalism
		\[
			\SH\colon\Span(\Sch)\longrightarrow1\Pr.
		\]
		If $f\colon Y\to X$ factors as an open immersion $j\colon Y\hookrightarrow\overline Y$ followed by a proper map $p\colon\overline Y\to X$, then
		\[
			f_!\simeq p_*j_\sharp.
		\]
	\end{corollary}
	\begin{proof}
		Apply \Cref{thm:Construction_6FF} with $I$ the open immersions, $P$ the proper maps, and $E=\Sch$. Both classes are stable under base change and left cancellable. By Nagata compactification, every morphism in $\Sch$ factors as an open immersion followed by a proper map, and every morphism in $I\cap P$ is $0$-truncated. The smooth functoriality recalled above supplies the left adjoints, Beck--Chevalley isomorphisms, and projection formula required for $I$. \Cref{thm:Ayoub_Cisinski_Deglise_Axiomatic_Proper_Base_Change} gives the dual data for $P$ and, applied to open immersions, gives the double Beck--Chevalley condition. The construction theorem therefore produces the asserted three-functor formalism and the formula for $f_!$. Since all its values are presentable, the needed right adjoints exist, so it is a six-functor formalism by \Cref{rmk:Presentable_Targets}.
	\end{proof}

	\section{The motivic Gestalt is \texorpdfstring{$2$}{2}-affine}
	\label[section]{sec:Motivic_Gestalt_Affineness}

	Put $S:=\Spec(\Z)$ and write
	\[
		\SH^{(2)}(X):=\bbK_{\SH,X}
		\qquad(X\in\Sch).
	\]
	Thus $\SH^{(2)}(X)$ is the category of kernels of the motivic six-functor formalism on $\Sch_{/X}$. The following observation of Scholze is the first reason that the Gestalt attached to $\SH$ is unusually accessible.

	\begin{proposition}[{\citet[Proposition~9.9]{Scholze_Gestalten}}]
		\label[proposition]{prop:Motivic_Gestalt_Affine}
		The Gestalt $[S]_{\SH}$ is $2$-affine:
		\[
			[S]_{\SH} \; \simeq \; \Gest(\SH^{(2)}(S)).
		\]
		Moreover, $\SH^{(2)}(S)$ is generated, under colimits and $1\Pr$-tensors, by the corepresentable objects
		\[
			[\bbA^n_{\Z}]_{\SH^{(2)}(S)}
			\qquad(n\geq0).
		\]
		For every $X\in\Sch$, the structural map
		\[
			[X]_{\SH}\longrightarrow[S]_{\SH}
		\]
		is $1$-affine.
	\end{proposition}
	\begin{proof}
		Let $p\colon X\to S$ be the structural map. The adjunction of \Cref{prop:Adjoints_Category_Of_Kernels} gives a symmetric monoidal functor
		\[
			p_1^*\colon\SH^{(2)}(S)\longrightarrow\SH^{(2)}(X)
		\]
		with right adjoint $p_{1,*}\simeq p_{1,\sharp}$. Set
		\[
			[X/S]_1:=p_{1,*}(\unit_{\SH^{(2)}(X)})\in\CAlg\bigl(\SH^{(2)}(S)\bigr).
		\]
		This is the degree-$1$ relative algebra of the map $[X]_{\SH}\to[S]_{\SH}$. The functor $p_{1,*}$ induces a functor $\SH^{(2)}(X) \to \Mod_{[X/S]_1}(\SH^{(2)}(S))$, which admits a symmetric monoidal left adjoint
		\[
			\Phi_X\colon\Mod_{[X/S]_1}\bigl(\SH^{(2)}(S)\bigr)\longrightarrow \SH^{(2)}(X)
		\]
		carrying the free module $[X/S]_1 \otimes M$ to $p_1^*M$. It follows from the projection formula that the unit map $[X/S]_1 \otimes M \to p_{1,*}p_1^*M$ is an isomorphism on all free modules, and since the source of $\Phi_X$ is generated by free modules under colimits and $1\Pr$-tensors it follows that $\Phi_X$ is fully faithful.

		It remains to prove that the image of $p_1^*$ generates $\SH^{(2)}(X)$. The corepresentables associated to objects $Y\to X$ of $\Sch_{/X}$ generate $\SH^{(2)}(X)$. By Zariski descent, we may cover $Y$ by affine opens and reduce to the case that $Y$ is affine. Choose a closed immersion $Y\hookrightarrow\bbA^n_{\Z}$. Since $X$ is separated over $S$, the graph of $Y\to X$ is a closed immersion; composing it with the base change of $Y\hookrightarrow\bbA^n_{\Z}$ gives a closed immersion
		\[
			i\colon Y\hookrightarrow\bbA^n_X.
		\]
		Since \(i\) is a closed immersion, the recollement for the closed--open decomposition gives a fully faithful right adjoint \(i_*\) to \(i^*\) in the 2-category $\SH^{(2)}(X)$. Hence the composite
		\[
			[Y]\xrightarrow{i_*}[\mathbb A^n_X]\xrightarrow{i^*}[Y]
		\]
		is the identity, so \([Y]\) is a retract of \([\mathbb A^n_X]\).
		 But
		\[
			\bbA^n_X=X\times_S\bbA^n_{\Z},
		\]
		and base change identifies the latter corepresentable with the image under $p_1^*$ of $[\bbA^n_{\Z}]_{\SH^{(2)}(S)}$. We conclude that the image of $p_1^*$ generates $\SH^{(2)}(X)$, hence $\Phi_X$ is an equivalence. Taking $S = X$, this also proves the stated assertion about generators, and it identifies
		\[
			\SH^{(2)}(X)\simeq\Mod_{[X/S]_1}\bigl(\SH^{(2)}(S)\bigr)
		\]
		for every $X$.

		The projection formula and base change also give
		\[
			[X\times_S Y/S]_1 \simeq [X/S]_1\otimes [Y/S]_1.
		\]
		Hence the preceding equivalence yields, for all $X,Y\in\Sch$,
		\[
			\SH^{(2)}(X)\otimes_{\SH^{(2)}(S)}\SH^{(2)}(Y)
			\simeq
			\Mod_{[X/S]_1 \otimes [Y/S]_1}\bigl(\SH^{(2)}(S)\bigr)
			\simeq
			\SH^{(2)}(X\times_SY).
		\]
		In other words, the $2$-categorical coefficient system $X\mapsto\SH^{(2)}(X)$ satisfies the categorical Künneth formula. The argument was relative in the base: applying it in the slice over any $T\in\Sch$ shows that $\SH^{(2)}$ satisfies the relative Künneth formula required in \Cref{cor:Kunneth_Transmutation_Affine}. That corollary, applied to $\D=\SH$ and $n=2$, proves the $2$-affineness of $[S]_{\SH}$.

		The same corollary applied in the slice over $X$ shows that $[X]_{\SH}$ is also $2$-affine. Hence the structural map $[X]_{\SH}\to[S]_{\SH}$ is automatically $2$-affine: in Stefanich-ring language, this follows from \Cref{prop:Affineness_Properties}(4). To improve this to $1$-affineness, only the first affine comparison remains. But this comparison is exactly the equivalence
		\[
			\SH^{(2)}(X)\simeq\Mod_{[X/S]_1}\bigl(\SH^{(2)}(S)\bigr)
		\]
		proved above. Hence $[X]_{\SH}\to[S]_{\SH}$ is $1$-affine.
	\end{proof}

	\section{The universal property of the motivic 2-category}
	\label[section]{sec:Aoki_Universal_Property}

	The result of the previous section shows that the Gestalt $[\Spec(\Z)]_{\SH}$ is the 2-affine Gestalt associated to the presentably symmetric monoidal 2-category $\SH^{(2)}(S)$. We now discuss the universal property of the functor $\Sch_S\catop \to \SH^{(2)}$ established by \citet{Aoki_PhD_Thesis}.

	In this section $S$ denotes a noetherian finite-dimensional base scheme, and $\Sch_S$ denotes the category of separated finite type $S$-schemes. For the rest of this chapter one may keep $S=\Spec(\Z)$ in mind.

	Let $\Cc$ be a stable presentably symmetric monoidal $2$-category. We use the same word \emph{motivic} for the $2$-categorical analogue of \Cref{def:Motivic_Coefficient_System}. Thus a symmetric monoidal functor
	\[
		F\colon\Sch_S\catop\longrightarrow\Cc
	\]
	is motivic if smooth pullback admits left adjoints satisfying smooth base change, closed--open decompositions give recollements, the projection $\bbA^1_X\to X$ is homotopy invariant, and the Tate operator $\pi_\sharp s_*$ is invertible. Morphisms between motivic functors are symmetric monoidal natural transformations which are compatible with the left adjoints for smooth maps.

	\begin{lemma}
		\label[lemma]{lem:Motivic_2_Functor_Extends_To_Spans}
		Let $\Cc$ be a stable presentably symmetric monoidal $2$-category, and let
		\[
			F\colon\Sch_S\catop\longrightarrow\Cc
		\]
		be motivic. Then $F$ is biadjointable with respect to open immersions and proper maps. Consequently, by \Cref{thm:Span_2_Universal}, it extends canonically to a symmetric monoidal $2$-functor
		\[
			\SpanTwoP{\Sch_S}{\Sch_S}{\mathrm{open}}{\mathrm{proper}}
			\longrightarrow\Cc.
		\]
	\end{lemma}
	\begin{proof}
		The motivic axioms already give the required left adjoints and Beck--Chevalley equivalences for open immersions, since open immersions are smooth. For proper maps, use \Cref{rmk:Nonmonoidal_ACD_Proper} via the represented functors $\Hom_{\Cc}(E,-)$ to obtain right adjoints and proper base change. The same input gives smooth--proper Beck--Chevalley; since open immersions are smooth, this is precisely the double Beck--Chevalley condition required in \Cref{def:Biadjointable_Functor}.

		It remains to explain the projection formulas. They are a formal consequence of symmetric monoidality together with the relevant Beck--Chevalley equivalence. We spell this out for a proper map $p\colon X\to Y$; the argument for open immersions and left adjoints is the same, with left adjoints in place of right adjoints. Consider the pullback square
		\[
		\begin{tikzcd}
			X \rar["\Gamma_p"] \dar["p"'] \drar[pullback] & Y\times X \dar["\id_Y\times p"] \\
			Y \rar["\Delta_Y"'] & Y\times Y,
		\end{tikzcd}
		\]
		where $\Gamma_p$ is the graph of $p$. By symmetric monoidality, pullback along $\id_Y\times p$ is the morphism $\id\otimes p^*$, so its right adjoint is $\id\otimes p_*$. For $N\in F(Y)$ and $M\in F(X)$, Beck--Chevalley for the displayed square gives
		\[
			\Delta_Y^*(\id\otimes p_*)(N\boxtimes M)
			\simeq
			p_*\Gamma_p^*(N\boxtimes M).
		\]
		The left-hand side is $N\otimes p_*M$, while the right-hand side is $p_*(p^*N\otimes M)$. This is precisely the projection formula for $p_*$. For an open immersion $j$, one uses instead the pullback square
		\[
		\begin{tikzcd}
			U \rar["\Gamma_j"] \dar["j"'] \drar[pullback] & U\times X \dar["j\times\id_X"] \\
			X \rar["\Delta_X"'] & X\times X
		\end{tikzcd}
		\]
		and the left Beck--Chevalley equivalence to obtain
		\[
			j_\sharp(M\otimes j^*N)\simeq j_\sharp(M)\otimes N.
		\]

		Thus $F$ is biadjointable with respect to open immersions and proper maps, and the final assertion is precisely \Cref{thm:Span_2_Universal}.
	\end{proof}

	We will also use the following concrete form of the universal property of the kernel category. Write
	\[
		\mathcal P_S:=\Fun(\Span(\Sch_S),1\Pr_{\st})
	\]
	with its Day convolution symmetric monoidal structure. Then
	\[
		\SH^{(2)}(S)=\bbK_{\SH,S}
		=
		\Mod_{\SH}(\mathcal P_S).
	\]

	\begin{lemma}
		\label[lemma]{lem:Maps_Out_Of_SH2}
		Let $\Cc$ be a presentably symmetric monoidal $2$-category. To construct a morphism
		\[
			H\colon\SH^{(2)}(S)\longrightarrow\Cc
		\]
		in $\CAlg(2\Pr)$ it is equivalent to construct the following two pieces of data:
		\begin{enumerate}
			\item a symmetric monoidal functor
			\[
				T\colon\Span(\Sch_S)\longrightarrow\Cc;
			\]
			\item a morphism of commutative algebra objects in $\Fun(\Span(\Sch_S),1\Pr_{\st})$
			\[
				\alpha\colon\SH\longrightarrow\Hom_{\Cc}(\unit,T(-)).
			\]
		\end{enumerate}
		Under this correspondence, $H$ sends the corepresentable object $[X]\in\SH^{(2)}(S)$ to $T(X)$. Moreover, natural transformations between such functors $H$ are determined by the corresponding transformations of the two pieces of data above.
	\end{lemma}
	\begin{proof}
		This is just the definition of $\SH^{(2)}(S)$ as a category of modules, unwound. A symmetric monoidal colimit-preserving $1\Pr$-linear functor out of the Day-convolution category $\mathcal P_S$ is determined by its restriction to the representables, hence by a symmetric monoidal functor from the span category. Passing from $\mathcal P_S$ to $\Mod_{\SH}(\mathcal P_S)$ imposes the additional datum that the commutative algebra $\SH\in\mathcal P_S$ acts on the unit of $\Cc$ after applying the chosen functor, i.e.\ the resulting commutative algebra in $\Cc$ comes with an augmentation. By adjunction, such an augmentation is equivalently a morphism
		\[
			\SH\longrightarrow\Hom_{\Cc}(\unit,T(-))
		\]
		of commutative algebra objects in $\Fun(\Span(\Sch_S),1\Pr_{\st})$. Finally, $\SH^{(2)}(S)$ is generated under colimits by the free $\SH$-modules on representables, which are the corepresentables $[X]$; this gives both the formula $H([X])\simeq T(X)$ and the assertion about natural transformations.
	\end{proof}

	\begin{theorem}[{\citet[Theorem~A]{Aoki_PhD_Thesis}}]
		\label[theorem]{thm:Aoki_Theorem_A}
		The functor of corepresentables
		\[
			[-]\colon\Sch_S\catop\longrightarrow\SH^{(2)}(S),
			\qquad X\longmapsto[X],
		\]
		is universal among motivic functors. Equivalently, for every stable presentably symmetric monoidal $2$-category $\Cc$, restriction along $[-]$ induces an equivalence
		\[
			\Fun^{\otimes,\mathrm L}\bigl(\SH^{(2)}(S),\Cc\bigr)
			\simeq
			\Fun^{\otimes}_{\mathrm{mot}}\bigl(\Sch_S\catop,\Cc\bigr).
		\]
	\end{theorem}

	\begin{proof}
		Let $\Mot^{(2)}(S)$ denote the universal stable presentably symmetric monoidal $2$-category equipped with a motivic functor
		\[
			[-]_{\Mot}\colon\Sch_S\catop\longrightarrow\Mot^{(2)}(S).
		\]
		Since $X\mapsto[X]\in\SH^{(2)}(S)$ satisfies the four axioms by the properties of motivic spectra recalled in \Cref{sec:Motivic_Spectra}, the universal property of $\Mot^{(2)}(S)$ gives a symmetric monoidal colimit-preserving functor
		\[
			F\colon\Mot^{(2)}(S)\longrightarrow\SH^{(2)}(S)
		\]
		under $\Sch_S\catop$. We have to construct an inverse.

		Applying \Cref{lem:Motivic_2_Functor_Extends_To_Spans} to $[-]_{\Mot}$ gives a canonical symmetric monoidal extension
		\[
			\widetilde{[-]}_{\Mot}\colon
			\SpanTwoP{\Sch_S}{\Sch_S}{\mathrm{open}}{\mathrm{proper}}
			\longrightarrow\Mot^{(2)}(S).
		\]
		Restricting this extension to the ordinary span category gives the first input in \Cref{lem:Maps_Out_Of_SH2}; call it
		\[
			T\colon\Span(\Sch_S)\longrightarrow\Mot^{(2)}(S).
		\]

		For the second input, map out of the unit of $\Mot^{(2)}(S)$. This produces a lax symmetric monoidal coefficient system
		\[
			D(X):=\Hom_{\Mot^{(2)}(S)}\bigl(\unit,[X]_{\Mot}\bigr),
			\qquad
			D\colon\Span(\Sch_S)\longrightarrow1\Pr_{\st}.
		\]
		Here $D$ is lax symmetric monoidal because $T$ is symmetric monoidal and $\Hom_{\Mot^{(2)}(S)}(\unit,-)$ is lax symmetric monoidal. The contravariant part $X\mapsto D(X)$ is a motivic coefficient system in the sense of \Cref{def:Motivic_Coefficient_System}. Indeed, the smooth left adjoints, closed--open recollements, $\bbA^1$-invariance, and Tate invertibility are obtained by applying $\Hom_{\Mot^{(2)}(S)}(\unit,-)$ to the corresponding structures for $[-]_{\Mot}$.

		The only point worth spelling out is the projection formula. Let $f\colon Y\to X$ be smooth. The graph-square argument in the proof of \Cref{lem:Motivic_2_Functor_Extends_To_Spans}, applied now to the smooth left adjoint $f_\sharp$, gives an equivalence in $\Mot^{(2)}(S)$
		\[
			f_\sharp(M\otimes f^*N)
			\simeq
			f_\sharp(M)\otimes N
		\]
		for $M\colon\unit\to[Y]_{\Mot}$ and $N\colon\unit\to[X]_{\Mot}$. Applying $\Hom_{\Mot^{(2)}(S)}(\unit,-)$ gives precisely the projection formula for the induced left adjoint on $D(Y)\to D(X)$. Hence \Cref{thm:Drew_Gallauer_Universal_SH} gives a unique morphism of motivic coefficient systems
		\[
			\alpha^{\op}\colon\SH\longrightarrow D|_{\Sch_S\catop}.
		\]

		We must still upgrade this contravariant map to spans. By definition of morphisms of motivic coefficient systems, $\alpha^{\op}$ is compatible with the smooth left adjoints, hence in particular with the left adjoints for open immersions. By the final clause of \Cref{thm:Ayoub_Cisinski_Deglise_Axiomatic_Proper_Base_Change}, it is also compatible with the right adjoints for proper maps. Therefore the same extension criterion of \citet{CLL_Span2} used in \Cref{lem:Motivic_2_Functor_Extends_To_Spans} extends $\alpha^{\op}$ uniquely to a morphism of commutative algebra objects
		\[
			\alpha\colon\SH\longrightarrow\Hom_{\Mot^{(2)}(S)}(\unit,T(-))
		\]
		in $\CAlg(\Fun(\Span(\Sch_S),1\Pr_{\st}))$, which is the second input required in \Cref{lem:Maps_Out_Of_SH2}.

		Now \Cref{lem:Maps_Out_Of_SH2}, applied to the pair $(T,\alpha)$, gives a symmetric monoidal colimit-preserving functor
		\[
			G\colon\SH^{(2)}(S)\longrightarrow\Mot^{(2)}(S)
		\]
		which sends the corepresentable $[X]\in\SH^{(2)}(S)$ to $[X]_{\Mot}$.

		It remains only to check that $F$ and $G$ are inverse equivalences. The composite $G F$ is a symmetric monoidal colimit-preserving endofunctor of $\Mot^{(2)}(S)$ whose restriction to $\Sch_S\catop$ is the identity. By the universal property defining $\Mot^{(2)}(S)$, this identifies $G F$ with the identity. For the other composite, use the last assertion of \Cref{lem:Maps_Out_Of_SH2}. The endofunctor $F G$ of $\SH^{(2)}(S)$ induces the identity on the span-valued functor of corepresentables, and its induced endomorphism of the coefficient system $\SH$ is the identity by Drew--Gallauer's initiality. Thus the two pieces of data attached to $F G$ agree with those attached to the identity functor of $\SH^{(2)}(S)$, so $F G\simeq\id$. This proves that $F$ is an equivalence, which is precisely the asserted universal property of $\SH^{(2)}(S)$.
	\end{proof}

	\begin{remark}
		Aoki phrases the theorem for quasiprojective schemes over $S$. I \emph{think} that this distinction is not essential here (this depends on whether \Cref{rmk:Nonmonoidal_ACD_Proper} is correct; I'll soon ask Denis-Charles about this.) One may replace the quasiprojective category by separated finite type schemes; the same universal property is obtained, using the usual basis argument for motivic spectra and compactification for the span formalism.
	\end{remark}

	\section{Ring Gestalten}
	\label[section]{sec:Ring_Gestalten}

	We now translate Aoki's universal property into the language of ring Gestalten. The guiding idea is very simple: a sufficiently geometric transmutation
	\[
		X\longmapsto X_R
	\]
	of schemes is controlled by the single object $R=(\bbA^1)_R$, and the motivic axioms become intrinsic conditions on this ring Gestalt.

	The section is organized as follows.
	\begin{itemize}
		\item First we prove the affine-line dictionary: finite-limit-preserving transmutations on affine schemes are the same thing as $0$-truncated commutative ring objects.
		\item Local ring Gestalten are the ones for which this affine construction satisfies Zariski gluing, and hence extends from affine schemes to all schemes.
		\item Sutured ring Gestalten are those for which the basic decomposition $\bbA^1=\bbG_m\cup\{0\}$ gives a recollement. By descent, this single condition controls general closed--open decompositions.
		\item Homological triviality and stability are the ring-Gestalt shadows of $\bbA^1$-invariance and invertibility of the Tate twist.
		\item Finally, weak suaveness and smoothness encode the smooth-adjoint part of the motivic formalism. Aoki's key observation is that, in the sutured situation, weak suaveness plus stability is equivalent to smoothness.
	\end{itemize}
	The final result says that the motivic ring Gestalt $[\bbA^1_{\Z}]_{\SH}$ is universal among ring Gestalten satisfying these conditions.

	\subsection{Transmutations from ring Gestalten}
	\label[subsection]{subsec:Transmutation_From_Ring_Gestalten}

	Write
	\[
		\Aff:=\Aff_{\Z}^{\mathrm{f.t.}}\subseteq\Sch
	\]
	for the full subcategory of affine schemes. We begin in an arbitrary category $C$ with finite limits. We write $\Ring(C)_{\leq0}$ for the category of commutative ring objects in $C$ whose underlying object is $0$-truncated. When $C$ is a category of Gestalten, we will simply refer to such objects as $0$-truncated ring Gestalten in $C$; in particular, $\RingGest$ denotes the category of commutative ring objects in $\Gest$.

	\begin{construction}[Evaluation at the affine line]
		\label[construction]{con:Evaluation_At_Affine_Line}
		Let $F\colon\Aff\to C$ preserve finite limits. Since $\bbA^1$ is a commutative ring object and $F$ preserves finite products, $F(\bbA^1)$ is a commutative ring object in $C$. Moreover, $\Aff$ is an ordinary category, so all its objects are $0$-truncated, and a finite-limit-preserving functor preserves truncated objects. We therefore obtain a functor
		\[
			\ev_{\bbA^1}\colon
			\Fun^{\lex}(\Aff,C)
			\longrightarrow
			\Ring(C)_{\leq0},
			\qquad F\longmapsto F(\bbA^1).
		\]
	\end{construction}

	\begin{proposition}
		\label[proposition]{prop:Affine_Transmutations_Ring_Gestalten}
		The functor $\ev_{\bbA^1}$ of \Cref{con:Evaluation_At_Affine_Line} is an equivalence.
	\end{proposition}
	In the Gestalten examples, this says that affine transmutations are the same thing as $0$-truncated ring Gestalten. More explicitly, let $R\in\Ring(C)_{\leq0}$ and suppose that
	\[
		X=\Spec\bigl(\Z[x_1,\ldots,x_k]/(f_1,\ldots,f_l)\bigr).
	\]
	Writing $f=(f_1,\ldots,f_l)\colon\bbA^k\to\bbA^l$, the functor corresponding to $R$ sends $X$ to the pullback
	\[
		\begin{tikzcd}
			X_R \rar \dar & R^k \dar["f_R"] \\
			* \rar["0"'] & R^l,
		\end{tikzcd}
	\]
	where $f_R$ is obtained by evaluating the polynomials $f_1,\ldots,f_l$ in $R$.
	\begin{proof}
		Let $\Poly_{\Z}\subseteq\Aff$ be the full subcategory spanned by the affine spaces $\bbA^k$. A finite-product-preserving functor $\Poly_{\Z}\to C$ is precisely a commutative ring object in $C$: the polynomial maps between affine spaces encode the ring operations and all their identities.

		The category $\Aff$ is the finite-limit completion of this algebraic theory in ordinary categories. Indeed, because $\Z$ is Noetherian, every finite type $\Z$-algebra is finitely presented, and hence every object of $\Aff$ is a finite limit of affine spaces as in the displayed pullback. It follows that a finite-product-preserving functor on $\Poly_{\Z}$ admits a finite-limit-preserving extension to $\Aff$ precisely when its underlying ring object is $0$-truncated; in that case the extension is unique and is computed by these finite-limit presentations. This construction is inverse to evaluation at $\bbA^1$, proving the proposition.
	\end{proof}

	\subsection{Local ring Gestalten}

	For this subsection, suppose that $C$ is countably toposic. Thus $C$ has finite limits and countable colimits, and these satisfy the countable Giraud axioms: countable colimits are universal, coproducts are disjoint, and groupoid objects are effective. We have in mind especially the categories $\Gest_{/S}^{n\text{-}\et}$, which are countably toposic by \Cref{thm:Etale_Countably_Toposic}.

	We next formulate Zariski locality intrinsically in $C$. Let $R\in\Ring(C)_{\leq0}$. Its object of units $R^\times$ is defined by the cartesian square
	\[
		\begin{tikzcd}
			R^\times \rar \dar \drar[pullback] & R\times R \dar["\mu"] \\
			* \rar["1"'] & R.
		\end{tikzcd}
	\]
	Composing the upper horizontal map with the first projection defines $\iota\colon R^\times\to R$. Equivalently, $R^\times$ classifies a generalized element of $R$ together with its multiplicative inverse. Define the locus on which $1-a$ is invertible by
	\[
		R^\times_{1-}:=R\times_{\,1-(-),\,R,\,\iota}R^\times.
	\]
	The two projections give a canonical map
	\[
		R^\times\sqcup R^\times_{1-}\longrightarrow R.
	\]

	\begin{definition}
		\label[definition]{def:Local_Ring_Gestalt}
		A $0$-truncated commutative ring object $R$ in $C$ is \emph{local} if the map
		\[
			R^\times\sqcup R^\times_{1-}\longrightarrow R
		\]
		is a cover, i.e.\ if the geometric realization of its \v{C}ech nerve is $R$. We write $\Ring^{\loc}(C)_{\leq0}$ for the full subcategory of local ring objects. In the Gestalten examples, its objects are the local ring Gestalten.
	\end{definition}

	This is the usual internal local-ring condition. Indeed, a generalized element of $R$ over $T\in C$ is a map $a\colon T\to R$. Pulling back the displayed cover gives the two loci in $T$ on which $a$ and $1-a$, respectively, are invertible; these loci cover $T$ precisely when $R$ is local.

	Let $\Fun^{\lex,\mathrm{Zar}}(\Sch,C)$ denote the full subcategory of $\Fun^{\lex}(\Sch,C)$ spanned by the functors satisfying Zariski gluing: for every finite affine open cover $U=\coprod_iU_i\to X$, the canonical map
	\[
		\colim_{[p]\in\simp\catop}F(U_p)\longrightarrow F(X)
	\]
	is an isomorphism, where $U_\bullet$ denotes the \v{C}ech nerve of $U\to X$.

	\begin{proposition}
		\label[proposition]{prop:Transmutation_Determined_By_Affine_Line}
		Evaluation at the affine line induces an equivalence
		\[
			\ev_{\bbA^1}\colon
			\Fun^{\lex,\mathrm{Zar}}(\Sch,C)
			\xrightarrow{\ \sim\ }
			\Ring^{\loc}(C)_{\leq0}.
		\]
	\end{proposition}
	\begin{proof}
		We first identify the effect of Zariski gluing on the affine restriction. Suppose that $F\colon\Sch\to C$ preserves finite limits and satisfies Zariski gluing, and put $R:=F(\bbA^1)$. Applying $F$ to the standard cover
		\[
			\bbA^1=D(t)\cup D(1-t)
		\]
		identifies the two open subobjects with $R^\times$ and $R^\times_{1-}$. The gluing assumption says that their coproduct is a cover of $R$, so $R$ is local.

		Conversely, let $R\in\Ring^{\loc}(C)_{\leq0}$ and let $F_R\colon\Aff\to C$ be the finite-limit-preserving functor of \Cref{prop:Affine_Transmutations_Ring_Gestalten}. Consider a finite principal-open cover
		\[
			\Spec(A)=D(f_1)\cup\dots\cup D(f_r).
		\]
		Choose $g_1,\ldots,g_r\in A$ with $\sum_i g_if_i=1$. The internal local-ring condition, applied inductively to this equality, implies that the loci on which one of the elements $g_if_i$ is invertible cover $F_R(\Spec(A))$. On such a locus $f_i$ is invertible, so it factors through $F_R(D(f_i))$. Hence $F_R$ carries the given principal-open cover to a cover in $C$. Every finite affine Zariski cover admits a refinement by a finite principal-open cover, and effectiveness of groupoid objects in $C$ then shows that $F_R$ carries its \v{C}ech nerve to a colimit diagram.

		It remains to pass between affine schemes and arbitrary objects of $\Sch$. Affine opens form a basis for the Zariski topology, and because the schemes in $\Sch$ are separated, this basis is closed under finite intersections. The usual basis-extension argument therefore identifies finite-limit-preserving Zariski-gluing functors on $\Sch$ with finite-limit-preserving functors on $\Aff$ which carry affine Zariski covers to colimit diagrams. Explicitly, if $X=\bigcup_iU_i$ is a finite affine open cover, the extension is forced to satisfy
		\[
			F_R(X)\simeq\colim_{[p]\in\simp\catop}F_R(U_p).
		\]
		Universality of countable colimits in $C$ shows that this construction is independent of the cover, functorial in $X$, and preserves finite limits. Thus restriction along $\Aff\hookrightarrow\Sch$ gives an equivalence onto the affine functors satisfying Zariski gluing. Combining this with the preceding paragraphs and \Cref{prop:Affine_Transmutations_Ring_Gestalten} proves the proposition.
	\end{proof}

	For $R\in\Ring^{\loc}(C)_{\leq0}$, we denote the corresponding transmutation by
	\[
		(-)_R\colon\Sch\longrightarrow C,
		\qquad X\longmapsto X_R.
	\]

	\subsection{Sutured ring Gestalten}
	\label[subsection]{subsec:Sutured_Ring_Gestalten}

	The next condition is no longer only a condition on the underlying finite-limit-preserving functor. It uses the coefficient categories attached to the objects $X_R$. Aoki phrases this using stacks, by which he means commutative algebras in a symmetric monoidal $2$-category. For us it is more natural to keep the Gestalten language throughout. Thus, from now on fix a presentably symmetric monoidal $2$-category $\Cc$, and set
	\[
		\Gg_{\Cc}:=\Gest^{1\text{-}\et}_{/\Gest(\Cc)}.
	\]
	Here $\Gest(\Cc)$ denotes the $1$-affine Gestalt corresponding to $\Cc$. The basic input will be a commutative ring object \(R\in\Gg_{\Cc}\) whose underlying object \(R=(\bbA^1)_R\) is \(1\)-affine over \(\Gest(\Cc)\). Then the affine-line dictionary gives \(X_R\) for every affine \(X\), and these affine pieces are again \(1\)-affine over \(\Gest(\Cc)\). If \(R\) is local, this extends by Zariski gluing to a transmutation
	\[
		(-)_R\colon\Sch\longrightarrow\Gg_{\Cc},
		\qquad X\longmapsto X_R
	\]
	on all schemes; the resulting \(X_R\) need not be \(1\)-affine for non-affine \(X\). Whenever \(X_R\) is defined and \(n\geq0\), set
	\[
		\D^{(n)}_R(X):=
		\bigl(\Oo(X_R)/\Oo(\Gest(\Cc))\bigr)_n.
	\]
	For affine \(X\), the \(1\)-affineness of \(X_R\) gives
	\[
		\D^{(1)}_R(X)\in\CAlg(\Cc),
		\qquad
		\D^{(2)}_R(X)\simeq\Mod_{\D^{(1)}_R(X)}(\Cc)\in\Mod_{\Cc}(2\Pr),
	\]
	while for general \(X\) the object \(\D^{(2)}_R(X)\) is defined by descent from affine opens. This is the Gestalten version of Aoki's convention that a stack has an associated coefficient category.

	We will use the following intrinsic form of recollement. Let $\Ee$ be a pointed $2$-category with finite limits and finite colimits. A sequence
	\[
		A\longrightarrow B\longrightarrow C
	\]
	in $\Ee$ is a \emph{bifiber sequence} if the square
	\[
	\begin{tikzcd}
		A \rar \dar & B \dar \\
		0 \rar & C
	\end{tikzcd}
	\]
	is both cartesian and cocartesian.

	\begin{definition}[Recollements in a pointed $2$-category]
		\label[definition]{def:Recollement_Pointed_2_Category}
		Let $\Ee$ be as above. A diagram
		\[
			C_U\xleftarrow{\ j^*\ }C_X\xrightarrow{\ i^*\ }C_Z
		\]
		in $\Ee$ is a \emph{recollement} if the following two conditions hold:
		\begin{enumerate}
			\item $j^*$ admits a left adjoint $j_\sharp$, and
			\[
				C_U\xrightarrow{\ j_\sharp\ }C_X\xrightarrow{\ i^*\ }C_Z
			\]
			is a bifiber sequence.
			\item $i^*$ admits a right adjoint $i_*$, and
			\[
				C_Z\xrightarrow{\ i_*\ }C_X\xrightarrow{\ j^*\ }C_U
			\]
			is a bifiber sequence.
		\end{enumerate}
	\end{definition}

	Recall that \(R^\times\to R\) denotes the object of units in the ring object \(R\), and that the zero section is the map \(0\colon *\to R\). More precisely, if \(R\) is the ring Gestalt, then in the following formulas \(R\), \(R^\times\), and \(*\) stand for \((\bbA^1)_R\), \((\bbG_m)_R\), and \((\Spec(\Z))_R\), respectively.

	\begin{definition}[Sutured ring Gestalten]
		\label[definition]{def:Sutured_Ring_Gestalt}
		Such a ring object \(R\in\Gg_{\Cc}\) is \emph{sutured} if the map \(R^\times\to R\) is an open immersion, the zero section \(*\to R\) is a closed immersion, and the diagram
		\[
			R^\times\longrightarrow R\longleftarrow *
		\]
		induces a recollement
		\[
			\D^{(2)}_R(\bbG_m)
			\xleftarrow{\ j^*\ }
			\D^{(2)}_R(\bbA^1)
			\xrightarrow{\ i^*\ }
			\D^{(2)}_R(\Spec(\Z))
		\]
		in the pointed $2$-category $\Mod_{\Cc}(2\Pr)$, where $j\colon R^\times\to R$ is the inclusion and $i\colon *\to R$ is the zero section.
	\end{definition}

	The open and closed immersion hypotheses are part of the geometry of the word ``sutured''. In particular, they imply that a sutured ring Gestalt is \(0\)-truncated: the shear automorphism of \(R\times R\) identifies the diagonal of \(R\) with the base change of the closed immersion \(*\to R\), namely the inclusion \(R\times\{0\}\to R\times R\). Thus the diagonal is a monomorphism.

	Thus suturedness says that the basic decomposition
	\[
		\bbA^1=\bbG_m\cup\{0\}
	\]
	behaves like an open--closed decomposition after transmutation. The following observation explains how, once the two pieces are sufficiently geometric, such a recollement also remembers the corresponding cover.

	\begin{lemma}
		\label[lemma]{lem:Recollement_Gives_Cover}
		Let $j\colon U\to X$ and $i\colon Z\to X$ be $1$-\'etale, $1$-affine maps of Gestalten. Suppose that the diagram
		\[
			\Oo(U)_2
			\xleftarrow{\ j^*\ }
			\Oo(X)_2
			\xrightarrow{\ i^*\ }
			\Oo(Z)_2
		\]
		is a recollement. Then
		\[
			U\sqcup Z\longrightarrow X
		\]
		is a cover.
	\end{lemma}
	\begin{proof}
		Set $f\colon U\sqcup Z\to X$. Since $1$-\'etale maps are closed under coproducts in the slice over $X$, the map $f$ is $1$-\'etale and hence $1$-suave. By the cover criterion of \Cref{prop:Surjective_Criteria}, it is therefore enough to show that
		\[
			f_1^*\colon \Oo(X)_2\longrightarrow\Oo(U\sqcup Z)_2
		\]
		is conservative. Relative global sections carry coproducts of Gestalten to products, so
		\[
			\Oo(U\sqcup Z)_2\simeq\Oo(U)_2\times\Oo(Z)_2,
		\]
		and under this identification $f_1^*$ is the pair $(j^*,i^*)$.

		It remains to recall the basic conservativity property of a recollement. If an object or morphism of the middle term becomes trivial after applying both $j^*$ and $i^*$, then the two bifiber sequences in \Cref{def:Recollement_Pointed_2_Category} force it to be trivial already in the middle term. Equivalently, the pair $(j^*,i^*)$ is jointly conservative. Thus $f_1^*$ is conservative, and \Cref{prop:Surjective_Criteria} shows that $f$ is a cover.
	\end{proof}

	In particular, suppose that $R$ is sutured and that the two maps $R^\times\to R$ and $*\to R$ in the defining sutured diagram are $1$-\'etale and $1$-affine. Then \Cref{lem:Recollement_Gives_Cover} shows that $R^\times\sqcup *\to R$ is a cover. Since the zero section factors through $R^\times_{1-}\to R$, this cover refines the locality map
	\[
		R^\times\sqcup R^\times_{1-}\longrightarrow R.
	\]
	Thus, under these geometric hypotheses, suturedness implies the locality cover from \Cref{def:Local_Ring_Gestalt}.

	By base change, the single recollement in the definition of suturedness controls all principal open--closed decompositions of affine schemes.

	\begin{proposition}
		\label[proposition]{prop:Sutured_Principal_Recollement}
		Let $R$ be sutured. Let $X=\Spec(A)$ be affine of finite type over $\Z$, and let $a\in A$. Put
		\[
			U:=D(a)\subseteq X,
			\qquad
			Z:=V(a)=\Spec(A/a)\subseteq X.
		\]
		Then the diagram
		\[
			\D^{(2)}_R(U)
			\xleftarrow{\ j^*\ }
			\D^{(2)}_R(X)
			\xrightarrow{\ i^*\ }
			\D^{(2)}_R(Z)
		\]
		is a recollement in $\Mod_{\Cc}(2\Pr)$.
	\end{proposition}
	\begin{proof}
		The element $a\in A$ determines a morphism $a\colon X\to\bbA^1$. The principal open $D(a)$ and the principal closed subscheme $V(a)$ are the two pullbacks
		\[
		\begin{tikzcd}
			D(a) \rar \dar & \bbG_m \dar \\
			X \rar["a"'] & \bbA^1,
		\end{tikzcd}
		\qquad
		\begin{tikzcd}
			V(a) \rar \dar & \Spec(\Z) \dar["0"] \\
			X \rar["a"'] & \bbA^1.
		\end{tikzcd}
		\]
		Applying the transmutation attached to $R$, these become the corresponding pullbacks of
		\[
			R^\times\longrightarrow R\longleftarrow *.
		\]
		Suturedness gives the recollement for this universal diagram in degree $2$. Base-changing this recollement along the map $X_R\to R$ induced by $a$ gives precisely the displayed diagram
		\[
			\D^{(2)}_R(U)\longleftarrow\D^{(2)}_R(X)\longrightarrow\D^{(2)}_R(Z).
		\]
		Hence the pulled-back diagram for $D(a)\subseteq X\supseteq V(a)$ is a recollement.
	\end{proof}

	The previous proposition is the affine principal form of the localization axiom. The general closed--open case follows from it by two elementary operations: finite induction and Zariski descent.

	\begin{lemma}
		\label[lemma]{lem:Sutured_General_Recollement}
		Let $R$ be local and sutured, and write $X\mapsto X_R$ for the associated Zariski-gluing transmutation. Then for every closed immersion $i\colon Z\hookrightarrow X$ in $\Sch$ with open complement $j\colon U\hookrightarrow X$, the diagram
		\[
			\D^{(2)}_R(U)
			\xleftarrow{\ j^*\ }
			\D^{(2)}_R(X)
			\xrightarrow{\ i^*\ }
			\D^{(2)}_R(Z)
		\]
		is a recollement.
	\end{lemma}
	\begin{proof}
		First note that recollements may be checked after finite Zariski covers. If $V=\coprod_iV_i\to X$ is a finite Zariski cover with \v{C}ech nerve $V_\bullet$, then Zariski gluing gives
		\[
			X_R\simeq\colim_{[p]\in\simp\catop}(V_p)_R
		\]
		in $\Gg_{\Cc}=\Gest^{1\text{-}\et}_{/\Gest(\Cc)}$. Applying relative global sections over $\Gest(\Cc)$ and taking degree $2$ gives
		\[
			\D^{(2)}_R(X)\simeq\lim_{[p]\in\simp}\D^{(2)}_R(V_p)
		\]
		in $\Mod_{\Cc}(2\Pr)$, and the same assertion holds for $U$ and $Z$. Thus the projections to the finite intersections of the cover jointly reflect the adjunctions and bifiber squares appearing in \Cref{def:Recollement_Pointed_2_Category}.

		We may therefore work affine-locally on $X$. Suppose first that $X=\Spec(A)$ and that $Z=V(I)$ with $I=(a_1,\ldots,a_r)$. The cases $U=\emptyset$ and $Z=\emptyset$ are tautological, so assume $r\geq1$.
		We claim, by induction on $r$, that
		\[
			\D^{(2)}_R(\bbA^r\setminus\{0\})
			\longleftarrow
			\D^{(2)}_R(\bbA^r)
			\longrightarrow
			\D^{(2)}_R(\Spec(\Z))
		\]
		is a recollement. For $r=1$ this is precisely the definition of suturedness. For the induction step, decompose $\bbA^r$ first into the principal open $D(t_1)$ and the principal closed $\bbA^{r-1}=V(t_1)$. On the closed part, the complement of the origin is $\bbA^{r-1}\setminus\{0\}$. Hence the induction hypothesis gives the localization sequence inside $V(t_1)$, and \Cref{prop:Sutured_Principal_Recollement} gives the principal localization sequence for $D(t_1)\subseteq\bbA^r$. Pasting the two recollements, or equivalently pasting the corresponding bifiber squares in \Cref{def:Recollement_Pointed_2_Category}, gives the desired recollement for $\bbA^r\setminus\{0\}\subseteq\bbA^r$.

		Now the tuple $(a_1,\ldots,a_r)$ defines a map $a\colon X\to\bbA^r$, and the closed--open pair
		\[
			U=X\setminus Z,
			\qquad
			Z=V(I)
		\]
		is the pullback of
		\[
			\bbA^r\setminus\{0\}\subseteq\bbA^r\supseteq\{0\}.
		\]
		Base-changing the recollement just constructed along $X_R\to(\bbA^r)_R$ gives the asserted recollement for the affine scheme $X$.

		For general $X$, choose a finite affine Zariski cover. Since the schemes in $\Sch$ are separated, the finite intersections in this cover are again affine. On each such affine, the closed subscheme is cut out by finitely many functions, so the affine case applies. The initial descent observation then glues these affine recollements to the desired recollement on $X$.
	\end{proof}

	Thus suturedness is precisely the ring-Gestalt translation of the closed--open localization axiom in \Cref{def:Motivic_Coefficient_System}.

	\subsection{Homologically trivial and stable ring Gestalten}
	\label[subsection]{subsec:Homologically_Trivial_Stable_Ring_Gestalten}

	We now translate the two remaining motivic axioms which only involve the affine line: $\bbA^1$-invariance and invertibility of the Tate twist. Continue to work in \(\Gg_{\Cc}\). Let
	\[
		\pi\colon R=(\bbA^1)_R\longrightarrow *=(\Spec(\Z))_R
	\]
	be the projection, and let
	\[
		s\colon *\longrightarrow R
	\]
	be the zero section.

	\begin{definition}[Homologically trivial ring Gestalten]
		\label[definition]{def:Homologically_Trivial_Ring_Gestalt}
		A ring Gestalt $R$ is \emph{homologically trivial} if the pullback
		\[
			\pi^*\colon\D^{(2)}_R(\Spec(\Z))\longrightarrow\D^{(2)}_R(\bbA^1)
		\]
		admits a left adjoint $\pi_\sharp$, and the counit
		\[
			\pi_\sharp\pi^*\longrightarrow\id_{\D^{(2)}_R(\Spec(\Z))}
		\]
		is an equivalence.
	\end{definition}

	Thus homological triviality says that the affine line has the same homology as the point: pulling coefficients to $R$ is fully faithful.

	\begin{definition}[Stable ring Gestalten]
		\label[definition]{def:Stable_Ring_Gestalt}
		Let $R$ be sutured, so that the zero-section pullback
		\[
			s^*\colon\D^{(2)}_R(\bbA^1)\longrightarrow\D^{(2)}_R(\Spec(\Z))
		\]
		admits a right adjoint $s_*$. We call $R$ \emph{stable} if $\pi^*$ admits a left adjoint $\pi_\sharp$ and the composite
		\[
			\pi_\sharp s_*\colon\D^{(2)}_R(\Spec(\Z))\longrightarrow\D^{(2)}_R(\Spec(\Z))
		\]
		is an autoequivalence.
	\end{definition}

	The autoequivalence $\pi_\sharp s_*$ is the \emph{Tate twist}.

	\begin{proposition}
		\label[proposition]{prop:Homologically_Trivial_Stable_Motivic_Axioms}
		Let $R$ be local. Then:
		\begin{enumerate}
			\item If $R$ is homologically trivial, then the assignment $X\mapsto\D^{(2)}_R(X)$ satisfies $\bbA^1$-invariance: for every $X\in\Sch$, the projection
			\[
				\pi_X\colon(\bbA^1_X)_R\longrightarrow X_R
			\]
			has the property that $\pi_X^*$ admits a left adjoint $\pi_{X,\sharp}$ and the counit
			\[
				\pi_{X,\sharp}\pi_X^*\longrightarrow\id_{\D^{(2)}_R(X)}
			\]
			is an equivalence.
			\item If $R$ is sutured and stable, then the assignment $X\mapsto\D^{(2)}_R(X)$ satisfies Tate invertibility: for every $X\in\Sch$, if
			\[
				s_X\colon X_R\longrightarrow(\bbA^1_X)_R
			\]
			is the zero section, then
			\[
				\pi_{X,\sharp}s_{X,*}\colon\D^{(2)}_R(X)\longrightarrow\D^{(2)}_R(X)
			\]
			is an autoequivalence.
		\end{enumerate}
	\end{proposition}
	\begin{proof}
		The point is that all of these constructions are obtained from the corresponding construction over the point by base change. We first treat the case where \(X\) is affine. Then \(X_R\) is \(1\)-affine over \(\Gest(\Cc)\), so products over the base are computed on degree-\(2\) coefficient categories by relative tensor product. In particular,
		\[
			\D^{(2)}_R(\bbA^1_X)
			\simeq
			\D^{(2)}_R(X)\otimes_{\Cc}\D^{(2)}_R(\bbA^1),
		\]
		and under this identification the pullback $\pi_X^*$ is obtained from $\pi^*$ by tensoring with $\D^{(2)}_R(X)$.

		If $R$ is homologically trivial, tensoring the adjunction $\pi_\sharp\dashv\pi^*$ with $\D^{(2)}_R(X)$ gives an adjunction $\pi_{X,\sharp}\dashv\pi_X^*$. The counit is the tensor product of the counit $\pi_\sharp\pi^*\to\id_{\D^{(2)}_R(\Spec(\Z))}$ with the identity of $\D^{(2)}_R(X)$, hence is an equivalence. This proves (1).

		For (2), suturedness gives the right adjoint $s_*$. Base change along $X_R\to *$ gives the right adjoint $s_{X,*}$, and the same tensor-product description identifies
		\[
			\pi_{X,\sharp}s_{X,*}
			\simeq
			\id_{\D^{(2)}_R(X)}\otimes(\pi_\sharp s_*).
		\]
		If $R$ is stable, the second tensor factor is an autoequivalence, hence so is the displayed functor.

		For general \(X\), choose a finite affine Zariski cover and pass to its \v{C}ech nerve. Since the transmutation \(X\mapsto X_R\) was constructed by gluing in \(\Gg_{\Cc}\), the categories \(\D^{(2)}_R(-)\) are obtained as the corresponding limits over this nerve. Adjunctions, counits, and the assertion that a functor is an equivalence may therefore be checked after restriction to the finite intersections of the cover, where the affine argument applies.
	\end{proof}

	\subsection{Weakly suave and smooth ring Gestalten}
	\label[subsection]{subsec:Weakly_Suave_Smooth_Ring_Gestalten}

	The last condition in Aoki's theorem is smoothness. Here one has to be slightly careful about terminology. In the general theory of Stefanich rings, recalled in \Cref{sec:Suave_Maps}, maps come with a hierarchy of \(n\)-suaveness conditions. Aoki's word ``suave'' corresponds in our indexing to \(0\)-suaveness. He first uses a weaker condition for stacks: only the visible pullback functor has to admit a linear left adjoint. In our notation this is the following.

	\begin{definition}[Weakly suave ring Gestalten]
		\label[definition]{def:Weakly_Suave_Ring_Gestalt}
		A ring Gestalt $R$ is \emph{weakly suave} if, for the projection
		\[
			\pi\colon R=(\bbA^1)_R\longrightarrow *=(\Spec(\Z))_R,
		\]
		the pullback
		\[
			\pi^*\colon\D^{(2)}_R(\Spec(\Z))\longrightarrow\D^{(2)}_R(\bbA^1)
		\]
		admits a $\D^{(2)}_R(\Spec(\Z))$-linear left adjoint $\pi_\sharp$.
	\end{definition}

	Thus weak suaveness is the degree-$2$ shadow of \(0\)-suaveness for the map $\pi$. Full \(0\)-suaveness asks for the corresponding higher adjointability and dualizability data from \Cref{def:N_Suave}. The point of Aoki's argument is that, for sutured ring Gestalten, stability supplies exactly the missing duality.

	\begin{definition}[Smooth ring Gestalten]
		\label[definition]{def:Smooth_Ring_Gestalt}
		A sutured ring Gestalt $R$ is \emph{smooth} if the projection
		\[
			\pi\colon R\to*
		\]
		is \(0\)-suave in the sense of the Gestalten/Stefanich-ring formalism, and its dualizing object
		\[
			\omega_\pi:=\pi^!\unit
		\]
		is tensor-invertible.
	\end{definition}

	This is the direct analogue of the usual criterion for cohomological smoothness: \(0\)-suaveness gives the exceptional pullback, and invertibility of the dualizing object is the smoothness condition.

	\begin{proposition}[{\citet{Aoki_PhD_Thesis}}]
		\label[proposition]{prop:Aoki_Weakly_Suave_Stable_Smooth}
		Let $R$ be a sutured weakly suave ring Gestalt. Then $R$ is smooth if and only if it is stable.
	\end{proposition}
	\begin{proof}
		First, note that suturedness makes the diagonal of $R$ into a closed immersion. Indeed, the shearing automorphism
		\[
			R\times R\longrightarrow R\times R,
			\qquad
			(x,y)\longmapsto(x,x-y),
		\]
		identifies the diagonal $\Delta\colon R\to R\times R$ with the map $\id_R\times0\colon R\times *\to R\times R$. But the zero section is the closed part of the sutured recollement, and closed immersions are stable under base change.

		Now use the left adjoint $\pi_\sharp$ supplied by weak suaveness. The closed diagonal gives a right adjoint $\Delta_*$ to $\Delta^*$, and hence maps
		\[
			\eta\colon
			\D^{(2)}_R(\Spec(\Z))
			\xrightarrow{\ \pi^*\ }
			\D^{(2)}_R(\bbA^1)
			\xrightarrow{\ \Delta_*\ }
			\D^{(2)}_R(\bbA^1\times\bbA^1)
		\]
		and
		\[
			\epsilon\colon
			\D^{(2)}_R(\bbA^1\times\bbA^1)
			\xrightarrow{\ \Delta^*\ }
			\D^{(2)}_R(\bbA^1)
			\xrightarrow{\ \pi_\sharp\ }
			\D^{(2)}_R(\Spec(\Z)).
		\]
		These are the expected coevaluation and evaluation maps for self-duality of $R$ over the point. The triangle composites are not literally identities at first. After applying the same shearing trick as above, they identify with the Tate operator $\pi_\sharp s_*$, where $s\colon *\to R$ is the zero section. Stability says precisely that this operator is invertible. The elementary duality fact used by Aoki says that if the two triangle composites are equivalences, then one can twist the coevaluation to obtain an honest duality datum. Thus \(R\to *\) is \(0\)-suave: weak suaveness has been promoted to the dualizability required in the Gestalten sense.

		It remains to identify the invertibility of the dualizing object. Let $B R$ denote the classifying object of the additive group underlying $R$, and let
		\[
			t\colon *\longrightarrow B R
		\]
		be the universal point. The standard torsor square shows that
		\[
			\omega_\pi=\pi^!\unit
		\]
		is pulled back from the object
		\[
			L:=t^!\unit.
		\]
		Consequently, smoothness is equivalent to invertibility of $L$. On the other hand, using $s^*\pi^!\unit\simeq L$, the projection formula, and the fact that the composite $\pi\circ s$ is the identity, Aoki computes
		\[
			L
			\simeq
			s^*\pi^!\unit
			\simeq
			\pi_\sharp s_*\unit.
		\]
		Thus $L$ is invertible exactly when the Tate operator $\pi_\sharp s_*$ is invertible. This is stability, and the proposition follows.
	\end{proof}

	In slogan form: suturedness gives the closed diagonal, weak suaveness gives the left adjoint, and stability turns the resulting weak duality into an invertible one.

	\subsection{Universality of the motivic ring Gestalt}
	\label[subsection]{subsec:Universality_Motivic_Ring_Gestalt}

	We can now assemble the preceding translations into Aoki's ring-Gestalt form of the universal property. The definitions above did not require stability of the ambient $2$-category $\Cc$. In Aoki's formulation, the category generated by the universal ring stack is stable, but the ambient presentably symmetric monoidal $2$-category in which one places ring stacks need not itself be stable.

	By \Cref{prop:Motivic_Gestalt_Affine}, we may identify the base Gestalt $[\Spec(\Z)]_{\SH}$ with $\Gest(\SH^{(2)}(\Spec(\Z)))$. Under this identification the motivic affine line is a ring Gestalt
	\[
		R_{\SH}:=[\bbA^1_{\Z}]_{\SH}
		\in
		\Gg_{\SH^{(2)}(\Spec(\Z))},
	\]
	and this object is \(1\)-affine over \(\Gest(\SH^{(2)}(\Spec(\Z)))\).
	If $\Cc$ is a presentably symmetric monoidal $2$-category, let
	\[
		\RingGest_{\mot}(\Cc)
	\]
	denote the category of local, sutured, homologically trivial, and smooth ring Gestalten in \(\Gg_{\Cc}\); morphisms are morphisms of ring Gestalten.

	\begin{theorem}[Aoki's ring-Gestalt theorem]
		\label[theorem]{thm:Aoki_Ring_Gestalt_Universality}
		For every presentably symmetric monoidal $2$-category $\Cc$, evaluation on the motivic affine line induces an equivalence
		\[
			\Fun^{\otimes,\mathrm L}
			\bigl(\SH^{(2)}(\Spec(\Z)),\Cc\bigr)
			\xrightarrow{\ \sim\ }
			\RingGest_{\mot}(\Cc),
			\qquad
			H\longmapsto H(R_{\SH}).
		\]
		Equivalently, the stable presentably symmetric monoidal $2$-category $\SH^{(2)}(\Spec(\Z))$ is freely generated by a local, homologically trivial, smooth, sutured ring Gestalt.
	\end{theorem}
	\begin{proof}
		This is \citet[Theorem~B]{Aoki_PhD_Thesis}, rewritten in the Gestalten language used here. Let us spell out how the dictionary works.

		First suppose that
		\[
			H\colon\SH^{(2)}(\Spec(\Z))\longrightarrow\Cc
		\]
		is symmetric monoidal and colimit-preserving. The notation $H(R_{\SH})$ means the following: the $1$-affine Gestalt $R_{\SH}\to[\Spec(\Z)]_{\SH}$ is equivalently encoded by its algebra of functions, a commutative algebra object of $\SH^{(2)}(\Spec(\Z))$, and we apply $H$ to that commutative algebra. This gives a ring Gestalt over $\Gest(\Cc)$.

		By \Cref{thm:Aoki_Theorem_A}, the associated functor
		\[
			X\longmapsto H([X])
		\]
		is motivic, where $[X]$ denotes the corepresentable object of $\SH^{(2)}(\Spec(\Z))$. Evaluating at $\bbA^1$ therefore gives a ring Gestalt with exactly the properties isolated above. Zariski descent gives locality, the localization axiom for the decomposition $\bbA^1=\bbG_m\cup\{0\}$ gives suturedness, and homotopy invariance gives homological triviality. The smooth axiom gives weak suaveness of the projection $H(R_{\SH})\to*$, while the Tate axiom gives stability; by \Cref{prop:Aoki_Weakly_Suave_Stable_Smooth}, this is the same as smoothness.

		Conversely, let $R$ be a local, sutured, homologically trivial, smooth ring Gestalt over $\Gest(\Cc)$. Locality gives the transmutation
		\[
			(-)_R\colon\Sch\longrightarrow
			\Gg_{\Cc}
		\]
		by \Cref{prop:Transmutation_Determined_By_Affine_Line}. Taking degree-$1$ functions gives a symmetric monoidal functor
		\[
			F_R\colon\Sch\catop\longrightarrow\Cc,
			\qquad
			X\longmapsto\D^{(1)}_R(X),
		\]
		where we forget the commutative algebra structure on $\D^{(1)}_R(X)$ when viewing it as an object of $\Cc$.

		The preceding subsections translate the ring-Gestalt conditions into the motivic axioms for $F_R$. Suturedness gives closed--open recollements, by \Cref{lem:Sutured_General_Recollement}. Homological triviality gives $\bbA^1$-invariance, by \Cref{prop:Homologically_Trivial_Stable_Motivic_Axioms}. Since smoothness implies weak suaveness, \Cref{prop:Aoki_Weakly_Suave_Stable_Smooth} gives stability, and hence the Tate operator is invertible. Finally, smoothness of $R\to*$ supplies the smooth-adjoint part of the motivic axioms; for general smooth morphisms this follows by the usual reduction to affine space, together with stability of \(0\)-suaveness under products, base change, composition, and Zariski descent.

		Thus $F_R$ is a motivic functor. By \Cref{thm:Aoki_Theorem_A}, it is induced by a unique symmetric monoidal colimit-preserving functor
		\[
			H_R\colon\SH^{(2)}(\Spec(\Z))\longrightarrow\Cc.
		\]
		The affine-line dictionary of \Cref{prop:Affine_Transmutations_Ring_Gestalten} and \Cref{prop:Transmutation_Determined_By_Affine_Line} says that the resulting ring Gestalt $H_R(R_{\SH})$ is canonically identified with the original $R$.

		The same argument also identifies morphisms. A natural transformation between functors out of $\SH^{(2)}(\Spec(\Z))$ is determined by its effect on the corepresentable motivic functor, by \Cref{lem:Maps_Out_Of_SH2}; under the affine-line dictionary, this is the same as a morphism of the corresponding local ring Gestalten. Hence the two constructions above are inverse equivalences.
	\end{proof}

	For a Gestalt $X$, let $\RingGest_{\mot}(X)$ denote the category of local, sutured, homologically trivial, and smooth ring Gestalten over $X$, i.e.\ such ring Gestalten in $\Gest^{1\text{-}\et}_{/X}$ whose affine-line object is \(1\)-affine over \(X\). If $X=\Gest(\Cc)$, this recovers the notation $\RingGest_{\mot}(\Cc)$ above.

	\begin{corollary}
		\label[corollary]{cor:Motivic_Gestalt_Represents_Ring_Gestalten}
		The presheaf
		\[
			\Gest\catop\longrightarrow\An,
			\qquad
			X\longmapsto\RingGest_{\mot}(X)^\simeq
		\]
		is represented by the motivic Gestalt $[\Spec(\Z)]_{\SH}$. Equivalently, for every Gestalt $X$ there is a natural equivalence
		\[
			\Hom_{\Gest}\bigl(X,[\Spec(\Z)]_{\SH}\bigr)
			\simeq
			\RingGest_{\mot}(X)^\simeq.
		\]
	\end{corollary}
	\begin{proof}
		By \Cref{prop:Motivic_Gestalt_Affine}, the motivic Gestalt is $2$-affine:
		\[
			[\Spec(\Z)]_{\SH}
			\simeq
			\Gest\bigl(\SH^{(2)}(\Spec(\Z))\bigr).
		\]
		Therefore maps from $X$ to $[\Spec(\Z)]_{\SH}$ are the same thing as symmetric monoidal colimit-preserving functors
		\[
			\SH^{(2)}(\Spec(\Z))\longrightarrow\Oo(X)_2.
		\]
		Applying \Cref{thm:Aoki_Ring_Gestalt_Universality} with target $\Oo(X)_2$ identifies these functors with motivic ring Gestalten over $X$. Passing to maximal subgroupoids gives the displayed equivalence, naturally in $X$.
	\end{proof}

	Using \Cref{prop:Aoki_Weakly_Suave_Stable_Smooth}, one may equivalently describe $\RingGest_{\mot}$ as the moduli problem of local, sutured, homologically trivial, weakly suave, and stable ring Gestalten. This is often the more practical formulation: weak suaveness is the visible adjoint, while stability is the extra condition which upgrades it to smoothness.

	\chapter{Cartier duality}
	\label[chapter]{chap:Cartier_Duality}

	\emph{Talk given by Han-Ung.}

	Today we introduce the statement of a very general \emph{Cartier duality} in the topos of Gestalten, following \citet[Lecture~X]{Scholze_Gestalten}, and illustrate it with many examples. This duality unifies and generalizes a host of classical ``Fourier isomorphisms''. The proof is deferred to the following talks; here we concentrate on the statement, its corollaries, and the way in which the classical examples are recovered.

	\emph{Most of the material here has been taken verbatim from Scholze's notes.}

	\section{Fourier isomorphisms}
	\label[section]{sec:Cartier_Examples}

	Cartier duality is a form of Fourier duality, so we begin by recalling various incarnations of the latter.

	\begin{example}[Fourier]\label[example]{ex:Fourier_Circle}
		We work with $\C$-coefficients throughout. The most basic instance of a Fourier isomorphism is the isomorphism of Hilbert spaces
		\[
			L^2(S^1) \iso L^2(\Z),
		\]
		where we regard $S^1 = \R/\Z$ with coordinate $x$. In one direction a sequence $(a_n)_{n \in \Z}$ is sent to the function $\sum_{n} a_n \exp(2\pi i n x)$; in the other, a function $f$ is sent to the sequence
		\[
			n \mapsto \int_{S^1} f(x)\exp(-2\pi i n x)\, dx.
		\]
		If we instead regard $S^1 \subseteq \C^{\times}$ with coordinate $t$, the formulas become $(a_n)_n \mapsto \sum_n a_n t^n$ and $f \mapsto \bigl(n \mapsto \int_{S^1} f(t) t^{-n}\, \tfrac{dt}{t}\bigr)$.

		There are various ``decompleted'' refinements of this isomorphism; for instance
		\[
			C^{\infty}(S^1) \cong \{\text{sequences } (a_n)_n \text{ of rapid decay}\}, \qquad \mathcal{D}'(S^1) \cong \{\text{sequences } (a_n)_n \text{ of polynomial growth}\},
		\]
		where $\mathcal{D}'(S^1)$ is the nuclear Fréchet space of Schwartz distributions on $S^1$ (the dual of $C^{\infty}(S^1)$); the Hilbert space isomorphism is recovered by passing to completions.

		This phenomenon has a categorical shadow: continuous finite-dimensional representations of $S^1$ decompose into weight spaces, giving an equivalence
		\[
			\Rep^{\mathrm{fd},\mathrm{cont}}(S^1) \simeq \bigoplus_{n \in \Z} \Vect^{\mathrm{fd}}(\C),
		\]
		sending $V$ to $(V_n)_n$ with $V_n = \{v \in V \mid t \cdot v = t^n v \text{ for all } t \in S^1\}$. This is even a symmetric monoidal equivalence, for the pointwise tensor product on the left and the Day convolution on the right. Heuristically, the right-hand side is ``finite sheaves of $\C$-vector spaces on $\Z$'' and the left-hand side is ``finite sheaves of $\C$-vector spaces on $*/S^1$''.
	\end{example}

	\begin{example}[The multiplicative group]\label[example]{ex:Gm_Fourier}
		The previous example has an elementary algebraic analogue, valid over any (for the moment, ordinary) commutative ring $R$. Viewing $S^1 \subseteq \C^{\times}$ inside the multiplicative group, one can decomplete $L^2(S^1)$ all the way to Laurent polynomials. Namely, for the multiplicative group $\bbG_{m,R} = \Spec(R[t^{\pm 1}])$ we have
		\[
			\Oo(\bbG_{m,R}) \cong \bigoplus_{n \in \Z} R = \{\text{sequences } (a_n)_n \text{ of finite support}\}, \qquad (a_n)_n \mapsto \sum_{n \in \Z} a_n t^n.
		\]
		The right-hand side may be understood as the compactly supported global sections (or homology) of the structure sheaf of the scheme $\Z \times \Spec(R)$, a countable disjoint union of copies of $\Spec(R)$.

		There is a clean categorification. Since representations of $\bbG_{m,R}$ decompose according to weights, there is an equivalence
		\[
			\D_{\qc}(*/\bbG_{m,R}) \simeq \prod_{n \in \Z} \D_{\qc}(\Spec(R)) = \D_{\qc}(\Z \times \Spec(R)).
		\]
		Concretely, an $R$-module $M$ with $\bbG_{m,R}$-action decomposes as a $\Z$-graded module $M = \bigoplus_{n} M_n$ into weight spaces. Dually, there is an equivalence
		\[
			\D_{\qc}(\bbG_{m,R}) \simeq \D_{\qc}(*/\Z \times \Spec(R)),
		\]
		since the right-hand side is a complex of $R$-modules equipped with an automorphism, i.e.\ a complex of $R[t^{\pm 1}]$-modules.

		The equivalences are realized by categorified Fourier transforms. Abbreviate $\bbG_m = \bbG_{m,R}$ and $\Z = \Z \times \Spec(R)$, and consider the correspondence
		\[
			\begin{tikzcd}[row sep=large, column sep=large]
				& */\bbG_m \times \Z \dlar["p"'] \drar["q"] \rar["m"] & */\bbG_m \\
				*/\bbG_m & & \Z,
			\end{tikzcd}
		\]
		where $m$ is induced by the pairing $\bbG_m \times \Z \to \bbG_m$, $(t,n) \mapsto t^n$, and $p, q$ are the projections. Let $\Ll \in \D_{\qc}(*/\bbG_m)$ be the tautological line bundle, corresponding to the tautological representation of $\bbG_m$. Using $\Ll^{-1}$ as a Fourier kernel gives
		\[
			\D_{\qc}(*/\bbG_m) \iso \D_{\qc}(\Z), \qquad M \mapsto q_!(p^*M \otimes m^*\Ll^{-1}),
		\]
		with inverse $N \mapsto p_!(q^*N \otimes m^*\Ll)$. Here $q_!$ plays the role of integration over $*/\bbG_m$, while $p^*M \otimes m^*\Ll^{-1}$ is the analogue of $f(t)t^{-n}$. Note that $q$ is proper in the quasicoherent six-functor formalism, so $q_! = q_*$, while $p$ is étale, so $p^! = p^*$.

		Symmetrically, the correspondence
		\[
			\begin{tikzcd}[row sep=large, column sep=large]
				& \bbG_m \times */\Z \dlar["p"'] \drar["q"] \rar["m"] & */\bbG_m \\
				\bbG_m & & */\Z
			\end{tikzcd}
		\]
		gives the equivalence $\D_{\qc}(\bbG_m) \simeq \D_{\qc}(*/\Z)$. Under it the structure sheaf of $\bbG_m$ corresponds to the regular representation of $\Z$, and passing to endomorphisms recovers the Fourier isomorphism on the level of vector spaces. Working with analytic stacks over the liquid complex numbers, one obtains analogous equivalences for the full derived categories of quasicoherent sheaves on the various ($C^{\infty}$, real-analytic, \dots) versions of $S^1$, categorifying the isomorphisms for $C^{\infty}$-functions or distributions above.
	\end{example}

	\begin{example}[Cartier]\label[example]{ex:Cartier_Finite_Flat}
		Let $R$ be an ordinary commutative ring. Cartier introduced a duality on finite flat commutative group schemes $G = \Spec(\Oo(G))$ over $R$; here $\Oo(G)$ is a finite projective $R$-module carrying a commutative algebra structure together with a compatible cocommutative coalgebra structure encoding the addition on $G$. Basic examples are the constant group $\Z/n$, the group of $n$-th roots of unity
		\[
			\mu_n = \ker(\bbG_m \xrightarrow{t \mapsto t^n} \bbG_m) = \Spec(R[t]/(t^n-1)),
		\]
		and the $n$-torsion $A[n] = \ker(A \xrightarrow{\cdot n} A)$ of an abelian variety $A$. If $R$ has characteristic $0$, then any such $G$ is étale-locally isomorphic to a constant commutative group (for $\mu_n$, one adjoins a primitive $n$-th root of unity); in positive or mixed characteristic, finite flat group schemes are far more interesting.

		As $\Oo(G)$ is simultaneously a commutative algebra and a cocommutative coalgebra, the dual $R$-module $\Oo(G)^*$ is again both, with the roles of algebra and coalgebra exchanged. Thus
		\[
			G^* = \Spec(\Oo(G)^*)
		\]
		is another finite flat group scheme, the \emph{Cartier dual} of $G$. This has a geometric meaning: $G^* = \Hh om(G, \bbG_m)$, the internal hom in, say, fpqc sheaves. It is induced by a natural pairing $m\colon G \times G^* \to \bbG_m$, corresponding to the map $R[t^{\pm 1}] \to \Oo(G) \otimes_R \Oo(G)^*$ sending $t$ to the image of $1$ under the coevaluation $R \to \Oo(G) \otimes_R \Oo(G)^*$ (which turns out to be an invertible element). For instance $(\Z/n)^* = \mu_n$, while $A[n]^* = A^*[n]$ where $A^*$ is the dual abelian variety. Using the correspondence
		\[
			\begin{tikzcd}[row sep=large, column sep=large]
				& */G \times G^* \dlar["p"'] \drar["q"] \rar["m"] & */\bbG_m \\
				*/G & & G^*,
			\end{tikzcd}
		\]
		one again obtains a Fourier equivalence $\D_{\qc}(*/G) \simeq \D_{\qc}(G^*)$, and by symmetry the version exchanging $G$ and $G^*$.
	\end{example}

	\begin{example}[Mukai]\label[example]{ex:Mukai}
		A parallel story applies to abelian varieties themselves. For an abelian variety $A/R$ there is a dual abelian variety
		\[
			A^* = \Ee xt^1(A, \bbG_m),
		\]
		the internal $\Ee xt^1$ in, say, fpqc sheaves. Note that as $\bbG_m$ is affine while $\Oo(A) = R$, one has $\Hh om(A, \bbG_m) = 0$, so one must pass to $\Ee xt^1$ to see anything interesting. There is a tautological pairing $m\colon A \times A^* \to */\bbG_m$ and hence a line bundle
		\[
			\mathcal{P} = m^*\Ll \in \D_{\qc}(A \times A^*),
		\]
		the \emph{Poincaré line bundle}. Mukai observed that the correspondence
		\[
			\begin{tikzcd}[row sep=large, column sep=large]
				& A \times A^* \dlar["p"'] \drar["q"] \rar["m"] & */\bbG_m \\
				A & & A^*
			\end{tikzcd}
		\]
		induces a Fourier equivalence $\D_{\qc}(A) \simeq \D_{\qc}(A^*)$, $M \mapsto p_*(q^*M \otimes \mathcal{P})$; see \citet{FourierMukai2006Huybrechts}. In contrast to the previous examples, here one does \emph{not} pass to classifying spaces.
	\end{example}

	\begin{warning}[Breen]\label[warning]{warn:Breen}
		The preceding examples make it tempting to form duals in general by taking the derived internal $\mathcal{R}\!\Hh om(-, \bbG_m)$. However, the higher $\Ext$-groups are (probably) very sensitive to the choice of Grothendieck topology, hard to compute in practice, and the naively expected vanishing in degrees $> 0$ (resp.\ $> 1$) already fails in the examples above. This makes it difficult to extend Cartier duality naively to ``chain complexes of finite flat group schemes'' or of abelian varieties. The theory of Gestalten circumvents this difficulty.
	\end{warning}

	\begin{example}[Deligne]\label[example]{ex:Deligne_Fourier}
		Similar Fourier equivalences exist for other sheaf theories. A famous example is the $\ell$-adic Fourier transform. Work over a field $k$ of positive characteristic and fix a nontrivial additive character $\psi\colon \F_p \to \overline{\Q}_{\ell}^{\times}$. Via the Artin--Schreier sequence
		\[
			0 \to \F_p \to \bbA^1_k \xrightarrow{x \mapsto x^p - x} \bbA^1_k \to 0,
		\]
		$\psi$ yields a rank-one local system $\Ll_{\psi}$ on $\bbA^1_k$, the \emph{Artin--Schreier sheaf}. For a finite-dimensional $k$-vector space $V$ with dual $V^*$, regarded as affine schemes, the correspondence
		\[
			\begin{tikzcd}[row sep=large, column sep=large]
				& V \times V^* \dlar["p"'] \drar["q"] \rar["m"] & \bbA^1_k \\
				V & & V^*
			\end{tikzcd}
		\]
		(with $m$ the evaluation pairing) induces an equivalence $\D_{\et}(V, \overline{\Q}_{\ell}) \simeq \D_{\et}(V^*, \overline{\Q}_{\ell})$, $M \mapsto p_!(q^*M \otimes m^*\Ll_{\psi})$.
	\end{example}

	\begin{example}[$D$-modules]\label[example]{ex:D_Modules_Fourier}
		The previous example has an analogue for $D$-modules over a field $k$ of characteristic $0$. Recall that via Simpson's de Rham stack one may define $\mathrm{Dmod}(X) = \D_{\qc}(X_{\mathrm{dR}})$, importing the six functors from the quasicoherent formalism. The role of the Artin--Schreier sheaf is played by the \emph{exponential $D$-module} $\Ll$ on $\bbA^1_k$, corresponding to $k[T] \cdot e$ with $\partial(e) = e$ (``$e = \exp(t)$''). This yields a map of commutative group stacks
		\[
			\bbA^1_{\mathrm{dR}} \to */\bbG_m,
		\]
		which may also be described as the composite $\bbA^1_{\mathrm{dR}} = \bbA^1/\widehat{\bbG}_a \to */\widehat{\bbG}_a \xrightarrow{\exp} */\bbG_m$. For $V$ as above, the correspondence
		\[
			\begin{tikzcd}[row sep=large, column sep=large]
				& V_{\mathrm{dR}} \times V^*_{\mathrm{dR}} \dlar["p"'] \drar["q"] \rar["m"] & \bbA^1_{\mathrm{dR}} \rar & */\bbG_m \\
				V_{\mathrm{dR}} & & V^*_{\mathrm{dR}}
			\end{tikzcd}
		\]
		induces a Fourier equivalence $\D_{\qc}(V_{\mathrm{dR}}) \simeq \D_{\qc}(V^*_{\mathrm{dR}})$, $M \mapsto q_!(p^*M \otimes m^*\Ll)$. For $\bbA^1_k$ this has an elementary description in terms of actual $D$-modules (modules over the Weyl algebra $k\langle t, \partial\rangle/([\partial, t] = 1)$) as the automorphism sending $t \mapsto \partial$ and $\partial \mapsto -t$.
	\end{example}

	\begin{remark}\label[remark]{rem:More_Fourier_Examples}
		There are many more dualities of this shape:
		\begin{enumerate}
			\item Over a field $k$ of characteristic $0$, the affine line $\bbA^1$ is Cartier dual to the formal additive group $\widehat{\bbG}_a$;
			\item Over a nonarchimedean field $k$ of characteristic $0$, the analytic affine line $\bbA^{1,\mathrm{an}}$ is Cartier dual to the overconvergent neighborhood of $0$ in $\bbA^1$;
			\item Over $\Q_p$, the open unit disc $D_{\Q_p}$ is Cartier dual to the locally analytic $p$-adic numbers $\Z_p^{\mathrm{la}}$, a version of Amice duality;
			\item Working directly with symmetric monoidal $(\infty,1)$-categories, Lurie has a generalization of Cartier duality.
		\end{enumerate}
		All of these---except the very first, using $L^2$-functions---turn out to be special cases of the general Cartier duality for Gestalten established below.
	\end{remark}

	\section{Cartier duality for Gestalten}
	\label[section]{sec:Cartier_Gestalten}

	Fix a Gestalt $X$ as base. The objects of the duality are \emph{spectrum objects} in Gestalten over $X$. Working in a fully derived setting, commutative ($E_{\infty}$-)groups in $\Gest_{/X}$ give an analogue of the connective part $\D_{\geq 0}(\S) \subseteq \D(\S)$; it is better to work stably and consider
	\[
		\Sp(\Gest_{/X}) := \bigl\{(M_0, M_1, \dots) \mid M_k \in \CGrp(\Gest_{/X}),\ M_k \iso \Omega M_{k+1}\bigr\}.
	\]
	For each $n \geq 0$ we have the full subcategory $\Sp(\Gest^{n\text{-}\et}_{/X}) \subseteq \Sp(\Gest_{/X})$ on those sequences with all $M_k$ being $n$-étale over $X$. By ambidexterity (\Cref{prop:Ambidexterity_Stefanich_Rings}) an $n$-étale map is $(n+1)$-proper and vice versa, and so
	\[
		\bigcup_n \Sp(\Gest^{n\text{-}\et}_{/X}) = \bigcup_n \Sp(\Gest^{n\text{-prop}}_{/X}).
	\]

	\begin{remark}
		Equivalently, $\Sp(\Gest^{n\text{-}\et}_{/X})$ is the category of $\aleph_1$-compact objects in the category of sheaves of $\S$-modules on the $(\infty,1)$-topos $\Ind_{\aleph_1}(\Gest^{n\text{-}\et}_{/X})$.
	\end{remark}

	\begin{construction}\label[construction]{con:GL1}
		There is a canonical object $\GL_1 \in \Sp(\Gest_{/X})$ representing the functor
		\[
			\Gest_{/X}\catop \to \Sp, \qquad Y \mapsto (\Oo(Y)^{\times}_0, \Oo(Y)^{\times}_1, \Oo(Y)^{\times}_2, \dots),
		\]
		where $\Oo(Y)^{\times}_n \subseteq \Oo(Y)_n$ is the subanima of $\otimes$-invertible objects, a commutative group under tensor product. This is indeed a spectrum object: $\Omega(\Oo(Y)^{\times}_{n+1})$ is the subanima of $\Hom_{\Oo(Y)_{n+1}}(\unit, \unit) \simeq \Oo(Y)_n$ spanned by the invertible endomorphisms of $\unit$, and the monoidal structure on $\Oo(Y)_n$ coming from composition agrees with the one inherited from $\Oo(Y)_{n+1}$, so an endomorphism is invertible precisely when the corresponding object of $\Oo(Y)_n$ is $\otimes$-invertible. We will in fact see that $\GL_1 \in \Sp(\Gest^{0\text{-prop}}_{/X})$.
	\end{construction}

	\begin{theorem}[Scholze--Stefanich]\label[theorem]{thm:Cartier_Duality}
		For every $M \in \bigcup_n \Sp(\Gest^{n\text{-}\et}_{/X})$, the functor
		\[
			\Sp(\Gest_{/X})\catop \to \An, \qquad N \mapsto \Hom(N \otimes M, \GL_1)
		\]
		is representable by an object
		\[
			M^* \; := \; \Hh om(M, \GL_1) \qin \bigcup_n \Sp(\Gest^{n\text{-}\et}_{/X}).
		\]
		For every $n \geq 0$, the assignment $M \mapsto M^*$ induces an equivalence
		\[
			\Sp(\Gest^{n\text{-}\et}_{/X}) \; \simeq \; \Sp(\Gest^{n\text{-prop}}_{/X})\catop;
		\]
		since $n$-proper maps are $(n+1)$-étale and vice versa, $(-)^*$ is thus an anti-self-equivalence of $\bigcup_n \Sp(\Gest^{n\text{-}\et}_{/X})$. Moreover, for every such $M$, all $m \geq n$, and all $0 \leq k \leq m$, there is an equivalence
		\[
			(\Oo(M_k)/\Oo(X))^!_m \; \simeq \; (\Oo(M^*_{m-k})/\Oo(X))_m,
		\]
		given by an explicit Fourier transform. (If $m > n$, the left-hand side is just $(\Oo(M_k)/\Oo(X))_m$ by ambidexterity.)
	\end{theorem}

	\begin{remark}
		All the dualities of \Cref{sec:Cartier_Examples}, except the $L^2$-example, are, suitably interpreted, instances of this theorem. For example, the connective spectrum objects $\Z$ and $\bbG_m = \Gest_X(\S[t^{\pm 1}])$ over $X = \Gest(\D(\S))$ are Cartier dual.
	\end{remark}

	\begin{remark}
		For $m > n$ the object $(\Oo(M_k)/\Oo(X))_m$ is dualizable, and the equivalence of \Cref{thm:Cartier_Duality} is induced by the canonical invertible object in
		\[
			(\Oo(M^*_{m-k})/\Oo(X))_m \otimes (\Oo(M_k)/\Oo(X))_m \cong (\Oo(M^*_{m-k} \times M_k)/\Oo(X))_m
		\]
		coming from the pairing $M^*_{m-k} \otimes M_k \to (\GL_1)_m$.
	\end{remark}

	The proof of the theorem is given in the next talk. We now discuss some examples and consequences.

	\begin{corollary}\label[corollary]{cor:Cartier_Biduality}
		For every $M \in \bigcup_n \Sp(\Gest^{n\text{-}\et}_{/X})$, the biduality map $M \to (M^*)^*$ is an isomorphism.
	\end{corollary}
	\begin{proof}
		Both sides lie in $\bigcup_n \Sp(\Gest^{n\text{-}\et}_{/X})$. For any $N$ in this union,
		\[
			\Hom(N, (M^*)^*) \simeq \Hom(N \otimes M^*, \GL_1) \simeq \Hom(M^*, N^*),
		\]
		using the symmetry of the tensor pairing. Under these identifications, the biduality map corresponds to the map $\Hom(N, M) \to \Hom(M^*, N^*)$ given by contravariant functoriality of $(-)^*$. As $(-)^*$ is an anti-self-equivalence of the union (\Cref{thm:Cartier_Duality}), this map is a bijection, so $M \to (M^*)^*$ is an isomorphism.
	\end{proof}

	Strikingly, \Cref{thm:Cartier_Duality} imposes essentially no finiteness hypothesis on $M$. The following example makes the resulting bidualities concrete.

	\begin{example}\label[example]{ex:Product_Sum_Fp}
		Work over $X = \Gest(\D(\C))$ (in fact $\Z[\tfrac1p, \zeta_p]$ would suffice) and fix a prime $p$. Then $\Z$ is Cartier dual to $\bbG_m = \Gest_X(\C[t^{\pm 1}])$. As $\Hh om(-, \GL_1)$ is exact, $\Z/p$ is Cartier dual to
		\[
			\mu_p = \ker(\bbG_m \xrightarrow{t \mapsto t^p} \bbG_m) \cong \Z/p
		\]
		(the isomorphism depending on a choice of $p$-th root of unity). On the derived category of sheaves of $\F_p$-vector spaces, Cartier duality becomes the naive duality $\Hh om_{\F_p}(-, \F_p)$. Using biduality (with all sums and products taken in the derived category) we find:
		\begin{enumerate}
			\item $\bigoplus_{\N} \F_p$ is $0$-étale, with dual the $0$-proper $\prod_{\N} \F_p$, which is therefore $1$-étale;
			\item $\bigoplus_{\N}\prod_{\N} \F_p$ is $1$-étale, with dual the $1$-proper (hence $2$-étale) $\prod_{\N}\bigoplus_{\N} \F_p$;
			\item $\bigoplus_{\N}\prod_{\N}\bigoplus_{\N}\F_p$ is $2$-étale, with dual the $2$-proper (hence $3$-étale) $\prod_{\N}\bigoplus_{\N}\prod_{\N}\F_p$; and so on.
		\end{enumerate}
		It is remarkable that any topos supports these bidualities. In naive $\F_p$-vector spaces the very first biduality for $\prod_{\N}\F_p$ already fails; it can be repaired in condensed $\F_p$-vector spaces, where it is consistent that the derived biduality holds for $\prod_{\N}\bigoplus_{\N}\F_p$---but this is independent of $\mathrm{ZFC}$ (it fails if the continuum is ``too small''), the third biduality certainly fails there, and everything fails for uncountable index sets (already for $\aleph_1$). Putting a Gestalt structure on these vector-space objects thus yields a surprisingly clean duality theory. The restriction to countable sums and products reflects our foundational choice of $\aleph_1$-presentable categories; a larger cutoff cardinal would permit larger index sets.
	\end{example}

	\section{Recovering the classical examples}
	\label[section]{sec:Cartier_Classical}

	We now indicate how the internal hom into $\GL_1$ recovers the classical duals. Consider $\GL_1 \in \Sp(\Gest_{/X})$. For each $n \geq 0$, we denote by $(\GL_1)_n \in \CGrp(\Gest_{/X})$ the commutative group representing the functor
	\[
		\Gest_{/X}\catop \to \CGrp(\An), \qquad X' \mapsto \Oo(X')^{\times}_n.
	\]
	We also denote by $(\bbA^1)_n \in \CMon(\Gest_{/X})$ the commutative monoid representing
	\[
		\Gest_{/X}\catop \to \CMon(\An), \qquad X' \mapsto \Oo(X')^{\simeq, \aleph_1}_n,
	\]
	the nonoid of $\aleph_1$-compact objects of $\Oo(X')_n$ under tensor product. Note that we have an inclusion $(\GL_1)_n \subseteq (\bbA^1)_n$ of commutative monoids. We then have the following result:

	\begin{theorem}\label[theorem]{thm:Internal_Hom_GL1}
		Let $X = \Gest(A)$ and let $G \in \CGrp(\Gest_{/X})$ be $n$-suave. Then the internal hom
		\[
			\Hh om_{\CGrp}(G, (\GL_1)_n)
		\]
		is representable by the $n$-affine Gestalt $\Gest\bigl((\Oo(G)/A)^!_n\bigr)$. More generally, for every $G \in \CMon(\Gest^{n\text{-suave}}_{/X})$ one has
		\[
			\Hh om_{\CMon(\Gest_{/X})}(G, (\bbA^1)_n) \simeq \Gest_X\bigl((\Oo(G)/A)^!_n\bigr).
		\]
	\end{theorem}

	Again, we will prove this result next week. Before continuing, let us provide more details about where the commutative algebra structure on $(\Oo(G)/A)^!_n$ comes from.
	
	\begin{construction}
		We construct a symmetric monoidal functor
		\[
			(\Gest^{n\text{-suave}}_{/X}, \times) \to A_{n+1}, \qquad Y \mapsto (\Oo(Y)/A)^!_n.
		\]
		Passing to commutative algebra objects then gives
		\[
			\CMon(\Gest^{n\text{-suave}}_{/X}) \simeq \CAlg(\Gest^{n\text{-suave}}_{/X}, \times) \to \CAlg(A_{n+1}),
		\]
		which endows $(\Oo(G)/A)^!_n$ with the commutative algebra structure induced by the group structure on $G$.

		To construct this functor, consider the lax monoidal right adjoint
		\[
			(\Oo(-)/A)_{n+1} \colon (\Gest^{n\text{-suave}}_{/X})\catop \to \CAlg(A_{n+2}) \to A_{n+2}
		\]
		to the $(n+1)$-affine inclusion. Passing to slices results in a lax symmetric monoidal functor
		\[
			(\Gest^{n\text{-suave}}_{/X})\catop \to (A_{n+2})_{\unit/}.
		\]
		This functor sends an $n$-suave map $f\colon Y \to X$ of Gestalten to the unit map $f_n^*\colon \unit \to (\Oo(Y)/A)_{n+1}$ in $A_{n+2}$, which by assumption admits a left adjoint $f_{n,\sharp}$. So this functor lands in the full subcategory $(A_{n+2})_{\unit/}^{\mathrm{adj}}$ of morphisms $\unit \to C$ admitting left adjoints. This subcategory admits a functor to
		\[
			\End_{A_{n+2}}(\unit) \simeq A_{n+1},
		\]
		by sending an arrow $u\colon \unit \to C$ with left adjoint $u_{\sharp}$ to the composite $u_{\sharp}u$; for a map $u \to v$ under $\unit$, say $v = ru$, the induced map $v_{\sharp}v \to u_{\sharp}u$ is obtained from the mate $v_{\sharp}r \to u_{\sharp}$ after whiskering with $u$. \red{[Q: How does one rigorously define this?]} In fact, since the tensor product on $A_{n+2}$ is a 2-functor in both variables, it preserves adjunctions, and the functor
		\[
			((A_{n+2})_{\unit/}^{\mathrm{adj}})\catop \to A_{n+1}
		\]
		is symmetric monoidal. \red{[Verify that that's an actual argument.]}
		Composing with our earlier functor gives the oplax monoidal functor
		\[
			\Gest^{n\text{-suave}}_{/X} \to A_{n+1}, \qquad (f\colon Y \to X) \mapsto (\Oo(Y)/A)^!_n = f_{n,\sharp}f_n^*\unit.
		\]
		For symmetric monoidality, it remains to check that the oplax comparison maps
		\[
			(\Oo(Y \times_X Z)/A)^!_n \to (\Oo(Y)/A)^!_n \otimes_{A_{n+1}} (\Oo(Z)/A)^!_n
		\]
		are isomorphisms. But this is an instance of the Künneth formula for suave maps \Cref{lem:Kunneth_Suave}.
	\end{construction}

	\begin{proposition}\label[proposition]{prop:Cartier_Dual_Suave_Group}
		Let $G \in \CGrp(\Gest_{/X})$ be $n$-suave with $BG = */G$ being $(n+1)$-affine. Then $G^* \in \Sp(\Gest_{/X})$ is $(-n)$-connective, and is determined by
		\[
			G^*_n = \Hh om_{\CGrp(\Gest_{/X})}(G, (\GL_1)_n) = \Gest\bigl((\Oo(G)/A)^!_n\bigr).
		\]
	\end{proposition}

	\begin{example}\label[example]{ex:Z_Gm_Cartier}
		Let $X = \Gest(\D(R))$ and let $G = \Z = \bigsqcup_{\Z} X$ with the commutative group structure inherited from $\Z$ (corresponding to the ring $\prod_{\Z} R$). This is $0$-étale, in particular $0$-suave, with
		\[
			(\Oo(G)/R)^!_0 = \bigoplus_{\Z} R \qin \D(R).
		\]
		Indeed, $n$-suave maps of Stefanich rings are stable under countable limits (\Cref{prop:Suave_Countable_Limits}), and the assignment $B \mapsto (B/A)^!_m$ carries limits to colimits, so the countable product $\prod_{\Z} R$ is sent to the countable coproduct $\bigoplus_{\Z} R$. The algebra structure is induced by the effect of $0\colon * \to \Z$ and $+\colon \Z^2 \to \Z$ on indexing sets,
		\[
			R = \bigoplus_{*} R \to \bigoplus_{\Z} R, \qquad \bigoplus_{\Z} R \otimes \bigoplus_{\Z} R \cong \bigoplus_{\Z^2} R \to \bigoplus_{\Z} R,
		\]
		so that $\Oo(\bbG_{m,R})_0 = R[t^{\pm 1}]$, as expected.
	\end{example}

	\begin{example}\label[example]{ex:Finite_Flat_Cartier}
		Let $G$ be a finite flat group scheme over $R$. Classically $\Oo(G)$ is a local complete intersection, hence $G$ is $0$-suave. Taking $n = 0$ and using that $\Oo(G)$ is finite projective,
		\[
			(\Oo(G)/A)^!_0 = \Oo(G)^* = \Oo(G^*),
		\]
		recovering the classical Cartier dual.
	\end{example}

	\section{Brown--Comenetz duality}
	\label[section]{sec:Brown_Comenetz}

	\begin{example}\label[example]{ex:GL1_Global_Sections}
		Let $X = \Gest(\D(\C))$ (the cyclotomic extension of $\Q$ would suffice for this discussion). The global sections of $\GL_1$ form the spectrum
		\[
			\GL_1(X) = (\C^{\times}, \D(\C)^{\times}, 1\Pr^{\times}_{\C}, \dots).
		\]
		Its homotopy groups are hard to compute in general, but:
		\begin{itemize}
			\item $\pi_0 = \C^{\times}$;
			\item $\pi_{-1} = \Z$, since any $\otimes$-invertible object of $\D(\C)$ is a shift $\C[n]$;
			\item $\pi_{-2} = 0$, by a theorem of Stefanich \citep{Stefanich_HigherQCoh};
			\item $\pi_{-3} = 0$ (working with ``small'' higher linear categories);
			\item $\pi_{-4}$ is the ``Witt group'', a highly nontrivial but known abelian group, countably but not finitely generated.
		\end{itemize}
		In the language of mathematical physics, these higher invertible linear categories give rise to \emph{invertible topological quantum field theories}, and there is an expectation (Kitaev and others) that these are classified by the Brown--Comenetz dual $I_{\C^{\times}}$ of the sphere spectrum, with
		\[
			\pi_{-n} I_{\C^{\times}} = \Hom(\pi_n \S, \C^{\times}).
		\]
		The rationale is that an $n$-dimensional invertible TQFT should assign to a framed closed $n$-manifold an invertible number in $\C^{\times}$, cobordism-invariantly, and stably framed cobordism classes of closed $n$-manifolds are given by $\pi_n \S$. 
		
		Note that this expectation differs drastically from the actual homotopy groups of $\GL_1(X)$ given above. However, looking not at global sections but at $\GL_1$ as a sheaf on Gestalten recovers the expected answer:
	\end{example}

	\begin{corollary}\label[corollary]{cor:GL1_Brown_Comenetz}
		The spectrum object $\GL_1$ in $\Gest_{/\D(\C)}$ is an internal Brown--Comenetz dual $I_{\bbG_m}$, in the sense that
		\[
			\Hh om(\Z, \GL_1) = \bbG_m.
		\]
		In particular, for all $n > 0$,
		\[
			\pi_{-n}\GL_1 \cong \Hom(\pi_n \S, \pi_0 \bbG_m) = \Hom(\pi_n \S, \C^{\times})
		\]
		as $0$-truncated abelian groups in $\Gest_{/\D(\C)}$, via the pairing $\pi_n \S \otimes \pi_{-n}\GL_1 \to \pi_0 \bbG_m$.
	\end{corollary}
	\begin{proof}
		The first assertion is the Cartier duality between $\Z$ and $\bbG_m$. For the negative homotopy groups, note that $\bbG_m$ is the connective cover of $\GL_1$, and its $1$-connective cover $\tau_{\geq 1}\bbG_m$ is rational: for any $A \in \CAlg(\D(\C))$ one has $\tau_{> 0}A^{\times} = \tau_{> 0}A$ via the logarithm, and $\tau_{> 0}A$ is rational. Hence
		\[
			\Hh om(\Z, \pi_0\GL_1) = \pi_0\bbG_m,
		\]
		with the right-hand side concentrated in degree $0$. The identification of the negative homotopy groups now follows from general Brown--Comenetz nonsense: in any countable $(\infty,1)$-topos, a coconnective spectrum object $G$ with $\Hh om(\Z, G) = A$ concentrated in degree $0$ and divisible satisfies $\pi_{-n}G = \Hom(\pi_n \S, A)$ via the pairing $\pi_n \S \otimes \pi_{-n}G \to \pi_0 G = A$. Since $\pi_n \S$ is a finite abelian group, so is $\Hom(\pi_n \S, \C^{\times})$; this generalizes the earlier computation $\pi_0 \Pic = \mu_2$ of \Cref{prop:Relation_Pic_BG_m}.
	\end{proof}

	\section{The motivic Fourier transform}
	\label[section]{sec:Motivic_Fourier}

	Most of the examples of \Cref{sec:Cartier_Examples} fit transparently into the setting of Gestalten. The subtlest is Deligne's $\ell$-adic Fourier transform (\Cref{ex:Deligne_Fourier}); to accommodate it one uses the transmutation of \Cref{chap:Transmutations}, turning a six-functor formalism into a Gestalt. One could work with the $\ell$-adic formalism directly, but a more refined discussion is available in the motivic six-functor formalism $\SH$. This example also highlights genuinely new geometry made possible by Gestalten.

	\begin{example}\label[example]{ex:Motivic_Deligne}
		Let $X = [\Spec(\Z)]_{\SH} \in \Gest$. Over it we have the transmutation $[\bbA^1_{\Z}]_{\SH}$ of the affine line, which carries a commutative group structure (via addition). In fact, it carries a commutative ring structure, and in particular a module structure over itself. Hence the Cartier dual $[\bbA^1_{\Z}]_{\SH}^*$ is again a $[\bbA^1_{\Z}]_{\SH}$-module. If it were isomorphic to $[\bbA^1_{\Z}]_{\SH}$ we would obtain a Fourier self-duality. This is not quite true, but it is true locally.
	\end{example}

	\begin{proposition}\label[proposition]{prop:Motivic_Fourier_Local}
		There is a cover $\widetilde{X} \to X$ such that
		\[
			[\bbA^1_{\Z}]^*_{\SH} \times_X \widetilde{X} \simeq \bigl([\bbA^1_{\Z}]_{\SH} \times_X \widetilde{X}\bigr)[-1]
		\]
		as $[\bbA^1_{\Z}]_{\SH} \times_X \widetilde{X}$-modules.
	\end{proposition}

	The universal such $\widetilde{X}$ is a $[\bbG_{m,\Z}]_{\SH}$-torsor $\widetilde{X} \to X$. Roughly, $\widetilde{X}$ is obtained by freely adjoining an exponential local system on $\bbA^1_{\Z}$ to the motivic six-functor formalism; the resulting six-functor formalism $S \mapsto \Oo([S]_{\SH} \times_X \widetilde{X})_1$ is a theory of \emph{exponential motives}. Over
	\[
		[\Spec(\F_p)]_{\SH} \times_{\Gest(\S)} \Gest(\S[\tfrac1p, \zeta_p]) \to X,
	\]
	one obtains an additive character $\psi$, and the Artin--Schreier cover yields an exponential local system on $\bbA^1_{\F_p}$ splitting this map to $\widetilde{X}$. Writing $[\bbA^1_{\F_p}]_{\SH}[p] := [\bbA^1_{\F_p}]_{\SH} \times_{\Gest(\S)} \Gest(\S[\tfrac1p, \zeta_p])$, one gets an isomorphism
	\[
		[\bbA^1_{\F_p}]_{\SH}^*[p] \simeq [\bbA^1_{\F_p}]_{\SH}[p][-1],
	\]
	a motivic version of Deligne's Fourier equivalence; base-changing to the Gestalt of the $\ell$-adic six-functor formalism recovers Deligne's equivalence itself.

	Finally, there is a Betti realization
	\[
		\Gest(\D(\S)) \to [\Spec(\Z)]_{\SH}
	\]
	corresponding to the ring Gestalt $\underline{\R}$ over $\Gest(\D(\S))$: the pullback of $[\bbA^1_{\Z}]_{\SH}$ is $\underline{\R} = \Gest(\Shv(\R, \D(\S)))$ over $\Gest(\D(\S))$. Under this realization the torsor $\widetilde{X} \to X$ yields an $\underline{\R}^{\times}$-torsor over $\Gest(\D(\S))$ which, up to a trivial $\{\pm 1\}$-torsor, is the $\underline{\R}_{> 0}$-torsor $\Gest(W) \to \Gest(\D(\S))$ of Witt vectors discussed by \citet{Scholze_Wild_Betti}.

	\chapter{The proof of Cartier duality}

	\emph{Talk given by Sil.}

	\appendix

	\titleformat{\chapter}{\normalfont\huge\bfseries}{\Alph{chapter}}{20pt}{\Huge}
	\titleformat{\section}{\Large\bfseries}{\Alph{chapter}.\arabic{section}}{10pt}{\Large}
	\titleformat{\subsection}{\large\bfseries}{\Alph{chapter}.\arabic{section}.\arabic{subsection}}{6pt}{\large}
	\titleformat{\subsubsection}{\large\bfseries}{}{6pt}{\large}
	\addtocontents{toc}{\vskip 15pt}
	\renewcommand{\thechapter}{\Alph{chapter}}
	\renewcommand{\thesection}{\Alph{chapter}.\arabic{section}}
	\renewcommand{\thesubsection}{\Alph{chapter}.\arabic{section}.\arabic{subsection}}
	\renewcommand{\thetheorem}{\Alph{chapter}.\arabic{theorem}}

	\chapter{\texorpdfstring{$\aleph_1$}{ℵ\_1}-compact objects}
	\label[appendix]{chap:Compact_Objects}

	The formalism developed in these notes heavily relies on manipulations with $\kappa$-compact objects for uncountable cardinals $\kappa$ (which we mostly fix to be $\aleph_1$). In this appendix, we record various basic properties of $\kappa$-compact objects.

	\emph{Bastiaan: I thank Peter Scholze for explaining the key steps in these arguments that I wasn't able to figure out.}

	\section{Compact algebras and modules}
	\label[section]{sec:Compact_Algebras}

	We state and prove various basic properties of $\kappa$-compact algebras in a $\kappa$-presentable category.

	\begin{lemma}\label[lemma]{lem:Kappa_Compact_Algebra_Iff_Module}
		Let $\kappa$ be an uncountable regular cardinal, let $C \in \CAlg(\PrL_{\kappa})$, let $R \in \CAlg(C)$, and let $B \in \CAlg_R(C)$. Then $B$ is $\kappa$-compact as an $R$-algebra if and only if its underlying $R$-module is $\kappa$-compact.
	\end{lemma}
	\begin{proof}
		The point is that the $E_{\infty}$-operad is countable. We write
		\[
			\operatorname{Sym}_R\colon \Mod_R(C) \rightleftarrows \CAlg_R(C) \colon U
		\]
		for the free-forgetful adjunction. We will use that $\kappa$-small colimits of $\kappa$-compact objects are again $\kappa$-compact: mapping out of such a colimit gives a $\kappa$-small limit of anima, and $\kappa$-filtered colimits of animae commute with $\kappa$-small limits.

		First let $M \in \Mod_R(C)$ be $\kappa$-compact. The underlying $R$-module of the free commutative $R$-algebra on $M$ is
		\[
			U\operatorname{Sym}_R(M)
			\simeq
			\coprod_{d \geq 0} (M^{\otimes_R d})_{h\Sigma_d}.
		\]
		Here each tensor power is $\kappa$-compact, the homotopy orbits are a countable colimit, and the coproduct is countable. Thus $U\operatorname{Sym}_R(M)$ is $\kappa$-compact. Also $\operatorname{Sym}_R(M)$ is $\kappa$-compact as an algebra, since $U$ preserves $\kappa$-filtered colimits.

		Since the monad $U\operatorname{Sym}_R$ is $\kappa$-accessible and preserves $\kappa$-compact objects by the formula above, the $\kappa$-compact objects in $\CAlg_R(C)$ are generated under $\kappa$-small colimits and retracts by free algebras on $\kappa$-compact modules. Indeed, for any $\kappa$-accessible monad on a $\kappa$-presentable category which preserves $\kappa$-compact objects, the free algebras on $\kappa$-compact objects are $\kappa$-compact and generate the algebra category under colimits, so the $\kappa$-compact algebras are the retracts of their $\kappa$-small colimits.

		It remains to check that the class of $R$-algebras with $\kappa$-compact underlying $R$-module is closed under these operations. Closure under retracts is clear. For $\kappa$-small colimits, write such a colimit as the geometric realization of its simplicial replacement; the $n$-simplices are $\kappa$-small coproducts of the algebras in the diagram. A $\kappa$-small coproduct of commutative $R$-algebras is the sifted colimit over its finite sub-coproducts, and the underlying $R$-module of a finite coproduct is the corresponding finite tensor product in $\Mod_R(C)$. Since $U$ preserves sifted colimits, and since $\kappa$-compact $R$-modules are closed under finite tensor products and $\kappa$-small colimits, the underlying module of each simplicial degree, and hence of the realization, is $\kappa$-compact. Thus a $\kappa$-compact $R$-algebra has $\kappa$-compact underlying $R$-module.

		Conversely, suppose that $U(B)$ is $\kappa$-compact. By monadicity, $B$ is the geometric realization of the standard simplicial resolution by free algebras,
		\[
			\cdots \rightrightarrows
			\operatorname{Sym}_R U\operatorname{Sym}_R U(B)
			\rightrightarrows
			\operatorname{Sym}_R U(B).
		\]
		By the free-algebra case above, every term in this simplicial object is $\kappa$-compact as an algebra. Since the geometric realization is a countable colimit, $B$ is $\kappa$-compact in $\CAlg_R(C)$.
	\end{proof}

	\begin{lemma}\label[lemma]{lem:Kappa_Compact_Modules_Change_Of_Scalars}
		Let $\kappa$ be an uncountable regular cardinal, let $C \in \CAlg(\PrL_{\kappa})$, and let $A \to B$ be $\kappa$-compact in $\CAlg_A(C)$. Then an object $M \in \Mod_B(C)$ is $\kappa$-compact as a $B$-module if and only if its underlying $A$-module is $\kappa$-compact.
	\end{lemma}
	\begin{proof}
		By \Cref{lem:Kappa_Compact_Algebra_Iff_Module}, the underlying $A$-module of $B$ is $\kappa$-compact. We write
		\[
			f^* = B \otimes_A - \colon \Mod_A(C) \rightleftarrows \Mod_B(C) \colon f_*
		\]
		for the extension-restriction adjunction. First suppose that $M$ is $\kappa$-compact as a $B$-module. The $\kappa$-compact objects in $\Mod_B(C)$ are generated under $\kappa$-small colimits and retracts by the free $B$-modules $B \otimes X$ with $X \in C^{\kappa}$. Since $f_*$ preserves colimits and sends such a free module to $B \otimes X$, which is $\kappa$-compact as an $A$-module because the $C$-action on $\Mod_A(C)$ preserves $\kappa$-compact objects, it follows that $f_*M$ is $\kappa$-compact.

		Conversely, suppose that $f_*M$ is $\kappa$-compact as an $A$-module. The standard monadic resolution
		\[
			M \simeq \left| \cdots \rightrightarrows f^*f_*f^*f_*M \rightrightarrows f^*f_*M \right|
		\]
		expresses $M$ as a countable colimit of $B$-modules of the form $f^*N$. Since $f_*f^*N \simeq B \otimes_A N$ and the tensor product on $\Mod_A(C)$ preserves $\kappa$-compact objects in both variables, the functor $f_*f^*$ preserves $\kappa$-compact $A$-modules. Hence every term in this resolution is of the form $f^*N$ with $N$ $\kappa$-compact as an $A$-module, and is therefore $\kappa$-compact as a $B$-module. Since the realization is a countable, hence $\kappa$-small, colimit, $M$ is $\kappa$-compact.
	\end{proof}

	\begin{corollary}\label[corollary]{cor:Restriction_Kappa_Compact_Algebra_Preserves_Compacts}
		Let $\kappa$ be an uncountable regular cardinal, let $C \in \CAlg(\PrL_{\kappa})$, and let $g\colon R \to S$ be $\kappa$-compact in $\CAlg_R(C)$. Then the restriction functor
		\[
			g_*\colon \Mod_S(C) \to \Mod_R(C)
		\]
		preserves $\kappa$-compact objects.
	\end{corollary}
	\begin{proof}
		This is one direction of \Cref{lem:Kappa_Compact_Modules_Change_Of_Scalars}.
	\end{proof}

	\begin{lemma}\label[lemma]{lem:Kappa_Compact_Algebra_Closure_Properties}
		Let $\kappa$ be an uncountable regular cardinal, let $C \in \CAlg(\PrL_{\kappa})$, and let $A \in \CAlg(C)$.
		\begin{enumerate}[(1)]
			\item The $\kappa$-compact objects in $\CAlg_A(C)$ are closed under $\kappa$-small colimits and retracts. In particular, they are closed under finite colimits.
			\item If $B \in \CAlg_A(C)$ is $\kappa$-compact and $A \to A'$ is any map in $\CAlg(C)$, then $A' \otimes_A B \in \CAlg_{A'}(C)$ is $\kappa$-compact.
			\item If $B \in \CAlg_A(C)$ is $\kappa$-compact and $D \in \CAlg_B(C)$ is $\kappa$-compact, then $D$ is $\kappa$-compact as an object of $\CAlg_A(C)$.
			\item If $B,D \in \CAlg_A(C)$ are $\kappa$-compact and $B \to D$ is a map in $\CAlg_A(C)$, then $D$ is $\kappa$-compact as an object of $\CAlg_B(C)$.
		\end{enumerate}
	\end{lemma}
	\begin{proof}
		Part (1) is the general closure property of $\kappa$-compact objects in a $\kappa$-presentable category. For (2), extension of scalars
		\[
			A' \otimes_A - \colon \CAlg_A(C) \to \CAlg_{A'}(C)
		\]
		is left adjoint to the forgetful functor, and this forgetful functor preserves $\kappa$-filtered colimits.

		For (3), \Cref{lem:Kappa_Compact_Algebra_Iff_Module} shows that $D$ is $\kappa$-compact as a $B$-module. By \Cref{lem:Kappa_Compact_Modules_Change_Of_Scalars}, its underlying $A$-module is $\kappa$-compact, so \Cref{lem:Kappa_Compact_Algebra_Iff_Module} again shows that $D$ is $\kappa$-compact as an $A$-algebra.

		For (4), the underlying $A$-module of $D$ is $\kappa$-compact by \Cref{lem:Kappa_Compact_Algebra_Iff_Module}. Since $B$ is $\kappa$-compact over $A$, \Cref{lem:Kappa_Compact_Modules_Change_Of_Scalars} implies that $D$ is $\kappa$-compact as a $B$-module. Applying \Cref{lem:Kappa_Compact_Algebra_Iff_Module} once more, $D$ is $\kappa$-compact as a $B$-algebra.
	\end{proof}

	\begin{lemma}\label[lemma]{lem:Kappa_Compact_Algebra_Section}
		Let $\kappa$ be an uncountable regular cardinal, let $C \in \CAlg(\PrL_{\kappa})$, and let $A \to B$ be $\kappa$-compact in $\CAlg_A(C)$. If $A \to B$ admits an $A$-algebra retraction $B \to A$, then $A$ is $\kappa$-compact as an object of $\CAlg_B(C)$.
	\end{lemma}
	\begin{proof}
		Apply \Cref{lem:Kappa_Compact_Algebra_Closure_Properties}(4) to the map $B \to A$ in $\CAlg_A(C)$: both $B$ and the initial $A$-algebra $A$ are $\kappa$-compact in $\CAlg_A(C)$.
	\end{proof}

	\section{Compact objects in sheaf categories}
	\label[section]{sec:Compacts_In_Sheaves}

	We now collect the compactness facts for presheaves and left exact
	localizations which are used in the main text.

	\begin{lemma}\label[lemma]{lem:Finite_Limits_Kappa_Coproducts_Representables}
		Let $\kappa$ be an uncountable regular cardinal, and let $C$ be a
		small category with finite limits. Then finite limits of
		$\kappa$-small coproducts of representables in $\PSh(C)$ are again
		$\kappa$-small coproducts of representables.
	\end{lemma}
	\begin{proof}
		The terminal object is representable, so it is enough to prove
		stability under pullbacks. Let
		\[
			F=\coprod_{a\in A} y(c_a),
			\qquad
			G=\coprod_{b\in B} y(d_b),
			\qquad
			H=\coprod_{e\in E} y(h_e)
		\]
		with $A,B,E$ $\kappa$-small. Consider maps $F\to H$ and $G\to H$.
		For each $a\in A$, the restriction $y(c_a)\to H$ chooses a summand
		$e(a)\in E$ and a map $c_a\to h_{e(a)}$ in $C$. Similarly, for each
		$b\in B$, the restriction $y(d_b)\to H$ chooses a summand
		$e'(b)\in E$ and a map $d_b\to h_{e'(b)}$.

		Coproducts in a presheaf $\infty$-topos are disjoint and universal.
		Therefore the pullback $F\times_H G$ decomposes as
		\[
			F\times_H G
			\simeq
			\coprod_{\substack{(a,b)\in A\times B\\ e(a)=e'(b)}}
			\left(y(c_a)\times_{y(h_{e(a)})}y(d_b)\right).
		\]
		Since the Yoneda embedding preserves finite limits, the summand is
		represented by $c_a\times_{h_{e(a)}}d_b$. The displayed coproduct is
		$\kappa$-small, because $\kappa$ is infinite and $A$ and $B$ are
		$\kappa$-small. This proves stability under pullbacks, and hence under
		all finite limits.
	\end{proof}

	\begin{lemma}\label[lemma]{lem:Ex_Infinity_Countable_Control}
		Let $C$ be a small category, and let
		$Y_{\bullet}\in\Fun(\simp\catop,\PSh(C))$ be a simplicial presheaf.
		For a finite simplicial set $K$, write
		\[
			Y(K):=\lim_{\Delta^n\to K}Y_n
		\]
		for the right Kan extension of $Y_{\bullet}$ from simplices to finite
		simplicial sets. For $r\geq 0$, set
		\[
			(\Ex^{r}(Y))_q:=Y(\sd^r\Delta^q),
			\qquad
			\Ex^{\infty}(Y):=\colim_{r\in\N}\Ex^{r}(Y).
		\]
		Then:
		\begin{enumerate}[(1)]
			\item Each $(\Ex^{r}(Y))_q$ is a finite limit of values of
			$Y_{\bullet}$.
			\item The natural map $Y_{\bullet}\to\Ex^{\infty}(Y)$ induces an
			isomorphism on geometric realizations, and
			$(\Ex^{\infty}(Y))_0\simeq Y_0$.
			\item The simplicial presheaf $\Ex^{\infty}(Y)$ is a Kan complex
			in $\PSh(C)$, i.e.\ the maps
			\[
				\Ex^{\infty}(Y)(\Delta^n)\to
				\Ex^{\infty}(Y)(\partial\Delta^n)
			\]
			are effective epimorphisms, equivalently pointwise surjective on
			$\pi_0$.
			\item If every $Y_n$ is a $\kappa$-small coproduct of
			representables, then each
			\[
				(\Ex^{\infty}(Y))_{q+1}
				\times_{(\Ex^{\infty}(Y))_q}
				(\Ex^{\infty}(Y))_{q+1}
			\]
			is the countable colimit, over $r$, of $\kappa$-small coproducts of
			representables.
		\end{enumerate}
	\end{lemma}
	\begin{proof}
		The simplicial set $\sd^r\Delta^q$ is finite. Hence its simplex
		category is finite, and the formula defining $(\Ex^{r}(Y))_q$ expresses it
		as a finite limit of the objects $Y_n$.

		The remaining assertions are the usual internal properties of the
		$\Ex^{\infty}$ fibrant replacement in an $\infty$-topos.\footnote{See
		the nLab entry on Kan fibrant replacement,
		\url{https://ncatlab.org/nlab/show/Kan+fibrant+replacement}.}
		The degree-zero assertion follows directly from
		$\sd^r\Delta^0=\Delta^0$.

		For (4), write $Y^{(r)}=\Ex^{r}(Y)$. By (1) and
		\Cref{lem:Finite_Limits_Kappa_Coproducts_Representables}, every
		$Y^{(r)}_q$ is a $\kappa$-small coproduct of representables. Since
		filtered colimits commute with finite limits in presheaves of animae,
		we have
		\[
		\begin{aligned}
			&(\Ex^{\infty}(Y))_{q+1}
			\times_{(\Ex^{\infty}(Y))_q}
			(\Ex^{\infty}(Y))_{q+1} \\
			&\qquad\simeq
			\colim_{r\in\N}
			\left(
			Y^{(r)}_{q+1}\times_{Y^{(r)}_q}Y^{(r)}_{q+1}
			\right).
		\end{aligned}
		\]
		Each term in the final colimit is again a $\kappa$-small coproduct of
		representables by
		\Cref{lem:Finite_Limits_Kappa_Coproducts_Representables}.
	\end{proof}

	\begin{proposition}\label[proposition]{prop:Two_Sided_Bar_Relation_Object}
		Let $C$ be a small category, and let
		$Y_{\bullet}\in\Fun(\simp\catop,\PSh(C))$ be a simplicial presheaf.
		There is a simplicial object $B_{\bullet}(Y)$, natural in
		$Y_{\bullet}$, with a natural isomorphism
		\[
			\abs{B_{\bullet}(Y)}
			\simeq
			Y_0\times_{\abs{Y_{\bullet}}}Y_0.
		\]
		Moreover, if $Y^{\fib}_{\bullet}:=\Ex^{\infty}(Y)$, then
		\[
			B_q(Y)\simeq
			Y^{\fib}_{q+1}\times_{Y^{\fib}_q}Y^{\fib}_{q+1},
		\]
		where the two maps $Y^{\fib}_{q+1}\to Y^{\fib}_q$ are the last and
		first face maps.
	\end{proposition}
	\begin{proof}
		Set $Y^{\fib}_{\bullet}:=\Ex^{\infty}(Y)$. By
		\Cref{lem:Ex_Infinity_Countable_Control}, the map
		$Y_{\bullet}\to Y^{\fib}_{\bullet}$ induces an isomorphism on
		geometric realizations and identifies degree zero. It is therefore
		enough to compute
		\[
			Y^{\fib}_0\times_{\abs{Y^{\fib}_{\bullet}}}Y^{\fib}_0.
		\]

		Let $\mathsf{Dec}_{\bot}(Y^{\fib})$ and
		$\mathsf{Dec}_{\top}(Y^{\fib})$ be the two decalages. Thus both have
		$q$-simplices $Y^{\fib}_{q+1}$: the first forgets the last face and
		degeneracy, and the second forgets the first face and degeneracy. They
		carry natural maps to $Y^{\fib}_{\bullet}$, given respectively by the
		last and first face maps. The augmented simplicial animae
		\[
			\mathsf{Dec}_{\bot}(Y^{\fib})\to Y^{\fib}_0,
			\qquad
			\mathsf{Dec}_{\top}(Y^{\fib})\to Y^{\fib}_0
		\]
		have extra degeneracies, hence their realizations are both equivalent
		to $Y^{\fib}_0$.

		Since $Y^{\fib}_{\bullet}$ is Kan in $\PSh(C)$, the maps
		$\mathsf{Dec}_{\bot}(Y^{\fib})\to Y^{\fib}_{\bullet}$ and
		$\mathsf{Dec}_{\top}(Y^{\fib})\to Y^{\fib}_{\bullet}$ are Kan
		fibrations in $\PSh(C)$.\footnote{For the decalage path fibration, see the
		nLab entry on decalage, \url{https://ncatlab.org/nlab/show/decalage};
		for preservation of homotopy pullbacks along such realization
		fibrations, compare Rezk, \emph{When are homotopy colimits compatible
		with homotopy base change?},
		\url{https://ncatlab.org/nlab/files/RezkHomotopyColimitsBaseChange.pdf}.}
		Thus geometric realization preserves the pullback defining
		\[
			B_{\bullet}(Y):=
			\mathsf{Dec}_{\bot}(Y^{\fib})
			\times_{Y^{\fib}_{\bullet}}
			\mathsf{Dec}_{\top}(Y^{\fib}).
		\]
		We obtain
		\[
			\abs{B_{\bullet}(Y)}
			\simeq
			Y^{\fib}_0\times_{\abs{Y^{\fib}_{\bullet}}}Y^{\fib}_0
			\simeq
			Y_0\times_{\abs{Y_{\bullet}}}Y_0.
		\]
		The formula for $B_q(Y)$ is immediate from the definition of the two
		decalages.
	\end{proof}

	\begin{corollary}\label[corollary]{cor:Relation_Object_Kappa_Small_Colimit}
		Let $\kappa$ be an uncountable regular cardinal, let $C$ be a small
		category with finite limits, and let $Y_{\bullet}$ be a
		simplicial object of $\PSh(C)$ such that every $Y_n$ is a
		$\kappa$-small coproduct of representables. Then
		\[
			Y_0\times_{\abs{Y_{\bullet}}}Y_0
		\]
		is a $\kappa$-small colimit of representables.
		More precisely, as an object over $Y_0\times Y_0$, it is the
		geometric realization of a simplicial object whose $q$-simplices are
		countable colimits of $\kappa$-small coproducts of representables.
	\end{corollary}
	\begin{proof}
		This follows from \Cref{prop:Two_Sided_Bar_Relation_Object} and
		part (4) of \Cref{lem:Ex_Infinity_Countable_Control}. Since $\kappa$ is
		uncountable regular, the countable colimits in the simplicial degrees
		and the countable geometric realization together form a
		$\kappa$-small colimit of representables.
	\end{proof}

	\begin{proposition}\label[proposition]{prop:Kappa_Compact_Presheaves_Finite_Limits}
		Let $\kappa$ be an uncountable regular cardinal, and let $C$ be a small
		category with finite limits. Then the $\kappa$-compact objects
		of $\PSh(C)$ are closed under finite limits.
	\end{proposition}
	\begin{proof}
		The category $\PSh(C)$ is $\kappa$-compactly generated by
		representables, and its $\kappa$-compact objects are $\kappa$-small
		colimits of representables. Indeed, the only point is to absorb
		retracts: if $X$ is a retract of $Y$ in a category with countable
		colimits, represented by an idempotent $e\colon Y\to Y$, then
		\[
			X\simeq\colim\left(Y\xrightarrow{e}Y\xrightarrow{e}Y
			\xrightarrow{e}\cdots\right).
		\]

		The terminal object of $\PSh(C)$ is representable, hence
		$\kappa$-compact. We are thus reduced to pullbacks. Let
		\[
			X \to Y \leftarrow Z
		\]
		be a diagram in which all three objects are $\kappa$-small colimits of
		representables. We first reduce to the case where $X$ and $Z$ are
		representable. Write $X \simeq \colim_i X_i$ and
		$Z \simeq \colim_j Z_j$ as $\kappa$-small colimits of representables,
		with the induced maps to $Y$. Since colimits in a presheaf topos are
		universal, pullback in the slice over $Y$ preserves colimits. Hence
		\[
			X \times_Y Z
			\simeq
			\colim_{i,j}(X_i \times_Y Z_j).
		\]
		The indexing category is still $\kappa$-small, so it remains to handle
		the case where $X$ and $Z$ are representable.

		Write $Y$ as a $\kappa$-small colimit of representables. The simplicial
		replacement expresses this colimit as the geometric realization of a
		simplicial object $Y_{\bullet}$ of $\PSh(C)$ such that each $Y_n$ is a
		$\kappa$-small coproduct of representables. The map
		$Y_0\to\abs{Y_{\bullet}}$ is an effective epimorphism. Since effective
		epimorphisms in $\An$ are the $\pi_0$-surjections, maps from
		representables lift in the slice along effective epimorphisms after
		evaluation. Thus the two maps $X,Z\to\abs{Y_{\bullet}}$ admit lifts to
		$Y_0$.

		Set
		\[
			R := Y_0 \times_{\abs{Y_{\bullet}}} Y_0.
		\]
		Then
		\[
			X \times_{\abs{Y_{\bullet}}} Z
			\simeq
			(X \times Z) \times_{Y_0 \times Y_0} R,
		\]
		where the map $X\times Z\to Y_0\times Y_0$ is induced by the chosen
		lifts. By \Cref{cor:Relation_Object_Kappa_Small_Colimit}, the map
		$R\to Y_0\times Y_0$ is the geometric realization, in the slice over
		$Y_0\times Y_0$, of a simplicial object $R_{\bullet}$ such that every
		$R_q$ is a countable colimit $\colim_s R_{q,s}$ of $\kappa$-small
		coproducts of representables, compatibly over $Y_0\times Y_0$. Since
		colimits in a presheaf topos are universal,
		\[
			(X \times Z) \times_{Y_0 \times Y_0} R
			\simeq
			\abs{(X \times Z) \times_{Y_0 \times Y_0} R_{\bullet}}.
		\]
		The $q$-simplices on the right are
		\[
			\colim_s
			\left((X \times Z)\times_{Y_0\times Y_0}R_{q,s}\right).
		\]
		Here $X\times Z$ is representable, and $Y_0\times Y_0$ is a
		$\kappa$-small coproduct of representables by
		\Cref{lem:Finite_Limits_Kappa_Coproducts_Representables}. Thus each
		term $(X \times Z)\times_{Y_0\times Y_0}R_{q,s}$ is a finite limit of
		$\kappa$-small coproducts of representables, hence is itself a
		$\kappa$-small coproduct of representables. Since $\simp\catop$ and
		$\N$ are countable and $\kappa$ is uncountable regular, the total
		realization is a $\kappa$-small colimit of representables. Therefore
		$X\times_Y Z$ is a $\kappa$-small colimit of representables.
	\end{proof}

	\begin{lemma}\label[lemma]{lem:Compact_Finite_Limits_Localization}
		Let $\kappa$ be an uncountable regular cardinal, and let
		\[
			L\colon C \rightleftarrows D : \iota
		\]
		be an accessible localization of presentable categories. Assume that
		$C$ is $\kappa$-compactly generated, that $C^{\kappa}$ is closed under
		finite limits, that $L$ preserves finite limits, and that $\iota$
		preserves $\kappa$-filtered colimits. Then $D$ is
		$\kappa$-compactly generated, and $D^{\kappa}$ is closed under finite
		limits.
	\end{lemma}
	\begin{proof}
		Since $\iota$ preserves $\kappa$-filtered colimits, the left adjoint
		$L$ preserves $\kappa$-compact objects. Thus $LK$ is
		$\kappa$-compact for every $K \in C^{\kappa}$.

		We next show that these objects generate $D$ under $\kappa$-filtered
		colimits. In fact, we will exhibit every object $d \in D$ as an
		explicit $\kappa$-filtered colimit of objects of this form. Observe
		that the category
		\[
			(C^{\kappa})_{/\iota d}
		\]
		is $\kappa$-filtered: a $\kappa$-small diagram of maps
		$K_{\alpha}\to\iota d$ is dominated by the map
		$\colim_{\alpha}K_{\alpha}\to\iota d$, and $C^{\kappa}$ is closed
		under $\kappa$-small colimits. Since $C$ is $\kappa$-compactly
		generated, the canonical map
		\[
			\colim_{(K \to \iota d)} K \to \iota d
		\]
		is an isomorphism. Applying $L$ gives
		\[
			\colim_{(K \to \iota d)} LK \simeq d.
		\]
		Thus $D$ is $\kappa$-compactly generated by the objects $LK$. Moreover,
		if $d$ is $\kappa$-compact, then the identity of $d$ factors through
		one of the structure maps $LK\to d$, so $d$ is a retract of some $LK$.

		Now let $F\colon J\to D^{\kappa}$ be a finite diagram. We apply the
		previous paragraph not to $D$ itself, but to the induced localization
		\[
			\Fun(J,C) \rightleftarrows \Fun(J,D).
		\]
		Since $J$ is finite, \Cref{lem:Functor_Category_Presentable} shows that
		$\Fun(J,C)$ is again $\kappa$-compactly generated, with
		$\kappa$-compact generators given by diagrams $J\to C^{\kappa}$;
		similarly, $F$ is $\kappa$-compact as an object of $\Fun(J,D)$. The
		right adjoint preserves $\kappa$-filtered colimits pointwise. Hence
		$F$ is a retract, as a diagram, of $LK_{\bullet}$ for some diagram
		$K_{\bullet}\colon J\to C^{\kappa}$.

		Taking limits, $\lim_J F$ is a retract of
		\[
			\lim_J LK_{\bullet} \simeq L\left(\lim_J K_{\bullet}\right),
		\]
		where we used that $L$ preserves finite limits. Since $C^{\kappa}$ is
		closed under finite limits, $\lim_J K_{\bullet}$ lies in $C^{\kappa}$.
		Therefore $L(\lim_J K_{\bullet})$ is $\kappa$-compact in $D$, and so
		its retract $\lim_J F$ is $\kappa$-compact.
	\end{proof}

	\begin{lemma}\label[lemma]{lem:Kappa_Compact_Covering_Sieves_Filtered_Colimits}
		Let $\kappa$ be an uncountable regular cardinal, let $C$ be a small
		category with finite limits, and let $\tau$ be a Grothendieck topology
		on $C$. Suppose that the $\tau$-sheaf condition can be tested on a
		collection of covering sieves
		\[
			R\subseteq y(c)
		\]
		which are $\kappa$-compact as objects of $\PSh(C)$. Then
		$\tau$-sheaves are closed under $\kappa$-filtered colimits in
		$\PSh(C)$.

		In particular, this applies if $\tau$ is generated by a pretopology
		whose covering families have fewer than $\kappa$ morphisms. More
		generally, it applies whenever there is such a testing collection in
		which every sieve is generated by fewer than $\kappa$ morphisms into
		the target.
	\end{lemma}
	\begin{proof}
		Let $F\simeq \colim_i F_i$ be a $\kappa$-filtered colimit in
		$\PSh(C)$, where every $F_i$ is a $\tau$-sheaf. Let
		$R\subseteq y(c)$ be one of the covering sieves in the testing
		collection. Since $y(c)$ and $R$ are $\kappa$-compact in $\PSh(C)$,
		we have
		\[
		\begin{aligned}
			F(c)
			&\simeq \Hom_{\PSh(C)}(y(c),F) \\
			&\simeq \colim_i \Hom_{\PSh(C)}(y(c),F_i) \\
			&\simeq \colim_i \Hom_{\PSh(C)}(R,F_i) \\
			&\simeq \Hom_{\PSh(C)}(R,F).
		\end{aligned}
		\]
		The middle isomorphism uses that each $F_i$ satisfies descent for the
		covering sieve $R$. Hence $F$ satisfies descent for every sieve in the
		testing collection, and is therefore a $\tau$-sheaf.

		It remains to justify the final assertion. Let $R\subseteq y(c)$ be
		the sieve generated by a family
		\[
			\{c_a\to c\}_{a\in A}
		\]
		with $A$ $\kappa$-small. Set $U:=\coprod_{a\in A}y(c_a)$. The map
		$U\to y(c)$ factors through an effective epimorphism $U\to R$, and
		$R$ is the geometric realization of its \v{C}ech nerve. Since
		$R\to y(c)$ is a monomorphism, the terms of this \v{C}ech nerve are the
		corresponding fiber products over $y(c)$:
		\[
			U_n
			\simeq
			\coprod_{(a_0,\ldots,a_n)\in A^{n+1}}
			y(c_{a_0}\times_c\cdots\times_c c_{a_n}).
		\]
		Here the indicated finite limits exist in $C$. Each $U_n$ is a
		$\kappa$-small coproduct of representables, and the simplicial
		indexing category is countable. Since $\kappa$ is uncountable regular,
		$R\simeq \abs{U_{\bullet}}$ is a $\kappa$-small colimit of
		representables. Such an object is $\kappa$-compact: mapping out of it
		gives a $\kappa$-small limit of evaluations, and $\kappa$-filtered
		colimits of animae commute with $\kappa$-small limits.
	\end{proof}

	\begin{corollary}\label[corollary]{cor:Kappa_Compact_Sheaves_Finite_Limits}
		Let $\kappa$ be an uncountable regular cardinal, let $C$ be a small
		category with finite limits, and let $\tau$ be a Grothendieck topology
		on $C$. Suppose that $\tau$-sheaves are closed under
		$\kappa$-filtered colimits in $\PSh(C)$ (e.g.\ in the situation of \Cref{lem:Kappa_Compact_Covering_Sieves_Filtered_Colimits}). Then
		$\Shv_{\tau}(C)$ is $\kappa$-compactly generated, and the
		$\kappa$-compact objects of $\Shv_{\tau}(C)$ are closed under finite
		limits.

		Moreover, sheafification sends $\kappa$-compact presheaves to
		$\kappa$-compact sheaves, and every $\kappa$-compact sheaf is a
		retract of a $\kappa$-small colimit of sheafified representables.
	\end{corollary}
	\begin{proof}
		Let
		\[
			L\colon \PSh(C)\rightleftarrows \Shv_{\tau}(C) : \iota
		\]
		be sheafification and the inclusion. This is an accessible left exact
		localization, so $L$ preserves finite limits. By assumption, $\iota$
		preserves $\kappa$-filtered colimits. The category $\PSh(C)$ is
		$\kappa$-compactly generated by representables, and
		\Cref{prop:Kappa_Compact_Presheaves_Finite_Limits} shows that its
		$\kappa$-compact objects are closed under finite limits. Applying
		\Cref{lem:Compact_Finite_Limits_Localization} gives the first claim.

		The same adjunction shows directly that $L$ sends $\kappa$-compact
		presheaves to $\kappa$-compact sheaves. The proof of
		\Cref{lem:Compact_Finite_Limits_Localization} also shows that every
		$\kappa$-compact sheaf is a retract of $LK$ for some
		$\kappa$-compact presheaf $K$. Since $\kappa$-compact presheaves are
		retracts of $\kappa$-small colimits of representables, and since such
		retracts may be absorbed by the usual countable telescope, every
		$\kappa$-compact sheaf is a retract of a $\kappa$-small colimit of the
		objects $Ly(c)$.
	\end{proof}

	\chapter{Six-functor formalisms}
	\label[appendix]{chap:Six_Functor_Formalisms}

	\emph{Some of the contents of this appendix were written by Ou Liu in preparation for Talk \ref{chap:Transmutations}.}

	In this appendix, we record the relevant definitions in the formalism of \emph{six-functor formalisms} that have appered throughout these notes. Following Mann \citep{mann2022sixFunctor} and Scholze \citep{Scholze_Six_Functors}, the six operations and all of their coherences are encoded by a single lax symmetric monoidal functor out of a category of spans. We then recall, following Mann \citep[Proposition~A.5.10]{mann2022sixFunctor}, a criterion that produces such a formalism from the more accessible datum of a functor $C\catop\to\CAlg(\Cat)$ equipped with adjoints satisfying base change and the projection formula; we also record the 2-categorical enhancement from \citet{CLL_Span2}. Finally, we recall the definitions of cohomologically smooth/suave/prim/étale/proper morphisms with respect to a six-functor formalism.

	\section{The span category}
	\label[section]{sec:Span_Category}

	We begin by recalling the construction of the category of spans. Throughout, $C$ denotes a category.

	\begin{definition}
		\label[definition]{def:Closed_Under_Base_Change}
		Let $E\subseteq C$ be a wide subcategory, i.e.\ a subcategory containing every object of $C$. We say that $E$ is \emph{closed under base change} if for every morphism $g\colon X\to Y$ in $E$ and every morphism $f\colon Y'\to Y$ in $C$, the pullback $Y'\times_Y X$ exists in $C$ and the projection $Y'\times_Y X\to Y'$ again lies in $E$.
	\end{definition}

	\begin{construction}[The span category]
		\label[construction]{con:Span_Category}
		Let $E\subseteq C$ be a wide subcategory closed under base change. The \emph{span category} (or \emph{category of correspondences}) $\Span(C,E)$ is the category with the same objects as $C$, whose morphisms from $X$ to $Y$ are \emph{spans}
		\[
		\begin{tikzcd}[column sep=small, row sep=small]
			& U \dlar[swap]{f} \drar{g} & \\
			X & & Y
		\end{tikzcd}
		\]
		in which the right leg $g$ lies in $E$ and the left leg $f$ is arbitrary. The identity of $X$ is the span $X\xleftarrow{=}X\xrightarrow{=}X$, and the composite of $X\xleftarrow{f}U\xrightarrow{g}Y$ with $Y\xleftarrow{f'}V\xrightarrow{g'}Z$ is the outer span of
		\[
		\begin{tikzcd}[column sep=small, row sep=small]
			& & U\times_Y V \dlar \drar & & \\
			& U \dlar[swap]{f}\drar{g} & & V \dlar[swap]{f'}\drar{g'} & \\
			X & & Y & & Z\rlap,
		\end{tikzcd}
		\]
		formed by the pullback $U\times_Y V$. This pullback exists, and its right leg $U\times_Y V\to Z$ lies in $E$ since $E$ is closed under base change and composition.
	\end{construction}

	\begin{remark}
		\label[remark]{rmk:Span_Coherence}
		Making \Cref{con:Span_Category} into a genuine $\infty$-category requires more care: one realizes $\Span(C,E)$ as a complete Segal anima whose anima of $n$-simplices is the anima of functors $\Tw[n]\to C$ out of the twisted arrow category sending forward maps into $E$ and forming the relevant pullbacks. We refer to \citet{Barwick2017SpectralMackey} and \citet{HHLN2022TwoVariable} for the details.
	\end{remark}

	\begin{construction}[Monoidal structure on spans]
		\label[construction]{con:Monoidal_Span}
		Suppose now that $C$ has finite products. Then $\Span(C,E)$ inherits a symmetric monoidal structure whose unit is the terminal object of $C$ and whose underlying tensor product sends a pair $(X,Y)$ to the product $X\times Y$.

		Indeed, for $g\colon X\to Y$ and $g'\colon X'\to Y'$ in $E$ the product $g\times g'\colon X\times X'\to Y\times Y'$ again lies in $E$, since it factors as
		\[
			X\times X'\xrightarrow{\;g\times\id\;}Y\times X'\xrightarrow{\;\id\times g'\;}Y\times Y'
		\]
		and each factor is a base change of $g$ resp.\ $g'$ along a projection. Hence the cartesian product equips the triple $(C,C,E)$ with the structure of a commutative monoid among adequate triples. As the functor $\Span(-)$ preserves finite products \citep{HHLN2022TwoVariable}, this results in a symmetric monoidal structure on $\Span(C,E)$.
	\end{construction}

	\section{Six-functor formalisms}
	\label[section]{sec:Six_Functor_Formalisms}

	Let $C$ be a category with finite products and $E\subseteq C$ a wide subcategory closed under base change. We endow $\Span(C,E)$ with the symmetric monoidal structure of \Cref{con:Monoidal_Span} and $\Cat$ with the cartesian one.

	\begin{definition}[Three-functor formalism]
		\label[definition]{def:Three_Functor_Formalism}
		A \emph{three-functor formalism} on $(C,E)$ is a lax symmetric monoidal functor
		\[
			\D\colon \Span(C,E)\to\Cat.
		\]
	\end{definition}

	\begin{discussion}[The six operations]
		\label[discussion]{disc:Six_Operations}
		Let us unwind the data of a three-functor formalism $\D$.
		\begin{enumerate}
			\item Each object $X$ of $C$ is canonically a commutative algebra in $(\Span(C,E),\times)$, with multiplication and unit given by the diagonal $X\to X\times X$ and the projection $X\to\ast$ read as backward maps. Since lax symmetric monoidal functors preserve commutative algebras, every $\D(X)$ is itself a symmetric monoidal category, with tensor product $\otimes$ and unit $\unit_X$.
			\item Every morphism $f\colon X\to Y$ in $C$ defines a backward span $Y\xleftarrow{f}X\xrightarrow{=}X$, and hence a \emph{pullback} functor $f^*\coloneqq\D(Y\xleftarrow{f}X\xrightarrow{=}X)\colon\D(Y)\to\D(X)$, which is symmetric monoidal and depends functorially on $f$.
			\item Every morphism $g\colon X\to Y$ in $E$ defines a forward span $X\xrightarrow{=}X\xrightarrow{g}Y$, and hence an \emph{exceptional pushforward} $g_!\coloneqq\D(X\xrightarrow{=}X\xrightarrow{g}Y)\colon\D(X)\to\D(Y)$, depending functorially on $g$. The lax structure makes $g_!$ into a $\D(Y)$-linear functor, which is the content of the \emph{projection formula} $g_!(g^*(A) \otimes B) \cong A \otimes g_!(B)$.
			\item The two factorizations of a span $X\xleftarrow{f}U\xrightarrow{g}Y$ as a backward span followed by a forward span yield the \emph{base change} isomorphism $g_!f^*\simeq f'^* g'_!$ associated to a pullback square, i.e.\ the assignment $(X\xleftarrow{f}U\xrightarrow{g}Y)\mapsto g_!f^*$ is functorial.
		\end{enumerate}
		These are the three functors $(\otimes,\ f^*,\ f_!)$. When the relevant adjoints exist, one obtains the remaining three operations as their right adjoints: the \emph{pushforward} $f_*$ right adjoint to $f^*$, the \emph{exceptional pullback} $g^!$ right adjoint to $g_!$, and the \emph{internal hom} $\iHom(M,-)$ right adjoint to $M\otimes-$. This accounts for the name: the six operations are $(f^*,\ f_*,\ f_!,\ f^!,\ \otimes,\ \iHom)$.
	\end{discussion}

	\begin{definition}[Six-functor formalism]
		\label[definition]{def:Six_Functor_Formalism}
		A three-functor formalism $\D\colon\Span(C,E)\to\Cat$ is a \emph{six-functor formalism} if the right adjoints of \Cref{disc:Six_Operations} all exist; that is, if each $\D(X)$ is closed symmetric monoidal, each pullback $f^*$ admits a right adjoint $f_*$, and each exceptional pushforward $g_!$ admits a right adjoint $g^!$.
	\end{definition}

	\begin{remark}
		\label[remark]{rmk:Presentable_Targets}
		The right adjoints in \Cref{def:Six_Functor_Formalism} are a \emph{property} of the three-functor formalism, not extra structure. In practice one frequently considers formalisms valued in presentable categories, e.g.\ a lax symmetric monoidal functor $\Span(C,E)\to\PrL$ as in the constructions $\Spec(A)\mapsto n\Pr_A$ considered earlier in these notes; there the existence of these adjoints is automatic by the adjoint functor theorem, so that every such three-functor formalism is already a six-functor formalism.
	\end{remark}

	\section{The construction of six-functor formalisms}
	\label[section]{sec:Construction_6FF}

	Three-functor formalisms encode their data in a highly coherent fashion, which makes them powerful but correspondingly difficult to construct directly. In practice one starts from the much more accessible datum of the pullbacks and tensor products alone, i.e.\ a functor
	\[
		\D_0\colon C\catop\to\CAlg(\Cat),\qquad X\mapsto(\D_0(X),\otimes),\quad f\mapsto f^*,
	\]
	and asks for conditions guaranteeing that it extends to a three-functor formalism. The fundamental result in this direction is due to Liu--Zheng and Mann \citep[Proposition~A.5.10]{mann2022sixFunctor}. The hypotheses concern a factorization of the maps in $E$ together with the interaction of the resulting adjoints. We first recall the Beck--Chevalley conditions and the projection formula.

	\begin{definition}[Beck--Chevalley condition]
		\label[definition]{def:Beck_Chevalley}
		Let $\D_0\colon C\catop\to\Cat$ and consider a pullback square
		\[
		\begin{tikzcd}
			X' \rar{g'} \dar[swap]{f'} \drar[pullback] & X \dar{f} \\
			Y' \rar[swap]{g} & Y
		\end{tikzcd}
		\]
		in $C$, giving a commutative square $f'^*g^*\simeq g'^*f^*$ of functors $\D_0(Y)\to\D_0(X')$.
		\begin{enumerate}
			\item Suppose $f^*$ and $f'^*$ admit left adjoints $f_\sharp$ and $f'_\sharp$. The square is \emph{left adjointable}, or satisfies the \emph{Beck--Chevalley condition}, if the mate
			\[
				\BC_\sharp\colon f'_\sharp\, g'^*\to g^*\, f_\sharp
			\]
			is an isomorphism.
			\item Dually, suppose $f^*$ and $f'^*$ admit right adjoints $f_*$ and $f'_*$. The square is \emph{right adjointable}, or satisfies the \emph{dual Beck--Chevalley condition}, if the mate
			\[
				\BC_*\colon g^*\, f_*\to f'_*\, g'^*
			\]
			is an isomorphism.
		\end{enumerate}
	\end{definition}

	\begin{definition}[Projection formula]
		\label[definition]{def:Projection_Formula}
		Let $\D_0\colon C\catop\to\CAlg(\Cat)$, so that each $\D_0(X)$ is symmetric monoidal and each $f^*$ is symmetric monoidal.
		\begin{enumerate}
			\item If $f\colon X\to Y$ is such that $f^*$ admits a left adjoint $f_\sharp$, then for $M\in\D_0(X)$ and $N\in\D_0(Y)$ there is a natural \emph{projection map}
			\[
				f_\sharp\big(M\otimes f^*N\big)\to f_\sharp(M)\otimes N,
			\]
			adjoint to $M\otimes f^*N\to f^*f_\sharp(M)\otimes f^*N\simeq f^*\big(f_\sharp(M)\otimes N\big)$. We say $f_\sharp$ satisfies the \emph{projection formula} if this map is an isomorphism for all $M$ and $N$.
			\item Dually, if $p^*$ admits a right adjoint $p_*$, then the \emph{projection map}
			\[
				N\otimes p_*(M)\to p_*\big(p^*N\otimes M\big)
			\]
			is required to be an isomorphism; we then say $p_*$ satisfies the \emph{projection formula}.
		\end{enumerate}
	\end{definition}

	With these notions in place, we can state the construction theorem.

	\begin{theorem}[{Mann \citep[Proposition~A.5.10]{mann2022sixFunctor}}]
		\label[theorem]{thm:Construction_6FF}
		Let $C$ be a category with finite products, and let $I,P\subseteq E\subseteq C$ be wide subcategories closed under base change satisfying:
		\begin{enumerate}
			\item $I$ and $P$ are left cancellable (if $gf,g\in I$ then $f\in I$, and likewise for $P$);
			\item every map in $E$ factors as a map in $I$ followed by a map in $P$;
			\item every map in $I\cap P$ is truncated.
		\end{enumerate}
		Let $\D_0\colon C\catop\to\CAlg(\Cat)$ be a functor satisfying:
		\begin{enumerate}
			\item for every $i\colon X\to Y$ in $I$, the pullback $i^*$ admits a left adjoint $i_\sharp$ satisfying the Beck--Chevalley condition (\Cref{def:Beck_Chevalley}) and the projection formula (\Cref{def:Projection_Formula});
			\item for every $p\colon X\to Y$ in $P$, the pullback $p^*$ admits a right adjoint $p_*$ satisfying the dual Beck--Chevalley condition and the projection formula;
			\item for every pullback square
			\[
			\begin{tikzcd}
				X'\rar{i}\dar[swap]{p}\drar[pullback] & X\dar{q}\\
				Y'\rar[swap]{j} & Y
			\end{tikzcd}
			\]
			in $C$ with horizontal maps in $I$ and vertical maps in $P$, the \emph{double Beck--Chevalley map} $j_\sharp\, p_*\to q_*\, i_\sharp$ is an isomorphism.
		\end{enumerate}
		Then $\D_0$ extends to a three-functor formalism
		\[
			\D\colon\Span(C,E)\to\Cat
		\]
		with the property that for every $e\in E$ and every factorization $e=p\circ i$ with $i\in I$ and $p\in P$, the exceptional pushforward is given by $e_!\simeq p_*\,i_\sharp$.
	\end{theorem}

	\section{A 2-categorical enhancement}
	\label[section]{sec:2Cat_Enhancement}

	The construction of \Cref{thm:Construction_6FF} is the shadow of a finer statement, in which $\Span(C,E)$ is promoted to a $2$-category and the extension is pinned down by a universal property. This refinement is due to Gaitsgory--Rozenblyum \citep{GaitsgoryRozenbluym2017StudyDAG} in special cases and due to Cnossen--Lenz--Linskens \citep{CLL_Span2} in the present generality. It has two useful features that we will make use of: it characterizes the extension uniquely, and it works with an arbitrary symmetric monoidal $2$-category as target. We give an informal account.

	\begin{construction}[The span $2$-category]
		\label[construction]{con:Span_2_Category}
		Let $I,P\subseteq E\subseteq C$ be wide subcategories closed under base change. There is a $2$-category $\SpanTwoP{C}{E}{I}{P}$, a mild generalization of a construction of Haugseng, refining $\Span(C,E)$ as follows:
		\begin{itemize}
			\item Its objects are the objects of $C$;
			\item Its $1$-morphisms $X\to Y$ are the spans $X\xleftarrow{f}U\xrightarrow{e}Y$ with $e\in E$, exactly as in $\Span(C,E)$;
			\item A $2$-morphism from $X\xleftarrow{f}U\xrightarrow{e}Y$ to $X\xleftarrow{f'}V\xrightarrow{e'}Y$ is a \emph{span of spans}, i.e.\ an object $W$ together with maps $p\colon W\to U$ in $P$ and $i\colon W\to V$ in $I$ making
			\[
			\begin{tikzcd}[row sep=small, column sep=small]
				& U \dlar \drar{e} & \\
				X & W \uar{p} \dar{i} & Y \\
				& V \ular \urar[swap]{e'} &
			\end{tikzcd}
			\]
			commute.
		\end{itemize}
		Composition of $2$-cells is again by pullback, and the underlying $1$-category is $\Span(C,E)$. Making this precise (as a $2$-fold complete Segal anima) is the technical heart of \citet{CLL_Span2}, to which we refer for details.
	\end{construction}

	\begin{definition}[$(I,P)$-biadjointable lax symmetric monoidal functor]
		\label[definition]{def:Biadjointable_Functor}
		Let $\bbE$ be a symmetric monoidal $2$-category. A lax symmetric monoidal functor
		\[
			F\colon C\catop\to\bbE,\qquad f\longmapsto f^*\coloneqq F(f),
		\]
		is \emph{$(I,P)$-biadjointable} if it satisfies the following conditions, interpreted via adjunctions internal to $\bbE$:
		\begin{enumerate}
			\item For every $i\in I$, the morphism $i^*$ admits a left adjoint $i_\sharp$ satisfying the Beck--Chevalley condition and the projection formula.
			\item For every $p\in P$, the morphism $p^*$ admits a right adjoint $p_*$ satisfying the dual Beck--Chevalley condition and the projection formula.
			\item For every pullback square
			\[
			\begin{tikzcd}
				X'\rar{i}\dar[swap]{p}\drar[pullback] & X\dar{q}\\
				Y'\rar[swap]{j} & Y
			\end{tikzcd}
			\]
			with horizontal maps in $I$ and vertical maps in $P$, the double Beck--Chevalley map $j_\sharp\,p_*\to q_*\,i_\sharp$ is an isomorphism.
		\end{enumerate}
		For $\bbE=\Cat$, these are precisely the Beck--Chevalley and projection-formula conditions of \Cref{def:Beck_Chevalley,def:Projection_Formula}.
	\end{definition}

	\begin{theorem}[{\citet[Theorem~B]{CLL_Span2}}]
		\label[theorem]{thm:Span_2_Universal}
		Let $I,P\subseteq E\subseteq C$ satisfy the hypotheses of \Cref{thm:Construction_6FF}, and let $\bbE$ be a symmetric monoidal $2$-category. Restriction along the inclusion $h\colon C\catop\hookrightarrow\SpanTwoP{C}{E}{I}{P}$ induces an equivalence of $2$-categories
		\[
			h^*\colon\Fun^{\mathrm{lax}\otimes}\big(\SpanTwoP{C}{E}{I}{P},\,\bbE\big)
			\iso
			\Fun^{\mathrm{lax}\otimes}_{(I,P)\textup{-biadj}}\big(C\catop,\,\bbE\big),
		\]
		where the right-hand side is the locally full sub-$2$-category on the $(I,P)$-biadjointable lax symmetric monoidal functors and the lax monoidal natural transformations satisfying the evident Beck--Chevalley conditions. In particular, every such functor extends uniquely to a lax symmetric monoidal $2$-functor on $\SpanTwoP{C}{E}{I}{P}$. For every $e\in E$ and every factorization $e=p\circ i$ with $i\in I$ and $p\in P$, its exceptional pushforward is given by
		\[
			e_!\simeq p_*\,i_\sharp.
		\]
	\end{theorem}

	Specializing to $\bbE=\Cat$ recovers \Cref{thm:Construction_6FF}, now with the additional information that the extension is unique and characterized by this universal property.

	Finally, we record an observation that we will need later, and which follows directly by inspecting \Cref{con:Span_2_Category}.

	\begin{construction}[Endomorphisms of the unit]
		\label[construction]{con:Endo_Unit}
		Assume that the class $E$ contains all maps in $C$. We may then describe the endomorphism category of $\ast$ in $\SpanTwoP{C}{E}{I}{P}$ as follows: an object of $\Hom(\ast,\ast)$ is a span of the form $* \leftarrow X \rightarrow *$ corresponding to an object $X\in C$, while a morphism $X\to Y$ corersponds to a span $X\xleftarrow{p}W\xrightarrow{i}Y$ with $p\in P$ and $i\in I$. In other words, we have
		\[
			\Hom_{\SpanTwoP{C}{E}{I}{P}}(\ast,\ast)\;\simeq\; \Span(C,P,I),
		\]
		the span category of $C$ with backward maps in $P$ and forward maps in $I$.

		Now let $\bbE$ be symmetric monoidal and let $F\colon\SpanTwoP{C}{E}{I}{P}\to\bbE$ be a lax symmetric monoidal $2$-functor, e.g.\ one obtained from \Cref{thm:Span_2_Universal}. Then $F$ carries the unit to a commutative algebra $F(\ast)\in\CAlg(\bbE)$, and its action on the endomorphisms of the units yields a functor (in fact lax symmetric monoidal)
		\[
			\Span(C,P,I)\;\simeq\;\Hom(\ast,\ast)\;\xrightarrow{\ F\ }\;\Hom_{\bbE}\big(F(\ast),F(\ast)\big).
		\]
	\end{construction}

	\section{Geometric properties of morphisms}
	\label[section]{sec:Geometric_Properties}

	The six operations make it possible to express, in purely terms of the six-functor formalism, several geometric properties of a morphism: cohomological smoothness, suaveness, primness, \'etaleness, and properness. We explicitly record this vocabulary here, so that we can compare it to the analogous notions in the Gestalten formalism under the transmutation functor $[-]_{\D}$; see \Cref{sec:Translation}. Throughout, $\D$ denotes a six-functor formalism on $(C,E)$ in the sense of \Cref{def:Six_Functor_Formalism}, and morphisms are assumed to lie in $E$ unless stated otherwise.

	\begin{definition}[Cohomologically smooth morphism]
		\label[definition]{def:Cohomologically_Smooth}
		A morphism $f \colon X \to Y$ in $E$ is \emph{cohomologically smooth} (or \emph{$\D$-smooth}) if it satisfies:
		\begin{enumerate}
			\item\label{def:Cohom_Smooth_Proj} The natural transformation
			\[
				f^!(\unit_Y) \otimes f^*(-) \;\longrightarrow\; f^!(-)
			\]
			is an isomorphism; we write $\omega_f \coloneqq f^!(\unit_Y)$ and call it the \emph{dualizing complex} of $f$.
			\item\label{def:Cohom_Smooth_Inv} The dualizing complex $\omega_f$ is $\otimes$-invertible in $\D(X)$.
			\item\label{def:Cohom_Smooth_BC} For every pullback square as above with $g \in E$, the morphism $f'$ also satisfies \eqref{def:Cohom_Smooth_Proj}--\eqref{def:Cohom_Smooth_Inv}, and the canonical map $g'^* f^!(\unit_Y) \to f'^!(\unit_{Y'})$ is an isomorphism.
		\end{enumerate}
	\end{definition}

	Cohomological smoothness can be conveniently formulated in terms of the \emph{2-category of kernels} $K_{\D}$ associated to $\D$.

	\begin{definition}[$2$-category of kernels]
		\label[definition]{def:Kernel_2_Category}
		Let $Y \in C$. The \emph{$2$-category of kernels} $K_{\D, Y}$ is the $2$-category whose:
		\begin{itemize}
			\item \emph{Objects} are objects of $(C_E)_{/Y}$, i.e.\ morphisms $X \to Y$ in $E$.
			\item The \emph{hom category} from $X_1$ to $X_2$ is $\Hom_{K_{\D}}(X_1, X_2) \coloneqq \D(X_1 \times_Y X_2)$.
			\item The \emph{identity} on $X$ is $\Delta_{f!}(\unit_X) \in \D(X \times_Y X)$, where $\Delta_f \colon X \to X \times_Y X$ is the relative diagonal.
			\item \emph{Composition} of $M \in \Hom_{K_{\D}}(X_1, X_2)$ and $N \in \Hom_{K_{\D}}(X_2, X_3)$ is
			\[
				N \circ M \;\coloneqq\; (\pi_{13})_!\bigl(\pi_{12}^* M \otimes \pi_{23}^* N\bigr) \;\in\; \D(X_1 \times_Y X_3),
			\]
			where the $\pi_{ij}$ are the projections from $X_1 \times_Y X_2 \times_Y X_3$.
		\end{itemize}
		To define $K_{\D,Y}$, we consider the span category $\Span(E_{/Y})$ equipped with the closed symmetric monoidal structure coming from the cartersian monoidal structure on $E_{/Y}$, regard it as self-enriched via the internal homs (which are given by $[X,X'] = X \times_Y X'$ by self-duality of all objects in a span category) and perform transfer of enrichments along the lax symmetric monoidal functor $D_Y\colon \Span(E_{/Y}) \to \Cat$. 
		
		An alternative approach to the 2-category of kernels is presented in \Cref{def:Category_Of_Kernels}.
	\end{definition}

	The origin of the name \emph{$2$-category of kernels} is that the $1$-morphisms $X_1\to X_2$ of $K_{\D,Y}$ are exactly the \emph{kernels} $K\in\D(X_1\times_Y X_2)$ of Fourier--Mukai transforms
	\[
		F_K := (\pi_2)_!\bigl(\pi_1^*(-) \otimes K\bigr)\colon \D(X_1) \to \D(X_2),
	\]
	and the composition law of \Cref{def:Kernel_2_Category} corresponds to composing the associated transforms. In fact, there is a canonical $2$-functor
	\[
		\Psi_{\D, Y} \;=\; \Hom_{K_{\D}}(Y, -) \;\colon\; K_{\D, Y} \longrightarrow \Cat,
	\]
	sending an object $X$ to $\D(X)$ and a $1$-morphism $K \in \Hom_{K_{\D}}(X_1, X_2)$ to the Fourier--Mukai transform $F_K$.

	\begin{theorem}[Smoothness via kernels, {cf.\ \citet[\textsection 4.5]{Heyer_Mann}}]
		\label[theorem]{thm:Smoothness_Via_Kernels}
		A morphism $f \colon X \to Y$ in $E$ is $\D$-smooth if and only if
		\begin{enumerate}
			\item\label{thm:Smooth_Kernel_Adjoint} The $1$-morphism $\unit_X \in \Hom_{K_{\D}}(X, Y) = \D(X)$ is a left adjoint in $K_{\D, Y}$. Denote its right adjoint by $\omega_f \in \Hom_{K_{\D}}(Y, X) = \D(X)$.
			\item\label{thm:Smooth_Kernel_Inv} $\omega_f$ is $\otimes$-invertible in $\D(X)$.
		\end{enumerate}
		When these hold, $\omega_f \simeq f^!(\unit_Y)$.
	\end{theorem}

	In light of this result, it makes sense to study more general adjunctions in $K_{D, Y}$. This leads to the following definition.

	\begin{definition}[Suave and prim objects]
		\label[definition]{def:Suave_Prim_Objects}
		Let $f \colon X \to Y$ be in $E$ and $A \in \D(X)$.
		\begin{enumerate}
			\item $A$ is \emph{$f$-suave} if, viewed as a $1$-morphism $X \to Y$ in $K_{\D, Y}$, it is a left adjoint. Its right adjoint is the \emph{$f$-suave dual} $\SD_f(A) \in \D(X)$. Write $\Suave_f(X) \subseteq \D(X)$ for the full subcategory of $f$-suave objects.
			\item $A$ is \emph{$f$-prim} if, viewed as a $1$-morphism $Y \to X$ in $K_{\D, Y}$, it is a left adjoint. Its right adjoint is the \emph{$f$-prim dual} $\PD_f(A) \in \D(X)$. Write $\Prim_f(X) \subseteq \D(X)$ for the full subcategory of $f$-prim objects.
		\end{enumerate}
		(Equivalently, via the self-duality $K_{\D, Y} \simeq K_{\D, Y}\catop$, $A$ is $f$-prim iff it is a \emph{right} adjoint $1$-morphism $X \to Y$; this is the form used by Heyer--Mann \citep{Heyer_Mann}.)
	\end{definition}

	\begin{proposition}[{\citet[\textsection 4.4]{Heyer_Mann}}]
		\label[proposition]{prop:Reflexive_Duality}
		Let $\D$ be a six-functor formalism on $(C, E)$ and $f \colon X \to Y$ in $E$.
		\begin{enumerate}
			\item If $A \in \D(X)$ is $f$-suave, then $\SD_f(A)$ is also $f$-suave (with right adjoint $A$ in $K_{\D, Y}$), and there is a natural isomorphism
			\[
				\SD_f(A) \otimes f^*(-) \;\simeq\; \uHom\bigl(A, f^!(-)\bigr) \colon \D(Y) \to \D(X).
			\]
			Setting $- = \unit_Y$ yields $\SD_f(A) \simeq \uHom(A, f^!(\unit_Y))$, exhibiting $\SD_f(A)$ as a relative Verdier dual of $A$. The bidual identity $\SD_f(\SD_f(A)) \simeq A$ then holds, and $\SD_f$ defines a contravariant isomorphism
			\[
				\SD_f \colon \Suave_f(X)\catop \xrightarrow{\;\sim\;} \Suave_f(X).
			\]
			\item Dually, if $A \in \D(X)$ is $f$-prim, then $\PD_f(A)$ is also $f$-prim, with
			\[
				f_!\bigl(\PD_f(A) \otimes -\bigr) \;\simeq\; f_*\uHom(A, -) \colon \D(X) \to \D(Y),
			\]
			and the bidual $\PD_f(\PD_f(A)) \simeq A$ holds. Consequently $\PD_f$ defines a contravariant isomorphism
			\[
				\PD_f \colon \Prim_f(X)\catop \xrightarrow{\;\sim\;} \Prim_f(X).
			\]
		\end{enumerate}
	\end{proposition}

	We may specialize \Cref{def:Suave_Prim_Objects} to $A = \unit_X$ to obtain two classes of morphisms.

	\begin{definition}[Suave and prim morphisms]
		\label[definition]{def:Suave_Prim_Morphisms}
		Let $f \colon X \to Y$ be in $E$.
		\begin{enumerate}
			\item $f$ is \emph{$\D$-suave} if $\unit_X \in \Suave_f(X)$. As before we write $\omega_f \coloneqq \SD_f(\unit_X) \in \D(X)$ for its \emph{dualizing complex}.
			\item $f$ is \emph{$\D$-prim} if $\unit_X \in \Prim_f(X)$. We then write $\delta_f \coloneqq \PD_f(\unit_X) \in \D(X)$ for the \emph{co-dualizing complex} of $f$.
		\end{enumerate}
	\end{definition}

	\begin{remark}
		\label[remark]{rmk:Suave_Vs_Smooth}
		Combining \Cref{thm:Smoothness_Via_Kernels} with \Cref{def:Suave_Prim_Morphisms} gives
		\[
			f \text{ is $\D$-smooth} \;\Longleftrightarrow\; f \text{ is $\D$-suave and } \omega_f \text{ is invertible}.
		\]
		By \Cref{prop:Reflexive_Duality}, the dualizing complex is the relative Verdier dual $\omega_f \simeq f^!(\unit_Y)$. The co-dualizing complex may be computed by
		\[
			\delta_f \simeq \pi_{2*}\Delta_{f!} \unit_X,
		\]
		where $\pi_2 \colon X \times_Y X \to X$ is the projection onto the second factor.
	\end{remark}

	For $\D$-suave $f$, the formula $\omega_f \otimes f^*(-) \xrightarrow{\sim} f^!(-)$ is equivalent (by adjunction) to $f^*(-) \xrightarrow{\sim} \uHom(\omega_f, f^!(-))$. Plugging in $\unit_Y$ yields $\unit_X \xrightarrow{\sim} \uHom(\omega_f, \omega_f)$, and for dualizable $\omega_f$ this gives $\unit_X \simeq \omega_f^{\vee} \otimes \omega_f$. We record this as an observation.

	\begin{observation}
		\label[observation]{obs:Invertible_Vs_Dualizable}
		For a $\D$-suave morphism $f$, the dualizing complex $\omega_f$ is $\otimes$-invertible if and only if it is dualizable. Consequently the verification of $\D$-smoothness reduces to checking suaveness and dualizability of $\omega_f$.
	\end{observation}

	Suave and prim morphisms satisfy various base change properties.

	\begin{lemma}[Base change for suave and prim morphisms, {\citet[Lemma~4.5.13]{Heyer_Mann}}]
		\label[lemma]{lem:Suave_Prim_Base_Change}
		Let $\D$ be a six-functor formalism on $(C, E)$ and consider a pullback square
		\[\begin{tikzcd}
			X' \rar["f'"] \dar["g'"'] \drar[pullback] & Y' \dar["g"] \\
			X \rar["f"'] & Y
		\end{tikzcd}\]
		in $C_E$.
		\begin{enumerate}
			\item If $f$ is $\D$-suave, the natural transformations
			\[
				g'_* f'^* \simeq f^* g_*, \quad g'_! f'^! \simeq f^! g_!, \quad f'^* g^! \simeq g'^! f^*, \quad g'^* f^! \simeq f'^! g^*
			\]
			are all isomorphisms.
			\item If $f$ is $\D$-prim, the natural transformations
			\[
				g^* f_* \simeq f'_* g'^*, \quad g^! f_! \simeq f'_! g'^!, \quad f_! g'_* \simeq g_* f'_!, \quad f_* g'_! \simeq g_! f'_*
			\]
			are all isomorphisms.
		\end{enumerate}
	\end{lemma}

	Passing to iterated diagonals gives notions of étaleness and properness.

	\begin{definition}[Cohomologically \'etale and proper morphisms]
		\label[definition]{def:Cohomologically_Etale_Proper}
		Let $\D$ be a six-functor formalism on $(C, E)$, and let $f \colon X \to Y$ be a truncated morphism in $E$ (i.e.\ $n$-truncated for some $n \in \N$).
		\begin{itemize}
			\item $f$ is \emph{$\D$-\'etale} (or \emph{cohomologically \'etale}) if $f$ is $\D$-suave and $\Delta_f$ is either $\D$-\'etale or an isomorphism.
			\item $f$ is \emph{$\D$-proper} (or \emph{cohomologically proper}) if $f$ is $\D$-prim and $\Delta_f$ is either $\D$-proper or an isomorphism.
		\end{itemize}
		These definitions are well-founded because $n$-truncated morphisms have $(n - 1)$-truncated diagonals.
	\end{definition}

	For $\D$-étale maps, one may inductively construct isomorphisms $f^* \iso f^!$ as follows. If $f$ is an isomorphism, $f_!$ is inverse to $f^*$, and so $f^* \iso f^!$. Now assume that $\Delta_f$ is $\D$-étale, and consider the following diagram:
	\[\begin{tikzcd}
		X \ar[dr, "\Delta_f"] \ar[drr, bend left=20, "\id_X"] \ar[ddr, bend right=15, "\id_X"' near end] \\
		& X \times_Y X \rar["\pi_2"] \dar["\pi_1"'] \drar[pullback] & X \dar["f"] \\
		& X \rar["f"'] & Y.
	\end{tikzcd}\]
	We have $f \pi_1 = f \pi_2$ and $\id_X = \pi_2 \Delta_f$. Using the inductively constructed isomorphism $\Delta_f^* \iso \Delta_f^*$, we obtain a map
	\[
		f^* = \Delta_f^!\pi_1^!f^* \to \Delta_f^!\pi_1^*f^! \cong \Delta_f^*\pi_1^*f^! = f^!,
	\]
	where we inserted the double Beck-Chevalley map $\pi_1^!f^* \to \pi_1^*f^!$. The following proposition shows that $f$ is $D$-étale if and only if this map is an isomorphism. A dual inductive construction produces a map $f_! \to f_*$ whenever $\Delta_f$ is known to be $\D$-proper.

	\begin{proposition}[Characterization of \'etaleness and properness, {cf.\ \citet[\textsection 4.6]{Heyer_Mann}}]
		\label[proposition]{prop:Etale_Proper_Equivalence}
		Let $\D$ be a six-functor formalism on $(C, E)$ and $f \colon X \to Y$ in $E$.
		\begin{enumerate}
			\item Suppose $\Delta_f$ is $\D$-\'etale. Then the following are equivalent:
			\begin{enumerate}
				\item[(a)] $f$ is $\D$-\'etale;
				\item[(b)] $f^! \unit_Y \iso f^* \unit_Y$ in $\D(X)$;
				\item[(c)] $f^! \iso f^*$ as functors.
			\end{enumerate}
			\item Suppose $\Delta_f$ is $\D$-proper. Then the following are equivalent:
			\begin{enumerate}
				\item[(a$'$)] $f$ is $\D$-proper;
				\item[(b$'$)] $f_! \unit_X \simeq f_* \unit_X$ in $\D(Y)$;
				\item[(c$'$)] $f_! \xrightarrow{\;\sim\;} f_*$ as functors.
			\end{enumerate}
		\end{enumerate}
	\end{proposition}

	In \Cref{sec:Translation} we study how the notions of cohomological smoothness, suaveness, primness, étaleness and properness can be translated to the corresponding notions for Gestalten via the transmutation $[-]_{\D}\colon C \to \Gest$.

	\section{Recovering six-functor formalisms from slices}
	\label[section]{sec:Recovering_From_Slices}

	Given a six-functor formalism on $\Span(C,E)$, we may restrict it along the forgetful functor $E_{/X} \to C$ for every $X \in C$ to obtain a six-functor formalism on the slice $E_{/X}$. The point of this section is to ask about the converse: when can the original six-functor formalism be recovered from all of these restrictions?

	The reason it is useful to be able to do so is that in categorified situations it is sometimes easier to get a handle on the restricted formalisms, using the fact that $\Span(E_{/X})$ is the endomorphism category of the span 2-category $\SpanTwo(C,E)$ considered in \Cref{con:Span_2_Category}. The abstract reconstruction result below is used in \Cref{prop:6FF_Truncated_Gest} to assemble the slicewise constructions on Gestalten.

	Before giving a precise result, let us set the stage by making precise the restriction data and the compatibilities between them. Given $\D\colon \Span(C,E) \to \Cat$, we may consider for every morphism $W \to X$ in $C$ the functor
	\[
		\Span(E_{/X}) \to \Cat, \qquad (Z \to X) \mapsto \D(Z \times_X W).
	\]
	Collecting these functors for all $W \in C_{/X}$ results in a functor
	\[
		\D_{/X}\colon \Span(E_{/X}) \to \Fun(C_{/X}\catop,\Cat).
	\]
	Moreover, these functors are compatible with the two canonical functorialities in $X$ on both sides:
	\begin{itemize}
		\item On the one hand, both sides are contravariantly functorial in maps $f\colon Y \to X$ in $C$. For the left-hand side, we may consider the pullback functor $f^*\colon E_{/X} \to E_{/Y}$, which induces a functor on span categories $\Span(f^*)\colon \Span(E_{/X}) \to \Span(E_{/Y})$. For the right-hand side, we can restrict along the postcomposition functor $f_{\sharp}\colon C_{/Y} \to C_{/X}$ to obtain a functor $f^*\colon\Fun(C_{/X}\catop,\Cat)\to\Fun(C_{/Y}\catop,\Cat)$. We then have the relation
		\[
			f^* \circ \D_{/X}\simeq \D_{/Y}\circ\Span(f^*), \qquad (Z,W)\mapsto \D(Z \times_X W) \simeq \D((Z \times_X Y) \times_Y W)
		\]
		of functors $\Span(E_{/X})\to\Fun(C_{/Y}\catop,\Cat)$.
		\item On the other hand, both sides are covariantly functorial in maps $e\colon Y \to X$ in $E$. For the left-hand side, the postcomposition functor $e_{\sharp}\colon E_{/Y} \to E_{/X}$ preserves pullbacks and thus induces a functor on spans. For the right-hand side, restricting along the pullback functor $e^*\colon C_{/Y} \to C_{/X}$ gives a functor $e_*\colon \Fun(C_{/Y}\catop,\Cat)\to\Fun(C_{/X}\catop,\Cat)$. We then have the relation
		\[
			\D_{/X} \circ \Span(e_{\sharp}) \simeq e_*\circ \D_{/Y}, \qquad (Z,W) \mapsto \D(Z \times_X W) \simeq \D(Z \times_Y (Y \times_X W))
		\]
		of functors $\Span(E_{/Y}) \to \Fun(C_{/X}\catop,\Cat)$.
	\end{itemize}
	These two compatibilities are closely related: the functor $\Span(e_{\sharp})$ induced by the left adjoint $e_{\sharp}\colon E_{/Y} \to E_{/X}$ to $e^*$ is a \emph{right} adjoint to $\Span(e^*)$, and similarly the restriction functor $e_*\colon \Fun(C_{/Y}\catop,\Cat)\to\Fun(C_{/X}\catop,\Cat)$ is right adjoint to $e^*$. The contravariant compatibility is then the mate of the covariant compatibility. In particular, if we remember the covariant compatibility relations as coherent data, the contravariant compatibilities can be stated as properties.

	We can now state the main result.

	\begin{theorem}[Slicewise reconstruction, {\citet[\textsection 8]{CLL_Span}}]
		\label[theorem]{thm:Slicewise_Reconstruction}
		Let $C$ have finite limits and let $E\subseteq C$ be closed under base change and left cancellation. For any category $A$, there is an equivalence between the following two types of data:
		\begin{itemize}
			\item Functors $\D\colon\Span(C,E)\longrightarrow A$;
			\item Natural transformations
			\[
				\D_{/\bullet}\colon \Span(E_{/\bullet}) \to \Fun(C_{/\bullet}\catop,A)
			\]
			of functors $C\catop \to \Cat$ satisfying the property that for each morphism $e\colon Y \to X$ in $E$ the canonical map
			\[
				\D_{/X} \circ \Span(e_{\sharp}) \to e_*\circ \D_{/Y}
			\]
			is a natural equivalence of functors $\Span(E_{/Y}) \to \Fun(C_{/X}\catop,A)$.
		\end{itemize}
	\end{theorem}
	\begin{proof}
		The proof is given in \cite[\textsection 8]{CLL_Span}; there, the data in the second bullet point occurs under the name \emph{$E$-commutative monoid in the $C$-category $X \mapsto \Fun(C_{/X},A)$}.

		The rough idea of the proof is as follows. Natural transformations between functors $C\catop \to \widehat{\Cat}$ can be alternatively encoded as morphisms between cocartesian fibrations over $C\catop$. The cocartesian unstraightening of $X \mapsto \Span(E_{/X})$ is the span category of $\Ar_E(C)$, the full subcategory of the arrow category $\Ar(C)$ spanned by the morphisms $Z \to X$ in $E$. The span category $\Span(C,E)$ is a localization of this span category, and the class of maps at which one needs to localize is precisely the class of maps encoding the second compatibility relation.
	\end{proof}

	We have already seen how to produce the data of the functor $\D_{/X}$ from $\D$. In the other direction, we may recover $\D$ from $\D_{/\bullet}$ by setting
	\[
		\D(X) := \D_{/X}(\id_X)(\id_X).
	\]
	For an arbitrary map $f\colon U \to X$, the pullback functor $f^*\colon \D(X) \to \D(U)$ is obtained by applying the presheaf $\D_{/X}(\id_X)\colon C_{/X}\catop \to A$ to the morphism $f\colon U \to X$ in $C_{/X}$, using that $\D_{/X}(\id_X)(U) \simeq \D_{/U}(\id_U)(\id_U)$. For a map $e\colon U \to Y$ in $E$, the exceptional pushforward $e_!\colon \D(U) \to \D(Y)$ is obtained by applying $\D_{/Y}$ to the forward span
	\[
		e \xleftarrow{\id} e \xrightarrow{e} \id_Y
	\]
	in $E_{/Y}$ and then evaluating at $\id_Y$. A general span $X \xleftarrow{f} U \xrightarrow{e} Y$ is sent to the composite $e_!f^*$. The compatibility with the right adjoints $e_*$ is precisely what makes this construction functorial in composites on spans.

	\begin{remark}
		\label[remark]{rmk:Slicewise_Reconstruction_Monoidal}
		The theorem is stated only in the non-monoidal form, which is the form proved in \citet[\textsection 8]{CLL_Span}. It is highly expected that the same argument can be used to prove a symmetric monoidal refinement: if $C$ has finite products and $E$ is closed under products, then three-functor formalisms $\Span(C,E)\to \Cat$ should be equivalent to compatible families of slicewise three-functor formalisms
		\[
			\Span(E_{/X}) \to \Fun(C_{/X}\catop,\Cat).
		\]
		The same expectation applies with target any symmetric monoidal category, in particular $1\Pr$ or $\PrL$. We will not prove this refinement here, and will take it for granted in the applications below.
	\end{remark}

	\chapter{Higher Zariski geometry of Gestalten}
	\label[appendix]{chap:Zariski_Gestalten}

	\emph{This appendix is exploratory: its material was neither treated in Scholze's notes nor discussed in the seminar. It records what the higher Zariski geometry of \citet{ABCSS_Higher_Zariski} looks like when internalized into the Gestalten formalism.}

	\emph{This appendix was written entirely by Claude Fable 5, based on some exploratory prompts and the contents of \citep{ABCSS_Higher_Zariski} and these notes. Please treat it with care: Some details may be wrong.}

	The paper \citet{ABCSS_Higher_Zariski} develops a Zariski geometry for \emph{$2$-rings}, i.e.\ idempotent-complete stable symmetric monoidal categories. It equips the opposite of the category of compact $2$-rings with the structure of a geometry in the sense of \citet{Lurie_DAG_V}: the admissible morphisms are the principal Karoubi quotients $\Kk \to \Kk/\lra{a}$, and a finite family of these is a cover when every object killed by all of its members is tensor-nilpotent. From this data the machinery of \emph{loc.\ cit.}\ produces a spectrum, a category of locally $2$-ringed topoi, and a spectrum--global-sections adjunction; the underlying space of the spectrum recovers the Balmer spectrum \citep{Balmer_Spectrum}, via its description in terms of the poset of tensor-triangulated localizations due to \citet{Kock_Pitsch}, and the spectrum functor is fully faithful on rigid $2$-rings.

	The goal of this appendix is to describe what this geometry becomes over an arbitrary base Stefanich ring $A$ and at an arbitrary categorical height $n$. The translation rests on the classification results of \Cref{chap:Injections_Surjections_Gestalten}: idempotent algebras classify proper injections of Gestalten, and idempotent coalgebras classify étale injections. Once the covering condition is defined geometrically, via surjections of Gestalten, a perhaps surprising amount of the theory of \citet{ABCSS_Higher_Zariski} becomes formal:
	\begin{itemize}
		\item Zariski descent, both for the internal rings themselves and for their module categories, holds for geometric covers essentially by construction (\Cref{thm:Zariski_Covers_And_Descent}). In \emph{loc.\ cit.}\ the corresponding statements are the main descent theorems and require rigidity; internally to $A$, rigidity is instead the hypothesis under which the geometric covers admit a support-theoretic description (\Cref{thm:Rigid_Covers_Conservative}).
		\item The dichotomy between Karoubi quotients and smashing localizations organizes itself along the categorical height: general localizations at height $n$ become smashing, and in fact clopen, at height $n+1$ (\Cref{obs:Localization_Ladder}). The theorem of Miller that finite localizations of rigid $2$-rings are smashing acquires a two-line internal proof (\Cref{prop:Rigid_Localizations_Smashing}).
	\end{itemize}
	Statements in this appendix which rely on results from the main text are proved; statements extrapolating beyond them are explicitly labelled as expectations or questions.

	Throughout, we fix a Stefanich ring $A$ and write $X := \Gest(A)$.

	\section{Internal higher rings}
	\label[section]{sec:Internal_Higher_Rings}

	\begin{definition}\label[definition]{def:Internal_Higher_Rings}
		For $n \geq 0$, the category of \emph{internal $n$-rings} in $A$ is
		\[
			\Ring_n(A) := \CAlg(A_{n+1}),
		\]
		and its \emph{affine core} is the opposite of its subcategory of $\aleph_1$-compact objects,
		\[
			\Cc_n(A) := \bigl(\Ring_n(A)^{\aleph_1}\bigr)\catop.
		\]
	\end{definition}

	By the affine correspondence of \Cref{def:N_Affine}, in the notation of \Cref{chap:Injections_Surjections_Gestalten}, an internal $n$-ring $R \in \Ring_n(A)$ determines an $n$-affine Gestalt
	\[
		Y_R := \Gest_X(R) \longrightarrow X,
	\]
	and this assignment is a fully faithful embedding
	\[
		j_n\colon \Cc_n(A) \hookrightarrow \Gest_{/X}
	\]
	(defined on all of $\Ring_n(A)\catop$; the restriction to the affine core is what will carry the geometry). Since $(\Oo(Y_R)/A)_n = R$, \Cref{prop:Affineness_Properties}(2) provides an equivalence
	\[
		\Oo(Y_R)_{n+1} \simeq \Mod_{R}(A_{n+1}),
	\]
	and consequently
	\[
		\CAlg\bigl(\Oo(Y_R)_{n+1}\bigr) \simeq \Ring_n(A)_{R/}.
	\]

	\begin{remark}[The role of $\aleph_1$]\label[remark]{rem:Role_Of_Aleph1}
		The cutoff cardinal $\aleph_1$ plays the role that finite presentation plays in classical algebraic geometry and that $\omega$-compactness plays in \citet{ABCSS_Higher_Zariski}. The category $\Ring_n(A)$ is $\aleph_1$-presentable, and by \Cref{lem:Kappa_Compact_Algebra_Closure_Properties} its $\aleph_1$-compact objects form an essentially small, idempotent-complete subcategory closed under finite colimits; the underlying reason is that the $E_\infty$-operad is countable (\Cref{lem:Kappa_Compact_Algebra_Iff_Module}). Thus $\Cc_n(A)$ is an idempotent-complete small category with finite limits, as required for a geometry in the sense of Lurie. All constructions below make sense before imposing compactness; the affine core only enters when forming the site.
	\end{remark}

	\begin{discussion}[Calibration with the $2$-rings of ABCSS]\label[discussion]{disc:ABCSS_Calibration}
		Take $A = \Oo(\Spec(\S))$ and $n = 1$. Then $A_2 = 1\Pr_{\S}$, so $\Ring_1(A)$ is the category of stable presentably symmetric monoidal categories, with morphisms the symmetric monoidal left adjoints preserving $\aleph_1$-compact objects. Passing to $\Ind$-objects embeds the $2$-rings of \citet{ABCSS_Higher_Zariski} into this category: a small idempotent-complete stable symmetric monoidal category $\Kk$ yields $\Ind(\Kk) \in \Ring_1(A)$, and exact symmetric monoidal functors yield morphisms (a functor preserving colimits and $\omega$-compact objects preserves $\aleph_1$-compact objects, these being retracts of countable colimits of compact ones).

		We emphasize that the two finiteness regimes genuinely differ, in both directions. On morphisms, $\Ring_1(A)$ is more permissive: a morphism $\Ind(\Kk) \to \Ind(\Ll)$ need only preserve $\aleph_1$-compact objects, so the embedding above is not full. On objects, $\aleph_1$-compactness of algebras (``countable presentation'') is weaker than compact generation in the sense relevant to \emph{loc.\ cit.}; and one cannot blindly transport rigidity arguments, since countable colimits of dualizable objects need not be dualizable. We therefore treat the geometry of \citet{ABCSS_Higher_Zariski} as a guide for the internal theory, not as a special case of it; the precise comparison of the two sites is left open.
	\end{discussion}

	\section{Localizations and immersions}
	\label[section]{sec:Internal_Localizations}

	The following class of maps is the internal counterpart of the Karoubi quotients of \citet{ABCSS_Higher_Zariski}.

	\begin{definition}\label[definition]{def:Internal_Localization}
		A map $q\colon R \to S$ in $\Ring_n(A)$ is a \emph{localization} if the multiplication map
		\[
			S \otimes_R S \longrightarrow S
		\]
		is an isomorphism, i.e.\ if $S$ is an idempotent $R$-algebra.
	\end{definition}

	\begin{observation}\label[observation]{obs:Localizations_Are_Epimorphisms}
		Localizations of $R$ are precisely the epimorphisms out of $R$: the condition that $S \otimes_R S \iso S$ is equivalent to the mapping animae $\Hom_{\Ring_n(A)_{R/}}(S,T)$ being $(-1)$-truncated for all $T$. In particular the localizations of $R$ form a poset. Every symmetric monoidal Bousfield localization which exists in $\Ring_n(A)$ (i.e.\ whose localization functor preserves colimits and $\aleph_1$-compact objects) is a localization in this sense, since factoring through a reflective localization is a property; so are, at height $1$, the $\Ind$-extensions of Karoubi quotients, as well as completions such as $\D(\Z) \to \D(\Z)^{\wedge}_p$. The finiteness of the admissible maps of \citet{ABCSS_Higher_Zariski} is instead imposed by $\aleph_1$-compactness.
	\end{observation}

	The classification of injections from \Cref{chap:Injections_Surjections_Gestalten} identifies this class geometrically.

	\begin{proposition}\label[proposition]{prop:Localizations_Proper_Injections}
		Let $q\colon R \to S$ be a map in $\Ring_n(A)$. The following are equivalent:
		\begin{enumerate}
			\item $q$ is a localization;
			\item the map $Y_S \to Y_R$ is an injection of Gestalten.
		\end{enumerate}
		In this case $Y_S \to Y_R$ is an $n$-proper injection, and the assignment $q \mapsto (Y_S \to Y_R)$ is an equivalence between the poset of localizations of $R$ and the poset of $n$-proper injections into $Y_R$.
	\end{proposition}
	\begin{proof}
		Since $\Gest_X(-)$ is fully faithful and symmetric monoidal, it carries the codiagonal $S \otimes_R S \to S$ to the diagonal $Y_S \to Y_S \times_{Y_R} Y_S$, giving the equivalence of (1) and (2). Under the identification $\CAlg(\Oo(Y_R)_{n+1}) \simeq \Ring_n(A)_{R/}$, a localization of $R$ is exactly an idempotent algebra in $\Oo(Y_R)_{n+1}$, so \Cref{prop:Injective_Prim_Idempotent_Algebras} (applied over the base $Y_R$) shows that $Y_S \to Y_R$ is an $n$-proper injection and that this establishes an equivalence of posets.
	\end{proof}

	Note that no compactness is needed in \Cref{prop:Localizations_Proper_Injections}: idempotency alone already produces the internal adjunctions witnessing $n$-properness. Among all localizations, we single out those which are induced from one categorical level down.

	\begin{definition}\label[definition]{def:Smashing_Internal_Localization}
		Let $n \geq 1$. A localization $q\colon R \to S$ in $\Ring_n(A)$ is \emph{smashing} if the corresponding injection $Y_S \hookrightarrow Y_R$ is $(n-1)$-proper. Equivalently, by \Cref{prop:Injective_Prim_Idempotent_Algebras} applied one level lower, $q$ is classified by an idempotent algebra $E \in \Idem(\Oo(Y_R)_n)$ in the \emph{underlying $n$-category} $\Oo(Y_R)_n \in n\Pr$ of $R \in A_{n+1}$. At height $n = 0$ we declare every localization to be smashing.
	\end{definition}

	\begin{remark}[The height-one picture]\label[remark]{rem:Smashing_Height_One}
		Take $A = \Oo(\Spec(\S))$ and $n = 1$, and let $R \in \Ring_1(A)$ be a stable presentably symmetric monoidal category. Under the canonical equivalence $\Oo(Y_R)_1 \simeq R$, \Cref{def:Smashing_Internal_Localization} recovers the classical notion: a smashing localization of $R$ is one of the form $R \to \Mod_E(R)$ for an idempotent algebra object $E \in \CAlg(R)$ with $E \otimes E \iso E$. Note the bookkeeping of finiteness conditions: the relevant compactness is that of $E$ as an \emph{algebra in $R$}, not as an underlying object of $R$. For example, $B[1/f] \in \CAlg(\D(B))$ is a compact algebra, having the finite presentation $B[t]/(ft-1)$, while its underlying $B$-module is generally not compact. The precise comparison between $\aleph_1$-compactness of $E$ in $\CAlg(\Oo(Y_R)_n)$ and $\aleph_1$-compactness of the localization $R \to S$ in $\Ring_n(A)_{R/}$ is the algebra-to-module compactness bridge for the internal module-endomorphism adjunction, which we treat as an input here (cf.\ \Cref{sec:Zariski_Height_Variation}).
	\end{remark}

	The two classes of maps interleave along the categorical height, in a way that dissolves the dichotomy between Karoubi quotients and smashing localizations:

	\begin{observation}[The localization ladder]\label[observation]{obs:Localization_Ladder}
		Recall from \Cref{sec:Zariski_Height_Variation} below the height-shift embedding $\Ring_n(A) \hookrightarrow \Ring_{n+1}(A)$, $R \mapsto \Mod_R(A_{n+1})$, which does not change the associated Gestalt. Let $q\colon R \to S$ be a localization at height $n$.
		\begin{enumerate}
			\item Viewed at height $n+1$, the map $q$ is a \emph{smashing} localization: the injection $Y_S \hookrightarrow Y_R$ is $n$-proper (\Cref{prop:Localizations_Proper_Injections}), which is the defining condition for smashing at height $n+1$. Its classifying idempotent algebra in $\Oo(Y_R)_{n+1}$ is $S$ itself.
			\item By ambidexterity (\Cref{prop:Ambidexterity_Stefanich_Rings}(2)), the $n$-proper injection $Y_S \hookrightarrow Y_R$ is also an $(n+1)$-étale injection, hence is classified by an idempotent \emph{coalgebra} in $\Oo(Y_R)_{n+2}$ as well (\Cref{prop:Suave_Etale_Idempotent_Coalgebras}); indeed the module category $\Mod_S(A_{n+1})$ is simultaneously an idempotent algebra and an idempotent coalgebra, cf.\ \Cref{lem:Smashing_Idempotent_Coalgebras}.
		\end{enumerate}
		In words: \emph{every localization becomes smashing, and in fact clopen, one categorical level up}. One may view this as a categorified and internal form of the theorem of Miller that finite localizations of rigid $2$-rings are smashing, with rigidity traded for one level of categorification; see \Cref{prop:Rigid_Localizations_Smashing} for the statement at a fixed height.
	\end{observation}

	Over a general base there is also a genuinely different, ``open'' class of admissible maps, invisible in the stable setting of \citet{ABCSS_Higher_Zariski}.

	\begin{definition}\label[definition]{def:Open_Internal_Localization}
		An \emph{open immersion at height $n$} over $R \in \Ring_n(A)$ is an $n$-étale injection $U \hookrightarrow Y_R$. By \Cref{prop:Suave_Etale_Idempotent_Coalgebras}, these are classified by idempotent coalgebras $C \in \coIdem(\Oo(Y_R)_{n+1})$, the corresponding Gestalt being $U_C = \Gest_{Y_R}\bigl(\coMod_C(\Oo(Y_R)_{n+1})\bigr)$.
	\end{definition}

	\begin{remark}\label[remark]{rem:Open_Level_Shift}
		The open $U_C$ is $(n+1)$-affine over $Y_R$, but in general not $n$-affine: opens of height-$n$ affines live one height up. They are, however, again localizations there: an $n$-étale injection is an $(n+1)$-proper injection by \Cref{prop:Ambidexterity_Stefanich_Rings}(4), i.e.\ a localization at height $n+1$. So the single class of localizations at height $n+1$ contains both the closed (\Cref{def:Internal_Localization}) and the open (\Cref{def:Open_Internal_Localization}) admissible maps from height $n$, and the internal Zariski geometry defined below automatically sees both. Complements relate the two classes: the complement of an open immersion is the closed immersion on the complementary idempotent algebra $\cofib(C \to \unit)$ (\Cref{prop:Complement_Construction}), but by the warning following that proposition this passage is irreversible without stability, which can only occur at the lowest level. The appropriate finiteness condition on the open side is likewise not $\aleph_1$-compactness of $C$ as an object: already for a finite localization the acyclization coalgebra can fail to be compact, and one should instead ask for compact generation of $\coMod_C$, or for $\aleph_1$-compactness of the resulting localization at height $n+1$.
	\end{remark}

	\begin{remark}[Hochster duality]\label[remark]{rem:Hochster_Duality_Gestalten}
		The dictionary above inverts the topology familiar from tensor-triangular geometry. In \citet{ABCSS_Higher_Zariski}, the admissible maps $\Kk \to \Kk/\lra{a}$ correspond to the quasicompact \emph{opens} $U(a)$ of the Balmer spectrum; internally, the corresponding localizations are classified by idempotent algebras, i.e.\ by \emph{closed} immersions of Gestalten. Concretely, $\Gest(\Z[\tfrac 1p]) \hookrightarrow \Gest(\Z)$ is a closed immersion (\Cref{ex:Closed_Immersion_Zp}), while $\Spec(\Z[\tfrac 1p]) \subseteq \Spec(\Z)$ is Zariski open; dually, the $p$-complete Gestalt is the open complement (\Cref{ex:Complement_Zp}), while the support-theoretic side of the recollement is closed in the Balmer spectrum. Thus the internal Zariski topology on Gestalten is Hochster dual to the Balmer topology; compare the remark in \citet{ABCSS_Higher_Zariski} that their nomenclature is Hochster dual to the Zariski frame of \citet{Kock_Pitsch}. By \Cref{obs:Localization_Ladder}, the tension disappears one categorical level up, where both classes become clopen.
	\end{remark}

	\section{The internal Zariski geometry}
	\label[section]{sec:Internal_Zariski_Geometry}

	We first record the admissibility axioms.

	\begin{proposition}\label[proposition]{prop:Localization_Admissibility}
		Localizations in $\Ring_n(A)$ are stable under:
		\begin{enumerate}
			\item cobase change: if $R \to S$ is a localization and $R \to T$ is any map, then $T \to T \otimes_R S$ is a localization;
			\item composition;
			\item left cancellation: if $R \to S \to T$ are maps such that $R \to S$ and $R \to T$ are localizations, then $S \to T$ is a localization;
			\item retracts in the arrow category of $\Ring_n(A)$.
		\end{enumerate}
		Moreover, each of these operations preserves the property that all objects involved are $\aleph_1$-compact and that the localizations are $\aleph_1$-compact in the relevant slice.
	\end{proposition}
	\begin{proof}
		(1) Since the tensor product commutes with itself,
		\[
			(T \otimes_R S) \otimes_T (T \otimes_R S) \simeq T \otimes_R (S \otimes_R S) \simeq T \otimes_R S.
		\]
		(2) If $R \to S$ and $S \to T$ are localizations, then $T \otimes_R T \simeq (T \otimes_S T) \otimes_{S \otimes_R S} S \simeq T \otimes_S T \simeq T$.
		(3) Rewriting the tensor product over $S$ in terms of the one over $R$,
		\[
			T \otimes_S T \simeq (T \otimes_R T) \otimes_{S \otimes_R S} S \simeq (T \otimes_R T) \otimes_S S \simeq T \otimes_R T \simeq T,
		\]
		using $S \otimes_R S \simeq S$ and idempotency of $T$ over $R$.
		(4) Idempotency of a map $R \to S$ asserts that a canonical map, the multiplication $S \otimes_R S \to S$, is an isomorphism. A retract diagram in the arrow category induces a retract diagram of multiplication maps, and a retract of an isomorphism is an isomorphism.
		The statements about $\aleph_1$-compactness are \Cref{lem:Kappa_Compact_Algebra_Closure_Properties}: part (2) of that lemma for cobase change, part (3) for composition, part (4) for cancellation, and part (1) for retracts.
	\end{proof}

	Next we define the covers. This is the point where the internal theory departs from \citet{ABCSS_Higher_Zariski}: their covering condition is support-theoretic, whereas the Gestalt formalism suggests a geometric one.

	\begin{definition}\label[definition]{def:Geometric_Cover}
		A finite family $\{q_i\colon R \to S_i\}_{i \in I}$ of $\aleph_1$-compact localizations in $\Ring_n(A)^{\aleph_1}$ is a \emph{geometric cover} if the map
		\[
			\coprod_{i \in I} Y_{S_i} \longrightarrow Y_R
		\]
		is a surjection of Gestalten in the sense of \Cref{def:Surjective_Gestalten}. We write $\tau_n$ for the Grothendieck topology on $\Cc_n(A)$ generated by the geometric covers.
	\end{definition}

	Each $Y_{S_i} \to Y_R$ is an $n$-proper injection, hence $(n+1)$-étale by \Cref{prop:Ambidexterity_Stefanich_Rings}(2); by \Cref{cor:Etale_Countable_Limits} the coproduct is again $(n+1)$-étale over $Y_R$. The covering condition therefore takes place inside the countably toposic category $\Gest^{(n+1)\text{-}\et}_{/Y_R}$ of \Cref{thm:Etale_Countably_Toposic}, where surjections are exactly the covers in the sense of topos theory.

	\begin{theorem}\label[theorem]{thm:Internal_Zariski_Geometry}
		The triple
		\[
			\Gg^{\mathrm{Zar}}_n(A) := \bigl(\Cc_n(A),\ \{\text{$\aleph_1$-compact localizations}\},\ \tau_n\bigr)
		\]
		is a geometry in the sense of \citet{Lurie_DAG_V}. Moreover the geometric covers form a pretopology: they contain the isomorphisms and are stable under base change and composition of families.
	\end{theorem}
	\begin{proof}
		The category $\Cc_n(A)$ is small, idempotent-complete and has finite limits by \Cref{rem:Role_Of_Aleph1}, and the admissibility axioms for the localizations are \Cref{prop:Localization_Admissibility}. It remains to verify the pretopology assertions; the topology $\tau_n$ they generate is then by construction generated by admissible morphisms.

		A family containing an isomorphism is a cover, since isomorphisms are surjective. For stability under base change, let $\{R \to S_i\}$ be a geometric cover and $R \to T$ any map in $\Ring_n(A)^{\aleph_1}$. The base-changed family consists of $\aleph_1$-compact localizations by \Cref{prop:Localization_Admissibility}, and the corresponding map of Gestalten is the pullback of $\coprod_i Y_{S_i} \to Y_R$ along $Y_T \to Y_R$. As in the proof of \Cref{lem:Pullback_Along_Surjection_Conservative}, base change along an arbitrary map commutes with the simplicial colimit computing the image of an eventually suave map, by \Cref{prop:Suave_Countable_Limits}(2); hence pullbacks of surjections along maps of Gestalten are surjections.

		For composition, let $\{R \to S_i\}_i$ be a geometric cover and let $\{S_i \to T_{ij}\}_j$ be geometric covers for each $i$. All maps in sight are $(n+1)$-étale, hence so are the two coproduct maps
		\[
			\coprod_{i,j} Y_{T_{ij}} \longrightarrow \coprod_i Y_{S_i} \longrightarrow Y_R
		\]
		and their composite. By \Cref{prop:Surjective_Criteria}(4), applied at level $n+1$, surjectivity of an $(n+1)$-suave map is equivalent to conservativity of its pullback functor on levels $n+2$; the first map has this property because a finite product of conservative functors is conservative, the second by assumption, and conservativity is stable under composition. Hence the composite family is a geometric cover.
	\end{proof}

	\section{Covers, descent, and structure sheaves}
	\label[section]{sec:Zariski_Covers_Descent}

	The next theorem unwinds what a geometric cover is. It is the internal counterpart of the two main descent theorems of \citet{ABCSS_Higher_Zariski}, descent for the structure presheaf and Zariski descent for module categories, and simultaneously of the classical statement that a Zariski cover of an affine scheme is of effective descent. The point is that, with the geometric definition of covers, all of this is essentially tautological: a surjection of Gestalten is \emph{by definition} a map whose Čech nerve realizes to the target, i.e.\ an effective-descent map for the structure sheaf.

	\begin{theorem}[Covers and descent]\label[theorem]{thm:Zariski_Covers_And_Descent}
		Let $\{q_i\colon R \to S_i\}_{i \in I}$ be a finite family of $\aleph_1$-compact localizations in $\Ring_n(A)^{\aleph_1}$. The following are equivalent:
		\begin{enumerate}
			\item the family is a geometric cover;
			\item (effective descent) for every $m \geq 0$, the canonical map
			\[
				\Oo(Y_R)_m \longrightarrow \lim_{[p] \in \simp} \; \prod_{(i_0,\dots,i_p) \in I^{p+1}} \Oo\bigl(Y_{S_{i_0}} \times_{Y_R} \cdots \times_{Y_R} Y_{S_{i_p}}\bigr)_m
			\]
			is an equivalence;
			\item the pullback functors $\Oo(Y_R)_{n+2} \to \Oo(Y_{S_i})_{n+2}$, $i \in I$, are jointly conservative.
		\end{enumerate}
		In this case, the pullback functors are jointly conservative on \emph{every} level; in particular:
		\begin{enumerate}
			\item[(4)] (level $n$) the functors $q_i\colon \Oo(Y_R)_n \to \Oo(Y_{S_i})_n$ underlying the localizations are jointly conservative, and the Amitsur comparison
			\[
				R \iso \lim_{[p] \in \simp} \; \prod_{(i_0,\dots,i_p)} S_{i_0} \otimes_R \cdots \otimes_R S_{i_p}
			\]
			holds in $\Ring_n(A)$;
			\item[(5)] (level $n+1$) if $M \in \Mod_R(A_{n+1})$ satisfies $M \otimes_R S_i \simeq 0$ for all $i$, then $M \simeq 0$, and the categorified Amitsur comparison
			\[
				\Mod_R(A_{n+1}) \iso \lim_{[p] \in \simp} \; \prod_{(i_0,\dots,i_p)} \Mod_{S_{i_0} \otimes_R \cdots \otimes_R S_{i_p}}(A_{n+1})
			\]
			holds.
		\end{enumerate}
	\end{theorem}
	\begin{proof}
		Write $f\colon Y := \coprod_i Y_{S_i} \to Y_R$; as discussed after \Cref{def:Geometric_Cover}, this is a map in the countably toposic category $\Gest^{(n+1)\text{-}\et}_{/Y_R}$, whose countable colimits are computed in $\Gest_{/Y_R}$, i.e.\ as limits of Stefanich rings, by \Cref{cor:Etale_Countable_Limits}. Since countable colimits in a countably toposic category are universal (\Cref{prop:Countably_Toposic_Characterization}(4)), the Čech nerve of $f$ has $p$-simplices
		\[
			Y^{\times_{Y_R}(p+1)} \simeq \coprod_{(i_0,\dots,i_p)} Y_{S_{i_0}} \times_{Y_R} \cdots \times_{Y_R} Y_{S_{i_p}}.
		\]
		By \Cref{def:Surjective_Gestalten}, the family is a geometric cover if and only if the colimit of this simplicial object in $\Gest$ is $Y_R$, i.e.\ if and only if $\Oo(Y_R)$ is the limit of the corresponding cosimplicial diagram of Stefanich rings. Limits of Stefanich rings under $\Oo(Y_R)$ are computed levelwise (\Cref{prop:Slice_Description}), and levels take coproducts of Gestalten to products; this proves (1) $\Leftrightarrow$ (2). The equivalence with (3) is \Cref{prop:Surjective_Criteria}, (1) $\Leftrightarrow$ (4), applied to the $(n+1)$-suave map $f$ at level $n+1$.

		Now assume (2). Each level of the limit is a limit of categories, in which isomorphisms are detected on the entries of the diagram; restricting attention to the entries with $p = 0$ gives joint conservativity at every level. For the displayed formula in (4), it remains to identify the levels of the Čech objects: the maps $Y_{S_i} \to Y_R$ are $n$-prim (\Cref{prop:Localizations_Proper_Injections}), so the Künneth formula for prim maps (\Cref{lem:Kunneth_Prim}) computes the relative level-$n$ algebra of $Y_{S_{i_0}} \times_{Y_R} \cdots \times_{Y_R} Y_{S_{i_p}}$ over $Y_R$ as $S_{i_0} \otimes_R \cdots \otimes_R S_{i_p}$. The formula in (5) is the same identification one level up, using $\Oo(Y_{S})_{n+1} \simeq \Mod_S(A_{n+1})$ from \Cref{prop:Affineness_Properties}(2), and the vanishing statement follows since an object of a limit of categories vanishes if all its images do (its identity being detected there).
	\end{proof}

	\begin{remark}[Structure sheaves of internal $n$-rings]\label[remark]{rem:Structure_Sheaf_Internal}
		\Cref{thm:Zariski_Covers_And_Descent} says that the tautological presheaf
		\[
			\Oo\colon \bigl(\Cc_n(A)_{/R}\bigr)\catop \longrightarrow \Ring_n(A), \qquad S \longmapsto S,
		\]
		together with all of its categorified levels $S \mapsto \Oo(Y_S)_m$, is a sheaf for the topology $\tau_n$, without any rigidity hypothesis. This inverts the logic of \citet{ABCSS_Higher_Zariski}: there, covers are defined support-theoretically and the sheaf property of the structure presheaf on the spectrum of a rigid $2$-ring is the main theorem; here, covers are defined geometrically, the sheaf property is automatic, and the substance migrates into the question of \emph{which} support-theoretically defined families are geometric covers. Global sections over $R$ tautologically recover $R$. What is not yet constructed here is the internal analogue of the absolute spectrum and its adjunction with locally ringed objects; see \Cref{disc:Zariski_Outlook}.
	\end{remark}

	The support-theoretic covering conditions are the following.

	\begin{definition}\label[definition]{def:Conservative_Cover}
		A finite family $\{q_i\colon R \to S_i\}$ of localizations in $\Ring_n(A)$ is:
		\begin{enumerate}
			\item \emph{conservative} if the functors $\Oo(Y_R)_n \to \Oo(Y_{S_i})_n$ underlying the $q_i$ are jointly conservative;
			\item \emph{nil-conservative} if every map $u$ in $\Oo(Y_R)_n$ inverted by all $q_i$ is tensor-nilpotent, meaning that some iterated pushout product $u^{\mathbin{\square}m}$ is an isomorphism.
		\end{enumerate}
	\end{definition}

	\begin{remark}
		Condition (2) is the direct internalization of the covers of \citet{ABCSS_Higher_Zariski}: when $\Oo(Y_R)_n$ is stable one has $\cofib(u^{\mathbin{\square}m}) \simeq \cofib(u)^{\otimes m}$, so nil-conservativity says exactly that the common kernel of the $q_i$ consists of tensor-nilpotent objects. Every conservative family is nil-conservative; the converse holds whenever tensor-nilpotent maps are isomorphisms, e.g.\ in the rigid stable setting, where a tensor-nilpotent object with a dual vanishes by the usual retract argument. We do not know an internal version of this retract argument at higher height.
	\end{remark}

	By \Cref{thm:Zariski_Covers_And_Descent}, every geometric cover is conservative. The converse is the internal analogue of the Mayer--Vietoris descent theorem of \citet{ABCSS_Higher_Zariski} (in turn a coherent enhancement of gluing results of Balmer--Favi). At height $0$ it holds without any rigidity hypothesis:

	\begin{theorem}\label[theorem]{thm:Height_Zero_Covers}
		Let $X = \Gest(A)$ be a Gestalt over $\Spec(\S)$, so that $A_1$ is stable, and let $\{q_i\colon R \to S_i\}_{i=1}^{k}$ be a finite family of $\aleph_1$-compact localizations in $\Ring_0(A)^{\aleph_1}$. Then the family is a geometric cover if and only if it is conservative, i.e.\ iff the base-change functors $\Mod_R(A_1) \to \Mod_{S_i}(A_1)$ are jointly conservative.
	\end{theorem}
	\begin{proof}
		One direction is \Cref{thm:Zariski_Covers_And_Descent}. (At height $0$ the two readings of conservativity agree: level $n = 0$ conservativity concerns the anima-level functors, and the module-level statement (5) is the natural stable formulation; we prove the converse from the latter, which is the weaker hypothesis.)

		Conversely, assume the family is conservative, and work in the stable symmetric monoidal category $\Mod_R(A_1)$. Set $I_i := \fib(R \to S_i)$. Tensoring together the fiber sequences $I_i \to R \to S_i$ produces a $k$-cube $T \mapsto S_T := \bigotimes^R_{i \in T} S_i$ (with $S_{\emptyset} = R$) whose total fiber is $I_1 \otimes_R \cdots \otimes_R I_k$. Since $q_i$ is a localization, $I_i \otimes_R S_i = \fib(S_i \to S_i \otimes_R S_i) = 0$, so the total fiber dies after base change to every $S_i$ and hence vanishes by joint conservativity. Thus the cube is cartesian:
		\[
			R \iso \lim_{\emptyset \neq T \subseteq \{1,\dots,k\}} S_T.
		\]
		Each $S_T$ with $T \neq \emptyset$ is a module over $S := \prod_i S_i$, and any $S$-module $M$ is a retract of $S \otimes_R M$; since the limit above is finite and $\Mod_R(A_1)$ is stable, it follows that $R$ lies in the thick tensor ideal generated by $S$, i.e.\ that $S$ is descendable in the sense of \Cref{def:Descendable}.

		Finally, since $A_1$ is stable, the finite product $S = \prod_i S_i$ splits off complementary idempotents at every level, so that $Y_S \simeq \coprod_i Y_{S_i}$. The map $Y_S \to Y_R$ is $0$-affine and countably presented, hence $0$-prim and eventually suave (\Cref{prop:Prim_Properties}(1), \Cref{prop:Proper_Properties}(1), \Cref{prop:Ambidexterity_Stefanich_Rings}), and its pushforward of the unit is the descendable algebra $S$. \Cref{cor:Descendable_Implies_Surjective} shows that it is surjective, i.e.\ the family is a geometric cover.
	\end{proof}

	\begin{warning}\label[warning]{warn:Conservative_Not_Geometric}
		At height $n \geq 1$, conservative families need not be geometric covers, and rigidity genuinely enters. The counterexample is already in \citet{ABCSS_Higher_Zariski}: for the non-rigid $2$-ring $\Fun([1],\Sp)^{\omega}$, the source and target evaluations to $\Sp^{\omega}$ are Karoubi quotients forming a nil-conservative --- indeed conservative --- family, but descent fails, as it would identify $\Fun([1],\Sp)$ with $\Sp \times \Sp$. Viewed internally at height $1$ over $\Spec(\S)$, this is a conservative family of $\aleph_1$-compact localizations which is not a geometric cover, by \Cref{thm:Zariski_Covers_And_Descent}(4).
	\end{warning}

	\begin{question}\label[question]{quest:Higher_Height_Descent}
		At height $n \geq 1$, under which internal rigidity hypotheses does conservativity imply geometric covering? The proof of \Cref{thm:Height_Zero_Covers} breaks down because $A_{n+1}$ is not stable, so the total-fiber argument is unavailable; the levels $A_{m}$ for $m \geq 2$ are, however, semiadditive, and one may hope for an ambidexterity-based replacement for the fracture cube. \Cref{thm:Rigid_Covers_Conservative} below settles the module-level version of this question; what remains open is the comparison with conservativity at the level of the underlying objects $\Oo(Y_R)_n$, which at height $1$ is the Mayer--Vietoris theorem of \citet{ABCSS_Higher_Zariski}.
	\end{question}

	\section{Rigidity}
	\label[section]{sec:Zariski_Rigid}

	\begin{definition}\label[definition]{def:Rigid_Internal_N_Ring}
		For $n \geq 1$, an internal $n$-ring $R \in \Ring_n(A)$ is \emph{rigid} if the structure map $Y_R \to X$ is $(n-1)$-proper, i.e.\ if the Stefanich $A$-algebra $\Oo(Y_R)$ is $(n-1)$-proper over $A$. At height $n = 0$ we declare every internal $0$-ring to be rigid. We write $\Cc_n^{\mathrm{rig}}(A) \subseteq \Cc_n(A)$ for the full subcategory of rigid objects.
	\end{definition}

	\begin{discussion}[Shifted rigidity]\label[discussion]{disc:Shifted_Rigidity}
		Note the shift by one: for a countably presented $n$-affine algebra, $n$-properness is automatic (\Cref{prop:Proper_Properties}(1)), and $(n-1)$-properness is the genuine extra condition. At height $n = 1$, \Cref{prop:Proper_Rigid} identifies it with rigidity of $R \in \CAlg(A_2)$ in the sense of Gaitsgory--Rozenblyum: the unit and the multiplication of $R$ admit the relevant internal right adjoints. For $A = \Oo(\Spec(\S))$ and $R$ compactly generated this is closely related to the condition of \citet{ABCSS_Higher_Zariski} that every compact object of $R$ be dualizable. The proof of \Cref{prop:Proper_Rigid} is levelwise and suggests the analogous identification at every height, with $(n-1)$-properness corresponding to internal rigidity of $R$ as an object of $\CAlg(A_{n+1})$; we have not checked this and use $(n-1)$-properness as the definition, since it is intrinsic to the Gestalt formalism and supplies all iterated codiagonal conditions at once. The convention at height $0$ is consistent with this: $0$-properness of countably presented $0$-affines is automatic.
	\end{discussion}

	\begin{proposition}\label[proposition]{prop:Rigid_Maps_Proper}
		Let $n \geq 1$ and let $R, S \in \Ring_n(A)$ be rigid. Then every map $R \to S$ induces an $(n-1)$-proper map $Y_S \to Y_R$ of Gestalten.
	\end{proposition}
	\begin{proof}
		Both composites $Y_S \to X$ and $Y_R \to X$ are $(n-1)$-proper; apply the cancellation statement of \Cref{prop:Proper_Properties}(4).
	\end{proof}

	\begin{proposition}[Internal Miller theorem]\label[proposition]{prop:Rigid_Localizations_Smashing}
		Let $n \geq 1$. Every localization $q\colon R \to S$ between rigid internal $n$-rings is smashing.
	\end{proposition}
	\begin{proof}
		By \Cref{prop:Localizations_Proper_Injections}, the map $Y_S \to Y_R$ is an injection, and by \Cref{prop:Rigid_Maps_Proper} it is $(n-1)$-proper. An $(n-1)$-proper injection is a smashing localization by \Cref{def:Smashing_Internal_Localization}.
	\end{proof}

	\begin{remark}\label[remark]{rem:Projection_Formula_Reading}
		\Cref{prop:Rigid_Localizations_Smashing} internalizes the theorem of Miller that finite localizations of rigid $2$-rings are smashing, which in \citet{ABCSS_Higher_Zariski} enters through the criterion that a symmetric monoidal Bousfield localization is smashing if and only if it satisfies the projection formula. The internal counterpart of that criterion is visible here as well: an $(n-1)$-proper map is in particular $(n-1)$-prim, and the prim projection formula (\Cref{lem:Prim_Projection_Formula}) identifies $M \otimes_R q_*(N) \simeq q_*(q^*(M) \otimes_S N)$ at level $n-1$. Note that, unlike in \emph{loc.\ cit.}, no compactness hypothesis is needed: rigidity of the two endpoints suffices.
	\end{remark}

	Rigidity also lowers by one the level at which the covering condition can be tested, which is the internal counterpart of the rigid descent theorem of \citet{ABCSS_Higher_Zariski}:

	\begin{theorem}\label[theorem]{thm:Rigid_Covers_Conservative}
		Let $n \geq 1$ and let $\{q_i\colon R \to S_i\}_{i \in I}$ be a finite family of localizations in $\Ring_n(A)$ with $R$ and all $S_i$ rigid. The following are equivalent:
		\begin{enumerate}
			\item the map $\coprod_i Y_{S_i} \to Y_R$ is a surjection of Gestalten;
			\item the base-change functors $\Mod_R(A_{n+1}) \to \Mod_{S_i}(A_{n+1})$, $M \mapsto M \otimes_R S_i$, are jointly conservative.
		\end{enumerate}
	\end{theorem}
	\begin{proof}
		By \Cref{prop:Rigid_Maps_Proper} each $Y_{S_i} \to Y_R$ is $(n-1)$-proper, hence $n$-étale by \Cref{prop:Ambidexterity_Stefanich_Rings}(2), and by \Cref{cor:Etale_Countable_Limits} the coproduct map $f\colon \coprod_i Y_{S_i} \to Y_R$ is $n$-étale, in particular $n$-suave. Now \Cref{prop:Surjective_Criteria}, (1) $\Leftrightarrow$ (4), applied to $f$ at level $n$, states that $f$ is surjective if and only if $f_n^*\colon \Oo(Y_R)_{n+1} \to \prod_i \Oo(Y_{S_i})_{n+1}$ is conservative. Under the identifications $\Oo(Y_R)_{n+1} \simeq \Mod_R(A_{n+1})$ and $\Oo(Y_{S_i})_{n+1} \simeq \Mod_{S_i}(A_{n+1})$ of \Cref{prop:Affineness_Properties}(2), the components of $f_n^*$ are the base-change functors of the statement.
	\end{proof}

	\begin{remark}
		Compare with \Cref{thm:Zariski_Covers_And_Descent}(3), where the covering condition for a general family is detected one level higher, on $\Oo(-)_{n+2}$; rigidity moves it down to $\Oo(-)_{n+1}$, i.e.\ to modules over the internal $n$-rings themselves. The same proof works verbatim for countable families. What rigidity does \emph{not} formally provide is the further descent to level $n$, i.e.\ the comparison with conservativity of the underlying functors $q_i$ on $\Oo(Y_R)_n$ (\Cref{def:Conservative_Cover}); at height $1$ over $\Spec(\S)$ this last comparison is the rigid Mayer--Vietoris theorem of \citet{ABCSS_Higher_Zariski}, and we expect their argument to internalize, but we have not carried this out; cf.\ \Cref{quest:Higher_Height_Descent}.
	\end{remark}

	\begin{question}\label[question]{quest:Rigid_Closure}
		Is a localization of a rigid internal $n$-ring again rigid? For $2$-rings this is the statement that Karoubi quotients of rigid $2$-rings are rigid, proved in \citet{ABCSS_Higher_Zariski}; the internal statement would ensure that $\Cc_n^{\mathrm{rig}}(A)$, with the smashing localizations (\Cref{prop:Rigid_Localizations_Smashing}) and the geometric covers, is itself a geometry, sitting inside $\Gg^{\mathrm{Zar}}_n(A)$.
	\end{question}

	\begin{discussion}[Realization and density]\label[discussion]{disc:Realization_Density}
		Since maps between rigid objects are $(n-1)$-proper, hence $n$-étale, the affine embedding restricts to
		\[
			j_n^{\mathrm{rig}}\colon \Cc_n^{\mathrm{rig}}(A) \longrightarrow \Gest^{n\text{-}\et}_{/X},
		\]
		landing in the countably toposic category of $n$-étale Gestalten over $X$. By construction, $j_n^{\mathrm{rig}}$ sends the Čech nerves of geometric covers to colimit diagrams. Left Kan extension along the Yoneda embedding therefore yields a realization functor
		\[
			\Shv_{\tau_n}\bigl(\Cc_n^{\mathrm{rig}}(A)\bigr)^{\aleph_1} \longrightarrow \Gest^{n\text{-}\et}_{/X},
		\]
		whose existence and image are governed by \Cref{thm:Sheaf_Valued_In_Etale}. The internal analogue of the full faithfulness of the spectrum functor on rigid $2$-rings is then a density question: are the rigid $n$-affine objects dense in the subtopos of the $n$-étale topos over $X$ that they generate, and does every $n$-étale Gestalt in this subtopos admit an effective-epimorphic resolution by rigid affines? We leave this as an open question; a positive answer would identify sheaves on the internal Zariski site with a localic piece of $\Gest^{n\text{-}\et}_{/X}$, in parallel with the identification of the Zariski spectrum of a $2$-ring with the sheaves on its Balmer spectrum.
	\end{discussion}

	\section{Variation of the categorical height}
	\label[section]{sec:Zariski_Height_Variation}

	The internal higher rings at successive heights are related by the fully faithful internal module functors
	\[
		\Mod^A_{n+1}\colon \Ring_n(A) = \CAlg(A_{n+1}) \hookrightarrow \CAlg(A_{n+2}) = \Ring_{n+1}(A)
	\]
	of \Cref{sec:Recollections_And_Notations}. On affine Gestalten this transition is invisible: an $n$-affine object is $(n+1)$-affine, with level-$(n+1)$ algebra the module category, so
	\[
		j_{n+1} \circ \bigl(\Mod^A_{n+1}\bigr)\catop \simeq j_n.
	\]
	Consequently:
	\begin{enumerate}
		\item localizations at height $n$ are carried to smashing localizations at height $n+1$, with classifying idempotent algebra the localization itself (\Cref{obs:Localization_Ladder});
		\item rigid objects are carried to rigid objects, since $(n-1)$-proper implies $n$-proper;
		\item geometric covers are literally unchanged, being conditions on the underlying map of Gestalten.
	\end{enumerate}
	The one assertion which is not formal is that $\Mod^A_{n+1}$ preserves $\aleph_1$-compactness, i.e.\ that a countably presented internal $n$-ring is countably presented as an internal $(n+1)$-ring. This algebra-to-module compactness bridge is used by Scholze and we treat it as an input here, cf.\ \Cref{sec:Compact_Algebras}. Granting it, the transitions define morphisms of geometries
	\[
		\Gg^{\mathrm{Zar}}_n(A) \longrightarrow \Gg^{\mathrm{Zar}}_{n+1}(A),
	\]
	and the realization functors of \Cref{disc:Realization_Density} are compatible with them, since the underlying Gestalten do not change.

	\begin{discussion}[Outlook]\label[discussion]{disc:Zariski_Outlook}
		We conclude with the questions this picture suggests, beyond those already recorded.
		\begin{enumerate}
			\item \emph{The internal Zariski spectrum.} \Cref{thm:Internal_Zariski_Geometry} feeds the machinery of \citet{Lurie_DAG_V}: there results an absolute spectrum functor from $\Ind$-objects of $\Cc_n(A)\catop$ to topoi with local $\Gg^{\mathrm{Zar}}_n(A)$-structure. For $A = \Oo(\Spec(\S))$, $n = 1$, and rigid compactly generated inputs, the poset of localizations is governed, after the Hochster duality of \Cref{rem:Hochster_Duality_Gestalten}, by the lattice of thick tensor ideals, and one expects the underlying locale of the spectrum to recover the Balmer spectrum via \citet{Kock_Pitsch}; at $\aleph_1$ and at higher height it is an intrinsic replacement for it. It would also be interesting to identify the internal analogue of a \emph{local} $2$-ring: plausibly, an internal $n$-ring all of whose geometric covers split.
			\item \emph{Principality and telescope-type questions.} Call a localization \emph{principal} if it is universal among maps inverting a single $\aleph_1$-compact morphism of $\Oo(Y_R)_n$. Principal localizations should be $\aleph_1$-compact; the extent to which the converse fails is a coherence question of telescope type, and the stalk-locality results in the appendix of \citet{ABCSS_Higher_Zariski} suggest what its internal formulation should look like.
			\item \emph{Stabilization of the tower.} The presentation $\StRing_{A/} \simeq \colim_m \CAlg(A_m)$ exhibits the tower of geometries $\Gg^{\mathrm{Zar}}_0(A) \to \Gg^{\mathrm{Zar}}_1(A) \to \cdots$ as an exhaustion of an all-height geometry on countably presented Stefanich $A$-algebras. By \Cref{obs:Localization_Ladder} every admissible map becomes clopen after one step, so the colimit geometry is generated by clopen immersions; whether the associated spectra stabilize on objects coming from a fixed finite height, and what the resulting limit anima is for classical inputs, we do not know.
		\end{enumerate}
	\end{discussion}

	\addtocontents{toc}{\vskip 10pt}
%	\addcontentsline{toc}{chapter}{Bibliography}
	\printbibliography[heading=bibintoc]
	
\end{document}
